% Propietario conservado: /Users/ruben/Documents/New project/output/EXTREMA_MEDIA_RAZON_RECIPROCIDAD_ES_EN_20260918/ARTICULO/ES/source/nuclear/antecedentes_fisicos.tex
\chapter{Réalisation canonique de la covariance}
\label{x_reciprocidad:nuclear:covarianza-antecedente}

% BEGIN OWNER /Users/ruben/Documents/New project/output/EXTREMA_MEDIA_RAZON_RECIPROCIDAD_ES_EN_20260918/ARTICULO/ES/source/antecedentes_ii/sections/07c_covarianza_accion.tex
% Propuesta de inserción después de 07b_geometria_elipse.tex.
% RESULTADO_RECUPERADO: c52_06_01_incertidumbre.tex, líneas 10--91.
% Dependencias: c52_04_01_ley_areal.tex, líneas 137--198;
% c52_06_02_intercambio_constitutivo.tex, líneas 5--29.
% Explicitación local: dominio operatorio, prueba de covarianza,
% transporte recíproco y distinción dimensional posición--momento.
% Este fragmento no introduce una constante ni un generador adicional.

\section{Covariance canonique et réciprocité de l'ellipse d'action}
\label{x_reciprocidad:sec:ii-area-covarianza}

La construction précédente détermine deux invariants complémentaires :
le produit des demi-axes fixe l'aire d'action et leur quotient conserve
l'anisotropie du couple angulaire. Précisons maintenant la réalisation canonique
dans laquelle ces demi-axes valent deux écarts-types. La structure de
départ est l'ellipse déjà obtenue à partir des publications angulaires et
d'action d'APP--TRIT--TPK ; la réalisation opératoire est établie après
cette construction.

Dans ce paragraphe, on écrit \(\eta=\eta_{\rm el}
=\operatorname{artanh}(C^*/A)\), avec \(|C^*/A|<1\) et
\(\hbar>0\). Le quotient angulaire utilise la même unité pour ses deux
composantes. Dans les coordonnées équilibrées \((Q,P)\), toutes deux de
dimension racine carrée d'action, définissons
\begin{equation}
 V_\eta:=\frac14
 \begin{pmatrix}a_S^2&0\\0&b_S^2\end{pmatrix}
 =\frac{\hbar}{2}
 \begin{pmatrix}e^\eta&0\\0&e^{-\eta}\end{pmatrix}.
 \label{x_reciprocidad:eq:ii-covarianza-geometrica}
\end{equation}
Cette matrice positive possède des entrées de dimension action. Son interprétation
comme covariance quantique est spécifiée par la réalisation suivante.

\begin{proposition}[Réalisation gaussienne de la covariance d'action]
\label{x_reciprocidad:prop:ii-covarianza-gaussiana}
Considérons la représentation canonique dans \(L^2(\mathbb R,dQ)\), avec
\(\widehat Qf(Q)=Qf(Q)\) et
\(\widehat Pf(Q)=-i\hbar f'(Q)\), initialement sur l'espace de
Schwartz \(\mathcal S(\mathbb R)\). L'état normalisé
\begin{equation}
 \psi_\eta(Q)=\bigl(\pi\hbar e^\eta\bigr)^{-1/4}
 \exp\!\left(-\frac{Q^2}{2\hbar e^\eta}\right)
 \label{x_reciprocidad:eq:ii-gaussiana-accion}
\end{equation}
a des moyennes nulles et pour matrice de covariance \(V_\eta\). En particulier,
\begin{equation}
 (\Delta Q)^2=\frac{\hbar e^\eta}{2},\qquad
 (\Delta P)^2=\frac{\hbar e^{-\eta}}{2},\qquad
 \operatorname{Cov}(Q,P)=0,\qquad
 \det V_\eta=\frac{\hbar^2}{4}.
 \label{x_reciprocidad:eq:ii-varianzas-accion}
\end{equation}
Cette réalisation sature l'inégalité de Robertson--Schrödinger.
\end{proposition}

\begin{proof}
Les opérateurs précédents sont essentiellement autoadjoints sur
\(\mathcal S(\mathbb R)\), qui constitue un cœur commun invariant ;
sur cet espace, on a
\([\widehat Q,\widehat P]=i\hbar I\). Pour
\(\widehat Y=(\widehat Q,\widehat P)\), la covariance symétrique est
\[
 V_{jk}=\frac12\left\langle
 (\widehat Y_j-\langle\widehat Y_j\rangle)
 (\widehat Y_k-\langle\widehat Y_k\rangle)
 +(\widehat Y_k-\langle\widehat Y_k\rangle)
 (\widehat Y_j-\langle\widehat Y_j\rangle)
 \right\rangle.
\]
Avec \(s=\hbar e^\eta>0\), soit
\(I_s=\int_{\mathbb R}e^{-Q^2/s}\,dQ\). Le produit de deux
intégrales identiques, évalué en coordonnées polaires, donne
\(I_s^2=2\pi\int_0^\infty r e^{-r^2/s}\,dr=\pi s\).
En outre, l'intégrale de
\(\frac{d}{dQ}(Qe^{-Q^2/s})\) est nulle, par décroissance aux deux
extrémités ; d'où
\(\int_{\mathbb R}Q^2e^{-Q^2/s}\,dQ=sI_s/2\).
Ces identités et la parité démontrent la normalisation, la moyenne nulle et
\(\langle\widehat Q^2\rangle=s/2\). De plus,
\(\widehat P\psi_\eta=i e^{-\eta}Q\psi_\eta\).
Ainsi, \(\langle\widehat P\rangle=0\),
\(\|\widehat P\psi_\eta\|^2=\hbar e^{-\eta}/2\) et
\(\operatorname{Re}\langle\widehat Q\psi_\eta,
\widehat P\psi_\eta\rangle=0\).

Pour tout état normalisé de ce domaine, Cauchy--Schwarz appliquée
à \((\widehat Q-\langle\widehat Q\rangle)\psi\) et
\((\widehat P-\langle\widehat P\rangle)\psi\) donne
\[
 (\Delta Q)^2(\Delta P)^2
 \geq \operatorname{Cov}(Q,P)^2
       +\frac14\left|\langle[\widehat Q,\widehat P]\rangle\right|^2
 =\operatorname{Cov}(Q,P)^2+\frac{\hbar^2}{4}.
\]
Les variances calculées atteignent l'égalité. De manière équivalente,
l'opérateur
\[
 a_\eta=\frac{e^{-\eta/2}\widehat Q
                  +i e^{\eta/2}\widehat P}{\sqrt{2\hbar}}
\]
satisfait \([a_\eta,a_\eta^\dagger]=I\) et
\(a_\eta\psi_\eta=0\), par substitution directe.
\end{proof}

\begin{proposition}[Aire, dispersion et transport réciproque]
\label{x_reciprocidad:prop:ii-area-covarianza-reciprocidad}
L'ellipse d'action est exactement le sous-niveau de covariance
\begin{equation}
 \mathcal E_\eta
 =\left\{y=(Q,P)^{\mathsf T}:y^{\mathsf T}V_\eta^{-1}y\leq4\right\}
 =\left\{\frac{Q^2}{a_S^2}+\frac{P^2}{b_S^2}\leq1\right\}.
 \label{x_reciprocidad:eq:ii-elipse-covarianza}
\end{equation}
Son aire est \(h=2\pi\hbar\), ses demi-axes sont
\(a_S=2\Delta Q\), \(b_S=2\Delta P\) et
\begin{equation}
 \frac{\Delta P}{\Delta Q}=e^{-\eta}
 =\frac{b_S}{a_S}=R_{\rm el}^2.
 \label{x_reciprocidad:eq:ii-covarianza-forma}
\end{equation}
La rotation canonique \(J_2\) de la section~\ref{x_reciprocidad:iv:ellipse:section} transporte
\(V_\eta\) vers \(V_{-\eta}\) et conserve la forme symplectique.
\end{proposition}

\begin{proof}
L'inversion de la matrice diagonale
\eqref{x_reciprocidad:eq:ii-covarianza-geometrica} donne
\eqref{x_reciprocidad:eq:ii-elipse-covarianza}. Ses demi-axes sont les racines positives de
quatre fois ses variances, et
\[
 \operatorname{Area}(\mathcal E_\eta)
 =4\pi\sqrt{\det V_\eta}=2\pi\hbar=h,
\]
en accord avec la loi aréale~\eqref{x_reciprocidad:iv:ellipse:area-fase}.
Le rapport des racines positives démontre
\eqref{x_reciprocidad:eq:ii-covarianza-forma}.
Si \(M_\eta=\operatorname{diag}(e^{\eta/2},e^{-\eta/2})\), alors
\[
 V_\eta=M_\eta V_0M_\eta^{\mathsf T},\qquad
 M_\eta^{\mathsf T}J_2M_\eta=J_2,\qquad
 J_2V_\eta J_2^{\mathsf T}=V_{-\eta}.
\]
Ces identités s'obtiennent par multiplication. En termes opératoires,
\((\widehat Q',\widehat P')=(-\widehat P,\widehat Q)\) conserve
\([\widehat Q',\widehat P']=i\hbar I\). La permutation
\(S_2:(Q,P)\mapsto(P,Q)\) échange également les variances, mais inverse
le signe du commutateur et de la forme symplectique. Le transport de
covariance respecte ainsi la distinction entre réciprocité canonique et inversion
d'orientation déjà établie à la section~\ref{x_reciprocidad:iv:ellipse:section}.
\end{proof}

\subsubsection{Position, quantité de mouvement et transduction dimensionnelle.}
Soit \(\lambda_{xp}>0\) un facteur de la carte dimensionnelle, de dimension
quantité de mouvement par unité de longueur. La transformation
\[
 Q=\sqrt{\lambda_{xp}}\,x,\qquad
 P=\frac{p_x}{\sqrt{\lambda_{xp}}}
\]
réalise le couple équilibré dans une position \(x\) et sa quantité de mouvement conjuguée
\(p_x\) ; elle conserve \(dQ\wedge dP=dx\wedge dp_x\) et
\([\widehat x,\widehat p_x]=i\hbar I\). Par transformation de la
covariance,
\begin{equation}
 V_\eta^{(x,p_x)}=\frac{\hbar}{2}
 \begin{pmatrix}e^\eta/\lambda_{xp}&0\\
                 0&\lambda_{xp}e^{-\eta}\end{pmatrix},\qquad
 \Delta x\,\Delta p_x=\frac{\hbar}{2},\qquad
 \frac{\Delta p_x}{\Delta x}=\lambda_{xp}e^{-\eta}.
 \label{x_reciprocidad:eq:ii-covarianza-dimensiones}
\end{equation}
Le quotient de dispersion acquiert ici la dimension quantité de mouvement par longueur.
En revanche, \eqref{x_reciprocidad:eq:ii-covarianza-forma} appartient à la carte équilibrée.
Pour la réalisation inductive--capacitive déjà définie,
\eqref{x_reciprocidad:hmtII:eq:ii-elipse-impedancia} apporte l'identité adimensionnelle
\(Z_\eta/Z_*=e^{-\eta}=\Delta P/\Delta Q\).
Chaque lecture utilise sa transduction déclarée et conserve le même
déterminant d'action.

\subsubsection{Énergie du contour et énergie moyenne de l'état.}
La réalisation inductive--capacitive possède le hamiltonien
\[
 H_\eta(Q,P)=\frac{\omega}{2}
 \left(e^{-\eta}Q^2+e^\eta P^2\right),\qquad \omega>0.
\]
Sur \(\partial\mathcal E_\eta\), sa valeur est
\(H_\eta=\hbar\omega\), d'après les équations des demi-axes. Dans la
réalisation quantique de la proposition~\ref{x_reciprocidad:prop:ii-covarianza-gaussiana},
\begin{equation}
 \langle\widehat H_\eta\rangle_{\psi_\eta}
 =\frac{\omega}{2}
 \left(e^{-\eta}\frac{\hbar e^\eta}{2}
       +e^\eta\frac{\hbar e^{-\eta}}{2}\right)
 =\frac{\hbar\omega}{2}.
 \label{x_reciprocidad:eq:ii-energia-contorno-covarianza}
\end{equation}
Le contour de deux écarts-types possède donc le double de
l'énergie moyenne de cet état gaussien. La loi aréale fixe la normalisation
d'action ; le couple angulaire fixe sa distribution entre les deux coordonnées ;
la représentation canonique et l'état explicite réalisent cette distribution
comme covariance. La saturation de l'incertitude est ainsi établie dans
la réalisation spécifiée par ses opérateurs et son état.

% END OWNER


% BEGIN OWNER /Users/ruben/Documents/New project/output/EXTREMA_MEDIA_RAZON_RECIPROCIDAD_ES_EN_20260918/ARTICULO/ES/source/antecedentes_iii/sections/03_hojas_y_orientacion.tex
\chapter[Retour de feuillet, incidence aréale--volumétrique et orientation constitutive]{Retour de feuillet, incidence aréale--volumétrique\\et orientation constitutive}
\label{x_reciprocidad:iii:sec:hojas}
\label{x_reciprocidad:nuclear:constitutiva-antecedente}

La structure des feuillets intervient avant l’évaluation des grandeurs
électromagnétiques. Sa fonction consiste à distinguer la fermeture d’une projection,
la restitution d’une orientation et l’évolution du registre transporté.
Ces opérations agissent sur des objets liés, mais de types différents : une
trajectoire projetée peut s’être refermée tandis que son relèvement reste
sur le feuillet opposé ; l’orientation peut avoir été restituée tandis que la mémoire
conserve les deux tours effectués. La construction constitutive doit recevoir
cette information avec la structure qui la détermine.

\section{Retour projeté et restitution du revêtement orienté}

Soit $\ell$ une boucle de retour d’une extension survivante et soit
$\varepsilon:\langle\ell\rangle\to\{+1,-1\}$ le caractère de son
relèvement, avec $\varepsilon(\ell)=-1$. Le fibré d’intervalles associé
a pour présentation
\begin{equation}
 \mathcal M=([0,1]\times[-1,1])/{(0,u)\sim(1,-u)}.
 \label{x_reciprocidad:iii:eq:mobius}
\end{equation}
La donnée décisive est le signe du transport sur la boucle spécifiée.
L’identification des extrémités réalise géométriquement ce caractère.

\begin{proposition}[Deux ordres de clôture]
Un tour de la boucle ferme la base et inverse l’orientation de la fibre. Deux
tours restituent l’orientation. Dans la paramétrisation angulaire du lecteur de
bord, ces retours correspondent respectivement à la clôture aréale $2\pi$ et
à la clôture volumétrique $4\pi$ du revêtement orienté.
\end{proposition}
\begin{proof}
La relation d’équivalence de \eqref{x_reciprocidad:iii:eq:mobius} identifie la fibre terminale
à la fibre initiale au moyen de $u\mapsto-u$. La composition de deux transports est
$u\mapsto-(-u)=u$. Paramétrer un tour de la base par $2\pi$ attribue $4\pi$
à son carré. L’assertion concerne ce relèvement et cette lecture.
\end{proof}

La mémoire du TPK conserve en outre la trajectoire parcourue et les mises à jour du registre. Le carré du caractère de signe est l'identité, mais n'impose pas l'annulation de ces mises à jour. La représentation spinorielle, lorsqu'elle sera utilisée, devra transporter cette donnée au moyen de sa propre application.

\section{Descente de l’incidence depuis le calendrier tritique}

Dans le bloc primitif $(+1,-1,0)$, le TPK met à jour le mode arithmétique au moyen de
$+1:m\mapsto\Sigma$, $-1:m\mapsto\Pi$ et $0:m\mapsto m$. Son orbite cyclique
est $(\Sigma,\Pi,\Pi)$. Le lecteur de feuillets attribue $\Sigma$ au feuillet aréal et
$\Pi$ au feuillet volumétrique. Dans l’ordre des bases $(e_{\rm ar},e_{\rm vol})$,
l’incidence des arêtes est représentée par
\begin{equation}
 F_{\rm av}=\begin{pmatrix}0&1\\1&1\end{pmatrix}.
 \label{x_reciprocidad:iii:eq:fav}
\end{equation}
Les trois paires consécutives sont $\Sigma\to\Pi$, $\Pi\to\Pi$ et $\Pi\to\Sigma$, chacune de multiplicité un. Les calendriers de longueurs $9$, $27$ et $108$ contiennent $3$, $9$ et $36$ copies du bloc primitif. Leurs matrices d'incidence sont donc ces multiples de $F_{\rm av}$.

Cette construction conserve la différence entre fréquence chronologique et mesure
de l’espace des continuations. La première assigne des poids $(1/3,2/3)$ au cycle
des modes. La seconde se définit sur l’enveloppe symbolique minimale de mémoire
un qui admet exactement les couples précédents. Sa matrice d’adjacence est
$F_{\rm av}$ ; les dénombrements de trajectoires appartiennent à cette enveloppe.

\Needspace{10\baselineskip}
\begin{proposition}[Poids de retour des deux feuillets]
Soient $F_0=0$, $F_1=1$ et $F_{n+1}=F_n+F_{n-1}$. Pour $n\geq1$,
\begin{equation}
 F_{\rm av}^{n}=
 \begin{pmatrix}F_{n-1}&F_n\\F_n&F_{n+1}\end{pmatrix}.
 \label{x_reciprocidad:iii:eq:fav-powers}
\end{equation}
La mesure de Parry de cette enveloppe attribue aux feuillets des poids proportionnels à $(1,\varphi^2)$, et le rapport des décomptes diagonaux converge vers $\varphi^{-2}$.
\end{proposition}
\begin{proof}
La formule des puissances se prouve par induction en utilisant la récurrence.
Le carré de $F_{\rm av}$ est strictement positif. Son vecteur positif de
Perron est $(1,\varphi)$, de valeur propre $\varphi$, car
$\varphi^2=\varphi+1$. La matrice étant symétrique, les poids de Parry sont
proportionnels aux carrés des composantes de ce vecteur. L’autre racine
du polynôme caractéristique est $-\varphi^{-1}$ ; son module est strictement
inférieur à $\varphi$, ce qui donne la limite des quotients diagonaux.
\end{proof}

L’identification spectrale de cette incidence est postérieure à sa construction
à partir du calendrier. La coordonnée régionale $\varphi$ conserve sa généalogie
coinductive antérieure : la reconnaissance de la valeur propre ne sélectionne pas le calendrier.

\section{Marquage de la frontière et rapport surfacique--volumétrique}

Soit $\Pi_{\rm HMT}(X_\infty)$ la valeur du lecteur régional de clôture déjà construit. Avec les projecteurs diagonaux de feuillet, on définit
\begin{equation}
 D_\pi=\Pi_{\rm HMT}(X_\infty)P_{\rm ar}+P_{\rm vol},\qquad
 \Theta_n=
 \frac{\langle e_{\rm ar},D_\pi F_{\rm av}^{n}e_{\rm ar}\rangle}
 {\langle e_{\rm vol},F_{\rm av}^{n}e_{\rm vol}\rangle}.
 \label{x_reciprocidad:iii:eq:marked-return}
\end{equation}
La normalisation volumétrique est unitaire ; la clôture régionale marque une seule fois
le retour aréal. En appliquant \eqref{x_reciprocidad:iii:eq:fav-powers},
\begin{equation}
 \Theta_n=\Pi_{\rm HMT}(X_\infty)\frac{F_{n-1}}{F_{n+1}}
 \longrightarrow\frac{\pi}{\varphi^2}.
 \label{x_reciprocidad:iii:eq:av-limit}
\end{equation}
Le rapport de retour est déterminé par l'incidence, la mesure et le marquage. Le quotient $2\pi/4\pi$ décrit, en revanche, les deux ordres de clôture du revêtement. Les deux relations participent à la lecture des feuillets sans se substituer l'une à l'autre.

\section{Transport de l’orientation aux canaux constitutifs}

Dans une représentation à deux canaux, soit $\mathcal R$ une involution
autoadjointe et $\mathcal S$ un opérateur unitaire d’échange qui satisfait
$\mathcal S\mathcal R\mathcal S^{-1}=-\mathcal R$. Les coordonnées angulaires
déjà construites s’écrivent $x=A_{\rm rad}$ et $y=C^*_{\rm rad}$, avec $x>|y|$.
Alors
\begin{equation}
 B=xI+y\mathcal R,\qquad T=e^{-B},\qquad
 \mathcal ST(y)\mathcal S^{-1}=T(-y).
 \label{x_reciprocidad:iii:eq:channel-conjugation}
\end{equation}
En effet, la conjugaison transforme $B(y)$ en $B(-y)$ et commute avec sa série
exponentielle convergente. Les valeurs propres $q_\pm=e^{-x\mp y}$ sont échangées.
Pour fixer l’ordre des réponses, on utilise la chambre $x>y>0$, où
$0<q_+<q_-<1$. L’application $(x,y)\mapsto(x,-y)$ la transporte dans la
chambre opposée $x>-y>0$, où $0<q_-<q_+<1$. Toutes deux appartiennent au domaine
contractif $x>|y|$. À la frontière $y=0$, les deux canaux coïncident.
La conjugaison échange les chambres, elle ne conserve pas l’ordre strict de
leurs étiquettes ; elle agit sur les projecteurs au moyen de
$\mathcal SP_\pm\mathcal S^{-1}=P_\mp$, avec
$P_\pm=(I\pm\mathcal R)/2$.
Cette application concrète transmet l’orientation à l’opérateur constitutif
développé ci-après.

% END OWNER


% BEGIN OWNER /Users/ruben/Documents/New project/output/EXTREMA_MEDIA_RAZON_RECIPROCIDAD_ES_EN_20260918/ARTICULO/ES/source/antecedentes_iii/sections/03b_vacancias_prefactor.tex
% Ampliación propia de III; no modifica el núcleo común copiado de I.
% Residencia propuesta: antes de Retorno de hoja, incidencia areal--volumétrica.
% Procedencia y controles: VACANCIAS_PREFACTOR_PROCEDENCIA.md, en esta carpeta.
\chapter{Polarisation discrète et mode tritique des projections arithmétiques}
\label{x_reciprocidad:iii:sec:vacancias-prefactor}

Le raffinement et la comparaison entre les feuillets d’APP fournissent deux
observations complémentaires. La première suit l’écart accumulé
entre le nombre de transitions et la capacité acquise ; la seconde conserve
le mode du sélecteur tritique dans l’appariement des fibres additives et
multiplicatives. La polarisation discrète et le préfacteur électronique
s’obtiennent à partir de ces opérations respectives. Leur articulation avec le vide
requiert de maintenir tant l’information temporelle que les espaces sur
lesquels agissent les opérateurs.

\section{Discrépance du raffinement et convention des extrémités}

Les cardinaux du bloc ternaire et de sa représentation décimale fixent
\[
 \rho=\log_{1000}729=\log_{10}9,\qquad
 \delta=1-\rho=\log_{10}(10/9).
\]
On a $0<\delta<1$. En outre, $\delta$ est irrationnel : une égalité $\delta=p/q$, avec $q>0$, impliquerait $10^{q-p}=9^q$, incompatible avec la décomposition en facteurs premiers. Ce décalage intervient déjà dans le codage du raffinement. Sa lecture centrée est explicitée ici, en conservant l'origine de phase et la convention des extrémités.

\begin{definition}[Codages inférieur et supérieur]
Pour $\theta\in\mathbb R$ et $n\in\mathbb Z$, on définit
\begin{align}
 z_n^{\downarrow}(\theta)
 &=\lfloor(n+1)\delta+\theta\rfloor-\lfloor n\delta+\theta\rfloor,
 \label{x_reciprocidad:iii:vp:lower}\\
 z_n^{\uparrow}(\theta)
 &=\lceil(n+1)\delta+\theta\rceil-\lceil n\delta+\theta\rceil.
 \label{x_reciprocidad:iii:vp:upper}
\end{align}
Les deux prennent leurs valeurs dans $\{0,1\}$. La flèche distingue exclusivement la convention des extrémités des intervalles ; l'orientation tritique conserve sa définition précédente.
\end{definition}

Le codage inférieur coïncide avec celui utilisé dans le développement
du raffinement. Le codage supérieur permet de retrouver la formulation de la
polarisation avec la fonction plafond. Les deux expressions sont reliées par
un terme de frontière exactement déterminé.

\begin{proposition}[Changement de convention et terme de frontière]
Soit $b_n(\theta)=\mathbf1_{\mathbb Z}(n\delta+\theta)$. Alors
\begin{equation}
 z_n^{\uparrow}(\theta)-z_n^{\downarrow}(\theta)
 =b_n(\theta)-b_{n+1}(\theta).
 \label{x_reciprocidad:iii:vp:endpoint}
\end{equation}
En particulier, pour $N\geq0$,
\[
 \sum_{n=0}^{N-1}(z_n^{\uparrow}-z_n^{\downarrow})
 =b_0-b_N.
\]
\end{proposition}
\begin{proof}
Pour tout réel $x$, $\lceil x\rceil=\lfloor x\rfloor+1-
\mathbf1_{\mathbb Z}(x)$. Appliquer cette identité aux deux extrémités
de chaque accroissement donne \eqref{x_reciprocidad:iii:vp:endpoint} ; leur somme est télescopique.
L’irrationalité de $\delta$ implique que, $\theta$ étant fixée, au plus
un indice entier satisfait $n\delta+\theta\in\mathbb Z$.
\end{proof}

Avec $\theta=0$, le premier symbole supérieur est $1$ et l'inférieur est $0$ ; pour $n\geq1$, ils coïncident. Le terme initial appartient au registre et est conservé lors du changement de convention. Les densités et les bornes de discrépance sont transportées avec ce terme, sans modifier la mise à jour du TPK ni recommencer l'histoire à chaque retour de phase.

\section{Polarisation bornée et bilan temporel exact}

\begin{definition}[Polarisation discrète centrée]
La polarisation du codage supérieur, relativement à la densité
$\delta$, est
\begin{equation}
 P_N^{\uparrow}(\theta)
 =\sum_{n=0}^{N-1}\bigl(z_n^{\uparrow}(\theta)-\delta\bigr),
 \qquad P_0^{\uparrow}=0.
 \label{x_reciprocidad:iii:vp:polarization}
\end{equation}
Son accroissement temporel est
\begin{equation}
 J_N^{\uparrow}(\theta)
 =P_{N+1}^{\uparrow}(\theta)-P_N^{\uparrow}(\theta)
 =z_N^{\uparrow}(\theta)-\delta.
 \label{x_reciprocidad:iii:vp:current}
\end{equation}
\end{definition}

Le centrage sépare la croissance uniforme de densité $\delta$ de
l’information de phase. Ainsi, $P_N^{\uparrow}$ mesure un déséquilibre accumulé,
non le nombre total de vacances. L’accroissement $J_N^{\uparrow}$ conserve
l’impulsion locale et le fond par rapport auquel elle est évaluée.

\begin{theorem}[Borne uniforme et continuité discrète temporelle]
Pour $N\geq0$ et toute phase $\theta$,
\begin{equation}
 P_N^{\uparrow}(\theta)
 =\lceil N\delta+\theta\rceil-\lceil\theta\rceil-N\delta,
 \qquad |P_N^{\uparrow}(\theta)|<1.
 \label{x_reciprocidad:iii:vp:bound}
\end{equation}
Pour tous entiers $0\leq a<b$,
\begin{equation}
 \sum_{n=a}^{b-1}J_n^{\uparrow}(\theta)
 =P_b^{\uparrow}(\theta)-P_a^{\uparrow}(\theta),
 \qquad
 \left|\sum_{n=a}^{b-1}J_n^{\uparrow}(\theta)\right|<1.
 \label{x_reciprocidad:iii:vp:continuity}
\end{equation}
La densité de la suite des vacances est $\delta$.
\end{theorem}
\begin{proof}
La somme des accroissements de plafond de \eqref{x_reciprocidad:iii:vp:upper} produit la première identité. En écrivant $r(x)=\lceil x\rceil-x\in[0,1)$, on obtient $P_N^{\uparrow}=r(N\delta+\theta)-r(\theta)$ ; sa valeur absolue est inférieure à un. Sur l'intervalle $[a,b)$, on obtient de même $r(b\delta+\theta)-r(a\delta+\theta)$, ce qui prouve la seconde borne, plus précise que la somme de deux bornes individuelles. Enfin, $N^{-1}\sum_{n=0}^{N-1}z_n^{\uparrow}
=\delta+N^{-1}P_N^{\uparrow}\to\delta$.
\end{proof}

Le codage inférieur admet les mêmes démonstrations en remplaçant
$r(x)$ par $\lfloor x\rfloor-x\in(-1,0]$. En particulier,
$P_N^{\uparrow}-P_N^{\downarrow}=b_0-b_N$. L’écart borné
utilisé précédemment dans les lectures de Dirichlet et la polarisation
présente conservent donc une relation explicite d’origine et d’extrémités.
L’apériodicité n’introduit pas non plus de dérive séculaire de ce déséquilibre :
l’histoire peut se prolonger indéfiniment tandis que l’écart
centré reste uniformément borné.

\section{Réalisation dans les droites de charge et de temps}

Pour réaliser dimensionnellement le bilan, on spécifie une droite de charge
$\mathcal L_Q$ de section non nulle $u_Q$ et une suite croissante
d’instants $t_N$ dans une droite temporelle orientée. Avec
$\Delta t_N=t_{N+1}-t_N>0$, la charge centrée et le courant moyen de
chaque intervalle sont définis par
\begin{equation}
 Q_N=P_N^{\uparrow}u_Q,\qquad
 I_N=\frac{Q_{N+1}-Q_N}{\Delta t_N}
     =\frac{u_Q}{\Delta t_N}\bigl(z_N^{\uparrow}-\delta\bigr).
 \label{x_reciprocidad:iii:vp:dimensional}
\end{equation}
Ainsi, $Q_N\in\mathcal L_Q$ et $I_N\in\mathcal L_Q\otimes\mathcal L_T^{-1}$. La section $u_Q$ et les intervalles temporels sont des données de cette réalisation ; la construction arithmétique précédente a déjà déterminé $P_N^{\uparrow}$ et $J_N^{\uparrow}$.

\begin{proposition}[Conservation du bilan sous réalisation dimensionnelle]
Pour $a<b$,
\begin{equation}
 Q_b-Q_a=\sum_{n=a}^{b-1}\Delta t_n I_n,
 \qquad
 \left|\frac{Q_b-Q_a}{u_Q}\right|<1.
 \label{x_reciprocidad:iii:vp:dimensional-balance}
\end{equation}
Si tous les intervalles ont longueur $t_0$, les deux courants possibles
sont $-\delta u_Q/t_0$ et $(1-\delta)u_Q/t_0$.
\end{proposition}
\begin{proof}
Multiplier chaque identité de \eqref{x_reciprocidad:iii:vp:dimensional} par
$\Delta t_n$ et sommer fait télescoper les charges. La borne est
\eqref{x_reciprocidad:iii:vp:continuity} ; les deux valeurs locales s’obtiennent à partir de
$z_n^{\uparrow}\in\{0,1\}$.
\end{proof}

La continuité démontrée est temporelle et se rapporte à la charge centrée. Une réalisation spatiale incorpore en outre la localisation des charges, l'incidence entre cellules et l'identification des flux de frontière. Dans une description constitutive par canaux, le couplage est spécifié au moyen d'une application d'excitation $\iota:\mathcal L_Q\to\mathcal H_{\rm ch}\otimes\mathcal L_Q$ et d'une application d'observation $\ell:\mathcal H_{\rm ch}\otimes\mathcal L_Q\to\mathcal L_Q$, où $\mathcal H_{\rm ch}$ est une réalisation réelle de l'espace des canaux. Si les deux sont linéaires et indépendantes de l'indice temporel, la réponse d'un opérateur $\mathcal V$ fixe est transportée par
\[
 Q_N^{\rm resp}=\ell(\mathcal V\otimes I_{\mathcal L_Q})\iota(Q_N),
 \qquad
 I_N^{\rm resp}=\frac{Q_{N+1}^{\rm resp}-Q_N^{\rm resp}}{\Delta t_N}.
\]
La linéarité conserve le bilan télescopique. Cette composition
explicite les données d’une identification constitutive : la séquence
scalaire détermine l’impulsion, tandis que $\iota$, $\mathcal V$ et $\ell$
déterminent son excitation, son transport et son observation par canaux. Les
identités \eqref{x_reciprocidad:iii:vp:polarization}--\eqref{x_reciprocidad:iii:vp:dimensional-balance}
établissent la première construction indépendamment de ce choix de
réalisation. L’opérateur du vide est construit dans la
section~\ref{x_reciprocidad:iii:sec:constitutiva} sur ses canaux et ses projecteurs propres.

\section{Appariement des fibres additives et multiplicatives}

La seconde observation provient de la réalisation tridimensionnelle d’APP,
$\mathcal A_3=D_9^3$, avec une mesure uniforme et
$D_9=\{1,\ldots,9\}$. On conserve les observables
\[
 \Sigma(x)=\rho_9(x_1+x_2+x_3),\qquad
 \Pi(x)=\rho_9(x_1x_2x_3),
\]
où $\rho_9(n)=1+[(n-1)\bmod9]$ pour $n>0$. Dans
$\ell^2(\mathcal A_3)$, soient $P_{\Sigma,3}$ et $P_{\Pi,3}$ les
espérances conditionnelles sur leurs fibres. Par exemple,
\[
 (P_{\Sigma,3}f)(x)=
 \frac1{|\Sigma^{-1}(\Sigma(x))|}
 \sum_{y:\,\Sigma(y)=\Sigma(x)}f(y).
\]
Ce sont des projections orthogonales définies sur le même support. La matrice
d’incidences $N_{ab}=|\Sigma^{-1}(a)\cap\Pi^{-1}(b)|$ résulte de
l’énumération des $729$ triplets, en additionnant et en multipliant leurs coordonnées avant
d’appliquer $\rho_9$ :
\begin{equation}
 N=9\begin{pmatrix}
 0&1&1&0&1&1&0&1&4\\
 1&0&1&1&0&1&1&0&4\\
 1&0&2&0&1&2&0&0&3\\
 0&1&1&0&1&1&0&1&4\\
 1&0&1&1&0&1&1&0&4\\
 0&0&2&1&0&2&0&1&3\\
 0&1&1&0&1&1&0&1&4\\
 1&0&1&1&0&1&1&0&4\\
 0&1&2&0&0&2&1&0&3
 \end{pmatrix}.
 \label{x_reciprocidad:iii:vp:joint-matrix}
\end{equation}
Les lignes ont pour somme $81$ et les sommes des colonnes sont
\[
 (m_1,\ldots,m_9)=(36,36,108,36,36,108,36,36,297).
\]
L’appariement des indicatrices normalisées des fibres
est donc $C_{ab}=N_{ab}/\sqrt{81m_b}$. La normalisation recueille les
multiplicités produites par APP, non un choix de valeurs singulières.

\begin{proposition}[Mode singulier du sélecteur tritique]
Dans la carte canonique, le vecteur du sélecteur est $m_{\rm ph}=(1,-1,0)^3$. On a
\begin{equation}
 Cm_{\rm ph}=C^{\mathsf T}m_{\rm ph}=-\tfrac12m_{\rm ph}.
 \label{x_reciprocidad:iii:vp:trit-mode}
\end{equation}
Ainsi, la valeur singulière correspondante est $s=1/2$.
\end{proposition}
\begin{proof}
Dans les six colonnes où $m_{\rm ph}$ est non nul, $m_b=36$ et
$C_{ab}=N_{ab}/54$. Multiplier les lignes de
\eqref{x_reciprocidad:iii:vp:joint-matrix} par $(1,-1,0)^3$ produit
$-m_{\rm ph}/2$. Pour le produit transposé, les sommes signées
des colonnes $3,6,9$ sont nulles ; les six autres donnent les mêmes
valeurs $-m_{{\rm ph},b}/2$. Cela prouve les deux identités et, par
composition, $CC^{\mathsf T}m_{\rm ph}=m_{\rm ph}/4$.
\end{proof}

\section{Norme du commutateur et préfacteur électronique}

\begin{theorem}[Norme exacte d’incompatibilité arithmétique]
Les deux projections satisfont
\begin{equation}
 \bigl\|[P_{\Sigma,3},P_{\Pi,3}]\bigr\|=\frac{\sqrt3}{4}.
 \label{x_reciprocidad:iii:vp:prefactor}
\end{equation}
Le maximum est atteint dans le plan principal correspondant au mode
tritique de \eqref{x_reciprocidad:iii:vp:trit-mode}.
\end{theorem}
\begin{proof}
La matrice rationnelle $M=CC^{\mathsf T}$ a pour entrées $M_{ac}=\sum_bN_{ab}N_{cb}/(81m_b)$. Son spectre complet se vérifie au moyen des sous-espaces indépendants suivants : les vecteurs $(1,\ldots,1)$, $(1,-1,0)^3$ et $(1,1,-2)^3$ ont pour valeurs propres $1$, $1/4$ et $7/132$, respectivement ; le sous-espace des vecteurs à support dans $\{3,6,9\}$ dont la somme est nulle a pour dimension $2$ et pour valeur propre $1/18$ ; les vecteurs de somme nulle à support dans $\{1,4,7\}$ ou dans $\{2,5,8\}$ forment un noyau de dimension $4$. La substitution dans la matrice vérifie ces actions et les dimensions ont pour somme $9$. Par conséquent,
\[
 \operatorname{spec}M=
 \{1,1/4,7/132,1/18,1/18,0,0,0,0\}.
\]
Pour deux projections orthogonales, une valeur singulière $s\in(0,1)$ de
l’appariement détermine un plan principal dans lequel leurs matrices sont
\[
 P=\begin{pmatrix}1&0\\0&0\end{pmatrix},\qquad
 Q=\begin{pmatrix}s^2&s\sqrt{1-s^2}\\
 s\sqrt{1-s^2}&1-s^2\end{pmatrix}.
\]
Cette forme s’obtient en prenant un vecteur unitaire $u$ de la première
image avec $PQu=s^2u$ et en le complétant par
$(Qu-s^2u)/(s\sqrt{1-s^2})$. Leur commutateur est
$s\sqrt{1-s^2}\bigl(\begin{smallmatrix}0&1\\-1&0\end{smallmatrix}\bigr)$.
Les secteurs de valeur singulière $0$ ou $1$ contribuent par zéro.
Par conséquent, la norme totale est le maximum de
$\sqrt{\lambda(1-\lambda)}$ sur le spectre de $M$. Ce maximum
est $\sqrt{3/16}$, atteint en $\lambda=1/4$.
\end{proof}

La norme est le préfacteur qui intervient dans la composition électronique.
Son obtention conserve le support tridimensionnel, les fibres et le mode
du sélecteur ; l’évaluation scalaire ne remplace pas cette structure. En
particulier, connaître $\sqrt3/4$ permet de connaître la grandeur maximale de
l’incompatibilité, tandis que les projecteurs conservent les directions
dans lesquelles elle agit.

\clearpage
\section{Défaut quadratique et réponse des secteurs 90/120}

Le mode tritique détermine simultanément
\begin{equation}
 1-s^2=\frac34=\frac{90}{120},\qquad s=\frac12.
 \label{x_reciprocidad:iii:vp:defect-ratio}
\end{equation}
La première expression est le défaut quadratique du plan principal d'APP. La dernière conserve les longueurs des deux secteurs du calendrier. La réponse constitutive utilise ces secteurs sur un autre objet : un canal contractant $q\in(0,1)$. Son évaluation est
\[
 r(q)=\frac{1-q^{90}}{1-q^{120}}.
\]
La relation entre le quotient des secteurs et cette réponse peut être
précisée sans identifier leurs opérateurs.

\begin{proposition}[Limite du quotient de réponse]
Para $0<q<1$,
\begin{equation}
 \frac34<r(q)<1,\qquad
 \lim_{q\uparrow1}r(q)=\frac34=1-s^2.
 \label{x_reciprocidad:iii:vp:response-limit}
\end{equation}
\end{proposition}
\begin{proof}
Avec $u=q^{30}\in(0,1)$, la factorisation de $1-u^3$ et
$1-u^4$ donne
$r(q)=(1+u+u^2)/(1+u+u^2+u^3)$. La limite en $u=1$ est
$3/4$. De plus,
$4(1+u+u^2)-3(1+u+u^2+u^3)
=(1-u)(3u^2+2u+1)>0$, et le dénominateur dépasse le numérateur
de $u^3>0$. On obtient les deux inégalités.
\end{proof}

La coïncidence exacte concerne le défaut du mode singulier et la limite
de la réponse des secteurs. À l’intérieur contractif, la réponse
dépend de $q$ et est strictement supérieure à $3/4$. Le développement
constitutif conserve ses deux canaux $q_\pm$ et ses projecteurs
$P_\pm$ ; les projecteurs $P_{\Sigma,3}$ et $P_{\Pi,3}$ conservent,
pour leur part, les fibres arithmétiques qui ont déterminé le préfacteur.
Une équivalence opératorielle entre les deux réalisations exige une application
entre leurs espaces qui entrelace les actions correspondantes.
Les équations \eqref{x_reciprocidad:iii:vp:trit-mode}, \eqref{x_reciprocidad:iii:vp:prefactor} et
\eqref{x_reciprocidad:iii:vp:response-limit} précisent l’information partagée et
maintiennent visibles les données de chaque réalisation.

% END OWNER


% BEGIN OWNER /Users/ruben/Documents/New project/output/EXTREMA_MEDIA_RAZON_RECIPROCIDAD_ES_EN_20260918/ARTICULO/ES/source/antecedentes_iii/sections/04_respuesta_constitutiva.tex
\chapter{Opérateur constitutif et reconstruction des canaux du vide}
\label{x_reciprocidad:iii:sec:constitutiva}

Les grandeurs constitutives s'obtiennent au moyen d'une opération sur le transport angulaire. L'information de départ comprend les canaux, leurs projecteurs et l'orientation qui les distingue. Cette organisation permet de suivre séparément la construction de l'opérateur, l'évaluation de ses fonctions spectrales et sa réalisation dimensionnelle ultérieure.

\section{Réponse positive sur deux feuillets}

Dans la chambre orientée $x>y>0$ de la section~\ref{x_reciprocidad:iii:sec:hojas}, soient
$P_\pm=(I\pm\mathcal R)/2$ et
$T=q_+P_++q_-P_-$, avec $0<q_+<q_-<1$. Les incidences construites par le
calendrier fixent les secteurs $90$ et $120$. La réponse constitutive est
\begin{equation}
 \mathcal V=(I-T^{90})(I-T^{120})^{-1}.
 \label{x_reciprocidad:iii:eq:vacuum-op}
\end{equation}
Cette définition reçoit le transport provenant du TPK et conserve les
attributions de feuillet. La règle de réponse et son domaine font explicitement partie
de la construction.

Le quotient provient de la comparaison des deux incidences du calendrier sur un même transport, et non de l'ajustement séparé de quatre grandeurs. Si $S=T^{30}$ et $\Sigma_j(S)=I+S+\cdots+S^{j-1}$, la factorisation de $I-S^j$ donne
\begin{equation}
 \mathcal V\,\Sigma_4(S)=\Sigma_3(S),\qquad
 \mathcal V=\Sigma_3(S)\Sigma_4(S)^{-1}.
 \label{x_reciprocidad:iii:eq:completion-response}
\end{equation}
Les facteurs $\Sigma_3$ et $\Sigma_4$ expriment respectivement les
secteurs $90=3\cdot30$ et $120=4\cdot30$. Comme le spectre de $S$ est
contenu dans $(0,1)$, $\Sigma_4(S)$ est inversible : l’équation de
complétion détermine une unique réponse pour le transport fixé.
Cette unicité concerne la règle constitutive déclarée et ses
incidences. L’identification électromagnétique ajoute l’attribution
orientée des feuillets et les droites dimensionnelles précisées plus
bas.

\begin{proposition}[Décomposition et positivité]
La réponse est bien définie et
\begin{equation}
 \mathcal V=r_+P_++r_-P_-,\qquad
 r_\pm=\frac{1-q_\pm^{90}}{1-q_\pm^{120}}.
 \label{x_reciprocidad:iii:eq:rchannels}
\end{equation}
Ses valeurs propres satisfont $3/4<r_-<r_+<1$ dans cette chambre.
\end{proposition}
\begin{proof}
Chaque coefficient $1-q_\pm^{120}$ est positif. L'inverse de $I-T^{120}$ est la somme des projecteurs divisés par ces coefficients, ce qui prouve \eqref{x_reciprocidad:iii:eq:rchannels}. Avec $s=q^{30}$, on obtient
\begin{equation}
 f(s)=\frac{1+s+s^2}{1+s+s^2+s^3},\qquad
 f'(s)=-\frac{s^2(3+2s+s^2)}{(1+s+s^2+s^3)^2}<0.
 \label{x_reciprocidad:iii:eq:monotone}
\end{equation}
Les limites en $0$ et en $1$ sont $1$ et $3/4$, respectivement. La monotonie
et l’ordre des canaux concluent la preuve.
\end{proof}

\section{Quatre coordonnées d’une réponse commune}

On définit les évaluations normalisées
\begin{equation}
 \widehat\varepsilon=r_+^2,\qquad
 \widehat\mu=r_-^2,\qquad
 \widehat Z=\frac{r_-}{r_+},\qquad
 \widehat c=\frac1{r_+r_-}.
 \label{x_reciprocidad:iii:eq:four}
\end{equation}
Le premier couple conserve les réponses quadratiques assignées à chaque feuillet.
Le produit recueille le mode commun ; le quotient distingue leur orientation relative.
On vérifie exactement
\begin{equation}
 \widehat\mu\widehat\varepsilon=\widehat c^{-2},\qquad
 \frac{\widehat\mu}{\widehat\varepsilon}=\widehat Z^2,
 \qquad
 r_+=\frac1{\sqrt{\widehat Z\widehat c}},\quad
 r_-=\sqrt{\frac{\widehat Z}{\widehat c}}.
 \label{x_reciprocidad:iii:eq:four-identities}
\end{equation}
Les racines positives fixent l'inversion. Sous la conjugaison des feuillets, lue par rapport aux marques $P_\pm$ fixes, $\widehat\varepsilon$ et $\widehat\mu$ sont échangés, $\widehat Z$ est inversé et $\widehat c$ est conservé. La réponse est alors transportée vers la chambre $x>-y>0$ : les formules restent valables, avec les ordres $q_-<q_+$ et $r_+<r_-$ inversés. Par conséquent, le mode commun et l'orientation sont des données récupérables au moyen d'observables complémentaires.

\begin{proposition}[Détermination des évaluations quadratiques]
Soit $W=r_+P_++r_-P_-$ une réponse positive sur les deux feuillets.
Étant fixées les lectures de produit et de quotient
\[
 \widehat c^{-1}=\det W=r_+r_-,
 \qquad \widehat Z=r_-/r_+,
\]
il existe une unique paire positive $(\widehat\varepsilon,\widehat\mu)$
qui vérifie
$\widehat\mu\widehat\varepsilon=\widehat c^{-2}$ et
$\widehat\mu/\widehat\varepsilon=\widehat Z^2$.
C’est précisément $(r_+^2,r_-^2)$.
\end{proposition}
\begin{proof}
La positivité permet de prendre les racines sans ambiguïté :
\[
 \widehat\mu=\frac{\widehat Z}{\widehat c}
 =\frac{r_-}{r_+}r_+r_-=r_-^2,\qquad
 \widehat\varepsilon=\frac1{\widehat Z\widehat c}
 =\frac{r_+}{r_-}r_+r_-=r_+^2.
\]
Les expressions vérifient les deux relations et en sont les seules racines
positives.
\end{proof}

Ainsi, les carrés ne constituent pas deux choix indépendants ajoutés
aux canaux : ils sont déterminés par les lectures multiplicative et
orientée, une fois fixées les relations constitutives. Cette proposition
établit l’unicité au sein de ce système de règles. L’interprétation
des règles comme réponse électromagnétique conserve son contenu
supplémentaire de représentation physique.

La lecture déterminantale du transport et la réponse constitutive
doivent également conserver leur ordre. Pour $T=e^{-xI-y\mathcal R}$, on a
$\det T=e^{-2x}$, tandis que
\[
 \det\mathcal V=f(q_+^{30})f(q_-^{30})=\widehat c^{-1}.
\]
Les deux proviennent du même opérateur à deux feuillets, mais évaluent des fonctions différentes de son spectre. En général, appliquer d'abord le déterminant puis une fonction rationnelle n'équivaut pas à prendre le déterminant de cette fonction de l'opérateur. Maintenir les deux valeurs propres jusqu'à l'achèvement de la réponse conserve l'information d'orientation qu'un déterminant isolé ne distingue plus.

\section{Inversion cubique et récupération angulaire}

\begin{theorem}[Inversion constitutive dans le domaine contractif]
Pour chaque $r\in(3/4,1)$, il existe une unique $s\in(0,1)$ telle que $f(s)=r$.
Cette $s$ est l’unique racine dans cet intervalle de
\begin{equation}
 r s^3+(r-1)(1+s+s^2)=0.
 \label{x_reciprocidad:iii:eq:cubic}
\end{equation}
La réponse étiquetée reconstruit $s_\pm$, $q_\pm=s_\pm^{1/30}$ et
\begin{equation}
 x=-\frac12\log(q_+q_-),\qquad
 y=\frac12\log\frac{q_-}{q_+}.
 \label{x_reciprocidad:iii:eq:recover-angular}
\end{equation}
\end{theorem}
\begin{proof}
La dérivée et les limites de \eqref{x_reciprocidad:iii:eq:monotone} prouvent que $f$ est une
bijection décroissante entre les intervalles indiqués. Multiplier $f(s)=r$
par son dénominateur produit \eqref{x_reciprocidad:iii:eq:cubic}. La racine positive d’ordre
$30$ retrouve le canal. Enfin, la somme et la différence de
$\log q_\pm=-x\mp y$ donnent \eqref{x_reciprocidad:iii:eq:recover-angular}.
\end{proof}

La reconstruction conserve les étiquettes de feuillet. Avec l’opérateur complet
$\mathcal V$, on conserve aussi les projecteurs et on applique la fonction
inverse par calcul spectral. Récupérer deux coordonnées angulaires est une
opération définie sur cette représentation ; la trace historique complète
du TPK conserve des données supplémentaires.

L'intervalle $r\in(3/4,1)$ et la fonction inverse de ce théorème appartiennent à la normalisation interne de \eqref{x_reciprocidad:iii:eq:four}. Si l'on dispose de coordonnées exprimées dans une autre carte d'unités, on récupère d'abord cette normalisation au moyen de la transformation de l'équation~\eqref{x_reciprocidad:iii:eq:dimensions} ; on applique ensuite l'inversion cubique.

\section{Coordonnées logarithmiques et raffinement}

Soient $\mathcal E=\log(r_+r_-)$ et $\mathcal B=\log(r_-/r_+)$. Alors
\begin{equation}
 \log\mathcal V=\tfrac12\mathcal E I-\tfrac12\mathcal B\mathcal R.
 \label{x_reciprocidad:iii:eq:logvac}
\end{equation}
Cette décomposition exhibe la composante invariante et la composante orientée.
Dans l’amplification $\mathcal V_m=\mathcal V\otimes I_{9^m}$, les
espérances partielles normalisées vérifient
$E_m(\mathcal V_{m+1})=\mathcal V_m$. Pour tout polynôme $p$,
$p(\mathcal V_m)=p(\mathcal V)\otimes I_{9^m}$ ; d’où
$E_m(p(\mathcal V_{m+1}))=p(\mathcal V_m)$. La continuité uniforme
étend l’identité aux fonctions continues du spectre. Les projecteurs,
les puissances et le logarithme sont ainsi conservés dans la tour indiquée.

\section{Types dimensionnels des quatre évaluations}

Avec les droites d’action, de charge et de vitesse, munies de sections
$u_S,u_Q,u_C$, la réalisation dimensionnelle s’écrit
\begin{align}
 Z_{\rm vac}&=\widehat Z u_Su_Q^{-2},&
 c_{\rm vac}&=\widehat c u_C,\\
 \mu_{\rm vac}&=\widehat\mu u_Su_C^{-1}u_Q^{-2},&
 \varepsilon_{\rm vac}&=\widehat\varepsilon u_S^{-1}u_C^{-1}u_Q^2.
 \label{x_reciprocidad:iii:eq:dimensions}
\end{align}
En multipliant et en divisant ces expressions, on préserve $\mu_{\rm vac}\varepsilon_{\rm vac}=c_{\rm vac}^{-2}$ et $\mu_{\rm vac}/\varepsilon_{\rm vac}=Z_{\rm vac}^2$. La carte SI évalue ensuite les sections. La comparaison physique requiert d'exposer leur construction et leur normalisation avec l'évaluation ; les quatre coordonnées normalisées de \eqref{x_reciprocidad:iii:eq:four} conservent leur type sans dimension.

% END OWNER


% BEGIN OWNER /Users/ruben/Documents/New project/output/EXTREMA_MEDIA_RAZON_RECIPROCIDAD_ES_EN_20260918/ARTICULO/ES/source/antecedentes_iii/sections/04b_reconstruccion_forma_lc.tex
% INSERCIÓN FOCAL AUTORIZADA: no introduce un capítulo ni el concepto de radión.
% Incluir dentro de la sección constitutiva, después de la inversión cúbica
% y de la declaración de las cartas dimensionales. Etiquetas fuera de grupos
% de prefijado automático. Conserva íntegramente el bloque geométrico anterior.
\section{Reconstruction de la forme d’action à partir de la réponse constitutive}
\label{x_reciprocidad:iii:subsec:puente-forma-LC}

La réponse constitutive et la réalisation inductive--capacitive reçoivent
le même couple angulaire construit à partir d’APP--TRIT--TPK. La première évalue
des fonctions rationnelles des deux canaux ; la seconde utilise le rapport de
leurs exposants pour déterminer la forme de l’ellipse d’action.
L’inversion cubique permet de composer les deux lectures sans identifier leurs
quotients numériques.

Soient $x=A_{\rm rad}$ et $y=C^*_{\rm rad}$, avec $x>|y|$, les coordonnées
relevées déjà construites, et conservons leurs étiquettes de feuillet. Nous écrivons
\begin{equation}
 q_\pm=e^{-(x\pm y)},\qquad s_\pm=q_\pm^{30},\qquad
 r_\pm=f(s_\pm),\qquad
 f(s)=\frac{1+s+s^2}{1+s+s^2+s^3}.
 \label{x_reciprocidad:iii:eq:forma-canales}
\end{equation}
Le domaine contractant donne $s_\pm\in(0,1)$ et $r_\pm\in(3/4,1)$. Dans la chambre $x>y>0$, on a $s_+<s_-$ et $r_+>r_-$ ; l'échange des feuillets transporte cette chambre vers l'orientation opposée.

L’image des deux coordonnées constitutives normalisées est le domaine
\begin{equation}
 \mathcal U=\left\{(z,v)\in(0,\infty)^2:
   (zv)^{-1/2}\in(3/4,1),\quad
   (z/v)^{1/2}\in(3/4,1)\right\}.
 \label{x_reciprocidad:iii:eq:forma-dominio}
\end{equation}
Ici, $(z,v)=(\widehat Z,\widehat c)$. Les restrictions expriment l'appartenance des canaux reconstruits au domaine du théorème d'inversion ; elles ne constituent pas un choix ultérieur de leurs valeurs.

\Needspace{23\baselineskip}
\begin{proposition}[Composition constitutive de la forme et de l’impédance LC]
\label{x_reciprocidad:iii:prop:forma-LC}
Pour $(\widehat Z,\widehat c)\in\mathcal U$, soient $s_\pm\in(0,1)$ les
racines uniques de
\begin{equation}
 r_\pm s_\pm^3+(r_\pm-1)(1+s_\pm+s_\pm^2)=0,
 \qquad
 r_+=\frac1{\sqrt{\widehat Z\widehat c}},\quad
 r_-=\sqrt{\frac{\widehat Z}{\widehat c}}.
 \label{x_reciprocidad:iii:eq:forma-inversion}
\end{equation}
Le rapport angulaire, l’anisotropie de l’ellipse d’action et l’impédance
LC relative sont reconstruits exactement par
\begin{align}
 \frac{y}{x}
 &=\frac{\log s_+-\log s_-}{\log s_++\log s_-},
 \label{x_reciprocidad:iii:eq:forma-razon}\\
 \eta_{\rm el}
 &=\frac12\log\frac{\log s_+}{\log s_-},
 \label{x_reciprocidad:iii:eq:forma-eta}\\
 \frac{Z_{\rm LC}}{Z_*}
 &=e^{-\eta_{\rm el}}
  =\sqrt{\frac{\log s_-}{\log s_+}}.
 \label{x_reciprocidad:iii:eq:forma-LC}
\end{align}
Dans la dernière identité, $Z_{\rm LC}=Z_{\eta_{\rm el}}$ est l’impédance
de la carte LC de \eqref{x_reciprocidad:hmtII:eq:ii-elipse-lc}.
\end{proposition}
\begin{proof}
L'inversion cubique de \eqref{x_reciprocidad:iii:eq:cubic} détermine $s_\pm$ de manière univoque, car $f$ est une bijection décroissante de $(0,1)$ sur $(3/4,1)$. Ses racines positives d'ordre $30$ récupèrent $q_\pm$. Les identités
\[
 \log s_+=-30(x+y),\qquad \log s_-=-30(x-y)
\]
prouvent \eqref{x_reciprocidad:iii:eq:forma-razon} par différence et somme. Les deux logarithmes
sont négatifs ; leur quotient est positif et leur somme n’est pas nulle. Par conséquent,
\[
 \frac12\log\frac{\log s_+}{\log s_-}
   =\frac12\log\frac{x+y}{x-y}
   =\operatorname{artanh}(y/x)=\eta_{\rm el}.
\]
La relation $Z_{\eta_{\rm el}}/Z_*=e^{-\eta_{\rm el}}$, démontrée dans
\eqref{x_reciprocidad:hmtII:eq:ii-elipse-impedancia}, donne \eqref{x_reciprocidad:iii:eq:forma-LC}.
\end{proof}

La composition conserve l’information nécessaire à chaque reconstruction.
Avec $\hbar$, on retrouve
$a_S=\sqrt{2\hbar}\,e^{\eta_{\rm el}/2}$ et
$b_S=\sqrt{2\hbar}\,e^{-\eta_{\rm el}/2}$ ; avec la carte
$(\omega,Z_*)$, on retrouve l’inductance et la capacité déjà définies.
L’action, l’horloge et la référence dimensionnelle accompagnent la forme comme
données typées. L’inversion de deux coordonnées adimensionnelles ne
reconstruit ni leurs unités ni toute la mémoire de l’état.

La normalisation intervient avant de résoudre les équations cubiques.
Si l’on reçoit $Z_{\rm vac}$ et $c_{\rm vac}$ dans une carte dimensionnelle,
les relations \eqref{x_reciprocidad:iii:eq:dimensions} déterminent d’abord
\[
 \widehat Z=\frac{Z_{\rm vac}}{u_Su_Q^{-2}},\qquad
 \widehat c=\frac{c_{\rm vac}}{u_C}.
\]
Ce sont les coordonnées auxquelles s'applique \eqref{x_reciprocidad:iii:eq:forma-inversion}. La référence $Z_*$ appartient à la réalisation LC et est expressément conservée dans \eqref{x_reciprocidad:iii:eq:forma-LC}.

Les deux quotients sont reliés au moyen des canaux, mais évaluent
des fonctions différentes :
\begin{equation}
 \widehat Z=\frac{f(s_-)}{f(s_+)},\qquad
 \frac{Z_{\rm LC}}{Z_*}
       =\sqrt{\frac{\log s_-}{\log s_+}}.
 \label{x_reciprocidad:iii:eq:forma-dos-cocientes}
\end{equation}
La reconstruction précédente fournit l'application entre les deux lectures. En particulier, échanger les feuillets inverse les deux quotients, remplace $\eta_{\rm el}$ par $-\eta_{\rm el}$ et conserve $\widehat c$. Dans la limite symétrique $y=0$, les canaux coïncident et les deux quotients valent $1$. En dehors de ce cas, l'égalité d'origine n'impose pas l'égalité des deux fonctions. La forme d'action est récupérable à partir de la réponse constitutive en conservant l'orientation et la normalisation.

% END OWNER


% BEGIN OWNER /Users/ruben/Documents/New project/output/EXTREMA_MEDIA_RAZON_RECIPROCIDAD_ES_EN_20260918/ARTICULO/ES/source/antecedentes_iii/sections/04c_velocidad_y_reloj.tex
% RESULTADO_RECUPERADO: funtor dimensional integral, eq:cZ-internos,
% eq:longitud-ruta; composición reunida en Artículo V REV02, sección 69.
% Incorporación autónoma a III: mismo transporte, mismas hojas y misma carta.
\section{Vitesse constitutive et réalisation cinématique de l’horloge}
\label{x_reciprocidad:iii:sec:enlace-velocidad-constitutiva}

La réponse constitutive et le lecteur cinématique réalisent une même
section de vitesse. Leur identification conserve l’ordre de construction :
les feuillets d’APP, leur orientation TRIT et le transport angulaire publié
par le TPK précèdent les lectures de vitesse et le choix des unités.
Dans la représentation à deux feuillets de la
section~\ref{x_reciprocidad:iii:sec:hojas}, avec $x>|y|$ et les projecteurs de rang un
$P_\pm$, le transport est
\[
 T=e^{-xI-y\mathcal R}=q_+P_++q_-P_-,
 \qquad q_\pm=e^{-x\mp y},\qquad 0<q_\pm<1.
\]
Ses blocs temporels de longueurs $90$ et $120$ produisent la réponse
$\mathcal V$ de \eqref{x_reciprocidad:iii:eq:vacuum-op}. Les coordonnées $x,y$ et les
canaux sont des sorties du développement angulaire préalable ; cette composition
applique ses lecteurs à ce transport déjà construit.

Le lecteur logarithmique symétrique du transport temporel est
\begin{equation}
 \mathcal E=
 \log\!\bigl[(1-q_+^{90})(1-q_-^{90})\bigr]
 -\log\!\bigl[(1-q_+^{120})(1-q_-^{120})\bigr].
 \label{x_reciprocidad:iii:eq:velocidad-lector-simetrico}
\end{equation}
Dans la droite dimensionnelle de vitesse, de section positive $u_C$, le
lecteur cinématique publie $c_{\mathrm{int}}=\exp(-\mathcal E)u_C$.
La réponse constitutive publie
$c_{\mathrm{vac}}=\widehat c u_C$ au moyen de \eqref{x_reciprocidad:iii:eq:four}.

\begin{proposition}[Identité constitutive et cinématique]
\label{x_reciprocidad:iii:prop:identidad-velocidad}
Dans la même carte dimensionnelle, on vérifie
\begin{equation}
 c_{\mathrm{int}}=c_{\mathrm{vac}}=\widehat c\,u_C,
 \qquad \widehat c=(r_+r_-)^{-1}=\exp(-\mathcal E).
 \label{x_reciprocidad:iii:eq:identidad-velocidad}
\end{equation}
L'identification conserve l'échange des feuillets.
\end{proposition}
\begin{proof}
Tous les facteurs de \eqref{x_reciprocidad:iii:eq:velocidad-lector-simetrico} sont
positifs. La loi du logarithme et \eqref{x_reciprocidad:iii:eq:rchannels} donnent
\[
 \mathcal E=\log(r_+r_-)=\log\det\mathcal V,
 \qquad e^{-\mathcal E}=(\det\mathcal V)^{-1}.
\]
Le déterminant correspond à la réalisation à deux feuillets de rang un.
Les deux définitions publient donc la même section. L’échange
des feuillets permute $r_+$ et $r_-$ et conserve leur produit. Le déterminant
$\det T=e^{-2x}$ correspond au transport élémentaire, tandis que
$\det\mathcal V=r_+r_-$ correspond à sa réponse temporelle ;
\eqref{x_reciprocidad:iii:eq:identidad-velocidad} utilise le second.
\end{proof}

\subsection{Longueur de trajectoire et durée élémentaire}

Le repère dimensionnel positif $(u_S,u_C,u_T,u_Q)$ induit la base de longueur $u_L=u_Cu_T$. Pour une trajectoire enrichie $\gamma$ de $n=|\gamma|$ transitions et à canal constant, les lecteurs additifs sont
\begin{equation}
 T_{\mathrm{int}}(\gamma)=nu_T,
 \qquad L_{\mathrm{int}}(\gamma)=n\widehat c\,u_L.
 \label{x_reciprocidad:iii:eq:velocidad-longitud-reloj}
\end{equation}
La concaténation additionne le nombre de transitions et leurs réalisations.
Les routes enrichies conservent en outre leur orientation et leur mémoire.

\begin{corollary}[Normalisation de la longueur élémentaire]
Pour $n>0$, $L_{\mathrm{int}}(\gamma)/T_{\mathrm{int}}(\gamma)
=c_{\mathrm{vac}}$. Dans la normalisation $t_0=u_T$,
\begin{equation}
 \ell_0=c_{\mathrm{int}}t_0=\widehat c\,u_L.
 \label{x_reciprocidad:iii:eq:longitud-elemental-constitutiva}
\end{equation}
Dans une carte temporelle $t_0=\tau_0u_T$, avec $\tau_0>0$, on obtient
$\ell_0=\tau_0\widehat c\,u_L$.
\end{corollary}
\begin{proof}
Diviser les deux expressions de \eqref{x_reciprocidad:iii:eq:velocidad-longitud-reloj} simplifie $n$ et utilise $u_L/u_T=u_C$. Multiplier la section de vitesse par $t_0$ prouve les deux formules élémentaires.
\end{proof}

Lorsque le canal dépend de l’arête, le lecteur conserve la forme additive
$L_{\mathrm{int}}(\gamma)=\sum_{e\in\gamma}\widehat c(e)u_L$.
Son quotient par $nu_T$ est la moyenne des vitesses de ces arêtes.
L’expression à canal constant est la spécialisation homogène de
cette lecture.

\subsection{Changement de carte et normalisation adaptée}

\begin{proposition}[Covariance de l’identification]
Soient $u_C'=d u_C$ et $u_T'=\eta u_T$, avec $d,\eta>0$. Les coordonnées des mêmes sections satisfont
\begin{equation}
 \widehat c'=d^{-1}\widehat c,
 \qquad \tau_0'=\eta^{-1}\tau_0,
 \qquad u_L'=d\eta u_L.
 \label{x_reciprocidad:iii:eq:velocidad-cambio-carta}
\end{equation}
Si une autre publication écrit $c'=\widehat c' u_C'$ avec $u_C'=d u_C$,
l’égalité des sections équivaut exactement à $d\widehat c'=\widehat c$.
\end{proposition}
\begin{proof}
L'égalité de chaque section dans les deux bases donne $\widehat c'u_C'=\widehat c u_C$ et $\tau_0'u_T'=\tau_0u_T$. En particulier,
\[
 \tau_0'\widehat c'u_L'
 =\frac{\tau_0}{\eta}\frac{\widehat c}{d}\,d\eta u_L
 =\tau_0\widehat c u_L=\ell_0.
\]
La dernière équivalence provient de la même substitution pour la vitesse.
\end{proof}

Le choix adapté $d=\widehat c$ donne $u_C'=c_{\mathrm{vac}}$ et
$\widehat c'=1$ : il exprime la section construite dans une base adaptée.
Le coefficient de la carte interne demeure $(r_+r_-)^{-1}$, et c’est
la normalisation utilisée par l’inversion cubique. En maintenant les
bases d’action et de charge, \eqref{x_reciprocidad:iii:eq:dimensions} produit
\[
 \widehat\mu'=d\widehat\mu,\qquad
 \widehat\varepsilon'=d\widehat\varepsilon,\qquad
 \widehat\mu'\widehat\varepsilon'=(\widehat c')^{-2}.
\]
Dans la carte adaptée, on a $\widehat\mu'=r_-/r_+$ et
$\widehat\varepsilon'=r_+/r_-$, par substitution de
$\widehat\mu=r_-^2$ et $\widehat\varepsilon=r_+^2$.

\subsection{Conservation du lecteur sous amplification}

Pour le raffinement $\mathcal V_m=\mathcal V\otimes I_{9^m}$, la trace normalisée $\tau_m=\operatorname{Tr}/(2\cdot9^m)$ satisfait
\begin{equation}
 \exp[-2\tau_m(\log\mathcal V_m)]
 =\exp[-\log r_+-\log r_-]=\widehat c.
 \label{x_reciprocidad:iii:eq:velocidad-traza-normalizada}
\end{equation}
En effet, chaque valeur propre apparaît $9^m$ fois ; cette multiplicité
s’annule avec le dénominateur de la trace normalisée. Le déterminant
ordinaire amplifié est $(r_+r_-)^{9^m}$. La normalisation à deux feuillets,
exprimée de manière équivalente par \eqref{x_reciprocidad:iii:eq:velocidad-traza-normalizada},
conserve le lecteur constitutif pendant l’amplification. Cette
réalisation conserve le transport enrichi dont elle est issue.

% END OWNER


% BEGIN OWNER /Users/ruben/Documents/New project/output/EXTREMA_MEDIA_RAZON_RECIPROCIDAD_ES_EN_20260918/ARTICULO/ES/source/antecedentes_iii/sections/07_carga_y_cuantos.tex
\chapter{Section de charge et relations de quantification électrique}
\label{x_reciprocidad:iii:sec:carga}

La réponse constitutive établit une relation entre les deux feuillets et une échelle relative de propagation. La section de charge s'obtient en composant cette réponse avec la constante de structure fine et l'action, toutes deux construites dans les chapitres précédents. L'ordre a une conséquence précise : la charge ne détermine pas rétrospectivement le registre dodécaphasique. C'est le registre, avec l'action et la réponse du vide, qui détermine la section de charge dans la droite correspondante.

\section{Détermination positive de la section de charge}

Considérons les sections positives $\hbar_a$, avec
$a\in\{\mathrm{pre},\mathrm{ret}\}$, et $h_a=2\pi\hbar_a$.
L’impédance $Z_{\rm vac}$ est déjà construite au moyen de
\eqref{x_reciprocidad:iii:eq:dimensions}. Le quotient $h_a/Z_{\rm vac}$ appartient à
la droite quadratique de charge. Cette compatibilité dimensionnelle permet
de formuler la définition suivante sans choisir encore une carte SI.

\begin{definition}[Section positive de charge]
La section de charge correspondant à l’orientation d’action $a$ est
\begin{equation}
 e_a=\left(\frac{2\alpha h_a}{Z_{\rm vac}}\right)^{1/2}.
 \label{x_reciprocidad:iii:eq:charge}
\end{equation}
La racine est prise dans la demi-droite positive de charge. Sa section conjuguée
est $-e_a$ ; le signe de cette dernière n’est pas identifié à l’indice de
retour de l’action.
\end{definition}

\begin{proposition}[Existence, unicité et compatibilité du couplage]
L’équation $Z_{\rm vac}e_a^2=2\alpha h_a$ détermine une unique section
positive. Dans la réalisation constitutive, on a aussi
\begin{equation}
 \alpha=\frac{Z_{\rm vac}e_a^2}{2h_a}
 =\frac{e_a^2}{4\pi\varepsilon_{\rm vac}\hbar_a c_{\rm vac}}.
 \label{x_reciprocidad:iii:eq:coupling}
\end{equation}
\end{proposition}
\begin{proof}
La positivité de $\alpha$, $h_a$ et $Z_{\rm vac}$ garantit une unique racine positive. Les identités constitutives donnent $\varepsilon_{\rm vac}c_{\rm vac}=Z_{\rm vac}^{-1}$. Substituer cette égalité et $h_a=2\pi\hbar_a$ dans la première expression de \eqref{x_reciprocidad:iii:eq:coupling} produit la seconde.
\end{proof}

La dernière expression est la forme habituelle du couplage électromagnétique
dans une carte physique. Elle apparaît ici comme identité ultérieure des sections
construites. La preuve ne recalcule pas $\alpha$ à partir d’une charge
mesurée : elle établit la cohérence algébrique du résultat précédent avec la
section qui vient d’être définie. L’identification physique exige en outre
de maintenir la carte des unités et le régime du couplage employés dans la
comparaison. Une identité entre grandeurs et une comparaison expérimentale
ont des fonctions différentes dans cette reconnaissance.

\section{Résistance, conductance, flux et fréquence}

Les quatre grandeurs suivantes utilisent la même section $e_a$. Leur
dépendance peut être entièrement résolue avant l’évaluation de la moindre décimale.

\begin{theorem}[Composition des quanta électriques]
Les sections
\begin{align}
 R_K&=\frac{h_a}{e_a^2}=\frac{Z_{\rm vac}}{2\alpha},
 &G_0&=\frac{2e_a^2}{h_a}=\frac{4\alpha}{Z_{\rm vac}},
 \label{x_reciprocidad:iii:eq:rk-g0}\\
 \Phi_{0,a}&=\frac{h_a}{2e_a}
 =\sqrt{\frac{h_aZ_{\rm vac}}{8\alpha}},
 &K_{J,a}&=\frac{2e_a}{h_a}=\Phi_{0,a}^{-1}
 \label{x_reciprocidad:iii:eq:flux-josephson}
\end{align}
satisfont exactement
\begin{equation}
 R_KG_0=2,\qquad G_0Z_{\rm vac}=4\alpha,\qquad
 K_{J,a}\Phi_{0,a}=1,\qquad
 K_{J,a}^{2}R_K=\frac4{h_a}.
 \label{x_reciprocidad:iii:eq:quanta-identities}
\end{equation}
\end{theorem}
\begin{proof}
La substitution de $e_a^2=2\alpha h_a/Z_{\rm vac}$ dans les deux premières
définitions annule $h_a$ et donne \eqref{x_reciprocidad:iii:eq:rk-g0}. En élevant au
carré la section positive $h_a/(2e_a)$, on obtient
$h_aZ_{\rm vac}/(8\alpha)$ ; la positivité fixe la racine dans
\eqref{x_reciprocidad:iii:eq:flux-josephson}. Son inverse est $2e_a/h_a$.
Les trois premiers produits de \eqref{x_reciprocidad:iii:eq:quanta-identities} se
réduisent, respectivement, à $2$, $4\alpha$ et $1$. Enfin,
$(4e_a^2/h_a^2)(h_a/e_a^2)=4/h_a$.
\end{proof}

La réalisation physique reconnaît dans ces sections la résistance de von Klitzing, le quantum de conductance, le quantum de flux et la constante de Josephson. Le facteur $2$ de la conductance appartient à la normalisation de la grandeur définie ici ; son identification à une multiplicité de canaux d'un système matériel exige de spécifier cette représentation.
Les quatre équations ne constituent pas quatre générateurs indépendants : elles partagent $\alpha$, l'action et la réponse du vide. C'est précisément cette dépendance qui produit les identités croisées et permet de vérifier une évaluation sans altérer les opérations qui l'ont engendrée.

\section{Retour de l’action et transport des sections}

Le quotient $R_{\rm act}=h_{\rm pre}/h_{\rm ret}$, obtenu avant le
couple angulaire, conserve la différence entre les deux sections d’action.
Il ne change pas les assignations de feuillet de la réponse constitutive. Ce sont deux
opérations distinctes : retour de l’action et échange des canaux.

\begin{proposition}[Covariance sous le retour]
Avec $\alpha$ et $Z_{\rm vac}$ communs, on a
\begin{align}
 \frac{e_{\rm pre}}{e_{\rm ret}}&=R_{\rm act}^{1/2},
 &R_K^{\rm pre}&=R_K^{\rm ret},
 &G_0^{\rm pre}&=G_0^{\rm ret},\label{x_reciprocidad:iii:eq:charge-return}\\
 \frac{\Phi_{0,\rm pre}}{\Phi_{0,\rm ret}}&=R_{\rm act}^{1/2},
 &\frac{K_{J,\rm pre}}{K_{J,\rm ret}}&=R_{\rm act}^{-1/2}.
 \label{x_reciprocidad:iii:eq:flux-return}
\end{align}
\end{proposition}
\begin{proof}
La première identité découle de la division des deux versions de
\eqref{x_reciprocidad:iii:eq:charge} et de la prise de la racine positive. Dans
\eqref{x_reciprocidad:iii:eq:rk-g0}, l’action s’est déjà annulée. Les deux autres
identités découlent des puissances $h_a^{1/2}$ et $h_a^{-1/2}$ dans
\eqref{x_reciprocidad:iii:eq:flux-josephson}.
\end{proof}

La résistance et la conductance conservent le quotient constitutif, tandis que le flux et la fréquence retiennent également l'orientation de l'action. Cette séparation donne une lecture structurelle de leurs dépendances : certaines grandeurs sont des invariants du quotient des feuillets et d'autres mesurent en outre la section d'action utilisée. La comparaison métrologique doit déclarer quelle section elle évalue ; la proximité numérique des deux actions n'autorise pas à les identifier.

% END OWNER


% BEGIN OWNER /Users/ruben/Documents/New project/output/EXTREMA_MEDIA_RAZON_RECIPROCIDAD_ES_EN_20260918/ARTICULO/ES/source/antecedentes_vii/sections/v_funcional_barbero.tex
% RESULTADO_RECUPERADO desde II REV09. Véase MAPA_PROCEDENCIA.json.
% Selección de líneas y cambios mecánicos explícitos; ninguna prueba nueva.
\chapter{Incidence hexadique--octadique et fonctionnelle de Barbero--Immirzi}
\label{x_reciprocidad:v:ant:sec:barbero}
Soit $H\hookrightarrow O$ l’inclusion marquée d’une hexade dans une octade,
$|H|=6$, $|O|=8$, construite au moyen du résidu binaire de la
section~\ref{x_reciprocidad:exc:residuo-binario}. On conserve l’hexade, la dodécade
et la carte régionale qui déterminent cette inclusion.
Les objets d’incidence
\[
\mathcal F_6=H\times\binom H2,\qquad
\mathcal F_8=O\times\binom H2
\]
ont pour cardinaux $90$ et $120$. Leur différence normalisée donne
\[
\frac{|\mathcal F_6|-|\mathcal F_8|}{24\binom62}=-\frac1{12}.
\]
\begin{figure}[htbp]
\centering
\begin{tikzpicture}[>=Latex,every node/.style={font=\small}]
\draw[rounded corners,draw=ink] (-1.8,-1.35) rectangle (1.8,1.35);
\foreach \a in {0,60,...,300}{
\fill[ink] (\a:0.83) circle(2pt);}
\node at (0,-1.7) {$H:\ 6$ puntos};
\draw[->,thick,ink] (2.1,0)--(3.6,0);
\node[above] at (2.85,1.5) {inclusion marquée};
\draw[rounded corners,draw=ink] (4,-1.35) rectangle (7.6,1.35);
\foreach \a in {0,60,...,300}{
\fill[ink] ($(5.8,0)+(\a:0.83)$) circle(2pt);}
\fill[accent] (4.5,0) circle(2pt);
\fill[accent] (7.1,0) circle(2pt);
\node at (5.8,-1.7) {$O:\ 8$ puntos};
\node at (0,-2.35) {$6\binom62=90$};
\node at (5.8,-2.35) {$8\binom62=120$};
\end{tikzpicture}
\caption{Incidence marquée de six à huit points. Le dessin représente
les ensembles et l'inclusion ; il ne prétend pas transformer les faces ou sommets
d'un cube en blocs de Steiner sans leur application d'incidence.}
\end{figure}

\section{Chaîne fondamentale dodécaphasique et coefficient de pôle}
Dans le module libre des événements orientés, on fixe la chaîne
fondamentale du tour
\[
c_{12}^{+}=
\sum_{j\in\Z/12\Z}[j,\mathcal F_6]
-[\partial_{\rm ret},\mathcal F_8].
\]
La multiplicité douze correspond aux secteurs et la multiplicité un
au raccordement orienté. L'involution change l'orientation de la chaîne
complète. Soit
\[
S_m(q)=\frac{q^m}{1-q^{3m}},\qquad 0<q<1.
\]
Le lecteur linéaire attribue $S_{90}$ à chaque secteur et $S_{120}$ au raccordement ;
par conséquent,
\[
\mathscr R(c_{12}^{+})=12S_{90}-S_{120}.
\]

\begin{definition}[Évaluation orientée de l'incidence hexadique--octadique]
\label{x_reciprocidad:paper:barbero-definicion}
Le paramètre \(\gamma_{\HMT}\) est l’évaluation conjuguée de la chaîne
\(c_{12}^{+}\), complétée par sa composante angulaire. Dans le domaine
\(0<q_\pm<1\), il est défini par
\[
\gamma_{\HMT}=\frac{\pi AC^*}{180}
+[\mathscr R(c_{12}^{+})](q_-)-[\mathscr R(c_{12}^{+})](q_+).
\]
L'incidence marquée détermine les espaces d'événements ; la chaîne orientée fixe les multiplicités, et les canaux de la section~\ref{x_reciprocidad:iv:angular:inversion} déterminent l'évaluation.
\end{definition}

\begin{theorem}[Coefficient du pôle et évaluation conjuguée]
L'image de la chaîne fondamentale a pour coefficient $1/24$
devant $(1-q)^{-1}$ et pour résidu $-1/24$ devant $(q-1)^{-1}$.
Son évaluation dans les canaux préalablement engendrés détermine
\begin{equation}
\gamma_{\HMT}=
\frac{\pi AC^*}{180}
+12\bigl[S_{90}(q_-)-S_{90}(q_+)\bigr]
-\bigl[S_{120}(q_-)-S_{120}(q_+)\bigr].
\label{x_reciprocidad:v:ant:eq:gamma}
\end{equation}
\end{theorem}
\begin{proof}
La factorisation $1-q^{3m}=(1-q)\sum_{j=0}^{3m-1}q^j$ implique
$\lim_{q\uparrow1}(1-q)S_m(q)=1/(3m)$. Par linéarité,
\[
\frac{12}{270}-\frac1{360}=\frac1{24}.
\]
La coordonnée $q-1$ inverse le signe du résidu.
L’évaluation conjuguée et la composante angulaire de la définition~\ref{x_reciprocidad:paper:barbero-definicion}
produisent \eqref{x_reciprocidad:v:ant:eq:gamma}.
\end{proof}

Le certificat diophantien $4a-3b=45$ admet les couples
$(12+3t,1+4t)$, $t\ge0$ ; il certifie la minimalité de $(12,1)$,
mais ne remplace pas l'origine des multiplicités dans la chaîne.
L'évaluation de \eqref{x_reciprocidad:v:ant:eq:gamma} est
\[
\gamma_{\HMT}
=0.27400767269445927474725526077398041940\ldots.
\]
Ce chiffre est une sortie de la fonctionnelle orientée ; ce n’est pas une valeur externe
utilisée pour sélectionner $A$, $C^*$ ou leurs coefficients.

\section{Parité et sensibilité structurelle}
Inverser $C^*$ échange $q_+$ et $q_-$ et change le signe
de la composante angulaire. Ainsi,
\[
\gamma_{\HMT}(A,-C^*)=-\gamma_{\HMT}(A,C^*).
\]
La fonctionnelle conserve l'orientation, tandis que d'autres lecteurs,
comme $q_+q_-$, sont pairs. Cette différence précède son interprétation physique.

% END OWNER


% BEGIN OWNER /Users/ruben/Documents/New project/output/EXTREMA_MEDIA_RAZON_RECIPROCIDAD_ES_EN_20260918/ARTICULO/ES/source/antecedentes_vii/sections/v_unidad_areal.tex
% RESULTADO_RECUPERADO desde II REV09. Véase MAPA_PROCEDENCIA.json.
% Selección de líneas y cambios mecánicos explícitos; ninguna prueba nueva.
\chapter{Action, gravité et information thermique}
\label{x_reciprocidad:v:ant:sec:gravity}
\section{Période dodécaphasique et réalisation circulaire de l'action}
Le calendrier déjà construit détermine $T_{108}=108t_0$.
La vitesse constitutive, la durée et la longueur de trajectoire sont
obtenues au moyen des identités de la
proposition~\ref{x_reciprocidad:iii:prop:identidad-velocidad} et des
équations~\eqref{x_reciprocidad:iii:eq:velocidad-longitud-reloj} et
\eqref{x_reciprocidad:iii:eq:longitud-elemental-constitutiva}.
Dans sa réalisation circulaire,
\[
\omega_{108}=\frac{2\pi}{108t_0},\qquad
R_{108}=\frac{c_{\rm int}}{\omega_{108}}
=\frac{54}{\pi}\ell_0,\qquad \ell_0=c_{\rm int}t_0.
\]
L'unité $\ell_0$ et les cartes de temps et de vitesse restent
explicites ; la carte SI ne sélectionne pas le cycle.
Pour chaque section orientée d'action $a\in\{\mathrm{pre},\mathrm{ret}\}$,
\[
E_{108,a}=\hbar_a\omega_{108},\qquad
m_{108,a}=\frac{E_{108,a}}{c_{\rm int}^2}.
\]
Les longueurs réduite et complète sont respectivement
\[
\frac{\hbar_a}{m_{108,a}c_{\rm int}}=R_{108},
\qquad
\frac{h_a}{m_{108,a}c_{\rm int}}=2\pi R_{108}.
\]
La même action et la même durée interviennent dans les deux lectures.

\begin{proposition}[Sélecteur de coïncidence des rayons]
Dans la famille dimensionnelle
\[
G_a(\vartheta)=\vartheta\frac{c_{\rm int}^5t_0^2}{\hbar_a},
\qquad \vartheta>0,
\]
la condition de la réalisation circulaire
$G_a(\vartheta)m_{108,a}/c_{\rm int}^2=R_{108}$
admet l'unique solution
\[
\vartheta^\circ_{108}=(54/\pi)^2.
\]
\end{proposition}
\begin{proof}
Le premier rayon est
$r_g(\vartheta)=\vartheta(\pi/54)\ell_0$.
Le comparer à $(54/\pi)\ell_0$, avec $\ell_0>0$, réduit
l'égalité à une équation linéaire strictement croissante.
\end{proof}
La condition de coïncidence est une partie explicite de cette réalisation ;
le théorème résout son coefficient sans utiliser de valeur métrologique de $G$.
Le rayon employé ici est $Gm/c^2$, de sorte que le facteur deux
de la convention $2Gm/c^2$ n'est pas transféré par inadvertance.

\begin{proposition}[Action et gravité réciproques, aire conservée]
Dans les deux cartes précédentes,
\[
\frac{G_{\rm pre}}{G_{\rm ret}}=R_{\rm act}^{-1},\qquad
\hbar_aG_a=\vartheta^\circ_{108}c_{\rm int}^5t_0^2,
\qquad
\ell_P^2=\frac{\hbar_aG_a}{c_{\rm int}^3}
=\vartheta^\circ_{108}\ell_0^2.
\]
\end{proposition}
\begin{proof}
La définition de $G_a$ contient $\hbar_a^{-1}$ et les autres facteurs
sont communs aux deux cartes. Prendre le quotient et le produit donne
les trois égalités.
\end{proof}
Cette réciprocité explique pourquoi l'information d'orientation peut
varier sans que l'aire change. On dispose ainsi d'une relation précise
entre les sections réciproques d'action, l'échelle de Planck et l'opérateur surfacique
de la section~\ref{x_reciprocidad:v:ant:sec:area}.

% END OWNER


% BEGIN OWNER /Users/ruben/Documents/New project/output/EXTREMA_MEDIA_RAZON_RECIPROCIDAD_ES_EN_20260918/ARTICULO/ES/source/antecedentes_vii/sections/v_area_informacion.tex
% RESULTADO_RECUPERADO desde II REV09. Véase MAPA_PROCEDENCIA.json.
% Selección de líneas y cambios mecánicos explícitos; ninguna prueba nueva.
\chapter{Opérateur d'aire et reconstruction de l'information sectorielle}
\label{x_reciprocidad:v:ant:sec:area}
Le regroupement de six trits construit une bijection entre
$(\Z/9\Z)^3$ et $(\Z/27\Z)^2$, tous deux de cardinal $729$.
La bijection n'est pas présentée comme un isomorphisme additif. Sur un bloc
de bord $9\times9$, la source décompose 36 liens en deux
ensembles de 18, étiquetés par les canaux $90$ et $120$.

\section{Regroupement tritique et incidence du bord}
\label{x_reciprocidad:v:ant:sec:ii-area-incidencia-explicita}

La correspondance conserve un mot ordonné, en plus de son cardinal.
Écrivons chaque coordonnée \(x_j\in\mathbb Z/9\mathbb Z\), pour
\(j=1,2,3\), au moyen de ses représentants positionnels uniques
\(x_j=r_{2j-1}+3r_{2j}\), avec \(r_k\in\{0,1,2\}\). L'application est
\begin{equation}
 \beta_{729}(x_1,x_2,x_3)
 =\bigl(r_1+3r_2+9r_3,\ r_4+3r_5+9r_6\bigr)
 \in(\mathbb Z/27\mathbb Z)^2.
 \label{x_reciprocidad:v:ant:eq:ii-area-beta729}
\end{equation}

\begin{proposition}[Correspondance positionnelle entre intérieur et bord]
\label{x_reciprocidad:v:ant:prop:ii-area-beta729}
L'application \(\beta_{729}\) est bijective et conserve l'ordre des
six trits. Son inverse extrait les deux triplets de chiffres et les regroupe
en trois paires.
\end{proposition}
\begin{proof}
Le développement positionnel des représentants \(0,\ldots,8\) et
\(0,\ldots,26\) est unique. Extraire les six chiffres, les regrouper puis
les extraire de nouveau restitue le même mot. L'opération inverse
rétablit ainsi chacune des trois coordonnées initiales.
\end{proof}

La retenue arithmétique conserve sa position dans le mot enrichi. Les opérations
additives des deux regroupements ont des frontières différentes : par exemple,
\(\beta_{729}(0,2,0)=(18,0)\) et
\(\beta_{729}(0,1,0)=(9,0)\), dont la somme dans le tore d'ordre
vingt-sept est \((0,0)\), tandis que
\(\beta_{729}(0,3,0)=(0,1)\). Cette distinction permet de réutiliser
la correspondance positionnelle en maintenant le typage des opérations.

Dans le tore \((\mathbb Z/27\mathbb Z)^2\), le bloc
\(B_\partial=\{0,\ldots,8\}^2\) possède les ensembles de liens
orientés
\begin{align}
 \partial B_{90}
 &=\{((-1,j),(0,j)),\ ((8,j),(9,j)):0\le j\le8\},\\
 \partial B_{120}
 &=\{((i,-1),(i,0)),\ ((i,8),(i,9)):0\le i\le8\}.
 \label{x_reciprocidad:v:ant:eq:ii-area-enlaces}
\end{align}
Toutes les coordonnées se lisent modulo vingt-sept. Chaque famille contient
dix-huit liens et elles sont disjointes. Leurs étiquettes conservent les
canaux d'incidence ; le cardinal des liens est \(18+18\).

\begin{theorem}[Capacité minimale de la coupe cubique]
\label{x_reciprocidad:v:ant:thm:ii-area-corte-minimo}
Dans le graphe cubique de sommets \(\{0,\ldots,N-1\}^3\), avec
\(N\ge2\), des liens entre voisins et une capacité unitaire par lien,
la coupe minimale entre deux faces opposées a pour capacité \(N^2\).
Pour \(N=9\), un état produit maximalement mélangé sur les
trits d'une couche minimale possède l'information
\begin{equation}
 \operatorname{cap}_{\min}=81,\qquad s_{\rm corte}=81\log3.
 \label{x_reciprocidad:v:ant:eq:ii-area-corte-informacion}
\end{equation}
\end{theorem}
\begin{proof}
Les \(N^2\) droites parallèles à la normale aux faces forment des chemins
disjoints par les liens. Toute coupe séparatrice intercepte chaque chemin, de
sorte qu'elle contient au moins \(N^2\) liens. Les liens traversant
une couche transversale constituent une coupe de ce cardinal. Chaque trit
uniforme apporte \(\log3\) nats ; l'entropie est additive pour l'état
produit des quatre-vingt-un trits de la couche.
\end{proof}

La capacité de la coupe cubique et l'état des trente-six liens
de \eqref{x_reciprocidad:v:ant:eq:ii-area-enlaces} appartiennent à des observables définies sur
des coupes différentes. La généalogie conserve les deux domaines et leurs mesures.

\section{Transport des modes collectifs d'incidence}
\label{x_reciprocidad:v:ant:sec:ii-area-transporte-colectivo}

Les espaces \(\mathcal F_6\) et \(\mathcal F_8\) de la
section~\ref{x_reciprocidad:v:ant:sec:barbero} fournissent les vecteurs uniformes
\[
 u_{90}=\frac1{\sqrt{90}}\sum_{f\in\mathcal F_6}(\delta_f,0),
 \qquad
 u_{120}=\frac1{\sqrt{120}}\sum_{f\in\mathcal F_8}(0,\delta_f)
\]
dans \(\ell^2(\mathcal F_6)\oplus\ell^2(\mathcal F_8)\).
Dans \(\ell^2(\partial B_{90})\oplus\ell^2(\partial B_{120})\),
on construit
\[
 v_{90}=\frac1{\sqrt{18}}\sum_{e\in\partial B_{90}}(\delta_e,0),
 \qquad
 v_{120}=\frac1{\sqrt{18}}\sum_{e\in\partial B_{120}}(0,\delta_e).
\]
Les deux paires sont orthonormales. Notons
\(\mathscr H_{\rm inc}^{\rm coll}\) et
\(\mathscr H_\partial^{\rm coll}\) les plans qu'elles engendrent respectivement.

\begin{theorem}[Entrelacement collectif des étiquettes sectorielles]
\label{x_reciprocidad:v:ant:thm:ii-area-entrelazamiento-colectivo}
Il existe une unique isométrie linéaire surjective entre ces plans telle que
\begin{equation}
 \mathcal J_{6\to8}u_m=v_m\quad(m=90,120).
 \label{x_reciprocidad:v:ant:eq:ii-area-j-colectivo}
\end{equation}
Pour les opérateurs
\[
 \begin{aligned}
 D_{\rm inc}&=90|u_{90}\rangle\langle u_{90}|
             +120|u_{120}\rangle\langle u_{120}|,\\
 D_\partial&=90|v_{90}\rangle\langle v_{90}|
             +120|v_{120}\rangle\langle v_{120}|,
 \end{aligned}
\]
on a
\begin{equation}
 \mathcal J_{6\to8}D_{\rm inc}
 =D_\partial\mathcal J_{6\to8}.
 \label{x_reciprocidad:v:ant:eq:ii-area-j-entrelaza}
\end{equation}
Les projecteurs sectoriels et leur combinaison orientée de coefficients
\(12:-1\) sont également transportés.
\end{theorem}
\begin{proof}
L'image d'une base fixe l'unique application linéaire. Les normes et les
produits scalaires sont conservés puisque les deux bases sont orthonormales.
L'identité \eqref{x_reciprocidad:v:ant:eq:ii-area-j-entrelaza} vaut sur chaque vecteur
\(u_m\), ses deux membres étant \(m v_m\). Les projecteurs
s'expriment comme \((120I-D)/30\) et \((D-90I)/30\) ; leur calcul fonctionnel
linéaire et toute combinaison des deux conservent l'entrelacement.
\end{proof}

Ce transport agit sur les deux modes uniformes et leurs étiquettes
spectrales. Les incidences individuelles et les fluctuations au sein de
chaque secteur demeurent dans leurs espaces complets. Le lien conserve
exactement la donnée collective qu'utilisera la pondération aréale.

\section{Caractère angulaire et opérateur d'aire}
\label{x_reciprocidad:v:ant:sec:ii-area-caracter}

Soit $N_e=\operatorname{diag}(0,1,2)$ sur chaque lien et
$N_{90},N_{120}$ les sommes par secteur. L'échelle étant fixée par
\[
\theta_{\rm clk}=\frac{C^*}{6}\frac{\pi}{180},\qquad
a_0=\frac{1+2\cos(10^\circ)}3\theta_{\rm clk},
\]
le facteur angulaire est le caractère réel normalisé du triplet
\begin{equation}
 \chi_{10}=\frac13\Tr\operatorname{diag}
       (1,e^{i\pi/18},e^{-i\pi/18})
 =\frac{1+2\cos(10^\circ)}3,
 \qquad a_0=\chi_{10}\theta_{\rm clk}.
 \label{x_reciprocidad:v:ant:eq:ii-area-caracter-triplete}
\end{equation}
La somme des phases conjuguées prouve l'égalité. Dans la branche positive
considérée, \(\theta_{\rm clk}>0\), \(\chi_{10}>0\) et
\(\gamma_{\HMT}>0\) ; ainsi, \(a_0\) et l'échelle aréale sont
positifs. Le triplet angulaire et sa normalisation précèdent la
distribution entropique sur les liens. Avec ces quantités,
l'opérateur aréal normalisé est
\[
\widehat{\mathcal A}/\ell_P^2
=\gamma_{\HMT}a_0(N_{90}+N_{120}).
\]
Son spectre a pour valeurs $n\gamma_{\HMT}a_0$, $0\le n\le72$,
avec les multiplicités données par $(1+z+z^2)^{36}$.
La preuve consiste à sommer les valeurs propres $0,1,2$ des
36 facteurs tensoriels.

\begin{theorem}[Récupération des secteurs de bord]
Soient
\[
N=N_{90}+N_{120},\qquad D=90N_{90}+120N_{120}.
\]
Alors
\begin{equation}
N_{90}=4N-\frac D{30},\qquad
N_{120}=\frac D{30}-3N.
\label{x_reciprocidad:v:ant:eq:recover}
\end{equation}
Ainsi, le total et le lecteur pondéré retrouvent conjointement
les deux occupations.
\end{theorem}
\begin{proof}
La matriz $\left(\begin{smallmatrix}1&1\\90&120\end{smallmatrix}\right)$
a pour déterminant $30$. Son inversion donne \eqref{x_reciprocidad:v:ant:eq:recover}.
Dans le domaine des occupations entières, l'image satisfait
$D\in30\Z$ et $90N\le D\le120N$.
\end{proof}

Le Hamiltonien modulaire de la source est
$K_A=-\log\rho_A=cI+\Delta_{\rm mod}D$.
Pour $\Delta_{\rm mod}\ne0$ et une normalisation connue,
$D=(K_A-cI)/\Delta_{\rm mod}$ publie la deuxième observable.
Cet opérateur décrit la structure spectrale de l'état réduit ;
il ne s'identifie pas, par la notation, au générateur de l'évolution temporelle.
L'aire conserve l'occupation totale et la paire ajoute la provenance.
Par exemple, $(N_{90},N_{120})=(1,0)$ et $(0,1)$ partagent $N=1$,
mais ont $D=90$ et $D=120$.


% BEGIN OWNER /Users/ruben/Documents/New project/output/EXTREMA_MEDIA_RAZON_RECIPROCIDAD_ES_EN_20260918/ARTICULO/ES/source/antecedentes_vii/sections/v_estado_modular.tex
% RESULTADO_RECUPERADO desde II REV09. Véase MAPA_PROCEDENCIA.json.
% Selección de líneas y cambios mecánicos explícitos; ninguna prueba nueva.
\section{Normalisation de l'état et signification du Hamiltonien modulaire}
La deuxième observable est explicitée par l'état qui la définit.
La décomposition tensorielle précédente est maintenue : dix-huit qutrits
par secteur, chacun muni de l'opérateur nombre $\operatorname{diag}(0,1,2)$ ;
$N_m$ est la somme de ces opérateurs et les secteurs agissent sur des facteurs
distincts. Soit
\[
 z_m=1+e^{-m\Delta_{\rm mod}}+e^{-2m\Delta_{\rm mod}},
 \qquad m\in\{90,120\}.
\]
La famille d'états de la source admet la réalisation produit
\[
 \rho_A=Z^{-1}e^{-\Delta_{\rm mod}(90N_{90}+120N_{120})},
 \qquad Z=z_{90}^{18}z_{120}^{18}.
\]
La commutativité des opérateurs nombre et la factorisation tensorielle
prouvent l'expression de $Z$. L'état est de rang plein pour
$\Delta_{\rm mod}$ réel et fini, et son logarithme spectral donne
\[
 K_A=-\log\rho_A=(\log Z)I+\Delta_{\rm mod}D.
\]
Pour $\Delta_{\rm mod}\ne0$, $K_A$ pondère différemment les excitations
des deux secteurs. La lecture conjointe avec $N$ retrouve les deux occupations
par l'équation~\eqref{x_reciprocidad:v:ant:eq:recover}.

La nécessité de la lecture conjointe apparaît dans deux exemples complémentaires.
Les états d'occupation $(1,0)$ et $(0,1)$ ont le même total $N=1$
et un poids modulaire différent. Inversement, $(4,0)$ et $(0,3)$ ont
le même $D=360$ et des totaux différents. La paire $(N,D)$ sépare les deux
cas. Cette récupération concerne les occupations sectorielles, et non
chaque vecteur microscopique au sein d'une dégénérescence de ces occupations.

Lorsque $\Delta_{\rm mod}=0$, l'état est proportionnel à l'identité
et $K_A$ est scalaire : cette lecture cesse de distinguer les secteurs.
Pour $\Delta_{\rm mod}\ne0$, la normalisation connue permet de retrouver
$D$ et de le composer avec l'aire. Le rôle du Hamiltonien
modulaire est ainsi défini : il caractérise l'état réduit et sa distribution d'information.
Son inclusion dans l'article est motivée par cette fonction mathématique,
tandis que le Hamiltonien d'évolution appartient au développement dynamique
de l'action.

% END OWNER


\section{État réduit, décomposition de Schmidt et entropie d'intrication}
Si $\rho_A=\sum_\nu p_\nu|\nu\rangle\langle\nu|$,
la purification
\[
|\Psi\rangle=\sum_\nu\sqrt{p_\nu}\,
|\nu\rangle_A|\nu\rangle_B
\]
satisfait $\Tr_B|\Psi\rangle\langle\Psi|=\rho_A$.
Ainsi, $-\Tr(\rho_A\log\rho_A)$ est l'entropie d'intrication
de cette bipartition pure précise. L'état élargi déclare
quel système porte cette interprétation.

Pour une involution à deux feuillets possédant un vecteur propre de chaque signe,
\[
\rho_y=\frac{e^{-yR}}{2\cosh y},\qquad R^2=I,
\]
la diagonalisation fournit
\[
S(\rho_y)=\log(2\cosh y)-y\tanh y,\qquad
D(\rho_y\Vert\rho_{-y})=2y\tanh y.
\]
L'entropie est paire en $y$ ; la distance compare explicitement
les deux orientations. Cette entropie à deux niveaux ne s'identifie pas
à la capacité $81\log3$ de la coupe triadique.

% END OWNER


% BEGIN OWNER /Users/ruben/Documents/New project/output/EXTREMA_MEDIA_RAZON_RECIPROCIDAD_ES_EN_20260918/ARTICULO/ES/source/antecedentes_vii/sections/v_transduccion_termica.tex
% RESULTADO_RECUPERADO desde II REV09. Véase MAPA_PROCEDENCIA.json.
% Selección de líneas y cambios mecánicos explícitos; ninguna prueba nueva.
% Fragmento aditivo. No sustituye 10b_boltzmann.tex.
% Procedencia y distinción entre recuperación y formalización:
% 10c_boltzmann_procedencia.md. Etiquetas locales prefijadas ii-kb-aut.

\chapter{Mémoire généalogique, réduction informationnelle et transduction thermique}
\label{x_reciprocidad:v:ant:sec:ii-kb-aut-desarrollo}

La construction thermique utilise deux publications du même développement APP--TRIT--TPK. La première conserve les états, les routes admissibles et l'information du registre ; la seconde fournit l'action orientée et la durée du retour. Leur composition détermine une énergie par décision tritique. Les coordonnées de température et d'énergie sont introduites ensuite au moyen d'applications dimensionnelles explicites. Cette séquence permet d'identifier quelle information est conservée, quelle information publie un lecteur et quelle grandeur convertit cette information en entropie dimensionnelle.

\section{Lecteur, mémoire et récupération de l'état}

Fixons un niveau fini du développement et notons \(\Omega\) l'
ensemble des états enrichis admissibles. Le registre conserve feuillet,
orientation, résidu, quotient, retenue, route, phase et mémoire dans les
domaines où intervient chaque coordonnée. Soient
\[
 q:\Omega\longrightarrow Y,\qquad
 M:\Omega\longrightarrow\mathcal M
\]
la publication visible et la mémoire retenue, respectivement. Pour une
distribution \(p\) sur \(\Omega\), nous écrivons
\(\mathsf H(p)=-\sum_xp_x\log p_x\), avec \(0\log0=0\).
Le logarithme s'applique à des probabilités adimensionnelles.

\begin{definition}[Récupération généalogique d'une publication]
Le couple \(\widehat q=(q,M)\) est un lecteur récupérable sur le support
\(\Omega_p\) lorsqu'il existe
\[
 d:\widehat q(\Omega_p)\longrightarrow\Omega_p,
 \qquad d\circ\widehat q=\operatorname{id}_{\Omega_p}.
\]
L'information de fibre de \(q\) est
\(\mathsf H_{\mathrm{fib}}(q;p)=\mathsf H(X\mid q(X))\), où
\(X\) a pour loi \(p\).
\end{definition}

\begin{proposition}[Identité de récupération et coût de la réduction]
Pour un lecteur récupérable,
\begin{align}
 \mathsf H(X\mid q(X),M(X))&=0,\\
 \mathsf H(X)&=\mathsf H(q(X))+\mathsf H(X\mid q(X)),\\
 \mathsf H(X\mid q(X))&=\mathsf H(M(X)\mid q(X)).
 \label{x_reciprocidad:v:ant:eq:ii-kb-aut-fibra}
\end{align}
\end{proposition}
\begin{proof}
L'application \(d\) restitue \(X\) à partir du couple \((q(X),M(X))\),
de sorte que la première entropie conditionnelle est nulle. Comme \(q\) est
une fonction de \(X\), la règle de la chaîne donne la seconde égalité.
Pour la troisième, on développe de deux manières
\(\mathsf H(X,M(X)\mid q(X))\). L'une fournit
\(\mathsf H(X\mid q(X))\), puisque \(M\) est une fonction de \(X\) ;
l'autre fournit
\(\mathsf H(M(X)\mid q(X))+\mathsf H(X\mid q(X),M(X))\).
\end{proof}

La division euclidienne des feuillets d'APP montre le mécanisme
arithmétique élémentaire. Avec \(i,j\in\{1,\ldots,9\}\), on conserve
\[
 i+j=9q^+_{ij}+R^+_{ij},\qquad
 ij=9q^\times_{ij}+R^\times_{ij},\qquad
 1\leq R^\sigma_{ij}\leq9\quad(\sigma\in\{+,\times\}).
\]
La paire résidu--quotient récupère la valeur entière de l'opération. En additionnant les différences sur les quatre-vingt-une positions, on obtient
\[
 2025-810=54+9\cdot129.
\]
La mémoire du quotient restaure la contribution que ne publie pas la
réduction résiduelle. La récupérabilité de l'état complet exige
en outre de conserver les coordonnées de position et de route pertinentes ;
le couple résidu--quotient identifie ici le résultat de l'opération.

\section{Mise à jour réversible, décision tritique et effacement}

Soit \(U:\Omega\to\Omega\) la mise à jour TPK dans un domaine où le transport enrichi possède une inverse. La loi transportée est \(p'_y=p_{U^{-1}y}\). La bijection conserve le multiensemble des probabilités et, par conséquent,
\begin{equation}
 \mathsf H(U(X))=\mathsf H(X),\qquad
 \mathsf H(X\mid U(X))=0.
 \label{x_reciprocidad:v:ant:eq:ii-kb-aut-reversible}
\end{equation}
Si la sortie observée est \(q(U(X))\), l'information non publiée
est quantifiée par \(\mathsf H(X\mid q(U(X)))\). L'état complet,
sa publication et la réinitialisation d'un registre sont ainsi trois applications distinctes.

Pour les trois états orientés du TRIT, soit
\(p=(p_{-1},p_0,p_{+1})\) une distribution normalisée. Son information est
\begin{equation}
 I_{\mathrm{TRIT}}(p)=-\sum_{a=-1}^{1}p_a\log p_a,
 \qquad 0\leq I_{\mathrm{TRIT}}(p)\leq\log3.
 \label{x_reciprocidad:v:ant:eq:ii-kb-aut-trit}
\end{equation}
L'extrémité supérieure correspond à \(p_a=1/3\). En effet,
\[
 D(p\Vert(1/3,1/3,1/3))
 =\sum_ap_a\log(3p_a)=\log3-I_{\mathrm{TRIT}}(p)\geq0;
\]
l'inégalité s'obtient à partir de la convexité de \(t\log t\), et l'égalité caractérise la distribution uniforme. L'extrémité inférieure correspond à une distribution concentrée sur un état.

La réalisation idéale de Landauer utilise une mémoire aux états logiques
énergétiquement dégénérés, un environnement isotherme à température \(\Theta\)
et un bilan qui conserve la sortie de l'application logique et comptabilise les
registres auxiliaires. Dans ce modèle, la diminution de l'entropie logique est
\(k_B[\mathsf H(X)-\mathsf H(F(X))]
=k_B\mathsf H(X\mid F(X))\) pour une application déterministe \(F\).
Le bilan entropique exige de transférer au milieu au moins cette quantité ;
son expression énergétique est \(Q\geq k_B\Theta\mathsf H(X\mid F(X))\).
La distinction entre évolution logique réversible et effacement physique
correspond à la réalisation de Landauer--Bennett \cite{x_reciprocidad:bennett2003}.
La mise à jour inversible
\eqref{x_reciprocidad:v:ant:eq:ii-kb-aut-reversible} a une contribution logique nulle ; la
réinitialisation d'un triplet uniforme exige le transfert de \(\log3\) nats
au milieu et satisfait \(Q\geq k_B\Theta\log3\).
L'énergie de contrôle et la dissipation d'une réalisation physique
concrète appartiennent à son bilan thermodynamique supplémentaire.

\section{Action, fréquence et énergie de décision}

Soient \(\mathcal L_{\rm act}\) et \(\mathcal L_T\) les droites
dimensionnelles d'action et de durée ; la droite énergétique est
\(\mathcal L_E=\mathcal L_{\rm act}\otimes\mathcal L_T^*\).
Le développement précédent fournit les sections positives
\(\hbar_a\in\mathcal L_{\rm act}\), avec
\(a\in\{\mathrm{pre},\mathrm{ret}\}\), et la durée élémentaire
\(t_0\in\mathcal L_T\). Le calendrier complet fournit
\[
 T_{108}=108t_0,\qquad
 \omega_{108}=\frac{2\pi}{T_{108}},\qquad
 h_a=2\pi\hbar_a.
\]
La contraction dimensionnelle entre action et fréquence définit
\begin{equation}
 \epsilon_a:=\hbar_a\omega_{108}
 =\frac{h_a}{108t_0}
 =\frac{\hbar_a}{t_P},\qquad
 t_P=\frac{54}{\pi}t_0.
 \label{x_reciprocidad:v:ant:eq:ii-kb-aut-energia}
\end{equation}
Le dernier membre emploie la même période circulaire réduite que la section gravitationnelle. Les trois expressions correspondent à une unique section d'énergie, avec une orientation \(a\) maintenue dans tous les facteurs. L'information normalisée d'une décision uniforme produit la section d'énergie par nat
\begin{equation}
 \varepsilon_a:=\frac{\epsilon_a}{\log3}\in\mathcal L_E^+.
 \label{x_reciprocidad:v:ant:eq:ii-kb-aut-energia-nat}
\end{equation}
L'action précédemment construite détermine l'échelle d'énergie lorsque
la durée est fixée ; le cardinal des alternatives du TRIT détermine sa
normalisation informationnelle. La valeur de \(\log3\) intervient par le
dénombrement et la mesure déclarés dans \eqref{x_reciprocidad:v:ant:eq:ii-kb-aut-trit}.

\section{Définition dimensionnelle de Boltzmann et sélection des coordonnées}

Soit \(\mathcal L_\Theta\) la droite orientée de température. L'espace dimensionnel d'entropie est \(\mathcal L_{\rm ent}=\mathcal L_E\otimes\mathcal L_\Theta^*\). Une application linéaire positive \(k_B:\mathcal L_\Theta\to\mathcal L_E\) est, de manière équivalente, un élément positif de \(\mathcal L_{\rm ent}\). C'est pourquoi le même objet peut agir sur une température ou multiplier une information sans dimension.

\begin{definition}[Normalisation thermique de la décision tritique]
Un couple thermique compatible avec la section orientée \(a\) se compose d'
un transducteur positif \(k_B\) et d'une section positive
\(\Theta_{{\rm clk},a}\in\mathcal L_\Theta\) tels que
\begin{equation}
 k_B(\Theta_{{\rm clk},a})=\varepsilon_a,
 \qquad
 k_B\Theta_{{\rm clk},a}\log3=\epsilon_a.
 \label{x_reciprocidad:v:ant:eq:ii-kb-aut-par-termico}
\end{equation}
La constante de Boltzmann remplit la fonction de transducteur
énergie--température de l'information tritique dans cette réalisation.
\end{definition}

\begin{theorem}[Unicité relative à la section thermique et covariance dimensionnelle]
Une section positive \(\Theta_{{\rm clk},a}\) étant fixée, il existe une unique application linéaire positive qui satisfait \eqref{x_reciprocidad:v:ant:eq:ii-kb-aut-par-termico}. Pour des bases positives \(u_E,u_\Theta\), si
\[
 \epsilon_a=\epsilon_a^{\rm num}u_E,\quad
 \Theta_{{\rm clk},a}=\theta_a u_\Theta,\quad
 k_B(u_\Theta)=b\,u_E,
\]
sa coordonnée est
\begin{equation}
 b=\frac{\epsilon_a^{\rm num}}{\theta_a\log3}.
 \label{x_reciprocidad:v:ant:eq:ii-kb-aut-coordenada}
\end{equation}
Sous \(u'_E=s_Eu_E\), \(u'_\Theta=s_\Theta u_\Theta\), avec
\(s_E,s_\Theta>0\), on a
\[
 b'=\frac{s_\Theta}{s_E}b,\qquad
 \theta'_a=\frac{\theta_a}{s_\Theta},\qquad
 (\epsilon_a^{\rm num})'=\frac{\epsilon_a^{\rm num}}{s_E}.
\]
\end{theorem}
\begin{proof}
Une section non nulle engendre une droite. Prescrire son image détermine
l'unique application linéaire ; la positivité des deux sections prouve
sa positivité. Substituer les coordonnées dans
\eqref{x_reciprocidad:v:ant:eq:ii-kb-aut-par-termico} donne
\(b\theta_a\log3=\epsilon_a^{\rm num}\). Les trois lois de
transformation s'obtiennent en exprimant les mêmes sections et la même
application dans les nouvelles bases.
\end{proof}

Une même constante de Boltzmann doit agir dans les deux orientations.
Puisque \(t_P\) est commun, la compatibilité simultanée équivaut à
\begin{equation}
 \frac{\Theta_{{\rm clk},{\rm pre}}}{\Theta_{{\rm clk},{\rm ret}}}
 =\frac{\epsilon_{\rm pre}}{\epsilon_{\rm ret}}
 =\frac{\hbar_{\rm pre}}{\hbar_{\rm ret}}.
 \label{x_reciprocidad:v:ant:eq:ii-kb-aut-orientaciones}
\end{equation}
En effet, un isomorphisme de droites conserve les quotients de sections. Réciproquement, si l'égalité précédente est satisfaite, l'isomorphisme déterminé dans une orientation satisfait également la normalisation de l'autre. Ainsi, le retour modifie la section thermique du cycle dans la même proportion que l'énergie, tandis que le transducteur reste commun aux deux lectures.

L'équation précédente explicite le choix requis pour publier une
valeur individuelle. Si la section thermique n'a pas encore été sélectionnée
par un lecteur indépendant, pour chaque \(\lambda>0\), le couple
\[
 (k_B,\Theta_{{\rm clk},a})
 \longmapsto(\lambda k_B,\lambda^{-1}\Theta_{{\rm clk},a})
\]
conserve \(\epsilon_a\). Par conséquent, la présentation d'une
coordonnée individuelle sélectionnée en interne exige de donner le lecteur de
\(\Theta_{{\rm clk},a}\), les bases et leur transport dimensionnel.
Les équations d'action et d'information restent déterminées avant
cette spécification. L'identification ultérieure des bases au
joule et au kelvin appartient à la publication métrologique du résultat.

\section{Entropie dimensionnelle et information de l'état réduit}

Pour une distribution de routes admises \(p\), ou pour un état réduit de matrice de densité \(\rho\), soient
\[
 s(p)=-\sum_xp_x\log p_x,\qquad
 s(\rho)=-\operatorname{Tr}(\rho\log\rho).
\]
Ses publications dimensionnelles sont
\begin{equation}
 S(p)=k_Bs(p),\qquad S(\rho)=k_Bs(\rho)
 \quad\text{dans }\mathcal L_{\rm ent}.
 \label{x_reciprocidad:v:ant:eq:ii-kb-aut-entropia}
\end{equation}
L'évaluation en \(\Theta\) produit une énergie
\(\Theta S=k_B\Theta\,s\). Pour le triplet uniforme à la température
du cycle, \eqref{x_reciprocidad:v:ant:eq:ii-kb-aut-par-termico} donne exactement
\[
 S_{\rm TRIT}=k_B\log3,\qquad
 \Theta_{{\rm clk},a}S_{\rm TRIT}=\epsilon_a.
\]
Si \(\rho_A\) est la réduction de la purification bipartite construite dans la section surfacique, \(s(\rho_A)\) mesure l'intrication de cette bipartition. Le transport des unités conserve ses valeurs propres, son orientation et les occupations sectorielles. L'échelle surfacique \(\gamma a_0\ell_P^2\) convertit les occupations en aire ; le transducteur \(k_B\) convertit l'information en entropie dimensionnelle. Les deux lectures utilisent le même registre au moyen d'opérateurs distincts.

\section{Normalisation spectrale des histoires thermiques}

Soit \(\mathcal G\) un graphe fini d'extensions admissibles déjà
construit, de matrice d'adjacence \(A\). Chaque arête \(e\) conserve
ses coordonnées de route et reçoit un incrément adimensionnel d'action
\(a_*(e)>0\) du lecteur correspondant. Une fois choisie la section d'énergie
\(\epsilon_*\), posons \(t=\beta\epsilon_*\), où
\(\beta=(k_B\Theta)^{-1}\). L'opérateur thermique est
\begin{equation}
 (L_tf)(y)=\sum_{e:x\to y}\exp[-t a_*(e)]f(x).
 \label{x_reciprocidad:v:ant:eq:ii-kb-aut-operador}
\end{equation}
Tous ses exposants sont sans dimension. La condition
\(\rho(L_t)=1\) normalise la croissance pondérée des histoires.

% REMATE_LOCAL_X_TERMICA
\Needspace{35\baselineskip}
\begin{proposition}[Existence et unicité de la pression normalisée]
Supposons que \(A\) est irréductible, \(\rho(A)>1\), et que
\(0<m\leq a_*(e)\leq M\) pour toutes les arêtes. Il existe un unique
\(t_*>0\) tel que \(\rho(L_{t_*})=1\), et
\begin{equation}
 \frac{\log\rho(A)}{M}\leq t_*\leq
 \frac{\log\rho(A)}{m}.
 \label{x_reciprocidad:v:ant:eq:ii-kb-aut-presion}
\end{equation}
\end{proposition}
\begin{proof}
La comparaison entrée par entrée fournit \(e^{-tM}A\leq L_t\leq e^{-tm}A\) pour \(t\geq0\). La monotonie du rayon spectral des matrices non négatives implique les bornes correspondantes pour \(\rho(L_t)\). Le rayon spectral est continu, vaut \(\rho(A)>1\) en zéro et tend vers zéro à l'infini. Pour \(s>0\), en outre, \(L_{t+s}\leq e^{-sm}L_t\), de sorte que \(\rho(L_{t+s})\leq e^{-sm}\rho(L_t)<\rho(L_t)\). La continuité et la décroissance stricte prouvent l'existence et l'unicité. Substituer la valeur un dans les bornes initiales démontre \eqref{x_reciprocidad:v:ant:eq:ii-kb-aut-presion}.
\end{proof}

La réalisation locale de trois alternatives équipondérées, représentée
par un sommet avec trois boucles et \(a_*(e)=1\), satisfait
\(\rho(L_t)=3e^{-t}\). Sa normalisation donne \(t_*=\log3\).
Avec \(\epsilon_*=\epsilon_a\), il en résulte
\(\beta_{\rm clk}=\log3/\epsilon_a\), exactement la relation
\eqref{x_reciprocidad:v:ant:eq:ii-kb-aut-par-termico}. Ce calcul identifie la
confluence locale tritique. Pour un graphe d'histoires avec des contraintes
supplémentaires, la confluence avec la même horloge exige de vérifier
\[
 \rho\!\left(L_{(\log3/\epsilon_a)\epsilon_*}\right)=1
\]
avec leurs incréments effectifs. La normalisation locale et la pression
du système complet conservent ainsi leurs domaines respectifs.

La relation obtenue réunit structure, valeur et signification. Le registre
détermine l'information récupérable ; les probabilités sur les
continuations déterminent sa lecture entropique ; l'action et la période
déterminent \(\epsilon_a\) ; le transducteur énergie--température
publie l'entropie dimensionnelle. L'aire et l'intrication sont
reliées à travers l'état réduit et ses opérateurs de bord,
en gardant visible la donnée qui produit chaque grandeur.

% END OWNER


% BEGIN OWNER /Users/ruben/Documents/New project/output/EXTREMA_MEDIA_RAZON_RECIPROCIDAD_ES_EN_20260918/ARTICULO/ES/source/antecedentes_vii/28_portador_tensorial.tex
% FORMALIZACION_REUNIDA: representación pentacomponente de II,
% con antecedente Paley y demostración directa de los momentos del marco.
% El retorno transporta la carta; no se identifica un ciclo de orden doce con A5.
\chapter{Incidence orientée et représentation tensorielle gravitationnelle}
\label{x_reciprocidad:vii:sec:pentacomponente}
La publication exceptionnelle de l'état conserve l'incidence qui accompagne
sa lecture arithmétique. La matrice ternaire \(A_W\) de
\eqref{x_reciprocidad:exc:aw} fournit ici un antécédent matériel précis. Son
relèvement équilibré \(C=\operatorname{bal}(A_W)\), avec les représentants
\(0,1,-1\), satisfait
\[
 C=C^{\mathsf T},\qquad C^2=5I_6,\qquad \operatorname{Tr}C=0.
\]
Les six positions ont pour ordre \((\infty,0,1,2,3,4)\). Cette incidence produit d'abord un espace tridimensionnel ; sa lecture quadratique produit ensuite l'espace à cinq composantes sur lequel agit le couplage gravitationnel.

\section{Revêtement signé et transport du calendrier}
On définit
\[
 E_C=\frac12(I_6+C/\sqrt5),\quad
 \mathcal H_C=\operatorname{im}E_C,\quad
 v_{u,\sigma}=\sigma\sqrt2E_Ce_u,\quad \sigma\in\{1,-1\}.
\]
\begin{proposition}[Douze positions orientées]
\label{x_reciprocidad:vii:prop:doce-orientadas}
Les vecteurs \(v_{u,\sigma}\) sont douze vecteurs unitaires distincts
dans un espace de dimension trois. Leur incidence est déterminée par
\(\langle v_{u,\sigma},v_{v,\tau}\rangle=1/\sqrt5\).
L'inversion de l'orientation agit par antipodie.
\end{proposition}
\begin{proof}
L'identité quadratique de \(C\) implique \(E_C^2=E_C=E_C^*\) ; sa trace est trois. Puisque la diagonale de \(C\) est nulle,
\[
 \langle v_{u,\sigma},v_{v,\tau}\rangle
 =\begin{cases}\sigma\tau,&u=v,\\
 \sigma\tau C_{uv}/\sqrt5,&u\ne v.
 \end{cases}
\]
Cela prouve la normalisation, la distinction des douze vecteurs et l'incidence.
Remplacer \(\sigma\) par \(-\sigma\) change le signe du vecteur.
\end{proof}

Une expression euclidienne explicite s'obtient en posant
\(\nu=\sqrt{1+\varphi^2}\) et
\[
 B=\frac1\nu
 \begin{pmatrix}
 0&-1&-\varphi&0&\varphi&1\\
 -1&-\varphi&0&1&0&-\varphi\\
 -\varphi&0&-1&-\varphi&-1&0
 \end{pmatrix}.
\]
La relation \(\varphi^2=\varphi+1\), pour la coordonnée déjà générée, donne \(B^*B=2E_C\). En conséquence, \(B/\sqrt2|_{\mathcal H_C}\) est une isométrie et envoie \(v_{u,\sigma}\) sur \(\sigma Be_u\). L'image est
\[
 \mathcal V=\frac1\nu\bigl(
 \{0\}\times\{\pm1\}\times\{\pm\varphi\}\ \cup\
 \{\pm1\}\times\{\pm\varphi\}\times\{0\}\ \cup\
 \{\pm\varphi\}\times\{0\}\times\{\pm1\}\bigr).
\]
La coordonnée interne, l'incidence et leurs
douze représentants géométriques sont ainsi reliés.

La coordonnée du calendrier s'écrit \((b,r)\in
\mathbb Z/12\mathbb Z\times\{0,\ldots,8\}\). Une route de longueur \(n\)
induit
\[
 k(r,n)=\left\lfloor\frac{r+n}{9}\right\rfloor,\qquad
 (b,r)\longmapsto
 \bigl(b+k(r,n)\bmod12,\ (r+n)\bmod9\bigr).
\]
L'alignement explicite de leurs positions \(p=b+1\) avec la carte signée est
\[
\begin{array}{c|rrrrrrrrrrrr}
p&1&2&3&4&5&6&7&8&9&10&11&12\\ \hline
u&\infty&0&0&\infty&1&3&3&1&2&4&4&2\\
\sigma&+&-&+&-&-&+&-&+&-&+&-&+
\end{array}
\]
Ce tableau fixe une trivialisation du porteur. Chaque changement de position
transporte conjointement les vecteurs et leur matrice d'incidence.

\begin{proposition}[Covariance du repère sous retour]
Si \(S e_p=e_{p+1\bmod12}\), la représentation matricielle d'un
opérateur \(P\) du porteur est transportée sous la forme \(P_k=S^kPS^{-k}\).
Ce transport conserve les produits scalaires, les incidences et la composition
des routes. Après douze retours, la carte revient, tandis que la mémoire
accumulée dans l'état source est conservée.
\end{proposition}
\begin{proof}
\(S\) est orthogonal et \(S^{12}=I\) ; les propriétés métriques et la périodicité de la carte en découlent. Pour la composition, avec \(r'=(r+n)\bmod9\), la division euclidienne donne \(k(r,n+m)=k(r,n)+k(r',m)\). Les exposants de \(S\) s'additionnent exactement avec la retenue. L'opération agit sur la coordonnée du calendrier ; les mises à jour de la mémoire de la route restent présentes dans son domaine enrichi.
\end{proof}

\section{Lecture quadratique et espace à cinq composantes}
Soient \(\mathbb R^{\mathcal V}\) l'espace des douze coordonnées,
\(\mathsf A e_v=e_{-v}\) son involution antipodale et
\(\mathsf J=\mathbf1\mathbf1^{\mathsf T}\). On définit
\[
 P_5=\frac12(I+\mathsf A)-\frac1{12}\mathsf J,\qquad
 Q_v=vv^{\mathsf T}-\frac13I_3,\qquad
 \Gamma_0x=\sum_{v\in\mathcal V}x_vQ_v.
\]
\begin{theorem}[Projecteur à cinq composantes et isométrie tensorielle]
\label{x_reciprocidad:vii:thm:portador-pentacomponente}
\(P_5\) est le projecteur orthogonal de rang cinq sur l'espace des
coordonnées paires sous l'antipodie et de somme nulle. De plus,
\[
 \Gamma_0^*\Gamma_0=\frac85P_5,\qquad
 \Gamma_0\Gamma_0^*=\frac85I_{\operatorname{Sym}^2_0(\mathbb R^3)}.
\]
Par conséquent
\[
 \Gamma_5=\sqrt{\frac58}\,\Gamma_0|_{\operatorname{im}P_5}
\]
est une isométrie sur l'espace des tenseurs symétriques sans trace.
\end{theorem}
\begin{proof}
L'espace pair possède une coordonnée pour chaque paire antipodale, donc
sa dimension est six. Le vecteur constant lui appartient ; soustraire son
projecteur de rang un donne \(P_5\), de rang cinq.

Comme
\(\langle Q_v,Q_w\rangle_F=(v^{\mathsf T}w)^2-1/3\), les entrées
de \(\Gamma_0^*\Gamma_0\) valent \(2/3\) pour \(w=\pm v\) et
\(-2/15\) aux dix autres emplacements de chaque ligne. Ce sont
exactement les entrées de \((8/5)P_5\).

On peut également vérifier l'identité dans l'espace tensoriel sans recourir à une classification des représentations. La liste des sommets donne
\[
 \sum_vv_i^4=\frac{12}{5},\qquad
 \sum_vv_i^2v_j^2=\frac45\quad(i\ne j),
\]
et les moments comportant une puissance impaire s'annulent par les changements de signe. En effet, \((1+\varphi^2)^2=5\varphi^2\) et \(1+\varphi^4=3\varphi^2\). Pour un tenseur symétrique \(Q\), ces sommes donnent
\[
 \sum_v(v^{\mathsf T}Qv)\,vv^{\mathsf T}
 =\frac85Q+\frac45(\operatorname{tr}Q)I.
\]
Avec \(\operatorname{tr}Q=0\), et en utilisant
\(\sum_vv^{\mathsf T}Qv=4\operatorname{tr}Q=0\), on obtient
\(\Gamma_0\Gamma_0^*Q=(8/5)Q\). Les deux identités prouvent
l'isométrie normalisée et sa surjectivité.
\end{proof}

La section gravitationnelle de \(\ref{x_reciprocidad:v:ant:sec:gravity}\) détermine
l’opérateur \(G_aP_5\). Son spectre a \(G_a\) de multiplicité
cinq et zéro de multiplicité sept. L’orientation est conservée dans
le porteur et dans son transport ; \(G_a\) fixe le couplage commun.

\section{Projection transversale et hélicités}
Pour \(n\in S^2\), posons \(\Pi=I-nn^{\mathsf T}\) et
\[
 \Lambda_nQ=\Pi Q\Pi-\frac12\operatorname{tr}(\Pi Q\Pi)\Pi.
\]
\begin{proposition}[Réduction transverse de rang deux]
\(\Lambda_n\) est le projecteur orthogonal sur \(\{Q\in\operatorname{Sym}^2_0(\mathbb R^3):Qn=0\}\), de dimension deux.
\end{proposition}
\begin{proof}
\(\Pi n=0\) donne la transversalité et \(\operatorname{tr}\Pi=2\) donne
une trace nulle. Un tenseur transverse sans trace reste fixe.
La formule est autoadjointe pour le produit de Frobenius.
Dans un repère \((e_1,e_2,n)\), l'image a pour base orthonormée
\[
 E_+=\frac{e_1e_1^{\mathsf T}-e_2e_2^{\mathsf T}}{\sqrt2},\qquad
 E_\times=\frac{e_1e_2^{\mathsf T}+e_2e_1^{\mathsf T}}{\sqrt2}.
\]
\end{proof}
Les trois autres vecteurs d'une base orthonormée de l'espace porteur sont
\[
 E_1=\frac{e_1n^{\mathsf T}+ne_1^{\mathsf T}}{\sqrt2},\quad
 E_2=\frac{e_2n^{\mathsf T}+ne_2^{\mathsf T}}{\sqrt2},\quad
 E_0=\frac{2nn^{\mathsf T}-e_1e_1^{\mathsf T}-e_2e_2^{\mathsf T}}{\sqrt6}.
\]
Lorsque le repère transversal tourne d'un angle \(\theta\), les paires \((E_+,E_\times)\) et \((E_1,E_2)\) tournent de \(2\theta\) et \(\theta\), respectivement, tandis que \(E_0\) reste fixe. La complexification possède ainsi les poids \(2,-2,1,-1,0\). \(\Lambda_n\) conserve les deux premiers et annule les autres. On établit ainsi la chaîne
\[
 \mathbb R^{12}\xrightarrow{P_5}\operatorname{im}P_5
 \xrightarrow{\Gamma_5}\operatorname{Sym}^2_0(\mathbb R^3)
 \xrightarrow{\Lambda_n}\mathcal H_{\rm TT}(n),
 \qquad 12\longrightarrow5\longrightarrow5\longrightarrow2.
\]
La représentation à cinq composantes et sa réduction sont des résultats
algébriques antérieurs au choix des équations de champ.
Dans la réalisation gravitationnelle sans masse, les contraintes de jauge
sélectionnent l'image transverse. Cette distinction conserve tant les
cinq coordonnées structurelles que les deux hélicités radiatives.


% END OWNER


% BEGIN OWNER /Users/ruben/Documents/New project/output/EXTREMA_MEDIA_RAZON_RECIPROCIDAD_ES_EN_20260918/ARTICULO/ES/source/antecedentes_vii/29_retorno_y_memoria_traslacional.tex
% Artículo VII, recuperación aditiva, revisión 04, 10-09-2026.
% RESULTADO_RECUPERADO: S15, 11c_gravedad_holonomia_rev6.tex:16--126.
% El cociclo canónico c27=(1,18) se conserva como antecedente del cómputo
% de rutas; este fragmento prueba su composición y representación.
% La memoria traslacional finita y la amplitud cosmológica s0 tienen
% tipos distintos. El cociclo no evalúa por sí solo una densidad local.

\chapter{Retour de phase et représentation de la mémoire accumulée}
\label{x_reciprocidad:vii:sec:retorno-memoria-gravedad}

La chronologie des routes enrichies distingue le retour de position,
le retour d'orientation et le retour de phase. La publication observable
identifie certains états d'arrivée, tandis que le registre cumulatif
conserve la séquence parcourue. La représentation géométrique doit
transporter les deux données conjointement.

\section{Cocycle d'inversion et changement de feuillet}
\label{x_reciprocidad:vii:subsec:cociclo-memoria}

Soit \(F_{\rm obs}\) la publication chronologique du transport TPK.
Ses blocs canoniques ont pour hiérarchie
\[
\begin{array}{c|l}
27&\text{retour positionnel avec inversion d'orientation},\\
54&\text{retour positionnel et orientationnel},\\
108&\text{retour complet de la phase observable}.
\end{array}
\]
Le registre compte les inversions d'orientation \(g(a)\) et les
changements effectifs de feuillet \(s(a)\). Pour une route,
\begin{equation}
 c(\gamma)=\sum_{a\subset\gamma}(g(a),s(a)),\qquad
 c(\gamma_2\gamma_1)=c(\gamma_2)+c(\gamma_1).
 \label{x_reciprocidad:vii:eq:cociclo-ruta}
\end{equation}
Le calcul canonique du bloc de longueur \(27\) publie
\(c_{27}=(1,18)\). La composition conserve ces incréments dans
les quatre blocs qui complètent la période observable.
Sur les routes orientées vers l'avant, le registre est cumulatif ;
son extension au groupoïde des routes utilise \(\mathbb Z^2\) et
\(c(\bar\gamma)=-c(\gamma)\).

\begin{proposition}[Mémoire du tour observable complet]
\label{x_reciprocidad:vii:prop:memoria-vuelta-completa}
En écrivant \(\mathcal K_{27}=F_{\rm obs}^{27}\), le relèvement
au registre, pour \(x\in\mathcal O_{108}\) et \(z\in\mathbb Z^2\), satisfait
\[
 \widetilde{\mathcal K}_{27}(x,z)
   =(\mathcal K_{27}x,z+(1,18)),\qquad
 \widetilde{\mathcal K}_{27}^{\,4}(x,z)
   =(x,z+(4,72)).
\]
Après \(N\) tours complets, l'incrément est \(N(4,72)\).
\end{proposition}

\begin{proof}
La projection observable de \(\mathcal K_{27}\) est d'ordre quatre. Le calcul canonique précédent fixe \(c_{27}=(1,18)\) de manière constante sur cette orbite observable. L'additivité de \eqref{x_reciprocidad:vii:eq:cociclo-ruta} additionne les quatre incréments et donne \(c_{108}=4(1,18)=(4,72)\). La répétition applique la même loi de concaténation et produit \(Nc_{108}\). Le retour de la première coordonnée conserve l'incrément de la seconde.
\end{proof}

\section{Représentation affine du cocycle}
\label{x_reciprocidad:vii:subsec:representacion-afin-cociclo}

La réalisation discrète utilise un élément \(R_4\) d'ordre quatre dans
\(\operatorname{Spin}^+(1,3)\), un vecteur temporel unitaire \(u\)
fixé par son action vectorielle et une unité de longueur \(\ell_0>0\).
L'échelle \(\ell_0\) conserve sa provenance dans la section dimensionnelle.
Nous définissons
\begin{equation}
 \rho_{\rm C}^{\rm disc}(\gamma)=
 \bigl(R_4^{c_g(\gamma)},\,\ell_0c_s(\gamma)u\bigr)
 \in\operatorname{Spin}^+(1,3)\ltimes\mathbb R^{1,3}.
 \label{x_reciprocidad:vii:eq:representacion-discreta-memoria}
\end{equation}
La puissance agit sur le vecteur au moyen de la représentation vectorielle
associée du groupe de spin.

\begin{theorem}[Composition et amplitude translationnelle du retour]
\label{x_reciprocidad:vii:thm:traslacion-memoria}
L'application \(\rho_{\rm C}^{\rm disc}\) conserve l'identité,
la concaténation et l'inversion. Dans les retours complets,
\begin{equation}
 \rho_{\rm C}^{\rm disc}(\gamma_{108})=(I,72\ell_0u),\qquad
 \rho_{\rm C}^{\rm disc}(\gamma_{108}^{\,N})=(I,72N\ell_0u).
 \label{x_reciprocidad:vii:eq:traslacion-retorno-108}
\end{equation}
L'inversion change le signe de la translation.
\end{theorem}

\begin{proof}
Le produit semi-direct est
\((R,a)(S,b)=(RS,a+Rb)\). Comme \(R_4u=u\),
\[
\begin{split}
 \rho_{\rm C}^{\rm disc}(\gamma_2)
 \rho_{\rm C}^{\rm disc}(\gamma_1)
 &=
 \bigl(R_4^{c_g(\gamma_2)+c_g(\gamma_1)},
 \ell_0(c_s(\gamma_2)+c_s(\gamma_1))u\bigr)\\
 &=\rho_{\rm C}^{\rm disc}(\gamma_2\gamma_1).
\end{split}
\]
Le cocycle de l'identité est nul et celui de l'inverse est son opposé.
Pour le tour complet, on substitue \(c_{108}=(4,72)\) et
\(R_4^4=I\). La concaténation de \(N\) tours maintient la partie
rotationnelle égale à l'identité et additionne leurs translations.
\end{proof}

L'amplitude \(72\ell_0\) est une longueur associée à une route finie. Une réalisation au moyen d'une cotétrade et d'une connexion la transporte vers son holonomie affine. La contribution locale aux équations gravitationnelles s'obtient ensuite au moyen de la courbure de cette connexion, de la variation de l'action matérielle et de la normalisation de sa densité. La représentation conserve ainsi le contenu cumulatif exact qui doit intervenir dans cette réalisation.

% END OWNER


% BEGIN OWNER /Users/ruben/Documents/New project/output/EXTREMA_MEDIA_RAZON_RECIPROCIDAD_ES_EN_20260918/ARTICULO/ES/source/antecedentes_vii/30_pantalla_y_respuesta.tex
% Artículo VII. Incorporación de resultados preexistentes, 10-09-2026.
% Propietarios íntegros preservados en fuentes_conservadas:
% 17_pantalla_cubica_soldadura.tex (P), 18_clausura_energetica.tex (E).
% P: incidencia, cociente, soldadura celular y refinamiento.
% E: respuesta, unicidad, extensión mínima y expectativa tracial.
% Dependencias anteriores de VII: núcleo APP--TRIT--TPK; realización
% cúbica de la ruta, cociclo de transporte y compresión q_vis=5/9.
% Este fragmento no constituye por sí solo el artículo autónomo.
% La identidad de pantalla se distingue de la contracción material de espín.

\chapter{Soudure de frontière et réponse énergétique}
\label{x_reciprocidad:vii:sec:pantalla-respuesta}

La réalisation cubique du transport conserve six incidences orientées.
Son quotient géométrique, la forme lorentzienne et la réponse énergétique
s'obtiennent à partir de la même matrice d'incidence. Le transport procède de
la route enrichie APP--TRIT--TPK : la face, le signe, la retenue et la
mémoire accompagnent chaque arête. Cette section développe
l'opération linéaire qui agit sur cette publication de frontière.

\section{Incidence cubique et quotient de rang quatre}
\label{x_reciprocidad:vii:sec:cociente-pantalla}

Soit \(H_F=\mathbb R^F\), avec
\(F=\{(i,\epsilon):1\leq i\leq3,\ \epsilon\in\{-1,1\}\}\),
son produit scalaire euclidien et sa base \(e_{i,\epsilon}\).
L'opposition des faces est l'involution orthogonale
\(Re_{i,\epsilon}=e_{i,-\epsilon}\). Nous définissons
\[
 u=\sum_{i,\epsilon}e_{i,\epsilon},\qquad
 P_u=\frac{uu^{\mathsf T}}6,\qquad P_-=\frac{I-R}{2},\qquad
 P_0=\frac{I+R}{2}-P_u,\qquad P_\square=P_u+P_-.
\]
Les sommets sont \(V=\{-1,1\}^3\). La matrice d'incidence \(M:\mathbb R^V\longrightarrow H_F\) est déterminée par
\[
 M_{(i,\epsilon),\nu}=\mathbf1_{\{\nu_i=\epsilon\}}.
\]

\begin{proposition}[Décomposition de l'incidence]
\label{x_reciprocidad:vii:prop:incidencia}
On a
\[
 MM^{\mathsf T}=12P_u+4P_-,
 \qquad
 \ker M^{\mathsf T}=P_0H_F.
\]
Par conséquent, le quotient
\(Q=H_F/P_0H_F\), identifié orthogonalement à
\(P_\square H_F\), est de dimension quatre.
\end{proposition}

\begin{proof}
Une face contient quatre sommets. Deux faces opposées n'ont pas de sommets
communs et deux faces d'axes différents ont deux sommets communs.
En écrivant \(\mathbf J=uu^{\mathsf T}\), on obtient
\[
 MM^{\mathsf T}=4I+2(\mathbf J-I-R)
 =2(I-R)+2\mathbf J=4P_-+12P_u.
\]
Les trois projecteurs \(P_u,P_-,P_0\) sont orthogonaux et leur somme vaut \(I\) ; leurs rangs sont respectivement un, trois et deux. L'identité \(\|M^{\mathsf T}x\|^2=\langle x,MM^{\mathsf T}x\rangle\) identifie le noyau annoncé. Le rang complémentaire est quatre.
\end{proof}

Sur \(Q\), la forme
\[
 B=P_--P_u
\]
a pour signature \((3,1)\). L'orientation temporelle est fixée par \(u\).
Avec
\[
 t=\frac{u}{\sqrt6},\qquad
 d_i=\frac{e_{i,+}-e_{i,-}}{\sqrt2},\qquad
 W=\sqrt6P_u+\sqrt2P_-,
\]
on a \(We_{i,\epsilon}=t+\epsilon d_i\). Ces six vecteurs
sont nuls pour \(B\), tandis que
\[
 3t+\nu_1d_1+\nu_2d_2+\nu_3d_3
\]
a pour norme quadratique \(-6\). La même incidence satisfait
\(MM^{\mathsf T}=2W^*W\) sur \(Q\). Les rotations propres du cube
commutent avec les projecteurs, préservent les deux formes et fixent \(t\).
Ces identités déterminent la métrique de la publication cubique.

\section{Soudure cellulaire et transport de mémoire}
\label{x_reciprocidad:vii:sec:soldadura-celular}

Pour chaque arête orientée \(a\), la publication de la route fournit
la face \(\lambda(\underline a)\), le signe
\(\sigma(a)\), les fibres d'origine et de destination et le transport
\(U(a):Q_{s(a)}\longrightarrow Q_{t(a)}\).
L'inversion vérifie \(U(\bar a)=U(a)^{-1}\).
Les demi-arêtes de frontière conservent leur inversion formelle et leur incidence
au bord. Avec \(\ell_m=\ell_0\,9^{-m}>0\), la soudure est
\begin{equation}
 \beta_m(a)=\ell_m\sigma_m(a)
              n_{\lambda(\underline a)}^{s(a)},\qquad
 n_{i,\epsilon}=t+\epsilon d_i.
 \label{x_reciprocidad:vii:eq:soldadura-arista}
\end{equation}
La face et le signe sont transportés avec la mémoire de la route. La convention d'inversion exige
\begin{equation}
 \beta_m(\bar a)=-U_m(a)\beta_m(a).
 \label{x_reciprocidad:vii:eq:inversion-soldadura}
\end{equation}
En appliquant deux fois cette relation, on retrouve \(\beta_m(a)\) ; le registre de retenue conserve séparément la composition ordonnée.

\begin{proposition}[Naturalité de la soudure par raffinement]
\label{x_reciprocidad:vii:prop:refinamiento-soldadura}
Soit \(a=a_1\cdots a_9\) un raffinement compatible d'une arête.
Supposons que les neuf incidences filles, transportées dans la fibre du
parent, conservent la face et le signe, et que
\(\ell_{m+1}=\ell_m/9\). Soit
\(\mathcal U_j:Q_{m,s(a)}\longrightarrow Q_{m+1,s(a_j)}\)
le transport cumulé qui inclut l'identification par raffinement
et le trajet jusqu'à l'origine de l'enfant \(j\). Alors
\[
 \beta_m(a)=
 \sum_{j=1}^{9}\mathcal U_j^{-1}\beta_{m+1}(a_j).
\]
L'égalité est compatible avec l'inversion et avec la concaténation du registre de mémoire.
\end{proposition}

\begin{proof}
Chaque terme, exprimé dans la fibre de l'origine du parent, est
\((\ell_m/9)\sigma_m(a)n_{\lambda(\underline a)}\).
La somme de neuf termes reproduit la soudure du parent.
L'inversion inverse l'ordre des transports et remplace chacun
par son inverse ; la relation \eqref{x_reciprocidad:vii:eq:inversion-soldadura} fournit
le signe global. La mémoire est concaténée dans ce même ordre, de sorte
que le raffinement conserve tant la publication vectorielle que son
registre enrichi.
\end{proof}

\section{Réponse positive, unicité et extension minimale}
\label{x_reciprocidad:vii:sec:respuesta-positiva}

La soudure linéaire d'écran est \(\bar\beta_m=\ell_mW|_Q\). Elle est inversible, avec
\[
 \bar\beta_m^{-1}
 =\ell_m^{-1}\left(\frac1{\sqrt6}P_u+\frac1{\sqrt2}P_-\right).
\]
La valeur propre \(5/9\) du secteur antisymétrique du bloc normalisé \(B_c/9\), construit dans la section~\ref{x_reciprocidad:vii:sec:bloque-compresion}, détermine le coefficient de compression \(q_{\mathrm{vis}}=5/9\). Nous réservons cette notation pour le distinguer du quotient euclidien \(q_9(n)\) du noyau arithmétique. Sa réalisation sur l'écran utilise une isométrie \(V_9:Q\longrightarrow\mathcal K\). La contribution visible \(C_9=q_{\mathrm{vis}}V_9\) et son complément \(D_9=\sqrt{1-q_{\mathrm{vis}}^2}\,V_9\) satisfont
\[
 C_9^*C_9+D_9^*D_9=I_Q,\qquad D_9^*D_9=\frac{56}{81}I_Q.
\]
Le coefficient \(56\) est \(9^2-5^2\) ; il appartient au défaut quadratique
de ce lecteur de compression.

\begin{theorem}[Réponse énergétique déterminée par la soudure]
\label{x_reciprocidad:vii:thm:respuesta-unica}
Il existe un unique opérateur \(\Lambda_m:Q\longrightarrow Q\) tel que
\(\bar\beta_m^*\Lambda_m\bar\beta_m=D_9^*D_9\). Il est défini positif et
\begin{equation}
 \Lambda_m=\frac{28}{243\ell_m^2}(P_u+3P_-)
 =\frac{112}{81\ell_m^2}
       (MM^{\mathsf T}|_Q)^{-1}.
 \label{x_reciprocidad:vii:eq:lambda-pantalla}
\end{equation}
Sur \(H_F\), l'identité correspondante a pour second membre
\((56/81)P_\square\).
\end{theorem}

\begin{proof}
La congruence par la soudure inversible détermine nécessairement
\[
 \Lambda_m=\bar\beta_m^{-*}\frac{56}{81}I_Q\bar\beta_m^{-1}.
\]
En substituant l'inverse, le coefficient de \(P_u\) est \(56/(486\ell_m^2)=28/(243\ell_m^2)\) ; celui de \(P_-\) est \(28/(81\ell_m^2)\). Tous deux sont positifs. Par la proposition \ref{x_reciprocidad:vii:prop:incidencia}, \((MM^{\mathsf T}|_Q)^{-1}=P_u/12+P_-/4\), ce qui donne la seconde expression. L'extension à six faces annule \(P_0H_F\) ; l'identité s'étend donc au moyen de \(P_\square\).
\end{proof}

\begin{proposition}[Caractérisation par extension minimale]
\label{x_reciprocidad:vii:prop:minima-extension}
Pour tout \(b\in Q\),
\[
 \langle b,\Lambda_m b\rangle
 =\frac{112}{81\ell_m^2}
   \min_{Mx=b}\|x\|^2.
\]
L'extension minimisante est unique :
\(x_{\min}=M^{\mathsf T}(MM^{\mathsf T}|_Q)^{-1}b\).
\end{proposition}

\begin{proof}
Le vecteur annoncé satisfait \(Mx_{\min}=b\) et appartient à
\(\operatorname{im}M^{\mathsf T}=(\ker M)^\perp\).
Toute autre solution est \(x_{\min}+z\), avec \(z\in\ker M\) ; par
orthogonalité, sa norme au carré est \(\|x_{\min}\|^2+\|z\|^2\).
De plus,
\[
 \|x_{\min}\|^2
 =\langle b,(MM^{\mathsf T}|_Q)^{-1}b\rangle.
\]
La formule de \(\Lambda_m\) conclut la preuve.
\end{proof}

L'opérateur \(MM^{\mathsf T}|_Q\) réalise la réponse
de Dirichlet--à--Neumann de cette formulation d'incidence ; son inverse
est la réponse de Neumann--à--Dirichlet. L'énergie précédente correspond
à cette dernière, avec la normalisation de la soudure et du défaut de
compression. Le complément de Schur de
\(\left(\begin{smallmatrix}I&M^{\mathsf T}\\M&0\end{smallmatrix}\right)\)
est \(-MM^{\mathsf T}\), ce qui offre la même élimination variationnelle.

\section{Conservation de la densité par raffinement}
\label{x_reciprocidad:vii:sec:densidad-refinamiento}

Considérons le secteur des transports cubiques, orthogonaux pour le produit scalaire énergétique et compatibles avec la forme \(B\). Nous fixons un espace commun d'étiquettes \(Q^{\mathrm{lab}}\), identifié isométriquement à l'écran du parent. Les soudures ont les types \(\bar\beta_m:Q^{\mathrm{lab}}\longrightarrow Q_m\) et \(\bar\beta_{m+1,j}:Q^{\mathrm{lab}}\longrightarrow Q_{m+1,j}\) ; les opérateurs \(\Lambda\) agissent dans leurs fibres d'arrivée respectives. Si \(U_j:Q_j\longrightarrow Q\) transporte la fibre fille vers celle du parent, alors
\[
 \bar\beta_{m+1,j}=\frac19U_j^{-1}\bar\beta_m,\qquad
 \Lambda_{m+1,j}=81U_j^{-1}\Lambda_mU_j.
\]
Le facteur \(81\) procède de \(\ell_{m+1}^{-2}=81\ell_m^{-2}\).

\begin{theorem}[Invariance de la densité normalisée]
\label{x_reciprocidad:vii:thm:densidad}
Chaque enfant conserve la réponse quadratique du parent :
\[
 \bar\beta_{m+1,j}^*\Lambda_{m+1,j}\bar\beta_{m+1,j}
 =\frac{56}{81}I_{Q^{\mathrm{lab}}}.
\]
Pour \(r\) raffinements, la réponse conjointe, exprimée dans les fibres
du parent, est
\[
 E_{m+r}=\frac{56}{81}I_{Q^{\mathrm{lab}}}\otimes I_{9^r}.
\]
L'espérance de trace normalisée
\(\tau_{9^r}(A)=9^{-r}\operatorname{Tr}(A)\) satisfait
\[
 (\operatorname{id}\otimes\tau_{9^r})(E_{m+r})
 =\frac{56}{81}I_{Q^{\mathrm{lab}}}
\]
à toute profondeur finie.
\end{theorem}

\begin{proof}
La congruence d'un enfant contient les facteurs \(1/9,81,1/9\) ;
leur produit vaut un. L'orthogonalité élimine les transports
intermédiaires. La somme directe de neuf réponses égales identifie
le premier raffinement à \((56/81)I_Q\otimes I_9\).
L'itération produit le tenseur avec \(I_{9^r}\).
Enfin, \(\tau_{9^r}(I_{9^r})=1\). Les espérances se composent
par l'identité
\(\tau_{9^{r+s}}=\tau_{9^r}\otimes\tau_{9^s}\).
\end{proof}

Les inclusions unitaires \(A\mapsto A\otimes I_9\) sont compatibles avec
cette famille de traces. Dans la limite inductive des algèbres matricielles,
la réponse est l'élément central
\((56/81)I_Q\otimes\mathbf1\). La somme extensive non normalisée produit,
en revanche, \(9^r(56/81)I_Q\). Cette distinction spécifie quelle grandeur
est conservée : la densité quadratique d'écran.

\section{Réalisation variationnelle de la réponse}
\label{x_reciprocidad:vii:sec:realizacion-respuesta}

La réalisation matérielle utilise une isométrie \(J_{\mathrm{scr}}:Q\longrightarrow\mathcal H_\partial\) compatible avec l'action de spin et avec le raffinement. La cotétrade de frontière est \(e=J_{\mathrm{scr}}\bar\beta_m\). Soit \(\Lambda_{\mathrm{CE}}\) la réponse de frontière obtenue par seconde variation de l'action matérielle et géométrique, avec ses conditions aux limites et l'élimination des variables intérieures. Le défaut de compatibilité est
\begin{equation}
 \mathcal R_{\mathrm{grav}}
 =\bar\beta_m^*
   (J_{\mathrm{scr}}^*\Lambda_{\mathrm{CE}}J_{\mathrm{scr}}-\Lambda_m)
   \bar\beta_m
 =e^*\Lambda_{\mathrm{CE}}e-D_9^*D_9.
 \label{x_reciprocidad:vii:eq:defecto-respuesta}
\end{equation}
Par l'inversibilité de \(\bar\beta_m\),
\[
 \mathcal R_{\mathrm{grav}}=0
 \quad\Longleftrightarrow\quad
 J_{\mathrm{scr}}^*\Lambda_{\mathrm{CE}}J_{\mathrm{scr}}=\Lambda_m.
\]
L'isométrie \(J_{\mathrm{scr}}\) transporte les états de frontière ;
le courant variationnel de spin est une forme de degré trois à valeurs
dans les bivecteurs. Chacun intervient avec son domaine propre dans la
réalisation de la réponse.

La construction détermine une forme énergétique positive, son extension
intérieure minimisante et la densité conservée par raffinement.
L'amplitude d'une contribution matérielle précise s'obtient en évaluant
le courant de cette réalisation dans sa forme quadratique variationnelle.
Sont ainsi distingués l'opérateur qui fixe la réponse, le courant
sur lequel il agit et l'observable gravitationnel qui publie l'évaluation.

% END OWNER


% BEGIN OWNER /Users/ruben/Documents/New project/output/EXTREMA_MEDIA_RAZON_RECIPROCIDAD_ES_EN_20260918/ARTICULO/ES/source/antecedentes_vii/30d_realizacion_geometrica.tex
% RESULTADO_RECUPERADO: S15, 11c_gravedad_holonomia_rev6.tex:128--212;
% S01, 02_dinamica_tres_hojas.tex:155--219; pantalla y respuesta de30.
% Pruebas de curvatura, covariancia y Bianchi reunidas para autonomía.
% La realización suave conserva sus condiciones de holonomía explícitas.

\chapter{Réalisation géométrique du transport et de la soudure}
\label{x_reciprocidad:vii:sec:realizacion-cartan}

La représentation discrète du retour et la soudure d'écran déterminent deux aspects du transport : l'accumulation de mémoire dans les routes et la comparaison de leurs fibres. Leur réalisation géométrique conserve les deux opérations. Dans cette section sont fixés l'espace d'arrivée, la normalisation dimensionnelle de la connexion et les identités qu'utilisera sa variation.

\section{Écran, cotétrade et connexion}

Le quotient d'écran \(Q_\square\), sa forme lorentzienne et la
soudure \(\beta_m\) procèdent de
\ref{x_reciprocidad:vii:sec:pantalla-respuesta}. Dans chaque cellule, une réalisation d'
écran \(J_{{\rm scr},m}\) transporte la forme et la soudure
vers la fibre géométrique correspondante. Son image cellulaire est
\[
 e_m=J_{{\rm scr},m}\beta_m.
\]
L'échelle \(\ell_m\) demeure à l'intérieur de la soudure. Les
changements de repère agissent sur \(J_{{\rm scr},m}\), la connexion
et le transport des cellules conjointement.

Une réalisation lisse de ces données sur une variété orientée \(\mathcal U\) de dimension quatre consiste en une cotétrade non dégénérée \(e=(e^I)\), une connexion de Lorentz \(\omega=(\omega^I{}_J)\) et une réalisation de routes \(\mathfrak R_{\rm C}\) qui satisfont
\begin{equation}
 \operatorname{Hol}_{\mathcal A_{\ell_0}}
       (\mathfrak R_{\rm C}(\gamma))
 =\rho_{\rm C}(\operatorname{Hol}_{\rm TPK}(\gamma)).
 \label{x_reciprocidad:vii:eq:compatibilidad-holonomia-geometrica}
\end{equation}
Dans cette égalité matricielle, la représentation est normalisée par
\[
 N_{\ell_0}(R,a)=(\Lambda(R),a/\ell_0),\qquad
 \rho_{\rm C}=N_{\ell_0}\circ\rho_{\rm C}^{\rm disc},
\]
où \(\Lambda:\operatorname{Spin}^+(3,1)\to SO^+(3,1)\)
est la représentation vectorielle. L'application \(N_{\ell_0}\)
préserve le produit semi-direct : diviser la translation
\(a+\Lambda(R)b\) par \(\ell_0\) équivaut à composer les
deux translations normalisées. Le tour complet publie donc
\(72u\) dans la connexion normalisée et \(72\ell_0u\) dans la
carte dimensionnelle. Le relèvement de spin conserve, séparément,
\(R_4^{c_g(\gamma)}\) ; la représentation vectorielle a pour
noyau central \(\{1,-1\}\), de sorte que ce
relèvement est retenu lors de l'utilisation de champs spinoriels.
L'identité, la concaténation, l'inversion et la subdivision sont maintenues dans la
réalisation correspondante. L'appariement d'
écran et la condition d'holonomie spécifient les données de la
réalisation utilisée dans les équations de champ.

Avec \(\eta=\operatorname{diag}(-1,1,1,1)\), on définit
\(g=\eta_{IJ}e^I\otimes e^J\). Pour une échelle élémentaire
positive \(\ell_0\), constante dans la carte, la connexion affine est
\begin{equation}
 \mathcal A_{\ell_0}
 =\begin{pmatrix}\omega&\ell_0^{-1}e\\0&0\end{pmatrix}.
 \label{x_reciprocidad:vii:eq:conexion-afin-realizada}
\end{equation}
La cotétrade a la dimension d'une longueur ; le facteur \(\ell_0^{-1}\) homogénéise sa composante translationnelle.

\begin{proposition}[Courbure et torsion de la connexion réalisée]
\label{x_reciprocidad:vii:prop:curvatura-afin}
La courbure de \eqref{x_reciprocidad:vii:eq:conexion-afin-realizada} est
\begin{equation}
 \mathcal F_{\ell_0}
 =\mathrm d\mathcal A_{\ell_0}
   +\mathcal A_{\ell_0}\wedge\mathcal A_{\ell_0}
 =\begin{pmatrix}F_\omega&\ell_0^{-1}T\\0&0\end{pmatrix},
 \quad F_\omega=\mathrm d\omega+\omega\wedge\omega,
 \quad T=\mathrm de+\omega\wedge e.
 \label{x_reciprocidad:vii:eq:curvatura-afin-realizada}
\end{equation}
\end{proposition}
\begin{proof}
Le produit de matrices de formes produit \(\omega\wedge\omega\) dans le bloc diagonal et \(\ell_0^{-1}\omega\wedge e\) dans le bloc translationnel. La dérivée extérieure apporte \(\mathrm d\omega\) et \(\ell_0^{-1}\mathrm de\), puisque \(\ell_0\) est constant. La somme démontre l'égalité.
\end{proof}

La courbure décrit le défaut de transport des repères et la torsion
celui de la cotétrade. Toutes deux appartiennent à la même connexion réalisée ;
leurs composantes conservent des domaines tensoriels distincts.

\begin{proposition}[Covariance et identités de Bianchi]
\label{x_reciprocidad:vii:prop:bianchi-cartan}
Sous un changement de repère \(g_L:\mathcal U\to SO^+(3,1)\),
\[
 e'=g_Le,\qquad
 \omega'=g_L\omega g_L^{-1}-\mathrm dg_L\,g_L^{-1},
\]
on a \(T'=g_LT\) et \(F_{\omega'}=g_LF_\omega g_L^{-1}\). En outre,
\begin{equation}
 D_\omega T=F_\omega\wedge e,\qquad D_\omega F_\omega=0.
 \label{x_reciprocidad:vii:eq:bianchi-realizacion}
\end{equation}
\end{proposition}
\begin{proof}
Dans \(\mathrm d(g_Le)+\omega'\wedge g_Le\), les termes
contenant \(\mathrm dg_L\wedge e\) se compensent.
La même substitution dans la courbure donne sa transformation par
conjugaison. Pour la première identité de Bianchi, on développe
\[
 D_\omega(\mathrm de+\omega\wedge e)
 =(\mathrm d\omega+\omega\wedge\omega)\wedge e.
\]
Dans \(\mathrm dF_\omega+\omega\wedge F_\omega-F_\omega\wedge\omega\), les termes contenant \(\mathrm d\omega\) et les produits cubiques s'annulent par paires. Cela prouve la seconde identité.
\end{proof}

\section{Les trois réalisations quadratiques de courbure}

Le feuillet \(\tau\in\{+1,0,-1\}\) du TRIT est représenté dans
l'algèbre de Cartan au moyen des générateurs \(J_{ab}=-J_{ba}\)
et \(P_a^{(\tau)}\), avec les relations
\begin{align*}
 [J_{ab},J_{cd}]&=
 \eta_{bc}J_{ad}-\eta_{ac}J_{bd}
 -\eta_{bd}J_{ac}+\eta_{ad}J_{bc},\\
 [J_{ab},P_c^{(\tau)}]&=
 \eta_{bc}P_a^{(\tau)}-\eta_{ac}P_b^{(\tau)},\\
 [P_a^{(\tau)},P_b^{(\tau)}]&=-\tau J_{ab}.
\end{align*}
Pour \(\ell>0\) constant, on écrit
\[
 \mathcal A_\tau^{\rm Cart}
 =\tfrac12\omega^{ab}J_{ab}+\ell^{-1}e^aP_a^{(\tau)}.
\]

\begin{proposition}[Courbure de la réalisation tritifiée]
\label{x_reciprocidad:vii:prop:curvatura-tritificada}
Avec les relations précédentes,
\begin{equation}
 \mathcal F_\tau^{\rm Cart}
 =\tfrac12(F_\omega^{ab}-\tau\ell^{-2}e^a\wedge e^b)J_{ab}
     +\ell^{-1}T^aP_a^{(\tau)}.
 \label{x_reciprocidad:vii:eq:curvatura-tritificada}
\end{equation}
\end{proposition}
\begin{proof}
On développe \(\mathrm d\mathcal A+\tfrac12[\mathcal A,\mathcal A]\). Le crochet du secteur de Lorentz produit \(F_\omega\), le crochet mixte produit \(D_\omega e\), et le crochet translationnel apporte \(-\tfrac12\tau\ell^{-2}e^a\wedge e^bJ_{ab}\).
\end{proof}

Dans le feuillet central, on retrouve la connexion affine ; les feuillets extrêmes
conservent les deux signes du terme de courbure constante. La
représentation maintient ainsi le régime tritique de la route avant
d'appliquer les équations variationnelles. L'égalité d'holonomie
\eqref{x_reciprocidad:vii:eq:compatibilidad-holonomia-geometrica} correspond au
feuillet central. Dans chaque feuillet extrême, on utilise la représentation
\(\rho_{{\rm C},\tau}\) dans son groupe de Cartan et sa condition
d'holonomie respective. Le courant et le tenseur métrique utilisent
les champs et la représentation de la même réalisation.

% END OWNER


% BEGIN OWNER /Users/ruben/Documents/New project/output/EXTREMA_MEDIA_RAZON_RECIPROCIDAD_ES_EN_20260918/ARTICULO/ES/source/antecedentes_vii/30b_variacion_energia_memoria.tex
% Artículo VII. Desarrollo aditivo, revisión 04, 10-09-2026.
% ARQUITECTURA_AUTORAL_PREEXISTENTE: transporte enriquecido, energía
% D_9^* W D_9 y compatibilidad por refinamiento, propietario S01,
% 02_dinamica_tres_hojas.tex, líneas 221--250.
% FORMALIZACION_NUEVA / CERTIFICADO_NUEVO: diferencial explícito de ese
% funcional, representante antihermítico, composición e inversión ponderada.
% Las variaciones pertenecen al espacio de realizaciones unitarias del
% transporte. No se presume que toda variación continua sea una ruta TPK.
% Su realización como corriente material de espín se trata por separado.
% \mathsf W_m es el peso energético; W de la sección anterior es soldadura.

\chapter{Réponse variationnelle de l'énergie de mémoire}
\label{x_reciprocidad:vii:sec:variacion-memoria}

Le transport d'une route identifie les fibres de ses extrémités et conserve l'orientation et la mémoire qui interviennent dans cette identification. L'énergie associée compare le champ final au champ initial transporté. Sa différentielle quantifie la réponse de cette comparaison à une variation du transport. Cette section obtient la réponse directement à partir de la fonctionnelle discrète et développe sa composition, sa covariance et son inversion.

\section{Fonctionnelle discrète et espace des variations}
\label{x_reciprocidad:vii:subsec:funcional-memoria}

À un niveau fini \(m\), soient \(\mathcal V_m\) les sommets de la
publication des routes et \(\mathcal E_m^+\) une orientation de chaque arête.
Chaque sommet porte une fibre hermitienne de dimension finie \(H_v\).
Nous définissons
\[
 \mathcal H_m=\bigoplus_{v\in\mathcal V_m}H_v,\qquad
 \mathcal K_m=\bigoplus_{a\in\mathcal E_m^+}H_{t(a)}.
\]
Le produit hermitien est antilinéaire dans la première variable.
Pour \(U_a:H_{s(a)}\to H_{t(a)}\) unitaire, l'opérateur de différences est
\begin{equation}
 (D_mf)_a=f_{t(a)}-U_af_{s(a)},\qquad
 E_m(f,U)=\langle D_mf,\mathsf W_mD_mf\rangle,\qquad
 \mathsf W_m=\mathsf W_m^*>0.
 \label{x_reciprocidad:vii:eq:energia-memoria-discreta}
\end{equation}
C'est l'énergie du terme \(D_m^*\mathsf W_mD_m\) du générateur discret.
Le poids \(\mathsf W_m\) peut coupler des arêtes distinctes. L'opérateur \(W\)
de la soudure de la section~\ref{x_reciprocidad:vii:sec:pantalla-respuesta} et ce poids
énergétique ont des domaines et des fonctions différents.

Le transport \(U\) est une publication de l'état enrichi. Le champ \(f\) et le poids \(\mathsf W_m\) restent explicites comme arguments de la fonctionnelle ; chaque application matérielle conserve les règles qui les déterminent. On considère des familles différentiables des réalisations unitaires du transport,
\[
 \dot U_a=A_aU_a,\qquad A_a^*=-A_a.
\]
Les variations physiquement utilisées forment le sous-espace qui conserve
les conditions de la réalisation correspondante. La différentielle est
calculée d'abord dans l'espace tangent unitaire, puis restreinte
à ce sous-espace.

\begin{proposition}[Covecteur de réponse au transport]
\label{x_reciprocidad:vii:prop:covector-memoria}
Écrivons
\[
 v_a=U_af_{s(a)},\qquad r=D_mf,\qquad q=\mathsf W_mr.
\]
À champ et poids fixés, la différentielle de l'énergie est
\begin{equation}
 \mathfrak j_m(A)=-2\operatorname{Re}
       \sum_{a\in\mathcal E_m^+}\langle q_a,A_av_a\rangle .
 \label{x_reciprocidad:vii:eq:covector-memoria}
\end{equation}
Son représentant pour le produit réel de Hilbert--Schmidt est
\begin{equation}
 \mathcal J_a=v_aq_a^*-q_av_a^*,\qquad
 \mathcal J_a^*=-\mathcal J_a,\qquad
 \mathfrak j_m(A)=
 \sum_a\operatorname{Re}\operatorname{Tr}(\mathcal J_a^*A_a).
 \label{x_reciprocidad:vii:eq:representante-corriente-discreta}
\end{equation}
Pour des variations simultanées du champ, du transport et du poids, on a
\begin{equation}
 \dot E_m=
 2\operatorname{Re}\sum_a
 \left\langle q_a,\dot f_{t(a)}-U_a\dot f_{s(a)}-A_av_a\right\rangle
 +\langle r,\dot{\mathsf W}_mr\rangle .
 \label{x_reciprocidad:vii:eq:variacion-completa-memoria}
\end{equation}
\end{proposition}

\begin{proof}
La mise à jour du premier ordre, écrite dans l'algèbre des nombres duaux avec \(\epsilon^2=0\), donne
\[
 r+\epsilon\dot r,\qquad
 \dot r_a=\dot f_{t(a)}-U_a\dot f_{s(a)}-A_av_a.
\]
Le coefficient de \(\epsilon\) dans l'énergie est
\[
 \langle\dot r,\mathsf W_mr\rangle+
 \langle r,\mathsf W_m\dot r\rangle+
 \langle r,\dot{\mathsf W}_mr\rangle.
\]
L'autoadjonction du poids produit
\eqref{x_reciprocidad:vii:eq:variacion-completa-memoria}. Lorsque le champ et le poids sont fixés, il reste
\eqref{x_reciprocidad:vii:eq:covector-memoria}. De plus,
\[
 \operatorname{Tr}(\mathcal J_a^*A_a)
 =v_a^*A_aq_a-q_a^*A_av_a.
\]
L'antihermiticité de \(A_a\) implique
\(v_a^*A_aq_a=-\overline{q_a^*A_av_a}\), d'où l'on obtient
le représentant annoncé.
\end{proof}

La réponse est ainsi évaluée : sur chaque arête, on transporte le champ, on calcule sa différence avec le champ final et on applique le poids énergétique. Le produit antisymétrisé de ces deux vecteurs détermine \(\mathcal J_a\in\mathfrak u(H_{t(a)})\). Cet opérateur représente, au moyen du produit réel de Hilbert--Schmidt, le covecteur de réponse sur les variations unitaires de la fibre de destination.

\section{Composition et transport de la réponse}
\label{x_reciprocidad:vii:subsec:composicion-respuesta}

\begin{proposition}[Règle de composition sur une route]
\label{x_reciprocidad:vii:prop:composicion-respuesta}
Pour une route formée de \(N\) transports, soient
\(U_\gamma=U_N\cdots U_1\) et
\(B_k=U_N\cdots U_{k+1}\), avec \(B_N=I\).
Si \(\dot U_k=A_kU_k\), alors
\begin{equation}
 \dot U_\gamma U_\gamma^{-1}
 =\sum_{k=1}^N B_kA_kB_k^{-1}.
 \label{x_reciprocidad:vii:eq:variacion-transporte-compuesto}
\end{equation}
Ainsi, un représentant \(\mathcal J_\gamma\) dans la fibre finale
induit dans le facteur \(k\) la contribution
\(B_k^*\mathcal J_\gamma B_k\). Lorsqu'un facteur appartient à plusieurs
routes de la fonctionnelle, ses contributions s'additionnent avec les incidences
qu'elle utilise.
\end{proposition}

\begin{proof}
La règle du produit donne
\[
 \dot U_\gamma=\sum_k
 U_N\cdots U_{k+1}A_kU_k\cdots U_1.
\]
Multiplier par \(U_\gamma^{-1}\) élimine les facteurs à droite
de \(A_k\) et produit \eqref{x_reciprocidad:vii:eq:variacion-transporte-compuesto}.
Par unitarité et cyclicité de la trace,
\[
 \operatorname{Re}\operatorname{Tr}
   (\mathcal J_\gamma^*B_kA_kB_k^{-1})
 =\operatorname{Re}\operatorname{Tr}
   ((B_k^*\mathcal J_\gamma B_k)^*A_k).
\]
La linéarité de la différentielle démontre la somme des contributions.
\end{proof}

Cette formule conserve la position de chaque opération dans la route.
Le transport postérieur à une incidence détermine comment sa contribution
est représentée dans la fibre finale ; la réponse locale est restituée
en la transportant vers la fibre de cette incidence.

\begin{proposition}[Covariance sous les changements unitaires de repère]
\label{x_reciprocidad:vii:prop:covariancia-corriente-memoria}
Soient \(g_v\) des changements unitaires de repère, indépendants du paramètre
de variation, et \(G_{\mathcal E}=\bigoplus_a g_{t(a)}\).
Sous les transformations
\[
 f'_v=g_vf_v,\quad U'_a=g_{t(a)}U_ag_{s(a)}^{-1},\quad
 \mathsf W'_m=G_{\mathcal E}\mathsf W_mG_{\mathcal E}^{-1},\quad
 A'_a=g_{t(a)}A_ag_{t(a)}^{-1},
\]
on a
\begin{equation}
 \mathcal J'_a=g_{t(a)}\mathcal J_ag_{t(a)}^{-1},
 \qquad \mathfrak j'_m(A')=\mathfrak j_m(A).
 \label{x_reciprocidad:vii:eq:covariancia-respuesta}
\end{equation}
\end{proposition}

\begin{proof}
On obtient \(r'=G_{\mathcal E}r\), \(q'=G_{\mathcal E}q\) et \(v'_a=g_{t(a)}v_a\). La substitution dans \eqref{x_reciprocidad:vii:eq:representante-corriente-discreta} démontre la première identité. L'invariance du produit hermitien, ou la cyclicité de la trace, démontre la seconde.
\end{proof}

Lorsque les changements de repère dépendent du paramètre, leurs dérivées
agissent également sur le champ et le poids. L'expression covariante
correspondante est alors la différentielle complète
\eqref{x_reciprocidad:vii:eq:variacion-completa-memoria}.

\section{Inversion orientée et transport du poids}
\label{x_reciprocidad:vii:subsec:inversion-respuesta}

L'inversion peut être vue explicitement dans le cas d'une énergie locale par arête, \(\mathsf W_m=\bigoplus_a W_a\). Pour \(a:s\to t\), son inverse transporte
\[
 U_{\bar a}=U_a^{-1},\qquad
 r_{\bar a}=-U_a^{-1}r_a,\qquad
 W_{\bar a}=U_a^{-1}W_aU_a.
\]
Ces règles conservent exactement l'énergie de l'arête.

\begin{proposition}[Compatibilité différentielle avec l'inversion]
\label{x_reciprocidad:vii:prop:inversion-respuesta}
Si \(W_a\) et les champs des deux extrémités restent fixes, alors
\[
 A_{\bar a}=-U_a^{-1}A_aU_a,\qquad
 \dot W_{\bar a}=U_a^{-1}[W_a,A_a]U_a.
\]
La variation complète de l'énergie de l'arête inversée coïncide avec
la variation originale. En particulier, si \([W_a,A_a]=0\), le covecteur partiel
de transport satisfait
\begin{equation}
 \mathfrak j_{\bar a}(U_a^{-1}A_aU_a)
 =-\mathfrak j_a(A_a).
 \label{x_reciprocidad:vii:eq:imparidad-transporte}
\end{equation}
\end{proposition}

\begin{proof}
La dérivée de \(U_a^{-1}\) est \(-U_a^{-1}A_a\), d'où les deux formules. En écrivant \(v_a=U_af_s\) et \(f_t=r_a+v_a\), la partie due exclusivement au transport inversé vaut
\[
 \mathfrak j_{\bar a}(A_{\bar a})
 =-2\operatorname{Re}\langle r_a,W_aA_af_t\rangle
 =\mathfrak j_a(A_a)
   -2\operatorname{Re}\langle r_a,W_aA_ar_a\rangle.
\]
La variation du poids apporte
\[
 \langle r_{\bar a},\dot W_{\bar a}r_{\bar a}\rangle
 =\langle r_a,[W_a,A_a]r_a\rangle
 =2\operatorname{Re}\langle r_a,W_aA_ar_a\rangle.
\]
Les deux termes se compensent. Si le commutateur est nul, cette contribution
disparaît ; la linéarité en \(A_a\) donne
\eqref{x_reciprocidad:vii:eq:imparidad-transporte}.
\end{proof}

L'énergie est sommée sur une orientation par arête. Une somme sur les deux
orientations porte le facteur \(1/2\). Cette convention conserve la même
fonctionnelle et donc la même réponse.

\section{Compatibilité avec le raffinement}
\label{x_reciprocidad:vii:subsec:refinamiento-respuesta}

\begin{theorem}[Transport de la différentielle par raffinement]
\label{x_reciprocidad:vii:thm:refinamiento-respuesta}
Soient \(J_m:\mathcal H_m\to\mathcal H_{m+1}\) et
\(K_m:\mathcal K_m\to\mathcal K_{m+1}\) des applications fixes.
Pour une famille \(C^1\), supposons les identités de raffinement
\begin{equation}
 D_{m+1}(t)J_m=K_mD_m(t),\qquad
 K_m^*\mathsf W_{m+1}(t)K_m=\mathsf W_m(t).
 \label{x_reciprocidad:vii:eq:familia-refinamiento-memoria}
\end{equation}
Alors \(E_{m+1}(J_mf)=E_m(f)\). Les covecteurs de transport coïncident sur les variations reliées par la première identité. Lorsque les poids varient, leurs contributions à la différentielle coïncident également.
\end{theorem}

\begin{proof}
Par substitution,
\[
 \langle D_{m+1}J_mf,\mathsf W_{m+1}D_{m+1}J_mf\rangle
 =\langle D_mf,K_m^*\mathsf W_{m+1}K_mD_mf\rangle=E_m(f).
\]
En différentiant la première identité de
\eqref{x_reciprocidad:vii:eq:familia-refinamiento-memoria}, on obtient
\(\dot D_{m+1}J_m=K_m\dot D_m\). Par conséquent,
\[
 2\operatorname{Re}
 \langle D_{m+1}J_mf,\mathsf W_{m+1}\dot D_{m+1}J_mf\rangle
 =2\operatorname{Re}\langle D_mf,\mathsf W_m\dot D_mf\rangle.
\]
La dérivée de la seconde identité est
\(K_m^*\dot{\mathsf W}_{m+1}K_m=\dot{\mathsf W}_m\), ce qui démontre
la coïncidence des termes de variation du poids.
\end{proof}

L'énoncé utilise la compatibilité de la famille de transports. Si \(J_m\) dépend du paramètre, la règle de dérivation en chaîne conserve également le terme
\[
 2\operatorname{Re}
 \langle D_{m+1}J_mf,\mathsf W_{m+1}D_{m+1}\dot J_mf\rangle.
\]
La naturalité obtenue transporte conjointement la fonctionnelle, ses
incidences et sa différentielle.

\section{Réalisation matérielle de la réponse}
\label{x_reciprocidad:vii:subsec:realizacion-respuesta-discreta}

Le résultat de cette section est le covecteur explicite
\(\mathfrak j_m\), avec son représentant \(\mathcal J_a\).
La réalisation matérielle compose ce covecteur avec la variation du
transport induite par la connexion et avec la normalisation de l'action.
Si \(\mathcal L_m\) désigne cette application différentielle, le covecteur sur
les variations de connexion est \(\mathcal L_m^*\mathfrak j_m\).
La règle de composition
\eqref{x_reciprocidad:vii:eq:variacion-transporte-compuesto} évalue sa partie
correspondant à des produits finis de transports.

Dans la section~\ref{x_reciprocidad:vii:sec:corriente-respuesta-cartan}, le courant
\(\sigma_{IJ}\) est une trois-forme définie par la variation de l'action
matérielle. L'identification entre les deux réalisations conserve le champ,
la représentation, le poids, la normalisation dimensionnelle et le passage à la
limite utilisés par cette action. Ces données permettent d'évaluer la
contraction qui intervient dans la densité torsionnelle. Le développement
présent apporte la différentielle discrète qui se compose dans ce calcul.

% END OWNER


% BEGIN OWNER /Users/ruben/Documents/New project/output/EXTREMA_MEDIA_RAZON_RECIPROCIDAD_ES_EN_20260918/ARTICULO/ES/source/antecedentes_vii/30c_composicion_corriente_conexion.tex
% Artículo VII, incorporación aditiva, revisión 05.
% Arquitectura autoral: energía de memoria APP--TRIT--TPK, S01.
% Antecedente recuperado: diferencial y refinamiento de 30b.
% Formalización reunida: nota COMPOSICION_VARIACIONAL_TRIT_CONEXION.md.
% Se conserva íntegra 30b; su representante antihermítico es la
% restricción unitaria del diferencial que aquí se desarrolla.
% Esta sección evalúa la corriente de una acción y una realización
% declaradas. No selecciona retrospectivamente su campo ni su peso.

\chapter{Réalisation du courant en coordonnées de connexion}
\label{x_reciprocidad:vii:sec:composicion-corriente-conexion}

L'énergie de mémoire attribue à chaque route une réponse qui conserve ses incidences et l'ordre des transports. Pour l'insérer dans une action géométrique, ses composantes par rapport aux variations de connexion sont nécessaires. Ce passage s'obtient en différentiant le transport qui réalise déjà la route : le courant de connexion est la composition de cette différentielle avec le covecteur énergétique de la section précédente.

Les régimes elliptique, parabolique et hyperbolique du TRIT requièrent de
conserver la différentielle complète avant de la restreindre à une
représentation particulière. Les rotations unitaires constituent un
domaine de cette construction. L'élargissement suivant conserve également
les directions de dilatation et de transport admises par la réalisation
géométrique utilisée.

\section{Différentielle sur les transports inversibles}

On maintient les fibres, les champs et le poids de
\eqref{x_reciprocidad:vii:eq:energia-memoria-discreta}. Considérons maintenant des transports
inversibles et des familles différentiables avec
\(\dot U_a=A_aU_a\), où \(A_a\) appartient à l'espace tangent
de la représentation choisie. Nous écrivons
\[
 v_a=U_af_{s(a)},\qquad r_a=f_{t(a)}-v_a,\qquad
 q=\mathsf W_mr,\qquad E_m=\langle r,\mathsf W_mr\rangle.
\]

\begin{proposition}[Représentant complet de la différentielle]
\label{x_reciprocidad:vii:prop:diferencial-invertible}
À champ et poids fixes, on a
\begin{equation}
 \mathrm d_UE_m(A)=-2\operatorname{Re}\sum_a
       \langle q_a,A_av_a\rangle
 =\sum_a\operatorname{Re}\operatorname{Tr}(\mathcal G_a^*A_a),
 \qquad \mathcal G_a=-2q_av_a^*.
 \label{x_reciprocidad:vii:eq:gradiente-transporte-completo}
\end{equation}
Les projections antihermitienne et hermitienne de ce représentant sont
\begin{equation}
 \frac{\mathcal G_a-\mathcal G_a^*}{2}
   =v_aq_a^*-q_av_a^*=\mathcal J_a,
 \qquad
 \frac{\mathcal G_a+\mathcal G_a^*}{2}
   =-(q_av_a^*+v_aq_a^*).
 \label{x_reciprocidad:vii:eq:proyecciones-gradiente-memoria}
\end{equation}
La restriction aux générateurs antihermitiens retrouve exactement
\eqref{x_reciprocidad:vii:eq:representante-corriente-discreta}.
\end{proposition}

\begin{proof}
L'identité \(\dot r_a=-A_av_a\) et le caractère autoadjoint du poids donnent \(\dot E_m=2\operatorname{Re}\langle q,\dot r\rangle\). En outre,
\[
 \operatorname{Tr}(\mathcal G_a^*A_a)
 =-2\operatorname{Tr}(v_aq_a^*A_a)=-2q_a^*A_av_a.
\]
Les deux projections s'obtiennent en prenant l'adjoint. Pour le produit
réel de Hilbert--Schmidt, les sous-espaces hermitien et antihermitien sont
orthogonaux ; la partie hermitienne a donc une contraction nulle avec
un générateur antihermitien.
\end{proof}

La formule permet des poids qui couplent les arêtes : leur action se calcule dans
\(q=\mathsf W_mr\) avant la séparation des contributions de chaque
incidence. Si le champ ou le poids varient, on conserve l'expression complète
\eqref{x_reciprocidad:vii:eq:variacion-completa-memoria}, également valable pour les
transports inversibles considérés ici.

\section{Représentations quadratiques du TRIT}

Dans une réalisation matricielle bidimensionnelle de \(\mathbb A_\tau\), on peut prendre
\[
 J_{+}=\begin{pmatrix}0&1\\-1&0\end{pmatrix},\qquad
 J_{-}=\begin{pmatrix}0&1\\1&0\end{pmatrix},\qquad
 J_{0}=\begin{pmatrix}0&1\\0&0\end{pmatrix}.
\]
Par multiplication directe,
\[
 J_+^2=-I,\qquad J_-^2=I,\qquad J_0^2=0.
\]
Leurs exponentielles réalisent respectivement les régimes elliptique, hyperbolique et parabolique. Le premier est orthogonal pour le produit positif fixé. Le second préserve la forme \(B=\operatorname{diag}(-1,1)\), puisque \(J_-^{\mathsf T}B+BJ_-=0\). Le troisième représente l'algèbre des nombres duaux. Ces représentations fixent des exemples des types quadratiques ; leur utilisation dans un écran d'une autre dimension conserve l'application de représentation correspondante.

Un calcul distingue les composantes de la différentielle. Sur une arête,
prenons \(U=I\), \(f_s=(1,1)^{\mathsf T}\),
\(f_t=2f_s\) et \(\mathsf W=I\). Alors \(q=v=f_s\),
\(\mathcal J=0\), tandis que
\begin{equation}
 \mathrm dE(J_-)=-4,\qquad \mathrm dE(J_0)=-2.
 \label{x_reciprocidad:vii:eq:control-regimenes-corriente}
\end{equation}
Le représentant complet conserve les deux réponses. La distinction
identifie l'espace des variations avant d'effectuer la contraction
matérielle.

\section{Composition des routes et inversion orientée}

\begin{proposition}[Transport contragrédient de la différentielle]
\label{x_reciprocidad:vii:prop:pullback-invertible}
Pour \(U_\gamma=U_N\cdots U_1\), soient
\(B_k=U_N\cdots U_{k+1}\) et \(B_N=I\). On a
\[
 \dot U_\gamma U_\gamma^{-1}
   =\sum_kB_kA_kB_k^{-1}.
\]
Le représentant \(\mathcal G_\gamma\) dans la fibre finale induit
à l'incidence \(k\)
\begin{equation}
 \mathcal G_k=B_k^*\mathcal G_\gamma(B_k^{-1})^*.
 \label{x_reciprocidad:vii:eq:transporte-contragrediente}
\end{equation}
\end{proposition}

\begin{proof}
La règle du produit, suivie de la multiplication par
\(U_\gamma^{-1}\), démontre la première égalité. La cyclicité
de la trace donne
\[
 \operatorname{Re}\operatorname{Tr}
  (\mathcal G_\gamma^*B_kA_kB_k^{-1})
 =\operatorname{Re}\operatorname{Tr}
  ((B_k^*\mathcal G_\gamma(B_k^{-1})^*)^*A_k).
\]
Pour \(B_k\) unitaire, on retrouve la conjugaison de
\ref{x_reciprocidad:vii:prop:composicion-respuesta}.
\end{proof}

\begin{proposition}[Inversion par congruence énergétique]
\label{x_reciprocidad:vii:prop:inversion-congruencia}
Pour une énergie locale par arête, l'inversion est représentée par
\begin{equation}
 U_{\bar a}=U_a^{-1},\qquad
 r_{\bar a}=-U_a^{-1}r_a,\qquad
 W_{\bar a}=U_a^*W_aU_a.
 \label{x_reciprocidad:vii:eq:inversion-peso-general}
\end{equation}
Elle conserve l'énergie. Si \(W_a\) reste fixe,
\[
 A_{\bar a}=-U_a^{-1}A_aU_a,\qquad
 \dot W_{\bar a}=U_a^*(A_a^*W_a+W_aA_a)U_a.
\]
La variation complète de l'énergie inversée coïncide avec la variation originale.
\end{proposition}

\begin{proof}
La substitution dans \(r_{\bar a}^*W_{\bar a}r_{\bar a}\) produit \(r_a^*W_ar_a\). La dérivée de l'inverse est \(\dot U_a^{-1}=-U_a^{-1}A_a\) ; celle du poids s'obtient par la règle du produit. L'égalité des énergies vaut pour toute la famille, de sorte que leurs dérivées coïncident, y compris le terme \(\langle r_{\bar a},\dot W_{\bar a}r_{\bar a}\rangle\).
\end{proof}

La congruence conserve la positivité pour tout transport inversible.
Dans la représentation unitaire, elle se réduit à la règle d'inversion de
\ref{x_reciprocidad:vii:prop:inversion-respuesta}. La somme utilise une orientation
par arête, ou la moitié de la somme sur les deux orientations.

\section{Évaluation des composantes de courant}

Soit \(\theta=(\theta_b)\) une variation réelle des coordonnées
de connexion de la réalisation utilisée. La différentielle effective du
transport définit les opérateurs
\begin{equation}
 A_a(\theta)=\dot U_aU_a^{-1}
     =\sum_bB_{ab}\theta_b.
 \label{x_reciprocidad:vii:eq:diferencial-conexion-transporte}
\end{equation}
Pour une composition finie, leurs contributions s'obtiennent par
\eqref{x_reciprocidad:vii:eq:transporte-contragrediente}. Les opérateurs
\(B_{ab}\) sont les dérivées de la représentation de connexion qui
réalise la route.

\begin{theorem}[Courant de connexion de l'action de mémoire]
\label{x_reciprocidad:vii:thm:corriente-conexion-evaluada}
Pour un terme matériel \(S_{\rm mem}=\mathfrak a_m E_m\), avec \(\mathfrak a_m\) sa normalisation dans la droite d'action, à champ, poids et normalisation fixes, on a
\begin{equation}
 \delta S_{\rm mem}=-\frac12\sum_b\theta_b s_b,\qquad
 s_b=4\mathfrak a_m\operatorname{Re}
        \sum_a\langle q_a,B_{ab}v_a\rangle.
 \label{x_reciprocidad:vii:eq:componentes-corriente-conexion}
\end{equation}
Si l'appariement cellulaire entre les bases de connexion et de courant
est une matrice inversible \(M\), définie par
\(\delta S_{\rm mem}=-\tfrac12\theta^{\mathsf T}M\sigma\),
les coordonnées du courant sont
\begin{equation}
 \sigma=M^{-1}s.
 \label{x_reciprocidad:vii:eq:coordenadas-corriente-material}
\end{equation}
\end{theorem}

\begin{proof}
On substitue \eqref{x_reciprocidad:vii:eq:diferencial-conexion-transporte} dans
\eqref{x_reciprocidad:vii:eq:gradiente-transporte-completo}, on multiplie par
\(\mathfrak a_m\) et on regroupe les coefficients de
\(\theta_b\). La convention variationnelle \(-1/2\) fixe le
facteur \(4\). Comparer les deux expressions pour toute variation
\(\theta\) donne \(M\sigma=s\) ; l'inversibilité de
l'appariement démontre l'existence et l'unicité de ces coordonnées.
\end{proof}

La formule réalise explicitement la composition mentionnée dans \ref{x_reciprocidad:vii:subsec:realizacion-respuesta-discreta}. Pour évaluer chaque composante, on transporte le champ, on calcule le résidu, on applique le poids et on contracte avec la variation de connexion de la même réalisation. L'accouplement \(M\) conserve l'orientation et le volume des cellules. Un accouplement diagonal donne la division correspondante par leurs volumes ; les autres bases conservent la matrice complète.

Si la connexion intervient aussi dans \(f\), \(\mathsf W_m\)
ou \(\mathfrak a_m\), on applique
\eqref{x_reciprocidad:vii:eq:variacion-completa-memoria} et l'on ajoute
\(E_m\delta\mathfrak a_m\). Les contributions de tous les
termes matériels sont additionnées avant de composer la réponse de Cartan.

\begin{corollary}[Compatibilité du courant avec le raffinement]
\label{x_reciprocidad:vii:cor:refinamiento-corriente-conexion}
Pour les familles compatibles de
\ref{x_reciprocidad:vii:thm:refinamiento-respuesta}, soit \(P_m\) l'application fixe
qui envoie les variations de connexion du niveau \(m\) au niveau
\(m+1\). Si les actions de mémoire coïncident sur ces champs
et transports raffinés, leurs covecteurs satisfont
\[
 s_m=P_m^{\mathsf T}s_{m+1},\qquad
 M_m\sigma_m=P_m^{\mathsf T}M_{m+1}\sigma_{m+1}.
\]
\end{corollary}

\begin{proof}
Différencier l'égalité des actions dans la direction
\(P_m\theta\) donne
\(-\tfrac12\theta^{\mathsf T}s_m
 =-\tfrac12\theta^{\mathsf T}P_m^{\mathsf T}s_{m+1}\).
Le caractère arbitraire de \(\theta\) et les représentations cellulaires
des deux covecteurs prouvent les deux identités. La composition de
ces applications itère l'égalité à tout nombre fini de raffinements.
\end{proof}

\section{Insertion dans l'équation de Cartan}

La section~\ref{x_reciprocidad:vii:sec:corriente-respuesta-cartan} normalise la trois-forme matérielle au moyen de \(\delta_\omega S_m=-\tfrac12\int\delta\omega^{IJ}\wedge
\sigma_{IJ}\). Le théorème précédent fournit ses composantes cellulaires pour l'action de mémoire et l'accouplement déclarés. La réponse de Cartan utilise ce courant avec la cotétrade, le paramètre de Barbero--Immirzi et le couplage gravitationnel de la même réalisation.

La dépendance matérielle par rapport à la contorsion détermine comment
cette équation se résout. Pour une matière affine, le courant est indépendant
de la contorsion et l’on applique l’inversion linéaire et l’élimination
quadratique démontrées dans la section suivante. Pour une action
d’holonomie, la différentielle est évaluée à la connexion actuelle et cette
dépendance est conservée lors de la résolution de l’équation. Par exemple, une arête avec
\(f_s=f_t=1\), un poids un et \(U(t)=e^{it}\) a
\[
 E(t)=2-2\cos t,\qquad E'(t)=2\sin t.
\]
Le courant dépend ici de \(t\). La formule complète maintient
cette dépendance au lieu de la remplacer par sa valeur dans une connexion
particulière.

Enfin, la densité gravitationnelle s'obtient à partir de la variation
métrique de l'action effective, avec son champ et sa normalisation.
La représentation du courant, la contorsion qu'il produit et la
densité résultante sont les opérations successives de cette composition.
La loi de dilution ultérieure transporte l'amplitude ainsi évaluée dans
la section matérielle correspondante.

% END OWNER


% BEGIN OWNER /Users/ruben/Documents/New project/output/EXTREMA_MEDIA_RAZON_RECIPROCIDAD_ES_EN_20260918/ARTICULO/ES/source/antecedentes_vii/30e_corriente_extension_minima.tex
% Artículo VII. Incorporación aditiva, 10-09-2026.
% RESULTADO_RECUPERADO: fuentes_conservadas/18_clausura_energetica.tex,
% líneas 357--485: Gramiano, peso, extensión mínima y eliminación de Schur;
% líneas 149--180 y 235--295: normalización y unicidad de la respuesta.
% Residencia anterior en VII: 30_pantalla_y_respuesta.tex,
% vii:prop:incidencia, vii:eq:lambda-pantalla, vii:prop:minima-extension.
% FORMALIZACION_NUEVA / CERTIFICADO_NUEVO: diferencial del mínimo con
% incidencia, frontera y normalización variables; incidencia transportada
% por bloques y composición con coordenadas de conexión.
% La familia de rutas que realiza cada incidencia se mantiene explícita.
% No se atribuyen al propietario histórico ni esa familia diferenciable
% completa ni la evaluación de una amplitud material de espín.
% Dependencias: 30, 30c y la convención variacional de 31. No modifica
% esos fragmentos ni constituye por sí solo una realización material.

\chapter{Différentielle de l'extension minimale de frontière}
\label{x_reciprocidad:vii:sec:corriente-extension-minima}

L'incidence cubique détermine simultanément le champ d'extension
minimale et le poids de sa réponse. Pour une donnée de frontière, tous deux sont
calculés au moyen de l'inverse du gramien d'incidence. On développe
maintenant la différentielle de cette même opération, en conservant les variations
de la frontière, du transport et de l'échelle de normalisation. De cette
manière, le champ qui intervient dans la réponse est sélectionné par
le problème variationnel d'écran.

\section{Champ minimisant et poids de réponse}
\label{x_reciprocidad:vii:subsec:extension-minima-parametrizada}

Soient \(H\) et \(Q\) des espaces hermitiens de dimension finie, à produits fixes et antilinéaires en la première variable. Le cas réel s'obtient en remplaçant les adjoints par les transposées. Dans un ouvert de paramètres réels \(\Omega\), considérons des familles de classe \(C^1\)
\[
 \mathsf C(\theta):H\longrightarrow Q,\qquad
 b(\theta)\in Q,\qquad \chi(\theta)>0,
 \qquad \operatorname{Ran}\mathsf C(\theta)=Q.
\]
La surjectivité équivaut à un rang en lignes maximal dans des bases orthonormées. Nous définissons le gramien, le potentiel conjugué et l'extension
\begin{equation}
 \mathsf L=\mathsf C\mathsf C^*,\qquad
 y=\mathsf L^{-1}b,\qquad f_{\min}=\mathsf C^*y,
 \qquad \Lambda=\chi\mathsf L^{-1}.
 \label{x_reciprocidad:vii:eq:gramiano-extension-minima}
\end{equation}
L'incidence \(\mathsf C\) et l'accouplement cellulaire \(M\) de \eqref{x_reciprocidad:vii:eq:coordenadas-corriente-material} sont des opérateurs différents.

\begin{proposition}[Sélection variationnelle de l'extension]
\label{x_reciprocidad:vii:prop:minimo-parametrizado}
L'extension \(f_{\min}\) est l'unique minimisante de la norme parmi
les champs \(f\in H\) qui satisfont \(\mathsf C f=b\). Son énergie est
\begin{equation}
 E(\theta)=\chi\min_{\mathsf C f=b}\|f\|^2
 =\chi\langle b,\mathsf L^{-1}b\rangle
 =\langle b,\Lambda b\rangle.
 \label{x_reciprocidad:vii:eq:energia-extension-minima}
\end{equation}
Le champ \(f_{\min}\), le potentiel \(y\) et l'énergie \(E\)
dépendent des paramètres de manière \(C^1\).
\end{proposition}

\begin{proof}
La surjectivité de \(\mathsf C\) rend \(\mathsf C^*\) injectif, de sorte que \(\langle z,\mathsf Lz\rangle=\|\mathsf C^*z\|^2>0\) pour \(z\ne0\). Ainsi, \(\mathsf L\) est inversible et \(\mathsf C f_{\min}=\mathsf L y=b\). Toute solution est \(f_{\min}+z\), avec \(z\in\ker\mathsf C\). Puisque \(f_{\min}\in\operatorname{Ran}\mathsf C^*=(\ker\mathsf C)^\perp\),
\[
 \|f_{\min}+z\|^2=\|f_{\min}\|^2+\|z\|^2,
 \qquad
 \|f_{\min}\|^2=\langle y,\mathsf L y\rangle
 =\langle b,\mathsf L^{-1}b\rangle.
\]
Ces identités prouvent le minimum, son unicité et sa valeur. L'inversion
est une opération différentiable sur l'ouvert des matrices inversibles ;
les expressions de \eqref{x_reciprocidad:vii:eq:gramiano-extension-minima} fournissent
la régularité affirmée.
\end{proof}

La réalisation cubique de \ref{x_reciprocidad:vii:prop:incidencia} correspond à
\(H=\mathbb R^{\{-1,1\}^3}\), au quotient d'écran
\(Q=P_\square\mathbb R^6\) et à l'incidence \(\mathsf C=M:H\to Q\).
Son gramien et sa normalisation sont
\begin{equation}
 \mathsf L_\square=12P_u+4P_-,\qquad
 \chi_m=\frac{112}{81\ell_m^2},\qquad
 \Lambda_m=\frac{28}{243\ell_m^2}(P_u+3P_-).
 \label{x_reciprocidad:vii:eq:normalizacion-minimo-cubico}
\end{equation}
Le facteur \(112/81\) provient de la soudure et du défaut quadratique de compression de \eqref{x_reciprocidad:vii:eq:lambda-pantalla}. Si \(b\) a la dimension d'une longueur, \(y\) et \(f_{\min}\) l'ont également, l'incidence est sans dimension et \(E\) est sans dimension. La conversion en action utilise la normalisation dans la droite d'action spécifiée plus loin.

\Needspace{15\baselineskip}
\section{Différentielle complète du minimum}
\label{x_reciprocidad:vii:subsec:diferencial-minimo}

\begin{theorem}[Réponse à la variation d'incidence et de frontière]
\label{x_reciprocidad:vii:thm:diferencial-minimo}
Pour toute direction tangente réelle \(\delta\theta\), on a
\begin{equation}
 \delta E
 =\delta\chi\,\langle b,y\rangle
   +2\chi\operatorname{Re}
       \langle y,\delta b-(\delta\mathsf C)f_{\min}\rangle.
 \label{x_reciprocidad:vii:eq:diferencial-minimo-completo}
\end{equation}
Si \(\chi=112/(81\ell_m^2)\), avec \(\ell_m>0\) variable,
\begin{equation}
 \delta E=-2E\frac{\delta\ell_m}{\ell_m}
 +2\chi_m\operatorname{Re}
       \langle y,\delta b-(\delta\mathsf C)f_{\min}\rangle.
 \label{x_reciprocidad:vii:eq:diferencial-minimo-escala}
\end{equation}
\end{theorem}

\begin{proof}
Différentier \(\mathsf L\mathsf L^{-1}=I_Q\) donne
\[
 \delta\mathsf L^{-1}
 =-\mathsf L^{-1}(\delta\mathsf L)\mathsf L^{-1},\qquad
 \delta\mathsf L=(\delta\mathsf C)\mathsf C^*
                      +\mathsf C(\delta\mathsf C)^*.
\]
En appliquant la règle du produit à \eqref{x_reciprocidad:vii:eq:energia-extension-minima} et en utilisant le caractère autoadjoint de \(\mathsf L^{-1}\), il vient
\[
 \delta E=\delta\chi\,\langle b,y\rangle
 +2\chi\operatorname{Re}\langle y,\delta b\rangle
 -\chi\langle y,(\delta\mathsf L)y\rangle.
\]
Les deux termes du dernier produit sont conjugués entre eux et
\(\mathsf C^*y=f_{\min}\) ; par conséquent,
\[
 \langle y,(\delta\mathsf L)y\rangle
 =2\operatorname{Re}\langle y,(\delta\mathsf C)f_{\min}\rangle.
\]
On obtient \eqref{x_reciprocidad:vii:eq:diferencial-minimo-completo}. Enfin, \(\delta\chi_m=-2\chi_m\delta\ell_m/\ell_m\) produit \eqref{x_reciprocidad:vii:eq:diferencial-minimo-escala}.
\end{proof}

La formule incorpore la variation du champ minimisant au moyen de l'
équation qui le détermine. Elle se vérifie aussi directement :
\(\mathsf C\delta f_{\min}=\delta b-(\delta\mathsf C)f_{\min}\)
et \(f_{\min}=\mathsf C^*y\) impliquent
\[
 \delta\|f_{\min}\|^2
 =2\operatorname{Re}\langle f_{\min},\delta f_{\min}\rangle
 =2\operatorname{Re}
     \langle y,\delta b-(\delta\mathsf C)f_{\min}\rangle.
\]
Par exemple, si \(\mathsf C\) est fixe et \(b=\ell_m b_0\), avec
\(b_0\) fixe, les deux termes de
\eqref{x_reciprocidad:vii:eq:diferencial-minimo-escala} se compensent et l'énergie
reste constante. Ce contrôle conserve conjointement la dimension
de frontière et sa normalisation.

\section{Quotient fixe et changements de repère}
\label{x_reciprocidad:vii:subsec:marcos-minimo}

Les dérivées précédentes sont prises dans des espaces \(H,Q\) et avec des produits hermitiens fixes. Dans la présentation à six faces, l'inverse agit toujours sur \(Q=P_\square H_F\). Pour une famille de quotients \(Q(\theta)\), on fixe d'abord une trivialisation isométrique \(V_Q(\theta):Q\to Q(\theta)\), et de même pour le domaine. Les dérivées de ces applications sont incluses lors de la différentiation des expressions transportées vers les espaces fixes. Les formules s'appliquent dans chaque ouvert où le rang est conservé. Les changements du produit énergétique exigent de différentier également ce produit ; une métrique lorentzienne et le produit positif utilisé pour minimiser conservent ici leurs types distincts.

\begin{proposition}[Covariance du minimum]
\label{x_reciprocidad:vii:prop:covariancia-minimo}
Soient \(g_H,g_Q\) unitaires. Sous
\[
 \mathsf C'=g_Q\mathsf C g_H^*,\qquad b'=g_Qb,\qquad \chi'=\chi,
\]
on a \(y'=g_Qy\), \(f'_{\min}=g_Hf_{\min}\) et \(E'=E\).
Lorsque les changements de cadre dépendent du paramètre, la différentielle
complète conserve cette égalité. En particulier, une variation pure de
cadre a \(\delta E=0\).
\end{proposition}

\begin{proof}
Le gramien transformé est \(\mathsf L'=g_Q\mathsf Lg_Q^*\),
d'où l'on obtient les transformations de \(y,f_{\min}\) et
l'égalité des énergies. Pour une variation pure de repère, en un point
avec des repères identité, soient \(A_Q^*=-A_Q\) et \(A_H^*=-A_H\)
leurs générateurs. Alors
\[
 \delta\mathsf C=A_Q\mathsf C-\mathsf C A_H,\qquad
 \delta b=A_Qb,\qquad
 \delta b-(\delta\mathsf C)f_{\min}=\mathsf C A_Hf_{\min}.
\]
Sa contribution à \eqref{x_reciprocidad:vii:eq:diferencial-minimo-completo} est
\(2\chi\operatorname{Re}\langle f_{\min},A_Hf_{\min}\rangle=0\).
La preuve inclut ainsi les termes dus aux cadres mobiles.
\end{proof}

\section{Incidence transportée et composantes de connexion}
\label{x_reciprocidad:vii:subsec:incidencia-conexion-minima}

La dépendance de connexion est spécifiée avant d'effectuer la
contraction. Une réalisation explicite de l'incidence par blocs
s'obtient comme suit. Soient \(\mathcal K\) une fibre hermitienne commune,
\(F=\{(i,\epsilon):1\leq i\leq3,\ \epsilon=\pm1\}\) et
\(V=\{-1,1\}^3\). Pour chaque incidence \(\nu_i=\epsilon\),
la réalisation utilisée déclare une route \(\gamma_{i,\epsilon,\nu}\)
de la fibre du sommet à celle de la face, avec le transport
\(U_{i,\epsilon,\nu}\). Les trivialisations identifient ces fibres
à \(\mathcal K\). Nous définissons
\begin{equation}
 \begin{split}
 (\mathsf B_Uf)_{i,\epsilon}
   &=\sum_{\nu_i=\epsilon}U_{i,\epsilon,\nu}f_\nu,\\
 \mathsf C_U&=(q_\square\otimes I_{\mathcal K})\mathsf B_U:
   \bigoplus_{\nu\in V}\mathcal K\longrightarrow Q\otimes\mathcal K,
 \qquad q_\square=P_\square:H_F\to Q.
 \end{split}
 \label{x_reciprocidad:vii:eq:incidencia-transportada-minima}
\end{equation}
Il s'agit d'une continuation variationnelle explicite de l'incidence cubique
dans la représentation des routes déclarée. Pour les transports identité,
\(\mathsf C_U=M\otimes I_{\mathcal K}\) et l'on retrouve
\eqref{x_reciprocidad:vii:eq:normalizacion-minimo-cubico}, tensorisée avec la fibre.
Le gramien reste positif dans un voisinage de cette réalisation,
car la positivité définie est une condition ouverte. Dans chaque domaine
de rang plein, le poids et le champ sont fixés par
\eqref{x_reciprocidad:vii:eq:gramiano-extension-minima}. Les routes qui réalisent les
incidences, leurs orientations et leurs mémoires sont des données de cette
réalisation du transport APP--TRIT--TPK et accompagnent la construction.

Si \(\omega=(\omega_j)\) sont des coordonnées réelles de connexion, la dérivée de \(\mathsf C_U\) se calcule en substituant dans chaque bloc
\begin{equation}
 \partial_j U_\gamma
 =\sum_{k=1}^{N}
    U_N\cdots U_{k+1}(\partial_jU_k)U_{k-1}\cdots U_1,
 \qquad U_\gamma=U_N\cdots U_1.
 \label{x_reciprocidad:vii:eq:derivada-incidencia-rutas}
\end{equation}
Les produits vides sont des identités. La formule provient de la règle du produit et conserve la place de chaque opération dans la route. Ainsi, \(\mathsf C_j:=\partial_j\mathsf C_U\) s'obtient à partir d'une somme finie de matrices connues dans la réalisation choisie, suivie de la projection \(q_\square\otimes I\). Si un quotient mobile est utilisé, sa dérivée est incorporée conformément à \ref{x_reciprocidad:vii:subsec:marcos-minimo}. La matrice cubique entière fixe correspond au cas \(\mathsf C_j=0\) ; la réponse à la connexion provient de la dépendance des transports effectivement déclarée dans \eqref{x_reciprocidad:vii:eq:incidencia-transportada-minima}.

% REMATE_LOCAL_X_EXTENSION_MINIMA
\Needspace{34\baselineskip}
\begin{theorem}[Composantes variationnelles de l'extension minimale]
\label{x_reciprocidad:vii:thm:componentes-minimo-conexion}
Pour l'action d'écran
\(S_{\min}=\mathfrak a_m E\), où \(\mathfrak a_m\) est sa
normalisation sur la droite d'action, écrivons
\[
 b_j=\partial_jb,\qquad \chi_j=\partial_j\chi,
 \qquad a_j=\partial_j\mathfrak a_m.
\]
Avec la convention \(\delta S_{\min}=-\tfrac12\sum_j s_j\delta\omega_j\),
les composantes sont
\begin{equation}
 \begin{split}
 s_j={}&4\mathfrak a_m\chi\operatorname{Re}
             \langle y,\mathsf C_jf_{\min}-b_j\rangle\\
      &-2\mathfrak a_m\chi_j\langle b,y\rangle-2a_jE.
 \end{split}
 \label{x_reciprocidad:vii:eq:componentes-minimo-conexion}
\end{equation}
Pour la normalisation cubique, le second terme équivaut à
\(4\mathfrak a_m E\,\partial_j\ell_m/\ell_m\).
Si l'appariement cellulaire de cette même action a une matrice
inversible \(\mathsf M_{\rm cel}\), ses coordonnées de courant
sont \(\sigma=\mathsf M_{\rm cel}^{-1}s\).
\end{theorem}

\begin{proof}
La règle du produit donne
\(\partial_jS_{\min}=a_jE+\mathfrak a_m\partial_jE\).
Substituer \eqref{x_reciprocidad:vii:eq:diferencial-minimo-completo} et multiplier
par \(-2\) prouve \eqref{x_reciprocidad:vii:eq:componentes-minimo-conexion}.
La dérivée de \(\chi_m\) produit l'expression d'échelle.
Enfin, comparer
\(\delta S_{\min}=-\tfrac12\delta\omega^{\mathsf T}s\)
avec \(-\tfrac12\delta\omega^{\mathsf T}\mathsf M_{\rm cel}\sigma\)
pour toute variation donne \(\mathsf M_{\rm cel}\sigma=s\),
et son inverse détermine les coordonnées.
\end{proof}

La construction retrouve un champ et un poids sélectionnés conjointement
par l'incidence, et évalue leur réponse aux variations de la même
réalisation. La donnée de frontière \(b\), la famille de routes, l'échelle
\(\ell_m\), la normalisation \(\mathfrak a_m\) et l'appariement
cellulaire conservent leurs règles de publication. L'interprétation comme
courant matériel utilise ces objets et l'action complète. La
réponse de Cartan et sa densité métrique ultérieure se composent avec
elle dans le domaine de matière correspondant, en conservant toute
dépendance du courant vis-à-vis de la connexion.

% END OWNER

\begin{HMTCompactSection}

% BEGIN OWNER /Users/ruben/Documents/New project/output/EXTREMA_MEDIA_RAZON_RECIPROCIDAD_ES_EN_20260918/ARTICULO/ES/source/antecedentes_vii/31_corriente_y_respuesta_cartan.tex
% Artículo VII. Fragmento de cuerpo: corriente y respuesta Cartan--Holst.
% RESULTADO_RECUPERADO: acción, corriente variacional, inversa de Holst,
% torsión algebraica, respuesta cuadrática y normalización métrica.
% Fuentes leídas íntegramente:
% [II10] output/ARTICULO_II_REV10_EDICION_INTEGRADA_20260910/manuscrito/
%        sections/10_gravedad.tex, líneas 109--155.
% [G11c] integral activo 2249, fuente/manuscrito/sections/md/
%        11c_gravedad_holonomia_rev6.tex, líneas 214--334.
% [B028] integral activo, fuente/manuscrito/integracion_83/parte_iv/body/
%        B028_ch68_realizacion_cartan_holst_barbero_7df2afd9f675_11_gravedad_cartan_holst.tex,
%        líneas 76--218; copia idéntica conservada en II10/fuentes/propietarios_II/G03/.
%        SHA-256 e1e1f3325399b3ef09c331ed8268e3883f7bcc561cd7a560cf80d1be7eb1e91c.
% CERTIFICADO_NUEVO de resultados recuperados: inversión explícita del mapa
% torsión -> D(e wedge e), prueba de unicidad de contorsión y eliminación
% homogénea del término cuadrático. Las dos inversas se comprobaron sobre
% las 24 bases correspondientes mediante aritmética racional exacta.
% Precisión de dominio: la respuesta lineal en corriente y su eliminación
% cuadrática se enuncian para materia afín en la contorsión. Se conserva
% la dependencia del coeficiente respecto de la acción material concreta.
%
% DEPENDENCIAS DE INTEGRACIÓN (no constituyen remisiones bibliográficas):
% 1. vii:sec:realizacion-cartan: cotetrada, conexión y realización del
%    transporte enriquecido; incluir su prueba en el cuerpo precedente.
% 2. vii:sec:acoplamientos: genealogía de gamma_HMT, G_int, c_int,
%    homogeneidad dimensional y elección de la carta, con pruebas.
% 3. Para evaluar una especie material: acción S_m concreta, representación
%    de sus campos y cálculo de su corriente sigma. Este fragmento recibe
%    esa acción declarada y demuestra su respuesta variacional.
% 4. La identificación de una corriente de pantalla con sigma, su
%    contracción física y su normalización material tienen residencia propia.
%    Aquí no se introduce ni se selecciona una amplitud s_0.
% 5. El balance métrico utiliza la covariancia de la acción material y
%    sus ecuaciones de movimiento; no supone conservación separada del
%    tensor de espín y del tensor material de Levi--Civita.

\chapter{Courant de spin et réponse algébrique de Cartan--Holst}
\label{x_reciprocidad:vii:sec:corriente-respuesta-cartan}

La réalisation du transport enrichi de la
section~\ref{x_reciprocidad:vii:sec:realizacion-cartan} fournit la cotétrade et la
connexion ; les sections dimensionnelles et le paramètre d'incidence des
sections~\ref{x_reciprocidad:v:ant:sec:barbero} et~\ref{x_reciprocidad:v:ant:sec:gravity} fixent leurs couplages. Leur composition
détermine un problème variationnel du premier ordre. Le courant matériel
y occupe une position précise : il est la dérivée de l'action par rapport à
la connexion et détermine la réponse torsionnelle avec la même normalisation.

\section{Domaine géométrique et normalisation du courant}
\label{x_reciprocidad:vii:subsec:dominio-corriente-cartan}

Nous travaillons sur une variété lisse orientée \(M\) de dimension quatre, munie d'un fibré lorentzien interne \(E\), d'une métrique \(\eta=\operatorname{diag}(-1,1,1,1)\) et d'une cotétrade non dégénérée \(e\in\Omega^1(M;E)\). Dans un repère local, nous écrivons \(e^I\), avec \(I=0,1,2,3\). La connexion métrique \(\omega^{IJ}=-\omega^{JI}\) agit au moyen de \(D_\omega\). Lorsque les champs \(\Psi\) sont spinoriels, une structure de spin et la représentation matérielle correspondante sont conservées. Nous définissons
\begin{equation}
 \Sigma^{IJ}=e^I\wedge e^J,\qquad
 F^{IJ}=\mathrm d\omega^{IJ}+\omega^I{}_{K}\wedge\omega^{KJ},
 \qquad T^I=D_\omega e^I.
 \label{x_reciprocidad:vii:eq:campos-cartan}
\end{equation}
L'opérateur \(\star\) agit sur les indices bivectoriels internes,
\begin{equation}
 (\star X)^{IJ}=\tfrac12\epsilon^{IJ}{}_{KL}X^{KL},
 \qquad \star^2=-I,\qquad
 \mathsf P_\gamma=\star+\gamma^{-1}I.
 \label{x_reciprocidad:vii:eq:operador-holst}
\end{equation}
On distingue ainsi cette dualité de la dualité extérieure des formes sur
\(M\). Les indices internes sont relevés et contractés avec \(\eta\).

Nous fixons une carte avec \(G_{\rm int}>0\), \(c_{\rm int}>0\),
\(\gamma=\gamma_{\rm HMT}\in\mathbb R\setminus\{0\}\) et
\(\Lambda_{\rm int}\) constants sur elle. Nous abrégeons
\(\kappa=8\pi G_{\rm int}/c_{\rm int}^4\) et \(c=c_{\rm int}\).
Ces sections demeurent fixes pendant la variation. L’action est
\begin{equation}
 \begin{split}
 S[e,\omega,\Psi]={}&
 \frac1{2\kappa c}\int_M\Sigma_{IJ}\wedge
                         (\mathsf P_\gamma F)^{IJ}
 -\frac{\Lambda_{\rm int}}{\kappa c}
                         \int_M\operatorname{vol}_e
 +S_{\rm m}[e,\omega,\Psi].
 \end{split}
 \label{x_reciprocidad:vii:eq:accion-cartan-holst}
\end{equation}
La convention de volume et d'orientation est celle qui donne
\(\Sigma_{IJ}\wedge(\star F)^{IJ}=R\operatorname{vol}_e\)
dans le secteur de Levi--Civita. Le type dimensionnel préalable fait de chaque
terme une section d'action.

Pour la variation de la connexion, on utilise des variations à support compact à l'intérieur, ou on fixe \(\delta\omega|_{\partial M}=0\). On peut également employer une complétion de bord dont le problème variationnel annule le même terme de frontière. À cotétrade et champs matériels fixes, le courant de spin est la forme de degré trois, à valeurs dans les bivecteurs internes duaux, déterminée par
\begin{equation}
 \delta_\omega S_{\rm m}
 =-\tfrac12\int_M\delta\omega^{IJ}\wedge\sigma_{IJ}.
 \label{x_reciprocidad:vii:eq:corriente-variacional}
\end{equation}
L'appariement avec une variation arbitraire à support compact fixe
\(\sigma\) de manière unique. Cette définition conserve l'unité
d'action du courant et le facteur de normalisation qui intervient dans
sa réponse géométrique.

\section{Variation de la connexion et inversion de Holst}
\label{x_reciprocidad:vii:subsec:variacion-holst}

\begin{theorem}[Équation variationnelle de Cartan--Holst]
\label{x_reciprocidad:vii:thm:ecuacion-cartan-holst}
Dans le domaine précédent, la stationnarité par rapport à la connexion
équivaut à
\begin{equation}
 D_\omega(\mathsf P_\gamma\Sigma)^{IJ}
       =\kappa c\,\sigma^{IJ}.
 \label{x_reciprocidad:vii:eq:ecuacion-cartan-holst}
\end{equation}
\end{theorem}
\begin{proof}
La variation de la courbure est
\(\delta F=D_\omega\delta\omega\). La dualité interne est parallèle
et autoadjointe pour la contraction bivectorielle ; comme \(\gamma\) est
constant, \(D_\omega\mathsf P_\gamma=\mathsf P_\gamma D_\omega\).
En posant \(B_{IJ}=(\mathsf P_\gamma\Sigma)_{IJ}\), on obtient
\[
 B_{IJ}\wedge D_\omega\delta\omega^{IJ}
 =\mathrm d(B_{IJ}\wedge\delta\omega^{IJ})
   -D_\omega B_{IJ}\wedge\delta\omega^{IJ}.
\]
La dernière forme change de signe lorsqu'on permute ses facteurs de degrés trois
et un. Les conditions au bord éliminent l'intégrale du terme
exact. Le terme cosmologique est indépendant de \(\omega\).
Par conséquent,
\[
 \delta_\omega S=
 \frac1{2\kappa c}\int_M\delta\omega^{IJ}\wedge
       \bigl[D_\omega B_{IJ}-\kappa c\,\sigma_{IJ}\bigr].
\]
Le caractère arbitraire de la variation prouve l'équation, y compris son signe et son coefficient.
\end{proof}

\begin{lemma}[Inverse réel de l'opérateur de Holst]
\label{x_reciprocidad:vii:lem:inversa-holst}
Pour \(\gamma\in\mathbb R\setminus\{0\}\),
\begin{equation}
 \mathsf P_\gamma^{-1}
 =\frac{\gamma^2}{1+\gamma^2}
                      (-\star+\gamma^{-1}I).
 \label{x_reciprocidad:vii:eq:inversa-holst}
\end{equation}
\end{lemma}
\begin{proof}
Les deux facteurs sont des polynômes en \(\star\) et commutent. Leur produit
avant normalisation est
\[
 (\star+\gamma^{-1}I)(-\star+\gamma^{-1}I)
 =-\star^2+\gamma^{-2}I=(1+\gamma^{-2})I.
\]
Le dénominateur \(1+\gamma^2\) est positif dans le domaine réel fixé. La multiplication par \(\gamma^2/(1+\gamma^2)\) donne l'identité.
\end{proof}

Le courant détermine ainsi la forme bivectorielle de degré trois
\begin{equation}
 Y^{IJ}:=\kappa c\,
   \frac{\gamma^2}{1+\gamma^2}
   (-\star+\gamma^{-1}I)\sigma^{IJ},
 \qquad D_\omega\Sigma^{IJ}=Y^{IJ}.
 \label{x_reciprocidad:vii:eq:corriente-dualizada}
\end{equation}

\section{Reconstruction ponctuelle de la torsion et de la contorsion}
\label{x_reciprocidad:vii:subsec:reconstruccion-torsion}

La règle de Leibniz donne
\begin{equation}
 D_\omega\Sigma^{IJ}
 =T^I\wedge e^J-e^I\wedge T^J
 =:\mathsf L_e(T)^{IJ}.
 \label{x_reciprocidad:vii:eq:mapa-torsion}
\end{equation}
En chaque point, \(\mathsf L_e\) agit de \(\Lambda^2T^*M\otimes E\) dans \(\Lambda^3T^*M\otimes\Lambda^2E\). Les deux espaces sont de dimension vingt-quatre. Soit \(E_I\) le repère dual de la cotétrade, de sorte que \(\iota_{E_I}e^J=\delta_I^J\).

\begin{theorem}[Inversion algébrique de l'équation torsionnelle]
\label{x_reciprocidad:vii:thm:inversion-torsion}
Pour une cotétrade non dégénérée, \(\mathsf L_e\) est un isomorphisme.
Son inverse sur la forme \(Y\) s'écrit
\begin{equation}
 B^I=\sum_J\iota_{E_J}Y^{IJ},\qquad
 b=\tfrac14\sum_I\iota_{E_I}B^I,\qquad
 T^I=B^I-e^I\wedge b.
 \label{x_reciprocidad:vii:eq:torsion-explicita}
\end{equation}
L'équation~\eqref{x_reciprocidad:vii:eq:ecuacion-cartan-holst} détermine donc
la torsion ponctuellement à partir de \(e\) et du courant qui y
figure.
\end{theorem}
\begin{proof}
Soit d'abord \(Y=\mathsf L_e(T)\) et posons
\(t=\sum_J\iota_{E_J}T^J\). La contraction de
\eqref{x_reciprocidad:vii:eq:mapa-torsion}, en utilisant le fait que \(T^I\) est de degré deux
et qu'il existe quatre éléments dans la cotétrade, donne
\[
 \sum_J\iota_{E_J}(T^I\wedge e^J)=2T^I,
 \qquad
 \sum_J\iota_{E_J}(e^I\wedge T^J)=T^I-e^I\wedge t.
\]
D'où \(B^I=T^I+e^I\wedge t\). Une seconde contraction produit
\[
 \sum_I\iota_{E_I}B^I=t+3t=4t.
\]
On retrouve \(t=b\) puis \(T^I=B^I-e^I\wedge b\).
En particulier, le noyau de \(\mathsf L_e\) est nul. L'égalité
des dimensions prouve sa surjectivité et rend la formule valable
pour toute forme \(Y\) du codomaine.
\end{proof}

La connexion se décompose de manière unique comme
\begin{equation}
 \omega^{IJ}=\mathring\omega^{IJ}(e)+K^{IJ},\qquad
 K^{IJ}=-K^{JI},\qquad T^I=K^I{}_{J}\wedge e^J,
 \label{x_reciprocidad:vii:eq:contorsion}
\end{equation}
où \(\mathring\omega(e)\) est la connexion de Levi--Civita et \(K\) la contorsion. Pour expliciter l'inverse, nous écrivons \(K_{IJ}=K_{IJK}e^K\) et \(T_I=\tfrac12T_{IJK}e^J\wedge e^K\). Alors
\begin{equation}
 T_{IJK}=K_{IKJ}-K_{IJK},\qquad
 K_{IJK}=\tfrac12(T_{KIJ}-T_{IJK}-T_{JKI}).
 \label{x_reciprocidad:vii:eq:contorsion-componentes}
\end{equation}
La seconde égalité s'obtient en combinant cycliquement la première et
en utilisant \(K_{IJK}=-K_{JIK}\) ; sa substitution restitue les composantes
de \(T\). La torsion et la contorsion sont donc deux
présentations équivalentes de la réponse algébrique à cotétrade fixée.

\begin{corollary}[Secteur sans courant et réponse linéaire minimale]
\label{x_reciprocidad:vii:cor:desacoplamiento-torsional}
Si \(\sigma=0\), on obtient \(T=0\) et \(\omega=\mathring\omega(e)\). Pour un couplage matériel dont le courant \(\sigma[e,\Psi]\) est indépendant de \(K\), les équations~\eqref{x_reciprocidad:vii:eq:corriente-dualizada}, \eqref{x_reciprocidad:vii:eq:torsion-explicita} et \eqref{x_reciprocidad:vii:eq:contorsion-componentes} fournissent une réponse unique \(K_*(e,\Psi)\), linéaire en \(\sigma\).
\end{corollary}
\begin{proof}
Les trois applications de reconstruction sont linéaires lorsque \(e\),
\(\gamma\) et \(\kappa c\) sont fixés. Le courant nul produit
\(Y=0\), puis \(T=0\) et enfin \(K=0\).
\end{proof}

Dans l'action minimale, l'équation de connexion contient cette réponse
ponctuelle. Son évolution s'obtient en faisant évoluer \(e\) et \(\Psi\) :
\(K_*=K_*(e,\Psi)\) change avec eux. Si une action matérielle conserve
une dépendance algébrique supplémentaire en \(K\), l'équation de Cartan
maintient sa forme variationnelle et doit être résolue avec ce courant ;
la réponse linéaire précédente correspond au domaine minimal spécifié.

\section{Élimination de la contorsion et contribution quadratique}
\label{x_reciprocidad:vii:subsec:respuesta-cuadratica}

Considérons maintenant le couplage minimal affine en \(K\), exprimé par
\begin{equation}
 S_{\rm m}[e,\mathring\omega+K,\Psi]
 =S_{\rm m}[e,\mathring\omega,\Psi]
   -\tfrac12\int_M K^{IJ}\wedge\sigma_{IJ}[e,\Psi].
 \label{x_reciprocidad:vii:eq:materia-afin-contorsion}
\end{equation}
Cette condition identifie l'action matérielle à laquelle s'applique
l'élimination quadratique ; le courant reste celui défini par
\eqref{x_reciprocidad:vii:eq:corriente-variacional}.

Le développement de la courbure est
\begin{equation}
 F(\mathring\omega+K)
 =F(\mathring\omega)+D_{\mathring\omega}K+K\wedge K,
 \qquad (K\wedge K)^{IJ}=K^I{}_{L}\wedge K^{LJ}.
 \label{x_reciprocidad:vii:eq:curvatura-contorsion}
\end{equation}
Comme \(D_{\mathring\omega}\Sigma=0\), le terme linéaire est la
dérivée extérieure de \(\Sigma_{IJ}\wedge(\mathsf P_\gamma K)^{IJ}\).
La contribution de bord correspondante reste explicite :
\begin{equation}
 S_\partial[e,K]=\frac1{2\kappa c}
     \int_{\partial M}\Sigma_{IJ}\wedge(\mathsf P_\gamma K)^{IJ}.
 \label{x_reciprocidad:vii:eq:borde-contorsion}
\end{equation}
Dans un domaine intérieur, ou avec la complétion de bord compatible,
la densité d'action qui dépend de la contorsion est
\begin{equation}
 \mathcal Q_{e,\gamma}(K)-\tfrac12K^{IJ}\wedge\sigma_{IJ},
 \qquad
 \mathcal Q_{e,\gamma}(K)
  =\frac1{2\kappa c}\Sigma_{IJ}\wedge
                         \mathsf P_\gamma(K\wedge K)^{IJ}.
 \label{x_reciprocidad:vii:eq:forma-cuadratica-contorsion}
\end{equation}
La densité \(\mathcal Q_{e,\gamma}\) est une forme différentielle de
degré quatre, homogène de degré deux en \(K\).

\begin{theorem}[Réponse effective quadratique en le courant]
\label{x_reciprocidad:vii:thm:respuesta-cuadratica}
Sous \eqref{x_reciprocidad:vii:eq:materia-afin-contorsion}, la substitution de la
solution algébrique \(K_*\) produit la densité effective de spin
\begin{equation}
 \mathcal L_{\rm spin}^{\rm eff}
 =-\tfrac14 K_*^{IJ}\wedge\sigma_{IJ}
 =-\mathcal Q_{e,\gamma}(K_*).
 \label{x_reciprocidad:vii:eq:accion-efectiva-spin}
\end{equation}
À cotétrade et couplages fixés, elle est homogène de degré deux en
le courant. L'action effective conserve en outre
\(S_\partial[e,K_*]\), ou sa complétion de bord déclarée.
\end{theorem}
\begin{proof}
D'après le corollaire~\ref{x_reciprocidad:vii:cor:desacoplamiento-torsional},
\(K_*\) est linéaire en \(\sigma\). L'équation de connexion est
précisément la stationnarité de
\eqref{x_reciprocidad:vii:eq:forma-cuadratica-contorsion}. Comme celle-ci est algébrique,
on applique ponctuellement la variation radiale \(\delta K=K_*\) ;
elle peut aussi être multipliée par une fonction à support compact.
L'identité d'Euler pour une forme quadratique donne
\[
 2\mathcal Q_{e,\gamma}(K_*)
             -\tfrac12K_*^{IJ}\wedge\sigma_{IJ}=0.
\]
La substitution dans la densité d'action prouve \eqref{x_reciprocidad:vii:eq:accion-efectiva-spin}. Si \(\sigma\mapsto\lambda\sigma\), alors \(K_*\mapsto\lambda K_*\) et la densité effective est multipliée par \(\lambda^2\). Le développement \eqref{x_reciprocidad:vii:eq:curvatura-contorsion} identifie séparément le terme de bord qui accompagne cette élimination.
\end{proof}

Le courant concret, sa contraction lorentzienne, la dualité et les
conventions de l'action matérielle déterminent le signe et le coefficient
de cette contribution. L'identité quadratique fournit la règle de
composition qui permet de les évaluer une fois ces données spécifiées.

\section{Source métrique effective et bilan covariant}
\label{x_reciprocidad:vii:subsec:fuente-metrica-efectiva}

La normalisation de la source métrique est fixée dans la même droite d'action. Après élimination de \(K\), nous définissons
\begin{align}
 \delta_g S_{\rm m}[e,\mathring\omega,\Psi]
 &=-\tfrac12\int_M\operatorname{vol}_e\,
                   \mathcal T^{S,\rm LC}_{\mu\nu}\,\delta g^{\mu\nu},
 \\
 \delta_g S_{\rm spin}^{\rm eff}
 &=-\tfrac12\int_M\operatorname{vol}_e\,
                   \mathcal U^{S,\rm spin}_{\mu\nu}\,\delta g^{\mu\nu}.
 \label{x_reciprocidad:vii:eq:fuentes-normalizadas-accion}
\end{align}
On considère ici des variations intérieures, avec les termes de bord
traités conformément au problème variationnel fixé. Pour le couplage
affine, la seconde source est la variation métrique de
\eqref{x_reciprocidad:vii:eq:accion-efectiva-spin}. La variation inclut la dépendance
du courant matériel à l'égard de \(e\) et de \(\Psi\).

\begin{proposition}[Équation métrique et conservation de la source totale]
\label{x_reciprocidad:vii:prop:fuente-metrica-balance}
Avec les conventions précédentes, les équations métriques sont
\begin{equation}
 G_{\mu\nu}[g]+\Lambda_{\rm int}g_{\mu\nu}
 =\kappa c\bigl(\mathcal T^{S,\rm LC}_{\mu\nu}
                    +\mathcal U^{S,\rm spin}_{\mu\nu}\bigr)
 =\kappa\bigl(T^{\rm phys,LC}_{\mu\nu}
                    +U^{\rm phys,spin}_{\mu\nu}\bigr),
 \label{x_reciprocidad:vii:eq:ecuacion-metrica-efectiva}
\end{equation}
où \(T^{\rm phys,LC}=c\mathcal T^{S,\rm LC}\) et \(U^{\rm phys,spin}=c\mathcal U^{S,\rm spin}\) lorsque la cotétrade prend ses valeurs dans la droite de longueur. Toute solution satisfait
\begin{equation}
 \mathring\nabla^\mu
   (T^{\rm phys,LC}_{\mu\nu}+U^{\rm phys,spin}_{\mu\nu})=0.
 \label{x_reciprocidad:vii:eq:balance-metrico-total}
\end{equation}
\end{proposition}
\begin{proof}
L'élimination de \(K\) conserve les équations des variables
restantes : dans la règle de dérivation en chaîne, le terme qui multiplie
\(\delta K_*\) s'annule par l'équation de connexion. Dans le secteur
de Levi--Civita, la première identité de Bianchi annule le terme
\(\Sigma_{IJ}\wedge F^{IJ}\), et la partie gravitationnelle s'écrit
\((2\kappa c)^{-1}\int_M(R-2\Lambda_{\rm int})\operatorname{vol}_e\).
Sa variation intérieure est
\[
 \frac1{2\kappa c}\int_M\operatorname{vol}_e\,
        (G_{\mu\nu}+\Lambda_{\rm int}g_{\mu\nu})\delta g^{\mu\nu}.
\]
L'addition avec \eqref{x_reciprocidad:vii:eq:fuentes-normalizadas-accion} donne \eqref{x_reciprocidad:vii:eq:ecuacion-metrica-efectiva}. L'identité de Bianchi contractée pour \(\mathring\nabla\), la compatibilité métrique et la constance de \(\kappa\) et \(\Lambda_{\rm int}\) dans la carte produisent \eqref{x_reciprocidad:vii:eq:balance-metrico-total}.
\end{proof}

Avant d'éliminer la connexion, les identités géométriques conservent
la forme
\begin{equation}
 D_\omega T^I=F^I{}_{J}\wedge e^J,\qquad D_\omega F^{IJ}=0.
 \label{x_reciprocidad:vii:eq:bianchi-cartan}
\end{equation}
La première est \(D_\omega^2e=F\wedge e\) ; la seconde s'obtient
en développant \(D_\omega(\mathrm d\omega+\omega\wedge\omega)\)
et en annulant les termes. Dans une action matérielle covariante, le
bilan métrique total coïncide avec l'identité de Noether sur les
équations matérielles. Géométrie et matière conservent ainsi une source
conjointe ; la répartition entre sa contribution de Levi--Civita et sa
contribution de spin est déterminée par l'action spécialisée.

% END OWNER

\end{HMTCompactSection}

% BEGIN OWNER /Users/ruben/Documents/New project/output/EXTREMA_MEDIA_RAZON_RECIPROCIDAD_ES_EN_20260918/ARTICULO/ES/source/antecedentes_vii/32_eliminacion_torsional_no_lineal.tex
% FORMALIZACION_NUEVA / CERTIFICADO_NUEVO, 10-09-2026.
% Composición de las operaciones efectivas ya reunidas en 30e y 31.
% No modifica la selección de rutas, la frontera, los acoplamientos ni
% la normalización de la acción; no evalúa una amplitud cosmológica s_0.
% Propietarios y alcance: desarrollo/COMPOSICION_METRICA_REV06.md.
% Control independiente: pruebas/verificar_eliminacion_torsional_no_lineal.py.

\chapter{Élimination torsionnelle et réponse métrique de l'extension minimale}
\label{x_reciprocidad:vii:sec:eliminacion-torsional-no-lineal}

L'extension minimale fournit un champ, un poids de réponse et un courant au moyen d'une seule opération variationnelle. Son incidence transportée dépend de la connexion à travers les produits de route de \eqref{x_reciprocidad:vii:eq:derivada-incidencia-rutas}. Cette dépendance est conservée lors de la composition du courant avec l'équation de Cartan--Holst. L'élimination de la connexion s'effectue donc sur l'action complète de cette extension. Le cas affine de \ref{x_reciprocidad:vii:thm:respuesta-cuadratica} est inclus comme une spécialisation exacte de la construction suivante.

\section{Composition variationnelle à un niveau cellulaire}
\label{x_reciprocidad:vii:subsec:composicion-torsional-celular}

Fixons un niveau fini, ses routes, la représentation de transport et
l'appariement cellulaire de
\ref{x_reciprocidad:vii:thm:componentes-minimo-conexion}. Soient \(p\) des coordonnées
réelles de la cotétrade et des champs matériels dans une trivialisation
locale déclarée, et \(k\in\mathbb R^d\) les coordonnées de la
contorsion \(K\). La famille \(p\mapsto e(p)\) conserve la
non-dégénérescence de la cotétrade. La connexion qui agit dans chaque route est
\begin{equation}
 \omega(p,k)=\mathring\omega(e(p))+K(k),\qquad
 F(p,k)=S_{\min}[e(p),\omega(p,k),\Psi(p)]
       =\mathfrak a(p,k)E(p,k).
 \label{x_reciprocidad:vii:eq:accion-minima-en-contorsion}
\end{equation}
Les espaces et les bases de cette expression sont transportés vers des espaces
fixes avant de dériver. On travaille dans un ouvert où l'incidence
\(\mathsf C(p,k)\) est surjective et les familles sont \(C^2\).
Cette régularité supplémentaire permet d'examiner le hessien de l'action.
Les couplages \(G_{\rm int},c_{\rm int},\gamma_{\rm HMT}\)
conservent la carte et la convention de variation de
\ref{x_reciprocidad:vii:subsec:dominio-corriente-cartan}.

Soit \(Q(p,k)\) l'intégrale, ou la réalisation cellulaire déclarée, de
la forme quadratique
\eqref{x_reciprocidad:vii:eq:forma-cuadratica-contorsion}. Dans des bases réelles,
\begin{equation}
 Q(p,k)=\tfrac12 k^{\mathsf T}\mathsf A(p)k,
 \qquad \mathsf A(p)^{\mathsf T}=\mathsf A(p).
 \label{x_reciprocidad:vii:eq:cuadratica-celular-contorsion}
\end{equation}
L'opérateur \(\mathsf A\) provient de l'action gravitationnelle et de son accouplement, tandis que le gramien \(\mathsf C\mathsf C^*\) provient du problème d'extension. La positivité du second garantit l'existence du champ minimisant ; la forme gravitationnelle conserve sa propre signature. Les termes de bord sont maintenus selon \eqref{x_reciprocidad:vii:eq:borde-contorsion} ; les formules intérieures suivantes s'appliquent avec des variations admissibles qui annulent leur contribution, ou avec la complétion de bord correspondante.

Définissons le covecteur de réponse par
\begin{equation}
 d_kF(p,k)[h]= -\tfrac12 h^{\mathsf T}s(p,k),
 \qquad s(p,k)=\mathsf M_{\rm cel}(p,k)\sigma(p,k).
 \label{x_reciprocidad:vii:eq:covector-torsional-minimo}
\end{equation}
Ici, \(s\) et \(\sigma\) conservent les types distincts de
\eqref{x_reciprocidad:vii:eq:componentes-minimo-conexion} : le premier contient les
coefficients de la différentielle et le second les coordonnées du
courant dans l'appariement matériel. Si l'on change de base dans
\(k\), cet appariement est également transporté. Lorsqu'on différencie
\(s=\mathsf M_{\rm cel}\sigma\), on inclut les dérivées des
deux facteurs lorsque l'appariement dépend du paramètre.

\begin{proposition}[Équation composée de l'extension et de la contorsion]
\label{x_reciprocidad:vii:prop:ecuacion-torsional-no-lineal}
La stationnarité intérieure de l'action par rapport à la contorsion est
\begin{equation}
 \mathcal R(p,k):=\mathsf A(p)k-\tfrac12s(p,k)=0.
 \label{x_reciprocidad:vii:eq:ecuacion-torsional-no-lineal}
\end{equation}
Chaque composante de \(s\) est calculée à partir de la même extension :
\begin{equation}
 \begin{split}
 s_j={}&4\mathfrak a\chi\operatorname{Re}
       \langle y,(\partial_{k_j}\mathsf C)f_{\min}
                         -\partial_{k_j}b\rangle\\
 &-2\mathfrak a(\partial_{k_j}\chi)\langle b,y\rangle
       -2(\partial_{k_j}\mathfrak a)E.
 \end{split}
 \label{x_reciprocidad:vii:eq:seleccion-variacional-corriente-k}
\end{equation}
Lorsque l'appariement cellulaire réalise l'équation de connexion de
\ref{x_reciprocidad:vii:thm:ecuacion-cartan-holst}, cette équation équivaut à composer
\(\sigma(p,k)\) avec l'inverse de Holst, l'inverse de
\(\mathsf L_e\) et la reconstruction de la contorsion :
\begin{equation}
 K=\mathsf C_e^{\rm tor}\,\mathsf L_e^{-1}
           \bigl(\kappa c\,\mathsf P_\gamma^{-1}
             \sigma[e,\mathring\omega(e)+K,\Psi]\bigr).
 \label{x_reciprocidad:vii:eq:composicion-implicita-cartan}
\end{equation}
Le symbole \(\mathsf C_e^{\rm tor}\) désigne exclusivement l'application \(T\mapsto K\) de \eqref{x_reciprocidad:vii:eq:contorsion-componentes}, à cotétrade fixe.
\end{proposition}

\begin{proof}
La différentielle de \(Q+F\), dans la direction arbitraire \(h\), est
\[
 d_k(Q+F)[h]
   =h^{\mathsf T}\bigl(\mathsf A k-\tfrac12s\bigr).
\]
Son annulation pour toute direction prouve \eqref{x_reciprocidad:vii:eq:ecuacion-torsional-no-lineal}. Le théorème~\ref{x_reciprocidad:vii:thm:componentes-minimo-conexion}, appliqué à \(\partial_{k_j}\), donne \eqref{x_reciprocidad:vii:eq:seleccion-variacional-corriente-k}. Avec \(p\) fixe, \(\delta\omega=\delta K\). L'équation de Cartan se transforme successivement au moyen de \eqref{x_reciprocidad:vii:eq:inversa-holst}, \eqref{x_reciprocidad:vii:eq:torsion-explicita} et \eqref{x_reciprocidad:vii:eq:contorsion-componentes} ; chacune de ces opérations est inversible dans son domaine. Cela prouve l'équivalence avec \eqref{x_reciprocidad:vii:eq:composicion-implicita-cartan} lorsque les deux expressions utilisent la même réalisation matérielle et cellulaire.
\end{proof}

La composition conserve la connexion dans les deux membres. Les routes
qui parcourent plusieurs cellules peuvent coupler les composantes du
courant entre elles. La reconstruction ponctuelle de \(T\) à partir
d'un courant donné et la résolution conjointe de
\eqref{x_reciprocidad:vii:eq:ecuacion-torsional-no-lineal} sont, dans ce cas,
des opérations successives. Cette formulation finie maintient explicite
la représentation des routes qui intervient dans la fonctionnelle.

\section{Existence locale d'une branche stationnaire}
\label{x_reciprocidad:vii:subsec:rama-torsional-local}

\begin{theorem}[Continuation locale de la solution torsionnelle]
\label{x_reciprocidad:vii:thm:rama-torsional-local}
Supposons que \((p_0,k_0)\) appartient à l'ouvert précédent,
\(\mathcal R(p_0,k_0)=0\), et que le hessien total
\begin{equation}
 \mathsf H_0
 =\mathsf A(p_0)+\partial_k^2F(p_0,k_0)
 =\mathsf A(p_0)-\tfrac12\partial_ks(p_0,k_0)
 \label{x_reciprocidad:vii:eq:hessiano-total-torsional}
\end{equation}
est inversible. Il existe des voisinages de \(p_0\) et de \(k_0\) dans lesquels
une unique branche \(C^1\), \(k=k_*(p)\), satisfait
\eqref{x_reciprocidad:vii:eq:ecuacion-torsional-no-lineal} et passe par \(k_0\).
Sur cette branche,
\begin{equation}
 \partial_{p_i}k_*
 =-\bigl[\mathsf A+\partial_k^2F\bigr]^{-1}
     \left[(\partial_{p_i}\mathsf A)k_*
                      +\partial_{p_i}\partial_kF\right].
 \label{x_reciprocidad:vii:eq:derivada-rama-torsional}
\end{equation}
\end{theorem}

\begin{proof}
L'inversion du gramien est \(C^2\) sur l'ouvert de rang maximal. Ainsi, \(F\) est \(C^2\) et \(\mathcal R\) est \(C^1\). Sa dérivée par rapport à \(k\), en \((p_0,k_0)\), est \(\mathsf H_0\). Le théorème des fonctions implicites fournit la branche, sa régularité et son unicité au voisinage du point. On peut également voir l'existence au moyen de l'application
\[
 \mathcal T_p(k)=k-\mathsf H_0^{-1}\mathcal R(p,k).
\]
Sa dérivée par rapport à \(k\) s'annule en \((p_0,k_0)\).
Dans une boule suffisamment petite, sa norme est au plus
\(q<1\). En réduisant le voisinage de \(p_0\), on obtient
\(\|\mathcal T_p(k_0)-k_0\|<(1-q)r\), où \(r\) est le
rayon de la boule. Ainsi, \(\mathcal T_p\) envoie la boule fermée
dans elle-même et est contractante. Son point fixe est la solution locale.
Différentier \(\mathcal R(p,k_*(p))=0\) et isoler
\(\partial_{p_i}k_*\) donne
\eqref{x_reciprocidad:vii:eq:derivada-rama-torsional}.
\end{proof}

L'hypothèse se vérifie sur l'action composée : le rang du
gramien assure le minimum d'écran et le hessien total contrôle
la branche torsionnelle. L'unicité affirmée est locale autour de la
solution spécifiée ; d'autres branches peuvent coexister dans le même
domaine de rang plein.

\section{Action effective et contribution non affine}
\label{x_reciprocidad:vii:subsec:accion-efectiva-no-afin}

\begin{theorem}[Identité exacte d'élimination]
\label{x_reciprocidad:vii:thm:eliminacion-torsional-no-lineal}
Soit \(k_*(p)\) une branche stationnaire. Supposons en outre que le
segment \(t\mapsto (p,tk_*)\), \(0\leq t\leq1\), demeure
dans l'ouvert où \(F\) est défini. La correction d'action
par rapport à la réalisation de Levi--Civita est
\begin{equation}
 \begin{split}
 \Delta S(p)
  &:=Q(p,k_*)+F(p,k_*)-F(p,0)\\
  &=\tfrac14 k_*^{\mathsf T}s(p,k_*)
     -\tfrac12\int_0^1
           k_*^{\mathsf T}s(p,tk_*)\,dt.
 \end{split}
 \label{x_reciprocidad:vii:eq:accion-efectiva-no-lineal}
\end{equation}
De manière équivalente,
\begin{equation}
 \Delta S=-\tfrac14 k_*^{\mathsf T}s(p,k_*)
  +\tfrac12\int_0^1 k_*^{\mathsf T}
       \bigl[s(p,k_*)-s(p,tk_*)\bigr]\,dt.
 \label{x_reciprocidad:vii:eq:correccion-no-afin-integrada}
\end{equation}
Ces identités se rapportent à l'action intégrée ou cellulaire. Dans une réalisation à densité locale, elles s'appliquent à la densité avec son accouplement de formes correspondant. L'action effective conserve séparément le terme de bord déclaré.
\end{theorem}

\begin{proof}
Par la règle de la chaîne et
\eqref{x_reciprocidad:vii:eq:covector-torsional-minimo},
\[
 \frac{d}{dt}F(p,tk_*)
      =-\tfrac12 k_*^{\mathsf T}s(p,tk_*).
\]
Intégrer entre \(0\) et \(1\) donne la différence des actions
matérielles. L'équation de stationnarité, évaluée dans la
direction \(k_*\), et l'homogénéité quadratique de \(Q\) donnent
\[
 2Q(p,k_*)=\tfrac12 k_*^{\mathsf T}s(p,k_*).
\]
La somme des deux identités démontre
\eqref{x_reciprocidad:vii:eq:accion-efectiva-no-lineal} ; ajouter et soustraire
\(\tfrac12k_*^{\mathsf T}s(p,k_*)\) fournit
\eqref{x_reciprocidad:vii:eq:correccion-no-afin-integrada}.
\end{proof}

\begin{corollary}[Spécialisation affine]
\label{x_reciprocidad:vii:cor:recuperacion-eliminacion-afin}
Si \(s(p,k)=s(p)\), alors
\[
 \Delta S=-\tfrac14 k_*^{\mathsf T}s(p)=-Q(p,k_*).
\]
C'est la réalisation intégrée de \eqref{x_reciprocidad:vii:eq:accion-efectiva-spin}.
\end{corollary}
\begin{proof}
La différence dans l'intégrale de
\eqref{x_reciprocidad:vii:eq:correccion-no-afin-integrada} s'annule. La seconde
égalité procède de l'identité radiale de stationnarité.
\end{proof}

\section{Différentielle métrique de l'action après élimination}
\label{x_reciprocidad:vii:subsec:diferencial-metrico-no-lineal}

\begin{theorem}[Réponse des variables géométriques]
\label{x_reciprocidad:vii:thm:envolvente-metrica}
Sur la branche du théorème~\ref{x_reciprocidad:vii:thm:rama-torsional-local},
\begin{equation}
 \partial_{p_i}\Delta S
 =\tfrac12 k_*^{\mathsf T}(\partial_{p_i}\mathsf A)k_*
  +(\partial_{p_i}F)(p,k_*)-(\partial_{p_i}F)(p,0),
 \label{x_reciprocidad:vii:eq:envolvente-metrica-no-lineal}
\end{equation}
où les dérivées partielles maintiennent fixes les coordonnées \(k\).
Dans des produits hermitiens fixes, chaque dérivée matérielle se calcule par
\begin{equation}
 \begin{split}
 \partial_{p_i}F={}&(\partial_{p_i}\mathfrak a)E
       +\mathfrak a(\partial_{p_i}\chi)\langle b,y\rangle\\
   &+2\mathfrak a\chi\operatorname{Re}
       \langle y,\partial_{p_i}b
                    -(\partial_{p_i}\mathsf C)f_{\min}\rangle.
 \end{split}
 \label{x_reciprocidad:vii:eq:diferencial-geometrico-minimo}
\end{equation}
Si \(p_i\) fait varier la cotétrade, \(\partial_{p_i}\mathsf C\) inclut la variation de \(\mathring\omega(e(p))+K(k)\), des routes réalisées et des trivialisations utilisées, selon leurs dépendances déclarées.
\end{theorem}

\begin{proof}
La dérivée totale de \(Q(p,k_*)+F(p,k_*)\) contient, outre
ses dérivées partielles, le terme
\[
 \bigl[\partial_kQ(p,k_*)+\partial_kF(p,k_*)\bigr]
                      \partial_{p_i}k_*.
\]
Le crochet est nul par l'équation de connexion. Soustraire la dérivée
de \(F(p,0)\) prouve
\eqref{x_reciprocidad:vii:eq:envolvente-metrica-no-lineal}.
La règle du produit et
\eqref{x_reciprocidad:vii:eq:diferencial-minimo-completo} prouvent
\eqref{x_reciprocidad:vii:eq:diferencial-geometrico-minimo}. La composition
\eqref{x_reciprocidad:vii:eq:accion-minima-en-contorsion} impose de différencier
la connexion de Levi--Civita lorsqu'on fait varier \(e\) à \(k\) fixé.
\end{proof}

Lorsque le produit positif d'extension varie également, sa
contribution est conservée explicitement. En coordonnées, soit
\(\mathsf h(p,k)>0\) sa matrice hermitienne et considérons le minimum
de \(f^*\mathsf h f\) sous \(\mathsf C f=b\). Alors
\begin{equation}
 \mathsf L=\mathsf C\mathsf h^{-1}\mathsf C^*,\quad
 y=\mathsf L^{-1}b,\quad
 f_{\min}=\mathsf h^{-1}\mathsf C^*y,\quad
 E=\chi b^*y.
 \label{x_reciprocidad:vii:eq:minimo-producto-variable}
\end{equation}
Ici, \(y\) représente le covecteur de la contrainte dans des bases fixes. La matrice positive \(\mathsf h\) conserve son rôle énergétique, distinct de celui de la métrique lorentzienne.

\begin{lemma}[Variation du produit d'extension]
\label{x_reciprocidad:vii:lem:producto-variable-extension}
Avec les conventions précédentes,
\begin{equation}
 \delta E=(\delta\chi)b^*y
   +2\chi\operatorname{Re}\bigl[y^*(\delta b-\delta\mathsf C f_{\min})\bigr]
   +\chi f_{\min}^*(\delta\mathsf h)f_{\min}.
 \label{x_reciprocidad:vii:eq:diferencial-producto-variable}
\end{equation}
\end{lemma}
\begin{proof}
La formule \(\delta\mathsf h^{-1}
=-\mathsf h^{-1}(\delta\mathsf h)\mathsf h^{-1}\) donne
\[
 \delta\mathsf L=(\delta\mathsf C)\mathsf h^{-1}\mathsf C^*
 +\mathsf C\mathsf h^{-1}(\delta\mathsf C)^*
 -\mathsf C\mathsf h^{-1}(\delta\mathsf h)
                              \mathsf h^{-1}\mathsf C^*.
\]
Lors de la différentiation de \(E=\chi b^*\mathsf L^{-1}b\), le terme \(-\chi y^*(\delta\mathsf L)y\) produit les deux termes avec \(\delta\mathsf C\) et le terme positif avec \(\delta\mathsf h\). Cela prouve l'identité complète.
\end{proof}

Dans une réalisation matérielle covariante dont la variation admet l'
appariement métrique de
\eqref{x_reciprocidad:vii:eq:fuentes-normalizadas-accion}, la correction détermine
\begin{equation}
 \delta_g\Delta S
 =-\tfrac12\int_M\operatorname{vol}_e\,
          \mathcal U^{S,\min}_{\mu\nu}\,\delta g^{\mu\nu},
 \qquad U^{\rm phys,\min}_{\mu\nu}
                      =c\mathcal U^{S,\min}_{\mu\nu}.
 \label{x_reciprocidad:vii:eq:fuente-metrica-minimo-no-lineal}
\end{equation}
Les équations \eqref{x_reciprocidad:vii:eq:envolvente-metrica-no-lineal} et
\eqref{x_reciprocidad:vii:eq:diferencial-geometrico-minimo}, complétées par
\eqref{x_reciprocidad:vii:eq:diferencial-producto-variable} le cas échéant,
évaluent leurs coefficients dans les variations métriques déclarées.
Le minimum d'écran, la solution torsionnelle et la dérivée métrique
sont trois opérations consécutives sur la même action.
L'action réduite conserve les éventuels couplages entre
cellules présents dans les routes originales.

La spécialisation homogène de
\ref{x_reciprocidad:vii:sec:dilucion-rebote} utilise en outre la réalisation du
tenseur obtenu : sa symétrie spatiale, le signe de la contraction
avec l'observateur et la loi d'échelle. Lorsque cette évaluation donne
\(\rho_{\rm corr}(a)=-s_0a^{-6}\), l'amplitude se lit dans l'
état de référence sous la forme
\(s_0=-a_{\rm ref}^{6}\rho_{\rm corr}(a_{\rm ref})\).
Cette lecture conserve l'état matériel et sa normalisation ; le
théorème d'élimination fournit l'action à contracter.

\section{Modèle exact de vérification non affine}
\label{x_reciprocidad:vii:subsec:control-no-afin}

Le modèle fini suivant vérifie les coefficients et la dépendance non affine des identités. Ses paramètres appartiennent à l'exemple algébrique et maintiennent séparé le problème de réalisation matérielle des routes.

\begin{proposition}[Élimination exacte d'une extension de rang un]
\label{x_reciprocidad:vii:prop:ejemplo-no-afin}
Soient \(H=\mathbb R^2\), \(Q=\mathbb R\), \(\mathsf C(k)=(1\ \ k)\), \(b=1\), \(\chi=1\), \(\mathfrak a=\lambda>0\) et \(Q_{\rm ex}(k)=k^2/2\). Alors
\begin{equation}
 f_{\min}=\frac{(1,k)^{\mathsf T}}{1+k^2},\quad
 F(\lambda,k)=\frac{\lambda}{1+k^2},\quad
 s(\lambda,k)=\frac{4\lambda k}{(1+k^2)^2},\quad
 \mathcal R(\lambda,k)
    =k\left(1-\frac{2\lambda}{(1+k^2)^2}\right).
 \label{x_reciprocidad:vii:eq:ejemplo-minimo-no-afin}
\end{equation}
Pour \(\lambda=25/2\), les solutions stationnaires sont
\(0,2,-2\). En \(k_*=2\),
\begin{equation}
 \mathsf H_*=\frac{16}{5},\quad s_*=4,\quad
 Q_{\rm ex}(k_*)=2,\quad F(\lambda,k_*)=\frac52,
 \quad \Delta S=-8.
 \label{x_reciprocidad:vii:eq:valores-ejemplo-no-afin}
\end{equation}
L'expression affine \(-k_*s_*/4=-2\) diffère de la correction
exacte. Sur la branche positive,
\begin{equation}
 \partial_\lambda\Delta S=-\frac45,
 \qquad \partial_\lambda k_* =\frac1{20}
 \quad\text{dans }\lambda=25/2.
 \label{x_reciprocidad:vii:eq:derivadas-ejemplo-no-afin}
\end{equation}
\end{proposition}

\begin{proof}
Le gramien est \(1+k^2>0\), de sorte que le minimum possède les expressions indiquées. Différentier \(F\) donne \(s=-2\partial_kF\), puis \(\mathcal R=k-s/2\). Pour \(\lambda=25/2\), l'équation est \(k[(1+k^2)^2-25]=0\) ; puisque \(1+k^2>0\), ses racines sont précisément \(0,\pm2\). La hessienne totale est
\[
 1-\frac{2\lambda(1-3k^2)}{(1+k^2)^3},
\]
qui vaut \(16/5\) aux deux racines non nulles. La correction est
\(2+5/2-25/2=-8\). Dans l'identité radiale,
\[
 \int_0^1 k_*s(\lambda,tk_*)\,dt
    =\int_0^1\frac{200t}{(1+4t^2)^2}\,dt=20,
\]
car une primitive est \(-25/(1+4t^2)\).
Ainsi, \(\Delta S=8/4-20/2=-8\), y compris le terme non affine.
Pour \(\lambda>1/2\), la branche positive satisfait
\(k_*^2=\sqrt{2\lambda}-1\) et
\(\Delta S=\sqrt{2\lambda}-1/2-\lambda\).
Leurs dérivées donnent \eqref{x_reciprocidad:vii:eq:derivadas-ejemplo-no-afin}.
La formule d'enveloppe fournit le même résultat :
\((\partial_\lambda F)(\lambda,k_*)
 -(\partial_\lambda F)(\lambda,0)=1/5-1=-4/5\).
\end{proof}

% END OWNER


% BEGIN OWNER /Users/ruben/Documents/New project/output/EXTREMA_MEDIA_RAZON_RECIPROCIDAD_ES_EN_20260918/ARTICULO/ES/source/antecedentes_vii/33_corriente_espinorial_y_densidad.tex
% Artículo VII. Incorporación aditiva autorizada, 10-09-2026.
% RESULTADO_RECUPERADO: representación real de Clifford y CAR:
% TEORIA_HOLOGRAFICA_INTEGRAL_MASA_HMT_MD_REV2_2026-08-20/fuente/modulos/
% 18_cuantica_estadistica_y_termodinamica.md, líneas 131--181 y 218--255;
% acción de primer orden: 14_topologia_exoticos_y_gravedad.md, 459--517;
% distinción corriente media/segundo momento y dilución:
% HMT_OBRAS/OBRA_II/chapters/11_cartan_bianchi_cosmologia.tex, 201--230.
% PDF REV2: b56bb38d10b9a4f4d0de47b7da8010fb31a6e8e8684a3bfd4237e37d7a3dddfa.
% FORMALIZACION_NUEVA: acción espinorial afín explícita, contracción
% normalizada mediante 31, observable cuártico CAR y funcional de amplitud.
% CERTIFICADO_NUEVO: cálculo racional de la forma cuadrática en las cuatro
% componentes de una trivector y evaluación de la realización CAR finita.
% Convención explicitada: las matrices g tienen firma (-+++); A_D=-i g^0
% y la acción cinética simétrica lleva hbar/2. La combinación equivalente
% i hbar Gamma^I usa Gamma^I=-i g^I y firma opuesta. No se mezclan ambas.
% Dependencias: cotetrada/conexión de 30d; constantes internas heredadas;
% corriente e inversas de 31. Estado material, celda y normal orden son
% datos de la realización declarada, no constantes universales extraídas.
% La acción afín de este fragmento y la acción de extensión mínima de 30e
% conservan sus problemas variacionales propios. No se identifican sus
% corrientes. No se cambia 50: su dominio de amplitud positiva se concreta.

\chapter{Courant spinoriel et amplitude torsionnelle de l'état matériel}
\label{x_reciprocidad:vii:sec:corriente-espinorial-densidad}

Le développement précédent et la réalisation spinorielle qui suit articulent
deux problèmes variationnels ayant des antécédents géométriques communs. L'
action d'extension minimale \(S_{\min}=\mathfrak a_m E\), construite dans
la section~\ref{x_reciprocidad:vii:sec:corriente-extension-minima}, incorpore la
dépendance de l'incidence transportée et de son minimum à l'égard de la
connexion ; son élimination torsionnelle conserve cette dépendance complète dans
la section~\ref{x_reciprocidad:vii:sec:eliminacion-torsional-no-lineal}. L'action
spinorielle \(S_D\) de \eqref{x_reciprocidad:vii:eq:accion-dirac-material} couple la
même connexion aux champs de sa représentation et est affine en la
contorsion lorsque la cotétrade et les champs sont fixés. Chaque courant s'obtient
en faisant varier sa propre action. Une préparation minimisante peut déterminer
l'état matériel initial de la réalisation spinorielle ; la substitution
de champs dépendant de la connexion conserve en outre les termes
de la différentielle composée. Cette distinction situe le calcul d'amplitude
suivant dans son action et dans son domaine variationnel spécifiques.

La représentation spinorielle du transport et la réponse algébrique de Cartan--Holst permettent de calculer l'intensité torsionnelle comme une fonctionnelle de l'état matériel. Son évaluation conserve trois données : la représentation des champs, le second moment de leur courant et le volume par rapport auquel la densité est définie. Le coefficient de couplage résulte de l'action et de son inversion ; l'état initial détermine l'amplitude matérielle.

\section{Représentation spinorielle et action matérielle affine}
\label{x_reciprocidad:vii:subsec:accion-espinorial-afin}

Nous conservons la cotétrade, l'orientation et les couplages de la
section~\ref{x_reciprocidad:vii:sec:corriente-respuesta-cartan}. La réalisation
spinorielle utilise le module complexe de dimension quatre obtenu en
complexifiant les matrices réelles
\begin{equation}
 \begin{gathered}
 J=\begin{pmatrix}0&-1\\1&0\end{pmatrix},\qquad
 R=\begin{pmatrix}0&1\\1&0\end{pmatrix},\qquad
 Z=\begin{pmatrix}1&0\\0&-1\end{pmatrix},\\
 g^0=J\otimes I_2,\qquad g^1=R\otimes I_2,\qquad
 g^2=Z\otimes R,\qquad g^3=Z\otimes Z.
 \end{gathered}
 \label{x_reciprocidad:vii:eq:clifford-material-explicito}
\end{equation}
La multiplication donne
\(\{g^I,g^J\}=2\eta^{IJ}I_4\), avec
\(\eta=\operatorname{diag}(-1,1,1,1)\). Les matrices se distinguent
du paramètre scalaire \(\gamma\) de l'opérateur de Holst. Nous définissons
\begin{equation}
 g_I=\eta_{IJ}g^J,\quad
 \mathbb B_{IJ}=\tfrac14[g_I,g_J],\quad
 A_D=-i g^0,\quad \bar\psi=\psi^\dagger A_D,\quad
 \Omega=\tfrac12\omega^{IJ}\mathbb B_{IJ}.
 \label{x_reciprocidad:vii:eq:adjunto-conexion-espinorial}
\end{equation}
Le produit de Clifford vérifie
\[
 [\mathbb B_{IJ},g^K]=\delta_J^K g_I-\delta_I^K g_J.
\]
Ainsi, la même connexion métrique agit sur les champs par
\(D\psi=\mathrm d\psi+\Omega\psi\) et
\(D\bar\psi=\mathrm d\bar\psi-\bar\psi\Omega\).
Le transport sur une route s'obtient en intégrant cette connexion dans
sa représentation spinorielle ; il conserve la connexion et la route de la
réalisation géométrique précédente.

Soit \(\eta_I=\iota_{E_I}\operatorname{vol}_e\), de sorte que \(e^J\wedge\eta_I=\delta_I^J\operatorname{vol}_e\). Pour un champ de type dimensionnel \([\psi]=L^{-3/2}\), nous prenons l'action matérielle minimale
\begin{equation}
 S_D=\int_M\left\{
 \frac{\hbar}{2}
 \bigl[\bar\psi g^I D\psi-(D\bar\psi)g^I\psi\bigr]\wedge\eta_I
 -mc\,\bar\psi\psi\operatorname{vol}_e\right\}.
 \label{x_reciprocidad:vii:eq:accion-dirac-material}
\end{equation}
Ici, \(\hbar=\hbar_{\rm int}\), \(c=c_{\rm int}\) et la masse \(m\) est la lecture de l'espèce matérielle dans la carte choisie. Ces valeurs sont des sorties antérieures à cette réalisation. La matrice \(A_D\) est hermitienne et \(A_Dg^I\) est antihermitienne. Le terme cinétique symétrique possède donc la convention réelle correspondant à \eqref{x_reciprocidad:vii:eq:clifford-material-explicito}. Dans cette signature, l'équation sans torsion a pour terme principal \(\hbar g^I D_I\) ; la notation équivalente \(i\hbar\Gamma^I D_I\), avec \(\Gamma^I=-i g^I\), emploie la relation de Clifford de signature opposée. Cette précision fixe conjointement l'adjoint, la signature et le facteur imaginaire de l'action.

Les champs matériels et la cotétrade restent fixes dans la variation
de la connexion. Une préparation issue de l'extension minimale de
la section~\ref{x_reciprocidad:vii:sec:corriente-extension-minima} peut déterminer
la donnée initiale du champ ou de l'état. Chaque action conserve sa
variation propre : si l'on substitue dans une action un champ minimisant
dépendant de la connexion, on applique en outre sa différentielle composée,
ou le théorème de l'enveloppe sous ses hypothèses de stationnarité.
L'action~\eqref{x_reciprocidad:vii:eq:accion-dirac-material} réalise ici un
couplage affine en la contorsion.

\begin{proposition}[Courant spinoriel à normalisation variationnelle]
\label{x_reciprocidad:vii:prop:corriente-dirac-explicita}
Soit
\(B^{IJK}=\bar\psi g^{[I}g^Jg^{K]}\psi\), où les crochets
désignent l'antisymétrisation normalisée. Le courant défini dans
\eqref{x_reciprocidad:vii:eq:corriente-variacional} par l'action
\eqref{x_reciprocidad:vii:eq:accion-dirac-material} est
\begin{equation}
 \sigma_{JK}
 =-\frac{\hbar}{2}\,
   \bar\psi\{g^I,\mathbb B_{JK}\}\psi\,\eta_I
 =-\frac{\hbar}{2}\,B^I{}_{JK}\eta_I.
 \label{x_reciprocidad:vii:eq:corriente-dirac-trivector}
\end{equation}
Il est indépendant de la contorsion lorsque \(e\) et \(\psi\) sont
fixés. L'inversion de la section~\ref{x_reciprocidad:vii:subsec:reconstruccion-torsion}
s'applique donc à ce courant.
\end{proposition}
\begin{proof}
La variation de \(\Omega\) est \(\delta\Omega=\frac12\delta\omega^{JK}\mathbb B_{JK}\). Les deux termes cinétiques donnent respectivement le produit avec \(g^I\) à gauche et à droite. En conséquence,
\[
 \delta_\omega S_D
 =\frac{\hbar}{4}\int_M\delta\omega^{JK}\wedge\eta_I\,
                 \bar\psi\{g^I,\mathbb B_{JK}\}\psi.
\]
La comparaison avec \eqref{x_reciprocidad:vii:eq:corriente-variacional} fixe le
facteur et le signe de \eqref{x_reciprocidad:vii:eq:corriente-dirac-trivector}.
En développant l'anticommutateur et en utilisant la relation de Clifford,
les contractions de degré un s'annulent, et il reste
\(\{g^I,\mathbb B_{JK}\}=g^{[I}g_Jg_{K]}\).
L'action est linéaire en \(\Omega\), ce qui prouve la dernière
affirmation.
\end{proof}

\section{Contraction effective et coefficient du second moment}
\label{x_reciprocidad:vii:subsec:contraccion-espinorial}

La forme quadratique lorentzienne du trivecteur s'écrit
\begin{equation}
 q(B)=\frac1{3!}B^{IJK}B_{IJK}
 =-(B^{012})^2-(B^{013})^2-(B^{023})^2+(B^{123})^2.
 \label{x_reciprocidad:vii:eq:cuadratica-trivector}
\end{equation}
Cette contraction conserve la signature du module spinoriel. Son second moment dépend de l'état et de la réalisation des produits de champs.

\begin{theorem}[Interaction effective du courant spinoriel]
\label{x_reciprocidad:vii:thm:coeficiente-torsional-espinorial}
Avec l'action matérielle~\eqref{x_reciprocidad:vii:eq:accion-dirac-material},
l'élimination algébrique de la contorsion donne
\begin{equation}
 \mathcal L_{\rm spin}^{\rm eff}
 =C_\gamma q(B)\operatorname{vol}_e,\qquad
 C_\gamma=\frac{3\kappa c\hbar^2}{16}
                      \frac{\gamma^2}{1+\gamma^2}.
 \label{x_reciprocidad:vii:eq:coeficiente-cuadratico-espinorial}
\end{equation}
Le coefficient s'obtient à partir du courant et de l'inverse de Holst
avec les normalisations de l'action précédente.
\end{theorem}
\begin{proof}
Par \eqref{x_reciprocidad:vii:eq:accion-efectiva-spin}, il suffit d'évaluer \(-\frac14K_*^{IJ}\wedge\sigma_{IJ}\), avec
\[
 Y^{IJ}=\kappa c\frac{\gamma^2}{1+\gamma^2}
           (-\star+\gamma^{-1}I)\sigma^{IJ},\qquad
 K_*=\mathsf C_e(Y).
\]
Ici, \(\mathsf C_e\) est la composition explicite de
\eqref{x_reciprocidad:vii:eq:torsion-explicita} avec
\eqref{x_reciprocidad:vii:eq:contorsion-componentes}. La linéarité de ces deux
inverses réduit le calcul aux parties \(-\star\) et \(I\).
Pour donner la contraction complète, nous ordonnons les composantes du
trivecteur sous la forme \((012,013,023,123)\), extrayons les facteurs
\(\kappa c\hbar^2\gamma^2/(1+\gamma^2)\) et substituons
\(\sigma_{JK}=-\frac12B^I{}_{JK}\eta_I\). Les deux matrices
de la forme quadratique restante sont
\begin{equation}
 \mathcal B_{-\star}
 =\frac3{16}\operatorname{diag}(-1,-1,-1,1),\qquad
 \mathcal B_I=0_{4\times4}.
 \label{x_reciprocidad:vii:eq:matrices-contraccion-espinorial}
\end{equation}
Ces matrices s'obtiennent en utilisant uniquement \(e^I\wedge\eta_J=\delta_J^I\operatorname{vol}_e\) : pour chaque composante \(B^{abc}\), on forme \(\sigma\), on calcule \(Y\), puis \(T^I=\sum_J\iota_{E_J}Y^{IJ}
 -\frac14e^I\wedge\sum_{L,J}\iota_{E_L}\iota_{E_J}Y^{LJ}\), et enfin \(K_{IJK}=(T_{KIJ}-T_{IJK}-T_{JKI})/2\). La contraction donne les quatre entrées diagonales présentées ; les six polarisations entre composantes distinctes sont nulles pour les deux opérateurs. C'est le calcul de tous les coefficients des deux formes quadratiques. La partie \(\gamma^{-1}I\) s'annule et la partie \(-\star\) produit exactement \eqref{x_reciprocidad:vii:eq:cuadratica-trivector}, avec le coefficient de \eqref{x_reciprocidad:vii:eq:coeficiente-cuadratico-espinorial}.
\end{proof}

\section{Réalisation fermionique finie et second moment de l'état}
\label{x_reciprocidad:vii:subsec:segundo-momento-car}

Nous fixons une cellule comobile périodique, un repère orthonormé adapté à son
observateur et le mode spatial homogène des quatre composantes
spinorielles. Cette réalisation finie a pour espace d'états
\(\mathcal F=\bigoplus_{n=0}^4\Lambda^n\mathbb C^4\).
Soient \(a_r,a_r^\dagger\), \(r=1,\ldots,4\), ses opérateurs
canoniques, avec
\(\{a_r,a_s^\dagger\}=\delta_{rs}\) et
\(\{a_r,a_s\}=0\). Le vide d'occupation fixe l'ordre
normal. Pour une matrice \(M\), nous écrivons
\[
 \mathrm d\Gamma(M)=\sum_{r,s}a_r^\dagger M_{rs}a_s,
 \qquad \widehat N=\mathrm d\Gamma(I_4).
\]
La lettre \(\Gamma\) de cette notation désigne la seconde
quantification ; l'holonomie nonadique conserve sa propre notation.

Les quatre matrices des bilinéaires sont calculées directement :
\begin{equation}
 \begin{array}{c|cccc}
 abc&012&013&023&123\\ \hline
 M_{abc}=A_Dg^ag^bg^c
       &iJ\otimes R&iJ\otimes Z&iI_2\otimes J&iZ\otimes J
 \end{array}
 \label{x_reciprocidad:vii:eq:matrices-bilineales-espin}
\end{equation}
Toutes sont hermitiennes, ont une trace nulle et satisfont \(M_{abc}^2=I_4\). L'observable quartique correspondant à \eqref{x_reciprocidad:vii:eq:cuadratica-trivector}, dans cet ordre normal déclaré, est
\begin{equation}
 \begin{split}
 \widehat{\mathscr Q}_{\rm sp}
 &=\sum_{a<b<c}\eta_{aa}\eta_{bb}\eta_{cc}
       :\!\mathrm d\Gamma(M_{abc})^2\!:\\
 &=\sum_{a<b<c}\eta_{aa}\eta_{bb}\eta_{cc}
       \mathrm d\Gamma(M_{abc})^2+2\widehat N.
 \end{split}
 \label{x_reciprocidad:vii:eq:observable-cuadratico-espin}
\end{equation}
L'ordre normal fait partie de la réalisation finie spécifiée ;
sa définition rend explicite la soustraction des contractions
d'une seule occupation.

\begin{proposition}[Évaluation par les relations canoniques]
\label{x_reciprocidad:vii:prop:operador-espin-car}
L'opérateur \(\widehat{\mathscr Q}_{\rm sp}\) est autoadjoint,
conserve chaque secteur d'occupation et a pour restrictions
\begin{equation}
 \widehat{\mathscr Q}_{\rm sp}|_{\Lambda^0\oplus\Lambda^1}=0,
 \qquad
 \widehat{\mathscr Q}_{\rm sp}|_{\Lambda^3}=4I_4,
 \qquad
 \widehat{\mathscr Q}_{\rm sp}|_{\Lambda^4}=8I_1.
 \label{x_reciprocidad:vii:eq:sectores-espin-car}
\end{equation}
Dans la base de \(\Lambda^2\mathbb C^4\) dont les sous-ensembles d'occupation sont \((12,13,23,14,24,34)\), sa matrice est
\begin{equation}
 \begin{pmatrix}
 0&0&0&0&0&-4\\
 0&2&0&0&6&0\\
 0&0&2&-6&0&0\\
 0&0&-6&2&0&0\\
 0&6&0&0&2&0\\
 -4&0&0&0&0&0
 \end{pmatrix}.
 \label{x_reciprocidad:vii:eq:sector-dos-espin-car}
\end{equation}
\end{proposition}
\begin{proof}
La relation \(a_s a_t^\dagger=\delta_{st}-a_t^\dagger a_s\)
appliquée à \(\mathrm d\Gamma(M)^2\) sépare sa contraction à
une occupation et donne
\[
 :\!\mathrm d\Gamma(M)^2\!:
 =\mathrm d\Gamma(M)^2-\mathrm d\Gamma(M^2).
\]
Comme les trois premiers signes de
\eqref{x_reciprocidad:vii:eq:matrices-bilineales-espin} sont négatifs et le dernier
positif, la somme des contractions est \(-2\widehat N\).
Cela prouve \eqref{x_reciprocidad:vii:eq:observable-cuadratico-espin}.
L'autoadjonction et la conservation de l'occupation découlent de
ces mêmes expressions.

Dans \(\Lambda^0\), tous les termes s'annulent. Dans \(\Lambda^1\), \(\mathrm d\Gamma(M)=M\), et les carrés ont pour somme \(-2I_4\), qui s'annule avec \(2\widehat N\). Dans \(\Lambda^4\), \(\mathrm d\Gamma(M)=\operatorname{tr}M=0\), ce qui laisse huit. L'identification par complément entre \(\Lambda^3\mathbb C^4\) et le dual de \(\mathbb C^4\) envoie \(\mathrm d\Gamma(M)\) sur \((\operatorname{tr}M)I-M^{\mathsf T}=-M^{\mathsf T}\). Leurs carrés ont à nouveau pour somme \(-2I_4\), tandis que \(2\widehat N=6I_4\) ; la restriction est donc \(4I_4\). Enfin, l'application des quatre matrices explicites aux six produits extérieurs de la base indiquée produit \eqref{x_reciprocidad:vii:eq:sector-dos-espin-car}. L'opérateur est ainsi évalué dans les seize dimensions de \(\mathcal F\).
\end{proof}

Pour un opérateur d'état \(\varrho\geq0\),
\(\operatorname{Tr}\varrho=1\), nous définissons
\(q_\varrho=\operatorname{Tr}(\varrho\widehat{\mathscr Q}_{\rm sp})\).
Sa dépendance à l'égard des moments d'ordre deux est explicite :
\begin{equation}
 q_\varrho=\sum_{a<b<c}\eta_{aa}\eta_{bb}\eta_{cc}
 \left\{
  \langle\widehat B_{abc}\rangle_\varrho^2
  +\operatorname{Var}_\varrho(\widehat B_{abc})
  -\langle\widehat N\rangle_\varrho
 \right\},\qquad
 \widehat B_{abc}=\mathrm d\Gamma(M_{abc}).
 \label{x_reciprocidad:vii:eq:covarianza-espin-material}
\end{equation}
La covariance et la soustraction normale sont conservées même lorsque
les courants moyens s'annulent. En particulier, si \(P_3\) est la
projection sur \(\Lambda^3\mathbb C^4\),
\begin{equation}
 \varrho_3=\tfrac14P_3,\qquad
 \langle\widehat B_{abc}\rangle_{\varrho_3}=0,\qquad
 q_{\varrho_3}=4>0.
 \label{x_reciprocidad:vii:eq:testigo-positivo-espin}
\end{equation}
C'est un état positif de trace un qui constitue un témoin explicite
du domaine positif de la fonctionnelle. Le secteur de deux occupations
contient également des valeurs négatives ; par exemple, le vecteur
\((|12\rangle+|34\rangle)/\sqrt2\) a pour valeur \(-4\).
Le signe de l'intensité procède ainsi de la contraction et de
l'état, et reste distinct de la sélection des conditions cosmologiques.

\section{Densité, pression et conservation de l'amplitude comobile}
\label{x_reciprocidad:vii:subsec:amplitud-espinorial-comovil}

Dans la réalisation spatialement homogène, nous fixons
\[
 e^0=N(t)c\,\mathrm dt,\qquad e^j=a(t)\,\mathrm dx^j,
 \qquad \mathcal V(t)=a(t)^3V_c,
\]
avec une fonction lapse \(N(t)>0\), un facteur d'échelle sans dimension \(a(t)>0\)
et un volume comobile \(V_c>0\). Le cadre et le mode périodique de la
sous-section~\ref{x_reciprocidad:vii:subsec:segundo-momento-car} restent déclarés.
La normalisation du champ homogène est
\(\widehat\psi=\mathcal V^{-1/2}(a_1,a_2,a_3,a_4)^{\mathsf T}\).
Chaque bilinéaire porte un facteur \(\mathcal V^{-1}\), et son produit
normalement ordonné porte \(\mathcal V^{-2}\). La forme temporelle
cinétique symétrique conserve cette normalisation canonique : les termes
provenant de la dérivation de \(\mathcal V^{-1/2}\) s'annulent entre
ses deux termes.

\begin{theorem}[Amplitude torsionnelle comme fonctionnelle de l'état]
\label{x_reciprocidad:vii:thm:amplitud-material-explicita}
En maintenant fixes l'état canonique et la réalisation du produit
pendant les variations métriques, le terme effectif de spin a pour
action réduite
\begin{equation}
 S_{\rm spin}^{\rm hom}
 =\int\!\mathrm dt\,N(t)
       \frac{cC_\gamma}{a(t)^3V_c}\,q_\varrho.
 \label{x_reciprocidad:vii:eq:accion-homogenea-espin}
\end{equation}
Sa densité d'énergie et sa pression isotrope sont
\begin{equation}
 \rho_{\rm spin}=p_{\rm spin}
 =-\frac{cC_\gamma}{V_c^2}\,q_\varrho a^{-6}.
 \label{x_reciprocidad:vii:eq:densidad-presion-espin}
\end{equation}
Dans toute évolution qui conserve \(q_\varrho\), l'amplitude constante correspondante est
\begin{equation}
 \boxed{\displaystyle
 s_0[\varrho,V_c]
 =\frac{3\kappa c^2\hbar^2}{16}
    \frac{\gamma^2}{1+\gamma^2}
    \frac{\operatorname{Tr}
       (\varrho\widehat{\mathscr Q}_{\rm sp})}{V_c^2}.}
 \label{x_reciprocidad:vii:eq:amplitud-s0-espinorial}
\end{equation}
\end{theorem}
\begin{proof}
L'intégration spatiale de
\eqref{x_reciprocidad:vii:eq:coeficiente-cuadratico-espinorial} multiplie le
facteur \(\mathcal V^{-2}\) par
\(Nc\mathcal V\,\mathrm dt\), ce qui donne
\eqref{x_reciprocidad:vii:eq:accion-homogenea-espin}. En variables canoniques,
la variation du lapse définit l'énergie par
\(\delta_N S_{\rm m}=-\int E\,\delta N\,\mathrm dt\).
Par conséquent,
\[
 E_{\rm spin}=-\frac{cC_\gamma q_\varrho}{a^3V_c},
 \qquad \rho_{\rm spin}=\frac{E_{\rm spin}}{a^3V_c}.
\]
La variation spatiale isotrope satisfait \(\delta_a S_{\rm m}=\int3Na^2V_c p\,\delta a\,\mathrm dt\). Appliquée à \eqref{x_reciprocidad:vii:eq:accion-homogenea-espin}, elle donne
\[
 3Na^2V_c p_{\rm spin}
 =-\frac{3NcC_\gamma q_\varrho}{a^4V_c}.
\]
Cela prouve simultanément le signe, la pression et la densité de \eqref{x_reciprocidad:vii:eq:densidad-presion-espin}. Puisque \(q_\varrho\) est conservé, le coefficient de \(-a^{-6}\) est constant et l'on obtient \eqref{x_reciprocidad:vii:eq:amplitud-s0-espinorial}. La dimensionnalité est cohérente : \([cC_\gamma]=\text{énergie}\,L^3\), de sorte que \(s_0\) a la dimension d'une densité d'énergie.
\end{proof}

\begin{corollary}[Domaine conservé d'amplitude positive]
\label{x_reciprocidad:vii:cor:dominio-positivo-conservado}
Soit une évolution unitaire de \(\mathcal F\) qui préserve
\(\Lambda^3\mathbb C^4\). Pour tout état initial de trace un
à support dans ce secteur, on a \(q_{\varrho(t)}=4\), et
\begin{equation}
 s_0=\frac{3\kappa c^2\hbar^2}{4V_c^2}
                   \frac{\gamma^2}{1+\gamma^2}>0.
 \label{x_reciprocidad:vii:eq:amplitud-positiva-sector-tres}
\end{equation}
\end{corollary}
\begin{proof}
L'invariance du secteur conserve le support de l'état ; dans celui-ci,
\eqref{x_reciprocidad:vii:eq:sectores-espin-car} donne
\(\widehat{\mathscr Q}_{\rm sp}=4I\). La conservation de la
trace fixe alors son espérance à quatre, indépendamment de
l'évolution interne de ses cohérences. La positivité des
couplages, \(V_c>0\) et \(\gamma\ne0\) prouvent la formule
et son signe.
\end{proof}

Pour des états généraux, la conservation de
\(\langle\widehat N\rangle\) détermine la dilution de la densité
d'occupation, \(n(t)=\langle\widehat N\rangle/(a^3V_c)\).
La conservation du moment quartique est vérifiée séparément.
Si \(i\hbar\dot\varrho=[H(t),\varrho]\), alors
\begin{equation}
 \frac{\mathrm d q_\varrho}{\mathrm dt}
 =\frac{i}{\hbar}\operatorname{Tr}
       \bigl(\varrho[H(t),\widehat{\mathscr Q}_{\rm sp}]\bigr).
 \label{x_reciprocidad:vii:eq:conservacion-momento-cuartico}
\end{equation}
La commutation, l'invariance d'un secteur où l'observable est scalaire, ou une conservation démontrée pour l'état concret fournissent les conditions pertinentes. Si le moment évolue, \eqref{x_reciprocidad:vii:eq:densidad-presion-espin} conserve son évaluation instantanée ; son bilan inclut
\[
 \dot\rho_{\rm spin}+6H\rho_{\rm spin}
 =-\frac{cC_\gamma}{V_c^2}\,a^{-6}\dot q_\varrho,
 \qquad H=\dot a/a
\]
en temps propre. Ce terme enregistre l'échange avec les
autres degrés de liberté matériels.

La composition est déterminée : représentation spinorielle du
transport, action matérielle, courant variationnel, inversion de
Cartan--Holst, second moment et densité normalisée. Le domaine
\(s_0>0\) de la section~\ref{x_reciprocidad:vii:sec:dilucion-rebote} dispose ainsi
d'une fonctionnelle évaluable et d'états témoins explicites. Son utilisation dans
la dynamique cosmologique conserve en outre les conditions matérielles,
géométriques et de conservation énoncées dans cette section.

% END OWNER


% BEGIN OWNER /Users/ruben/Documents/New project/output/EXTREMA_MEDIA_RAZON_RECIPROCIDAD_ES_EN_20260918/ARTICULO/ES/source/antecedentes_vii/50_dilucion_y_rebote.tex
% Artículo VII: pruebas preexistentes incorporadas al cuerpo.
% Propietario íntegro: fuentes_conservadas/17_rebote_einstein_cartan.tex.
% Incluye la ley de escala y las pruebas de rebote y persistencia.
% La amplitud s_0 es la salida de la realización material antecedente;
% este fragmento no certifica por sí solo la evaluación de esa realización.

\chapter{Dilution de la mémoire et persistance du rebond}
\label{x_reciprocidad:vii:sec:dilucion-rebote}

L'évolution de la mémoire et l'expansion spatiale sont deux lectures d'une même suite d'états. L'indice de raffinement peut être éliminé entre les deux ; cette élimination détermine l'exposant de dilution que la réalisation cosmologique transmet à sa contribution torsionnelle. L'amplitude initiale dépend de l'évaluation du courant matériel concret dans sa réponse variationnelle. La loi d'échelle et le théorème de rebond qui suivent conservent cette dépendance.

Le théorème~\ref{x_reciprocidad:vii:thm:amplitud-material-explicita} détermine cette
amplitude comme fonctionnelle de l'état matériel et du volume comobile.
Le corollaire~\ref{x_reciprocidad:vii:cor:dominio-positivo-conservado} fournit
un domaine explicite d'états d'amplitude positive et conservée.
Dans ce domaine, l'évaluation de \(s_0\) précède l'intégration
cosmologique développée ci-dessous. La densité ordinaire
\(\rho_0>0\) et sa loi barotropique sont spécifiées en outre dans la
section~\ref{x_reciprocidad:vii:sec:umbral-rebote} : le témoin de positivité de l'
amplitude n'évalue pas par lui-même cette densité matérielle.

\section{Loi de dilution \texorpdfstring{\(a^{-6}\)}{a élevé à -6}}
\label{x_reciprocidad:vii:sec:dilucion}

\begin{proposition}[Élimination de l'indice de mémoire]
\label{x_reciprocidad:vii:prop:dilucion}
Soient \(a_{\mathrm{ref}}>0\), \(\varphi>1\) et les lecteurs de la route
\[
 \frac{a_n}{a_{\mathrm{ref}}}=\varphi^{n/3},\qquad
 M_n=\varphi^{-2n}.
\]
Alors
\[
 M_n=\left(\frac{a_n}{a_{\mathrm{ref}}}\right)^{-6}
\]
pour tout indice de la suite. Si la réalisation de la densité
quadratique est \(s_n^2=s_*^2M_n\), avec \(s_*^2>0\), on obtient
\[
 s_n^2=s_*^2\left(\frac{a_n}{a_{\mathrm{ref}}}\right)^{-6}.
\]
\end{proposition}

\begin{proof}
En élevant la première égalité à la puissance \(-6\), on obtient \(\varphi^{-2n}\), qui est exactement le second lecteur. Multiplier par l'amplitude initiale reproduit la seconde affirmation. L'exposant est déterminé par le quotient des exposants des deux lecteurs ; l'amplitude est l'évaluation de l'état de référence.
\end{proof}

Dans la réalisation homogène continue, cette loi constitutive est conservée :
la densité torsionnelle s'écrit \(s_0a^{-6}\), où la normalisation
du facteur \(a\), le coefficient matériel et l'amplitude de référence
ont été absorbés dans \(s_0\). La loi d'échelle et l'évaluation de
\(s_0\) sont des opérations successives de la construction.

\section{Seuil homogène et transversalité}
\label{x_reciprocidad:vii:sec:umbral-rebote}

Considérons la carte homogène, spatialement plate, en unités \(c=1\),
avec \(G>0\) constant, des densités positives \(\rho_0,s_0\) et un paramètre
barotrope constant \(w<1\). La réalisation spécifiée satisfait
\(\Lambda_{\mathrm{eff}}=0\) et les équations
\begin{align}
 \rho_{\mathrm m}(a)&=\rho_0a^{-3(1+w)},&
 \rho_{\mathrm s}(a)&=s_0a^{-6},\\
 H^2&=\frac{8\pi G}{3}
        \bigl(\rho_{\mathrm m}-\rho_{\mathrm s}\bigr)
        =:F_0(a),&
 \dot H&=-4\pi G
        \bigl((1+w)\rho_{\mathrm m}-2\rho_{\mathrm s}\bigr).
 \label{x_reciprocidad:vii:eq:dinamica-homogenea}
\end{align}
Les couplages appartiennent à la section dimensionnelle déjà construite.
L'expression de \(\rho_{\mathrm s}\) intègre le signe de sa contribution
dans cette réalisation gravitationnelle.

\begin{theorem}[Existence et unicité du seuil transversal]
\label{x_reciprocidad:vii:thm:rebote}
Il existe un unique \(a_{\mathrm b}>0\) tel que \(F_0(a_{\mathrm b})=0\) :
\[
 a_{\mathrm b}^{\,3(1-w)}=\frac{s_0}{\rho_0}.
\]
La région permise par \(H^2\geq0\) est \(a\geq a_{\mathrm b}\).
Au seuil, avec
\(\rho_{\mathrm b}=\rho_{\mathrm m}(a_{\mathrm b})
                  =\rho_{\mathrm s}(a_{\mathrm b})>0\),
\[
 \dot H_{\mathrm b}=4\pi G(1-w)\rho_{\mathrm b}>0.
\]
La solution qui traverse \(H=0\) présente un minimum strict de \(a\), avec le développement local
\[
 a(t)=a_{\mathrm b}\left[
 1+2\pi G(1-w)\rho_{\mathrm b}(t-t_{\mathrm b})^2
 +O((t-t_{\mathrm b})^3)\right].
\]
\end{theorem}

\begin{proof}
La factorisation
\[
 F_0(a)=\frac{8\pi G}{3}a^{-6}
             \bigl(\rho_0a^{3(1-w)}-s_0\bigr)
\]
réduit ses zéros à celui d'une fonction strictement croissante de
\((0,\infty)\) sur \((-s_0,\infty)\). Cela démontre l'existence,
l'unicité et le signe de \(F_0\) de part et d'autre du seuil.
La seconde équation de \eqref{x_reciprocidad:vii:eq:dinamica-homogenea}, évaluée
à l'égalité des densités, donne la formule de \(\dot H_{\mathrm b}\).
Comme \(\dot a=aH\), au seuil
\(\ddot a_{\mathrm b}=a_{\mathrm b}\dot H_{\mathrm b}>0\).
Les équations sont lisses pour \(a>0\), de sorte que leur développement de
Taylor donne le développement indiqué. Le changement de signe de \(H\) est
transverse : il passe de la contraction à l'expansion.
\end{proof}

Pour \(w=1\), les deux termes ont le même exposant et l'égalité des
amplitudes détermine un cas dégénéré. Pour \(w>1\), le signe de
\(\dot H\) lors d'une coïncidence des densités est négatif.
Le domaine \(w<1\) du théorème exprime exactement l'ordre entre les
deux exposants de dilution.

\section{Persistance locale sous perturbations contrôlées}
\label{x_reciprocidad:vii:sec:persistencia}

La persistance du seuil utilise deux estimations indépendantes : le contrôle du signe aux extrémités et celui de la dérivée dans l'intervalle. Soit \(I=[a_-,a_+]\subset(0,\infty)\), avec \(a_-<a_{\mathrm b}<a_+\), choisi de sorte que
\[
 F_0(a_-)<0<F_0(a_+),\qquad
 m_I:=\inf_{a\in I}F_0'(a)>0.
\]
Un tel intervalle existe parce que
\[
 F_0'(a_{\mathrm b})
   =\frac{8\pi G}{a_{\mathrm b}}(1-w)\rho_{\mathrm b}>0.
\]
Nous définissons
\(\delta_I=\min\{-F_0(a_-),F_0(a_+)\}>0\).

\begin{theorem}[Persistance locale du rebond transversal]
\label{x_reciprocidad:vii:thm:persistencia}
Soit \(R\in C^1(I)\) une perturbation qui satisfait
\[
 \|R\|_{\infty,I}<\delta_I,\qquad
 \|R'\|_{\infty,I}<m_I.
\]
Alors \(F_R=F_0+R\) possède un unique zéro \(a_R\in(a_-,a_+)\) et \(F_R'(a_R)>0\). La réalisation \(H^2=F_R(a)\) admet le rebond local transversal, avec
\[
 \dot H_R=\frac{a_R}{2}F_R'(a_R)>0
\]
au seuil. Si \(R(a,\eta)\) est conjointement \(C^1\), les zéros
persistants dépendent localement de manière \(C^1\) du paramètre \(\eta\).
\end{theorem}

\begin{proof}
La première borne conserve \(F_R(a_-)<0<F_R(a_+)\) ; le théorème des valeurs intermédiaires fournit un zéro intérieur. La seconde donne \(F_R'(a)\geq m_I-\|R'\|_{\infty,I}>0\) dans \(I\), de sorte que le zéro est unique et simple.

En \(a>a_R\) proche du zéro,
\(F_R(a)=F_R'(a_R)(a-a_R)+o(a-a_R)\).
L'intégrale
\[
 t-t_{\mathrm b}=\int_{a_R}^{a}
          \frac{d\xi}{\xi\sqrt{F_R(\xi)}}
\]
est finie et strictement croissante. Son inverse définit la branche
en expansion ; la réflexion temporelle définit la branche en contraction.
Toutes deux se rejoignent avec \(\dot a(t_{\mathrm b})=0\) et
\[
 a(t)-a_R
 =\frac{a_R^2F_R'(a_R)}4(t-t_{\mathrm b})^2
   +o((t-t_{\mathrm b})^2).
\]
Hors du seuil, la dérivation de \(H^2=F_R(a)\) donne
\(\dot H=aF_R'(a)/2\) ; sa limite continue donne l'expression positive
en \(a_R\). La régularité \(C^1\) de \(F_R\) garantit la continuité
de l'accélération à cette jonction. Enfin,
\(\partial_aF_R(a_R)>0\) permet d'appliquer le théorème des fonctions
implicites lorsque la perturbation dépend conjointement de \((a,\eta)\).
\end{proof}

La première borne contrôle l'existence ; la seconde assure l'unicité locale et la transversalité. La stabilité démontrée correspond à ce seuil dans l'intervalle sélectionné. L'amplitude torsionnelle intervient dans sa localisation, tandis que la relation entre exposants détermine l'orientation temporelle du franchissement.

% END OWNER


% BEGIN OWNER /Users/ruben/Documents/New project/output/EXTREMA_MEDIA_RAZON_RECIPROCIDAD_ES_EN_20260918/ARTICULO/ES/source/antecedentes_vii/51_solucion_homogenea.tex
% FORMALIZACION_NUEVA de la realización homogénea preexistente.
% Prioridad de la fórmula integrada: procedencia histórica pendiente de cotejo.
% Conserva la ecuación y el dominio del capítulo de rebote; no selecciona s_0.
\chapter{Solution homogène explicite et complétude causale}
\label{x_reciprocidad:vii:sec:solucion-homogenea}
La détermination de l'amplitude permet de développer une réalisation
complète du rebond, en plus de l'argument local. On conserve la carte
spatialement plate de \eqref{x_reciprocidad:vii:eq:dinamica-homogenea}, en unités
\(c=1\), avec \(\rho_0,s_0,G>0\), et on spécialise le secteur matériel
à la poussière, \(w=0\), avec \(\Lambda_{\mathrm{eff}}=0\).
Le coefficient \(s_0\) est celui obtenu à partir de l'état
matériel dans la réalisation correspondante ; l'intégration suivante
utilise ce coefficient et conserve son domaine.

\begin{theorem}[Intégration exacte du rebond de poussière]
\label{x_reciprocidad:vii:thm:rebote-polvo}
Posons
\[
 a_{\rm b}^3=\frac{s_0}{\rho_0},\qquad
 \tau^2=\frac{s_0}{6\pi G\rho_0^2},\qquad u=t-t_{\rm b}.
\]
La solution du système homogène avec
\(a(t_{\rm b})=a_{\rm b}\) et \(H(t_{\rm b})=0\) est, pour tout
\(t\in\mathbb R\),
\begin{equation}
 a(t)=a_{\rm b}\left(1+\frac{u^2}{\tau^2}\right)^{1/3},
 \qquad H(t)=\frac{2u}{3(\tau^2+u^2)},\qquad
 \dot H(t)=\frac{2(\tau^2-u^2)}{3(\tau^2+u^2)^2}.
 \label{x_reciprocidad:vii:eq:solucion-polvo}
\end{equation}
Elle possède un unique minimum, en \(t_{\rm b}\) ; elle est contractante avant et expansive après. L'évolution est lisse et satisfait \(a(t)\geq a_{\rm b}>0\).
\end{theorem}
\begin{proof}
Soit \(y=a^3\). L'équation de Friedmann donne
\[
 \dot y^{\,2}=24\pi G(\rho_0y-s_0).
\]
La fonction \(y=s_0/\rho_0+6\pi G\rho_0u^2\) vérifie cette
identité et les conditions initiales. Ses dérivées donnent les
trois expressions de \eqref{x_reciprocidad:vii:eq:solucion-polvo}. Pour vérifier
également l'équation d'évolution, écrivons
\[
 \rho_{\rm b}=\frac{\rho_0^2}{s_0},\qquad
 x=\frac{u}{\tau}.
\]
Alors
\[
 \rho_{\rm m}=\frac{\rho_{\rm b}}{1+x^2},\qquad
 \rho_{\rm s}=\frac{\rho_{\rm b}}{(1+x^2)^2},\qquad
 4\pi G\rho_{\rm b}=\frac{2}{3\tau^2}.
\]
La substitution dans \(\dot H=-4\pi G(\rho_{\rm m}-2\rho_{\rm s})\) donne exactement la troisième formule. Le système complet \((\dot a,\dot H)\) est donc satisfait, et son champ est lisse pour \(a>0\). Son unicité locale, prolongée sur la solution explicite, fixe l'évolution par ces conditions initiales. Le dénominateur est positif pour tout \(t\), et les affirmations sur le minimum et le signe découlent de \(H\).
\end{proof}

\begin{figure}[htbp]
\centering
\begin{tikzpicture}[x=1.45cm,y=1.15cm,>=Latex]
\draw[->] (-3.25,0)--(3.3,0) node[right] {\(\displaystyle u/\tau\)};
\draw[->] (0,0)--(0,2.8) node[above] {\(\displaystyle a/a_{\rm b}\)};
\draw[densely dashed,gray] (-3.1,1)--(3.1,1);
\draw[very thick,ink,domain=-3:3,samples=121,smooth]
 plot (\x,{(1+\x*\x)^(1/3)});
\fill[accent] (0,1) circle (1.6pt);
\node[anchor=north east] at (-.08,.98) {\(1\)};
\node[ink] at (-2,2.65) {Contraction};
\node[ink] at (2,2.65) {Expansion};
\foreach \t in {-3,-2,-1,1,2,3}
{\draw (\t,.045)--(\t,-.045) node[below] {\(\t\)};}
\end{tikzpicture}
\caption{Profil adimensionnel de la solution démontrée :
\(a/a_{\rm b}=(1+(u/\tau)^2)^{1/3}\).
Les paramètres \(a_{\rm b}\) et \(\tau\) se calculent à partir de
\(\rho_0,s_0,G\) ; le dessin représente toute la famille remise à l'échelle.}
\label{x_reciprocidad:vii:fig:rebote-polvo}
\end{figure}

\section{Courbure finie et extension des géodésiques}
La métrique réalisée dans cette carte est
\(g=-dt^2+a(t)^2\,d\mathbf x^2\) sur
\(\mathbb R\times\mathbb R^3\). Avec la convention dans laquelle
le scalaire de courbure du cas plat est \(R=6(\dot H+2H^2)\), la solution donne
\[
 R(t)=\frac{4\tau^2+\frac43u^2}{(\tau^2+u^2)^2}.
\]
En particulier, \(R(t_{\rm b})=4/\tau^2\) est fini.
Dans un repère orthonormé, les composantes de Riemann dépendent de
\(H^2\) et de \(\dot H+H^2\), qui sont des fonctions bornées sur
\(\mathbb R\). Il en va de même des contractions polynomiales
de ces composantes ; par exemple,
\[
 R_{\mu\nu\rho\sigma}R^{\mu\nu\rho\sigma}
 =12\bigl[(\dot H+H^2)^2+H^4\bigr].
\]

\begin{proposition}[Complétude des géodésiques causales]
\label{x_reciprocidad:vii:prop:completitud-causal-polvo}
La réalisation métrique précédente est géodésiquement complète
pour les géodésiques temporelles et nulles de sa connexion de
Levi--Civita, dans les deux directions.
\end{proposition}
\begin{proof}
Les translations spatiales fournissent les constantes \(p_i=a^2\,dx^i/d\lambda\). La normalisation causale donne
\[
 \left|\frac{dt}{d\lambda}\right|
 =\sqrt{\epsilon+\frac{|\mathbf p|^2}{a(t)^2}},
 \qquad \epsilon=1\ \text{dans le cas temporel},\quad
 \epsilon=0\ \text{dans le cas nul}.
\]
Une géodésique nulle non constante a \(|\mathbf p|>0\). Puisque \(a\geq a_{\rm b}\), dans le cas temporel
\[
 \left|\frac{d\lambda}{dt}\right|
 \geq\left(1+\frac{|\mathbf p|^2}{a_{\rm b}^2}\right)^{-1/2}>0,
\]
et dans le cas nul
\(\left|d\lambda/dt\right|=a/|\mathbf p|
 \geq a_{\rm b}/|\mathbf p|>0\).
Ainsi, atteindre \(t=\pm\infty\) exige une longueur affine infinie.
Sur un intervalle temporel compact, \(a\) et ses dérivées sont
lisses et \(a\) est séparé de zéro. Les équations explicites
précédentes et \(dx^i/d\lambda=p_i/a^2\) permettent de prolonger la
géodésique à travers ses extrémités finies. Les deux possibilités
de prolongement sont couvertes.
\end{proof}

Le résultat de complétude appartient à cette réalisation homogène de
poussière, avec ses conditions d'espace, de matière et d'amplitude. Sa
persistance locale sous perturbations du seuil conserve les bornes
du théorème~\ref{x_reciprocidad:vii:thm:persistencia} ; cette propriété locale et la
complétude globale de la solution explicite sont des conclusions distinctes.
L'échelle minimale procède de la relation entre densité matérielle et
réponse torsionnelle. La mémoire intervient ainsi dans une évolution
géométrique vérifiable, avec une réalisation régulière de la contraction
et de l'expansion.

% END OWNER


% BEGIN OWNER /Users/ruben/Documents/New project/output/EXTREMA_MEDIA_RAZON_RECIPROCIDAD_ES_EN_20260918/ARTICULO/ES/source/antecedentes_vii/60_tensor_residual_y_balance.tex
% Artículo VII. Recuperación por dependencia, 10-09-2026.
% RESULTADO_RECUPERADO: integral activo 2249, fuente/colaboracion/
% partes_iv_v/tex/residencias/86_cosmologia_cerrada.tex (S04),
% líneas 2--44 (cartas seleccionadas), 174--248 (tensor y aceleración).
% Fuente S04 leída íntegramente, SHA-256:
% feb2582c3316c8b30829cfb0476bfa2f4e95889dabfba16fe43f2d835a226673.
% Se recuperan definición, unicidad, conservación, descomposición,
% intercambio y aceleración, con sus condiciones explícitas.
% FORMALIZACION_REUNIDA: relación con la fuente efectiva de 31;
% se distingue conservación total de conservación material separada.
% CERTIFICADO_NUEVO: expansión de las pruebas de traza, balance y signos;
% cálculo local de Ricci para las tres cartas ya seleccionadas por S04.
% No se modifica el rebote de 50 ni se determina aquí una amplitud s_0.
%
% DEPENDENCIAS DE INTEGRACIÓN:
% - vii:sec:realizacion-cartan: la métrica procede de la cotetrada realizada.
% - vii:sec:acoplamientos: G_int y el reloj t_0 del mismo estado HMT;
%   c_int=1 es una elección de unidades posterior a su construcción.
% - vii:eq:ecuacion-metrica-efectiva y vii:eq:balance-metrico-total, de 31:
%   se usan únicamente para identificar el residual sobre una solución.
% - T_b es el tensor publicado por el transductor material declarado;
%   su conservación separada es una condición expresamente enunciada.
% - El selector TRIT de la carta y t_0(x)>0 son antecedentes de la
%   especialización de Sitter, no datos cosmológicos observacionales.
% El presente fragmento no certifica la clausura del conjunto de VII.

\chapter{Tenseur résiduel, décomposition de trace et bilan covariant}
\label{x_reciprocidad:vii:sec:tensor-residual-balance}

La réalisation géométrique de l'état enrichi fournit une métrique, tandis que le transducteur matériel publie la source à laquelle se rapporte la lecture. Le tenseur résiduel conserve la différence entre les deux publications dans un même espace tensoriel. Sa décomposition permet de distinguer le secteur proportionnel à la métrique et le secteur de trace nulle, et détermine le courant d'échange correspondant à cette répartition.

\section{Données de la carte et construction du tenseur résiduel}
\label{x_reciprocidad:vii:subsec:datos-tensor-residual}

Fixons un état HMT \(x\) et une carte connexe \(\mathcal U\) de sa
réalisation géométrique. Sur \(\mathcal U\), la cotétrade non dégénérée
de la section~\ref{x_reciprocidad:vii:sec:realizacion-cartan} détermine une métrique lisse
\(g=g_x\) de signature \((-+++)\). Nous notons
\(\mathring\nabla\) sa connexion de Levi--Civita et écrivons
\[
 \mathsf G_{\mu\nu}[g]
 :=\operatorname{Ric}_{\mu\nu}[g]-\tfrac12R[g]g_{\mu\nu}.
\]
On distingue ainsi le tenseur d'Einstein \(\mathsf G[g]\) du couplage
gravitationnel \(G_{\rm int}\). La dérivée utilisée dans les bilans
de cette section est \(\mathring\nabla\), conformément à l'équation
métrique effective de la section~\ref{x_reciprocidad:vii:subsec:fuente-metrica-efectiva}.

Nous travaillons dans la carte dimensionnelle \(c_{\rm int}=1\). Nous supposons \(G_{\rm int}(x)>0\) et \(\Lambda_{\rm ref}(x)\) constants sur \(\mathcal U\), et posons \(\kappa_x=8\pi G_{\rm int}(x)\). Le symbole \(x\) identifie l'état qui fournit ces lecteurs ; la différentiation covariante agit sur les points de \(\mathcal U\). La généalogie des couplages demeure dans la section~\ref{x_reciprocidad:v:ant:sec:gravity}.

Soit \(T_b\) le tenseur symétrique lisse publié par le transducteur matériel
de la carte ; dans la lecture baryonique, c'est sa source baryonique. Nous définissons
\begin{equation}
 \mathcal E^{\rm ref}_{\mu\nu}
 :=\kappa_x^{-1}\bigl(\mathsf G_{\mu\nu}[g]
                            +\Lambda_{\rm ref}g_{\mu\nu}\bigr),
 \qquad
 T^{\rm res}_{\mu\nu}:=\mathcal E^{\rm ref}_{\mu\nu}-T_{b,\mu\nu}.
 \label{x_reciprocidad:vii:eq:residual-definicion}
\end{equation}
Le tenseur \(T^{\rm res}\) est le tenseur résiduel effectif, également
noté \(T_{\rm osc}\) lors de l'interprétation du secteur géométrique qui reste
à publier par rapport à \(T_b\). La construction conserve pour
arguments effectifs \(g_x,G_{\rm int},\Lambda_{\rm ref}\) et \(T_b\).

\begin{theorem}[Unicité du résiduel et bilan avec la source matérielle]
\label{x_reciprocidad:vii:thm:residual-unicidad-balance}
Pour les données précédentes, il existe un unique tenseur symétrique \(T^{\rm res}\) qui satisfait \(\mathcal E^{\rm ref}=T_b+T^{\rm res}\). Sa divergence est
\begin{equation}
 \mathring\nabla^\mu T^{\rm res}_{\mu\nu}
 =-\mathring\nabla^\mu T_{b,\mu\nu}.
 \label{x_reciprocidad:vii:eq:balance-residual-material}
\end{equation}
En particulier, si la source matérielle se conserve séparément,
\(\mathring\nabla^\mu T_{b,\mu\nu}=0\), \(T^{\rm res}\)
se conserve également.
\end{theorem}
\begin{proof}
La différence de deux tenseurs satisfaisant l'équation annoncée est nulle ; l'expression~\eqref{x_reciprocidad:vii:eq:residual-definicion} fournit son existence. L'identité de Bianchi contractée et la compatibilité métrique de Levi--Civita donnent
\[
 \mathring\nabla^\mu\mathsf G_{\mu\nu}=0,
 \qquad \mathring\nabla g=0.
\]
Comme \(\kappa_x\) et \(\Lambda_{\rm ref}\) sont constants sur la
carte, \(\mathring\nabla^\mu\mathcal E^{\rm ref}_{\mu\nu}=0\).
Dériver la différence qui définit \(T^{\rm res}\) prouve
\eqref{x_reciprocidad:vii:eq:balance-residual-material}, y compris l'affirmation sous
conservation matérielle séparée.
\end{proof}

La constance spatiale et temporelle des couplages dans la carte fixe
le domaine de ce bilan. Une famille d'états peut publier
des couplages différents dans des cartes différentes ; chaque identité précédente
s'applique à la carte et à l'état qui la spécifient.

\section{Décomposition canonique et courant d'échange}
\label{x_reciprocidad:vii:subsec:traza-corriente-residual}

Nous définissons la fonction de trace et les deux tenseurs
\begin{equation}
 \theta:=g^{\mu\nu}T^{\rm res}_{\mu\nu},\qquad
 T^{\Lambda}_{\mu\nu}:=\tfrac14\theta g_{\mu\nu},\qquad
 T^{0}_{\mu\nu}:=T^{\rm res}_{\mu\nu}-T^{\Lambda}_{\mu\nu}.
 \label{x_reciprocidad:vii:eq:descomposicion-traza-residual}
\end{equation}
L'exposant \(\Lambda\) identifie la composante proportionnelle à la
métrique ; \(\Lambda_{\rm ref}\) est le paramètre de référence de
\eqref{x_reciprocidad:vii:eq:residual-definicion}. Leurs fonctions sont distinctes.

\begin{proposition}[Projecteurs de trace et unicité de la décomposition]
\label{x_reciprocidad:vii:prop:proyectores-traza-residual}
Sur les tenseurs symétriques covariants d'ordre deux, les applications
\[
 P_{\rm tr}(S)=\tfrac14\operatorname{tr}_g(S)g,
 \qquad P_0(S)=S-P_{\rm tr}(S)
\]
satisfont
\begin{equation}
 P_{\rm tr}^2=P_{\rm tr},\quad P_0^2=P_0,\quad
 P_{\rm tr}P_0=P_0P_{\rm tr}=0,\quad P_{\rm tr}+P_0=I.
 \label{x_reciprocidad:vii:eq:proyectores-traza-residual}
\end{equation}
En particulier, \(T^{\rm res}=T^\Lambda+T^0\), \(\operatorname{tr}_gT^0=0\), et cette décomposition est unique.
\end{proposition}
\begin{proof}
En dimension quatre, \(g^{\mu\nu}g_{\mu\nu}=4\). Par conséquent,
\(\operatorname{tr}_g(P_{\rm tr}S)=\operatorname{tr}_gS\), ce qui
prouve l'idempotence de \(P_{\rm tr}\), l'annulation de la trace
de \(P_0S\) et, par développement, les autres identités.
Si \(S=fg+Q\) avec \(\operatorname{tr}_gQ=0\), prendre la trace
donne \(f=\operatorname{tr}_gS/4\). Cela détermine aussi
\(Q=S-fg\) et établit ainsi l'unicité.
\end{proof}

\begin{theorem}[Échange covariant entre les deux secteurs]
\label{x_reciprocidad:vii:thm:intercambio-traza-residual}
Dans le domaine de \eqref{x_reciprocidad:vii:eq:residual-definicion}, soit
\(b_\nu:=\mathring\nabla^\mu T_{b,\mu\nu}\). Le courant
d'échange de trace est la un-forme
\begin{equation}
 \mathcal J^{\rm tr}_\nu
 :=\tfrac14\partial_\nu\theta
 =\mathring\nabla^\mu T^\Lambda_{\mu\nu}.
 \label{x_reciprocidad:vii:eq:corriente-intercambio-traza}
\end{equation}
Les bilans complets sont
\begin{equation}
 \mathring\nabla^\mu T^0_{\mu\nu}
       =-b_\nu-\mathcal J^{\rm tr}_\nu,
 \qquad
 \mathring\nabla^\mu(T^\Lambda+T^0)_{\mu\nu}=-b_\nu.
 \label{x_reciprocidad:vii:eq:balances-sectores-residuales}
\end{equation}
Si \(b=0\), les deux secteurs échangent des courants exactement opposés. Dans ce cas, ils se conservent séparément si et seulement si \(\theta\) est constant dans la carte connexe.
\end{theorem}
\begin{proof}
La compatibilité métrique permet de calculer
\[
 \mathring\nabla^\mu\bigl(\tfrac14\theta g_{\mu\nu}\bigr)
 =\tfrac14g^{\mu\alpha}(\partial_\alpha\theta)g_{\mu\nu}
 =\tfrac14\partial_\nu\theta.
\]
Soustraire cette égalité à \eqref{x_reciprocidad:vii:eq:balance-residual-material} produit
\eqref{x_reciprocidad:vii:eq:balances-sectores-residuales}. Avec \(b=0\), la
conservation séparée équivaut à \(\mathrm d\theta=0\) ; une fonction
lisse de différentielle nulle est constante sur chaque composante connexe.
\end{proof}

La forme de degré un \(\mathcal J^{\rm tr}\) mesure l'échange de la répartition tensorielle. Elle conserve un type distinct de celui du covecteur de transport de la section~\ref{x_reciprocidad:vii:sec:variacion-memoria} et du courant de spin \(\sigma_{IJ}\in\Omega^3(\mathcal U;\Lambda^2E^*)\) de \eqref{x_reciprocidad:vii:eq:corriente-variacional}.

\section{Lien avec la source effective de Cartan--Holst}
\label{x_reciprocidad:vii:subsec:residual-fuente-cartan}

Sur une solution de \eqref{x_reciprocidad:vii:eq:ecuacion-metrica-efectiva}, dans les
mêmes unités \(c=1\), le résiduel a pour expression
\begin{equation}
 T^{\rm res}
 =T^{\rm phys,LC}+U^{\rm phys,spin}-T_b
   +\frac{\Lambda_{\rm ref}-\Lambda_{\rm int}}{\kappa_x}\,g.
 \label{x_reciprocidad:vii:eq:residual-composicion-cartan}
\end{equation}
En effet, substituer
\(\mathsf G+\Lambda_{\rm int}g
  =\kappa_x(T^{\rm phys,LC}+U^{\rm phys,spin})\)
dans \eqref{x_reciprocidad:vii:eq:residual-definicion} donne l'égalité terme à terme.
Lorsque le transducteur matériel publie précisément
\(T_b=T^{\rm phys,LC}\) et \(\Lambda_{\rm ref}=\Lambda_{\rm int}\),
on obtient \(T^{\rm res}=U^{\rm phys,spin}\).
La condition de conservation séparée de \(T_b\) du
théorème~\ref{x_reciprocidad:vii:thm:residual-unicidad-balance} conserve sa fonction
propre : \eqref{x_reciprocidad:vii:eq:balance-metrico-total} établit le bilan de la
source totale, et la répartition entre ses composantes dépend de l'action
matérielle utilisée.

\section{Secteur de trace et accélération des cartes sélectionnées}
\label{x_reciprocidad:vii:subsec:residual-aceleracion-ds}

Le lecteur TRIT de l'état \(x\) publie \(s=\tau(x)\in\{+1,0,-1\}\), et son horloge positive publie
\begin{equation}
 H_*=\frac{2\pi}{108t_0(x)}>0.
 \label{x_reciprocidad:vii:eq:reloj-carta-ds}
\end{equation}
Soit \(h_s\) une métrique tridimensionnelle de courbure sectionnelle constante
\(s\). Les trois cartes de la réalisation sélectionnée s'écrivent
\begin{equation}
 g_s=-\mathrm dt^2+a_s(t)^2h_s,\qquad
 a_+(t)=H_*^{-1}\cosh(H_*t),\quad
 a_0(t)=\exp(H_*t),\quad
 a_-(t)=H_*^{-1}\sinh(H_*t).
 \label{x_reciprocidad:vii:eq:cartas-ds-residual}
\end{equation}
Les cartes fermée et plate sont considérées pour \(t\in\mathbb R\) ; la carte ouverte, pour \(t>0\). Le signe \(s\) et l'horloge \(H_*\) maintiennent leurs antécédents dans l'état qui sélectionne la réalisation.

\begin{lemma}[Courbure des trois cartes de la même horloge]
\label{x_reciprocidad:vii:lem:curvatura-cartas-ds}
Pour les fonctions de \eqref{x_reciprocidad:vii:eq:cartas-ds-residual},
\begin{equation}
 \frac{\ddot a_s}{a_s}=H_*^2,\qquad
 \left(\frac{\dot a_s}{a_s}\right)^2+\frac{s}{a_s^2}=H_*^2,
 \qquad
 \operatorname{Ric}[g_s]=3H_*^2g_s.
 \label{x_reciprocidad:vii:eq:curvatura-cartas-ds}
\end{equation}
Par conséquent, \(R[g_s]=12H_*^2\) et
\(\mathsf G[g_s]=-3H_*^2g_s\).
\end{lemma}
\begin{proof}
Les deux premières égalités s'obtiennent en différentiant les trois fonctions :
dans les cartes fermée et ouverte, on utilise
\(\cosh^2z-\sinh^2z=1\), et dans la carte plate,
\(\dot a_0=H_*a_0\). Pour vérifier l'identification tensorielle,
les coefficients de connexion qui font intervenir la coordonnée temporelle sont
\[
 \Gamma^0{}_{ij}=a_s\dot a_s(h_s)_{ij},\qquad
 \Gamma^i{}_{0j}=(\dot a_s/a_s)\delta^i_j,
\]
et les composantes purement spatiales coïncident avec celles de \(h_s\). La contraction de la courbure donne
\[
 \operatorname{Ric}_{00}=-3\ddot a_s/a_s,\qquad
 \operatorname{Ric}_{0i}=0,\qquad
 \operatorname{Ric}_{ij}
  =(a_s\ddot a_s+2\dot a_s^2+2s)(h_s)_{ij}.
\]
Substituer les deux premières identités de
\eqref{x_reciprocidad:vii:eq:curvatura-cartas-ds} produit
\(\operatorname{Ric}_{00}=-3H_*^2\) et
\(\operatorname{Ric}_{ij}=3H_*^2a_s^2(h_s)_{ij}\).
Prendre la trace et former le tenseur d'Einstein prouve les autres
affirmations, avec la convention de courbure de
\eqref{x_reciprocidad:vii:eq:ecuacion-metrica-efectiva}.
\end{proof}

\begin{theorem}[Densité, pression et signe de l'accélération sélectionnée]
\label{x_reciprocidad:vii:thm:residual-ds-aceleracion}
Dans les cartes précédentes, pour \(\Lambda_{\rm ref}=0\) et \(T_b=0\), le tenseur résiduel est
\begin{equation}
 T^{\rm res}=-\rho_*g_s,\qquad
 \rho_*:=\frac{3H_*^2}{8\pi G_{\rm int}}>0,
 \qquad T^\Lambda=T^{\rm res},\quad T^0=0.
 \label{x_reciprocidad:vii:eq:residual-ds-puro}
\end{equation}
Pour un observateur comobile unitaire \(u=\partial_t\), la densité
et la pression isotrope sont \(\rho_*\) et \(p_*=-\rho_*\).
L'équation d'accélération s'évalue sous la forme
\begin{equation}
 \frac{\ddot a_s}{a_s}
 =-\frac{4\pi G_{\rm int}}3(\rho_*+3p_*)=H_*^2>0.
 \label{x_reciprocidad:vii:eq:aceleracion-residual-ds}
\end{equation}
\end{theorem}
\begin{proof}
Le lemme précédent et \eqref{x_reciprocidad:vii:eq:residual-definicion} donnent \(T^{\rm res}=-(3H_*^2/\kappa_x)g_s\). Puisque \(g_s(u,u)=-1\), \(T^{\rm res}_{\mu\nu}u^\mu u^\nu=\rho_*\). Le projecteur spatial \(q^{\mu\nu}=g_s^{\mu\nu}+u^\mu u^\nu\) satisfait \(q^{\mu\nu}(g_s)_{\mu\nu}=3\), et donc
\[
 p_*:=\tfrac13q^{\mu\nu}T^{\rm res}_{\mu\nu}=-\rho_*.
\]
La trace vaut \(-4\rho_*\), de sorte que
\eqref{x_reciprocidad:vii:eq:descomposicion-traza-residual} produit les deux secteurs
de \eqref{x_reciprocidad:vii:eq:residual-ds-puro}. Enfin,
\(\rho_*+3p_*=-2\rho_*\), et la substitution dans le second membre
de \eqref{x_reciprocidad:vii:eq:aceleracion-residual-ds} donne
\(\kappa_x\rho_*/3=H_*^2\), égal à la valeur géométrique déjà calculée.
La positivité utilise exactement \(G_{\rm int}>0\) et \(H_*>0\).
\end{proof}

Dans cette réalisation, le secteur de trace publie une densité constante et une pression opposée, et l'horloge interne fixe l'accélération positive. Dans une carte générale, \(T^\Lambda\) décrit la composante isotrope ponctuelle, tandis que \(T^0\) conserve le résiduel de trace nulle et tous deux obéissent à l'échange de \eqref{x_reciprocidad:vii:eq:balances-sectores-residuales}. L'identification de ces secteurs aux sources sombres observationnelles exige qu'une même réalisation reproduise conjointement expansion, lentilles, rotation, croissance des structures et fond cosmique. La présente construction fixe le tenseur, ses bilans et l'accélération des cartes sélectionnées qui interviennent dans cette comparaison.

% END OWNER


% BEGIN OWNER /Users/ruben/Documents/New project/output/EXTREMA_MEDIA_RAZON_RECIPROCIDAD_ES_EN_20260918/ARTICULO/ES/source/antecedentes_vii/70_horizontes_e_informacion.tex
% Recuperacion de 09c_entropia_sectorial.tex, 11_confluencia.tex y
% 85_agujero_negro_doble_hoja.tex; originales en fuentes_conservadas/
% horizontes_informacion. Procedencia: desarrollo/HORIZONTES_INFORMACION_FUENTES.md.
% Las referencias a antecedentes se acuerdan con el editor de VII.
% No incorpora una corriente material ni evalua la amplitud torsional.

\chapter{États bipartites purs, entropie réduite et lois d’aire de l’horizon}
\label{x_reciprocidad:vii:sec:horizontes-informacion}

L'état enrichi conserve les routes et leurs registres pendant
l'évolution ; une lecture de sous-système détermine quelle partie de cette information
reste accessible. La construction de cette section utilise les canaux
angulaires de la section~\ref{x_reciprocidad:iv:angular:section}, l'incidence et
l'opérateur d'aire de la section~\ref{x_reciprocidad:v:ant:sec:area}, ainsi que les unités
d'action, d'aire et de température des
sections~\ref{x_reciprocidad:v:ant:sec:gravity}
et~\ref{x_reciprocidad:v:ant:sec:ii-kb-aut-desarrollo}. Ce sont des publications antérieures de la
même construction APP--TRIT--TPK. Chacune transmet ici son domaine,
son orientation, sa normalisation et ses conditions de réalisation.

On distingue trois espaces : le registre binaire de feuillet, les facteurs
tritiques de l'incidence sectorielle et la réalisation géométrique des
extérieurs d'un horizon. La purification, le transport des observables
et la conversion dimensionnelle précisent leurs relations. Cette
distinction conserve la structure qui accompagne chaque valeur entropique.

\section{Registre de feuillet et purification quantique bipartite}
\label{x_reciprocidad:vii:subsec:hojas-purificacion}

Écrivons \(a=A_{\rm rad}\) et \(y=C^*_{\rm rad}\) pour les
coordonnées angulaires déjà construites. Les canaux et leurs probabilités sont
\begin{equation}
 \mathfrak q_\pm=e^{-a\mp y},\qquad
 p_\pm=\frac{\mathfrak q_\pm}{\mathfrak q_++\mathfrak q_-}
       =\frac{e^{\mp y}}{2\cosh y}.
 \label{x_reciprocidad:vii:eq:probabilidades-hoja}
\end{equation}
Les signes indexent ces deux canaux. Dans
\(\mathcal H_{\rm h}=\operatorname{span}\{|+\rangle,|-\rangle\}\),
nous définissons
\[
 \rho_{\rm h}(y)=p_+|+\rangle\langle+|+p_-|-\rangle\langle-|.
\]
L'espace auxiliaire \(\mathcal H_{\rm ar}=\operatorname{span}
\{|\mathrm{adv}\rangle,|\mathrm{ret}\rangle\}\) porte deux étiquettes orthonormées. Une phase \(\chi_0\) étant fixée, l'état conjoint est
\begin{equation}
 |\Psi_{\rm ar}\rangle
 =\sqrt{p_+}|+\rangle\otimes|\mathrm{adv}\rangle
  +e^{i\chi_0}\sqrt{p_-}|-\rangle\otimes|\mathrm{ret}\rangle.
 \label{x_reciprocidad:vii:eq:purificacion-hoja}
\end{equation}

\begin{proposition}[Information accessible dans le registre de feuillet]
\label{x_reciprocidad:vii:prop:purificacion-hoja}
Le vecteur~\eqref{x_reciprocidad:vii:eq:purificacion-hoja} est normalisé et sa trace
partielle sur \(\mathcal H_{\rm ar}\) est \(\rho_{\rm h}(y)\).
Son entropie d'intrication, exprimée en nats, est
\begin{equation}
 s_{\rm h}(y)=\log(2\cosh y)-y\tanh y.
 \label{x_reciprocidad:vii:eq:entropia-hoja}
\end{equation}
L'échange des étiquettes \(+\) et \(-\) conjugue
\(\rho_{\rm h}(y)\) à \(\rho_{\rm h}(-y)\) et conserve
\(s_{\rm h}\).
\end{proposition}
\begin{proof}
Le carré de la norme est \(p_++p_-=1\). L'orthogonalité des étiquettes auxiliaires élimine les termes croisés de la trace partielle. Les carrés des coefficients de Schmidt sont précisément \(p_\pm\). Par~\eqref{x_reciprocidad:vii:eq:probabilidades-hoja}, \(p_+-p_-=-\tanh y\) ; substituer \(\log p_\pm=\mp y-\log(2\cosh y)\) dans \(-\sum_\pm p_\pm\log p_\pm\) démontre la formule. L'échange des signes permute les deux valeurs propres et conserve l'entropie.
\end{proof}

Le facteur commun \(e^{-a}\) se simplifie lors de la normalisation du registre binaire.
La phase relative \(\chi_0\) demeure dans l'état conjoint bien que cette
réduction diagonale ne la publie pas. La réversibilité logique se rapporte
au transport du registre complet, avec ses étiquettes et ses corrélations ;
le retour de phase conserve l'avancée de la mémoire décrite dans la
section~\ref{x_reciprocidad:vii:sec:retorno-memoria-gravedad}. Une mise en œuvre thermique
a en outre sa propre loi d'énergie et de dissipation. En particulier,
\(s_{\rm h}\) mesure ce sous-système binaire et sa bipartition précise.

\section{Purification quantique des secteurs d'incidence}
\label{x_reciprocidad:vii:subsec:purificacion-sectorial}

L'incidence de bord construite précédemment distingue deux ensembles
de dix-huit liens, d'étiquettes sectorielles \(m=90,120\). Chaque lien
porte \(\mathbb C^3\), de base \(|0\rangle,|1\rangle,|2\rangle\)
et d'opérateur d'occupation \(\operatorname{diag}(0,1,2)\). Soient
\(\mathsf N_m\) leurs sommes sectorielles, et
\begin{equation}
 \mathsf N=\mathsf N_{90}+\mathsf N_{120},\qquad
 \mathsf D_\partial=90\mathsf N_{90}+120\mathsf N_{120}.
 \label{x_reciprocidad:vii:eq:ocupaciones-sectoriales}
\end{equation}
Dans la branche positive, \(\gamma_{\rm BI}\) désigne le paramètre de
Barbero--Immirzi déjà obtenu, et \(a_0>0\) son facteur aréal angulaire.
Cette notation distingue le paramètre de la connexion et de ses holonomies.
L'aire est publiée sous la forme
\begin{equation}
 \widehat{\mathcal A}_m=\ell_P^2\gamma_{\rm BI}a_0\mathsf N_m,
 \qquad
 \widehat{\mathcal A}=\widehat{\mathcal A}_{90}
                         +\widehat{\mathcal A}_{120}.
 \label{x_reciprocidad:vii:eq:area-sectores}
\end{equation}

Pour le paramètre fini \(\Delta=\Delta_{\rm mod}\in\mathbb R\) de l'état réduit, nous posons
\begin{align}
 x_m&=e^{-m\Delta},& z_m&=1+x_m+x_m^2,&
 Z(\Delta)&=z_{90}^{18}z_{120}^{18},\\
 \rho_\partial(\Delta)&=Z(\Delta)^{-1}e^{-\Delta\mathsf D_\partial}.
 \label{x_reciprocidad:vii:eq:estado-sectorial}
\end{align}
La trace dans cette section est la trace matricielle ordinaire et les états ont une trace égale à un. Le paramètre \(\Delta\) conserve la sélection et la normalisation de l'état antécédent.

\begin{proposition}[Purification tensorielle de la frontière]
\label{x_reciprocidad:vii:prop:purificacion-sectorial}
Pour un poids réel fini \(w\), soit
\begin{equation}
 |\operatorname{TFD}(w)\rangle
 =\frac{|0\rangle|0\rangle+e^{-w/2}|1\rangle|1\rangle
                   +e^{-w}|2\rangle|2\rangle}
        {\sqrt{1+e^{-w}+e^{-2w}}}.
 \label{x_reciprocidad:vii:eq:tfd-local}
\end{equation}
Le produit
\begin{equation}
 |\Psi_\Delta\rangle
 =|\operatorname{TFD}(90\Delta)\rangle^{\otimes18}
  \otimes|\operatorname{TFD}(120\Delta)\rangle^{\otimes18}
 \label{x_reciprocidad:vii:eq:tfd-sectorial}
\end{equation}
est une purification normalisée de \(\rho_\partial(\Delta)\).
\end{proposition}
\begin{proof}
Dans chaque lien, la somme des carrés des trois coefficients vaut un.
La trace partielle élimine les termes d'indices distincts et restitue
\(\operatorname{diag}(1,e^{-w},e^{-2w})/(1+e^{-w}+e^{-2w})\).
La trace partielle commute avec le produit tensoriel. Multiplier les
trente-six réductions donne l'exponentielle et la fonction de partition
de~\eqref{x_reciprocidad:vii:eq:estado-sectorial}.
\end{proof}

\section{Équation informationnelle et fluctuations surfaciques}
\label{x_reciprocidad:vii:subsec:ecuacion-informacional}

Les moyennes et les variances d'occupation par lien sont
\begin{equation}
 f_m=\frac{x_m+2x_m^2}{z_m},\qquad
 v_m=\frac{x_m+4x_m^2}{z_m}-f_m^2.
 \label{x_reciprocidad:vii:eq:momentos-sectoriales}
\end{equation}
Nous écrivons \(s(\rho)=-\operatorname{Tr}(\rho\log\rho)\), avec \(0\log0=0\). Dans la purification précédente, \(s(\rho_\partial)\) est son entropie d'intrication.

\begin{proposition}[Aire moyenne, entropie et réponse sectorielle]
\label{x_reciprocidad:vii:prop:respuesta-sectorial}
Pour \(\Delta\) réel et fini, on a
\begin{align}
 \langle\mathsf N\rangle_\Delta&=18(f_{90}+f_{120}),\\
 \langle\mathsf D_\partial\rangle_\Delta&=18(90f_{90}+120f_{120}),\\
 \langle\widehat{\mathcal A}\rangle_\Delta
 &=18\ell_P^2\gamma_{\rm BI}a_0(f_{90}+f_{120}),\\
 s(\Delta)&=\log Z(\Delta)+\Delta\langle\mathsf D_\partial\rangle_\Delta,
 \label{x_reciprocidad:vii:eq:entropia-sectorial}\\
 \frac{\mathrm d\langle\mathsf D_\partial\rangle_\Delta}{\mathrm d\Delta}
 &=-\operatorname{Var}_\Delta(\mathsf D_\partial)
  =-18(90^2v_{90}+120^2v_{120}),\\
 \frac{\mathrm ds}{\mathrm d\Delta}
 &=-\Delta\operatorname{Var}_\Delta(\mathsf D_\partial),\\
 \frac{\mathrm d\langle\widehat{\mathcal A}\rangle_\Delta}{\mathrm d\Delta}
 &=-18\ell_P^2\gamma_{\rm BI}a_0(90v_{90}+120v_{120}).
 \label{x_reciprocidad:vii:eq:respuestas-area-entropia}
\end{align}
\end{proposition}
\begin{proof}
Les poids locaux sont \((1,x_m,x_m^2)/z_m\), dont les deux premiers moments donnent \(f_m\) et \(v_m\). L'additivité des opérateurs d'occupation et la factorisation de l'état donnent les trois premières formules. Prendre la trace de \(-\rho_\partial\log\rho_\partial
=\rho_\partial(\log Z+\Delta\mathsf D_\partial)\) donne l'entropie. La différentiation des sommes finies produit \(\partial_\Delta\log Z=-\langle\mathsf D_\partial\rangle_\Delta\) et \(\partial_\Delta^2\log Z
=\operatorname{Var}_\Delta(\mathsf D_\partial)\). Les facteurs indépendants additionnent leurs variances. Dans la dérivée de \eqref{x_reciprocidad:vii:eq:entropia-sectorial}, les deux termes de premier moment s'annulent. Enfin, \(\partial_\Delta f_m=-mv_m\) donne la réponse de l'aire.
\end{proof}

La perte d'uniformité se quantifie en dimension \(d_\partial=3^{36}\) :
\begin{equation}
 I_\partial(\Delta)
 =36\log3-s(\Delta)
 =\mathcal D\bigl(\rho_\partial(\Delta)\Vert I/d_\partial\bigr)\ge0.
 \label{x_reciprocidad:vii:eq:deficit-informacional}
\end{equation}
Ici, \(\mathcal D\) désigne l'entropie relative, définie et justifiée
ci-dessous. Substituer \(\log(I/d_\partial)=-36\log3\,I\) prouve
l'identité. La convexité de \(t\log t\) donne la positivité ; l'égalité
correspond à la distribution uniforme, qui, dans cette famille, se produit
exactement en \(\Delta=0\). Le déficit mesure une information adimensionnelle ;
l'échelle \(\gamma_{\rm BI}a_0\ell_P^2\) mesure l'aire par occupation.

\begin{theorem}[Inversion des occupations et aire complémentaire]
\label{x_reciprocidad:vii:thm:inversion-sectorial}
Soit \(\mathcal R_\partial\) l'involution qui échange
\(|0\rangle\) et \(|2\rangle\) et fixe \(|1\rangle\) dans chaque
facteur. Alors
\begin{align}
 \mathcal R_\partial\mathsf N\mathcal R_\partial^*&=72I-\mathsf N,\\
 \mathcal R_\partial\mathsf D_\partial\mathcal R_\partial^*
 &=7560I-\mathsf D_\partial,\\
 \rho_\partial(-\Delta)&=\mathcal R_\partial\rho_\partial(\Delta)
                          \mathcal R_\partial^*,\\
 s(-\Delta)&=s(\Delta),\\
 \langle\widehat{\mathcal A}\rangle_{-\Delta}
 +\langle\widehat{\mathcal A}\rangle_\Delta
 &=72\gamma_{\rm BI}a_0\ell_P^2,\\
 \mathcal D\bigl(\rho_\partial(\Delta)\Vert\rho_\partial(-\Delta)\bigr)
 &=\Delta\bigl(7560-2\langle\mathsf D_\partial\rangle_\Delta\bigr).
 \label{x_reciprocidad:vii:eq:inversion-sectorial}
\end{align}
\end{theorem}
\begin{proof}
Chaque occupation se transforme par \(n\mapsto2-n\). Dans chaque secteur,
\(\mathsf N_m\mapsto36I-\mathsf N_m\), de sorte que le terme
constant du lecteur pondéré est \(36(90+120)=7560\).
De plus, \(z_m(-\Delta)=e^{2m\Delta}z_m(\Delta)\) et
\(Z(-\Delta)=e^{7560\Delta}Z(\Delta)\). La substitution dans l'
état normalisé prouve la conjugaison. Son spectre conserve l'entropie ;
la transformation de \(\mathsf N\) donne l'aire complémentaire. Soustraire
les logarithmes des deux états et prendre l'espérance dans
\(\rho_\partial(\Delta)\) donne la dernière égalité.
\end{proof}

En \(\Delta=0\), \(s=36\log3\). Lorsque \(\Delta\to+\infty\),
l'occupation nulle concentre l'état et l'entropie tend vers zéro ;
l'inversion envoie cette limite sur l'occupation maximale. Pour \(\Delta>0\),
l'entropie décroît strictement parce que le lecteur pondéré possède plus
d'une valeur propre et que l'état est de rang plein. La configuration
complète, l'aire et l'entropie conservent ainsi des fonctions différenciées.

\section{Loi modulaire et variations finies d'aire}
\label{x_reciprocidad:vii:subsec:ley-modular-finita}

Fixons \(\rho_0=\rho_\partial(\Delta_0)\), avec \(\Delta_0\) réel fini, et maintenons \(\Delta_0,\gamma_{\rm BI},a_0,\ell_P\) constants. Les aires normalisées \(\mathcal A_m^{\rm n}=\gamma_{\rm BI}a_0\mathsf N_m\) expriment l'opérateur modulaire comme
\begin{equation}
 \mathsf K_0=-\log\rho_0
 =c_0I+\sum_{m=90,120}\beta_m\mathcal A_m^{\rm n},
 \quad \beta_m=\frac{m\Delta_0}{\gamma_{\rm BI}a_0},
 \quad c_0=\log Z(\Delta_0).
 \label{x_reciprocidad:vii:eq:operador-modular-areal}
\end{equation}
On a \(\beta_{120}=\tfrac43\beta_{90}\), y compris lorsque les deux sont nuls. Le quotient est \(4/3\) lorsque \(\Delta_0\ne0\). \(\mathsf K_0\) est le logarithme de cet état de référence ; son type le distingue du transport temporel et de la connexion gravitationnelle.

\begin{theorem}[Première loi modulaire et correction finie]
\label{x_reciprocidad:vii:thm:ley-modular-finita}
Pour tout état normalisé \(\sigma\) du même espace, soit
\(\Delta_\sigma\langle B\rangle
=\operatorname{Tr}[(\sigma-\rho_0)B]\). Alors
\begin{equation}
 s(\sigma)-s(\rho_0)
 =\sum_{m=90,120}\beta_m\Delta_\sigma\langle\mathcal A_m^{\rm n}\rangle
  -\mathcal D(\sigma\Vert\rho_0),
 \label{x_reciprocidad:vii:eq:ley-modular-finita}
\end{equation}
où
\(\mathcal D(\sigma\Vert\rho_0)
=\operatorname{Tr}[\sigma(\log\sigma-\log\rho_0)]\)
est finie et non négative. Si \(\sigma(\varepsilon)\) est différentiable,
normalisée et \(\sigma(0)=\rho_0\), on obtient
\begin{equation}
 \left.\frac{\mathrm ds(\sigma(\varepsilon))}{\mathrm d\varepsilon}\right|_0
 =\sum_{m=90,120}\beta_m
   \left.\frac{\mathrm d\langle\mathcal A_m^{\rm n}\rangle_{
                 \sigma(\varepsilon)}}{\mathrm d\varepsilon}\right|_0.
 \label{x_reciprocidad:vii:eq:primera-ley-modular}
\end{equation}
\end{theorem}
\begin{proof}
Par définition,
\(\mathcal D(\sigma\Vert\rho_0)=-s(\sigma)
+\operatorname{Tr}(\sigma\mathsf K_0)\), tandis que
\(s(\rho_0)=\operatorname{Tr}(\rho_0\mathsf K_0)\).
La soustraction et~\eqref{x_reciprocidad:vii:eq:operador-modular-areal} prouvent l'identité
finie, car la normalisation annule le terme \(c_0I\).

Pour la positivité, diagonalisons \(\sigma\) avec les valeurs propres \(p_i\) et les vecteurs \(u_i\), et \(\rho_0\) avec les valeurs propres \(q_j>0\) et les vecteurs \(v_j\). Si \(r_i=\langle u_i,\rho_0u_i\rangle>0\), la concavité du logarithme donne
\[
 \sum_j|\langle u_i,v_j\rangle|^2\log q_j\le\log r_i.
\]
Par conséquent,
\(\mathcal D(\sigma\Vert\rho_0)
\ge\sum_i p_i\log(p_i/r_i)\ge0\).
La dernière inégalité est la convexité de \(t\log t\) avec les poids
\(r_i\), dont la somme vaut un. La fidélité de \(\rho_0\) maintient
ses logarithmes finis ; les valeurs propres nulles de \(\sigma\)
s'interprètent au moyen de \(0\log0=0\).

Pour la formule infinitésimale, au voisinage de l'état fidèle \(\rho_0\),
la dérivée de la trace de \(-\sigma\log\sigma\) est
\(-\operatorname{Tr}[\dot\sigma(\log\rho_0+I)]\).
Cette identité s'obtient à partir de la dérivée spectrale de la trace.
La condition \(\operatorname{Tr}\dot\sigma=0\) laisse
\(\operatorname{Tr}(\dot\sigma\mathsf K_0)\), et substituer
\eqref{x_reciprocidad:vii:eq:operador-modular-areal} conclut la preuve.
\end{proof}

Pour les aires physiques, les coefficients sont \(\beta_m/\ell_P^2\). L'application du transducteur commun \(k_B\) donne \(S=k_Bs\) et conserve intégralement le terme \(k_B\mathcal D\). La dérivée dans \eqref{x_reciprocidad:vii:eq:respuestas-area-entropia} fait varier \(\Delta\) ; la loi \eqref{x_reciprocidad:vii:eq:primera-ley-modular} maintient \(\Delta_0\) fixe et fait varier l'état autour de la référence. Ce sont des opérations distinctes.

\section{Action, énergie de décision et unité surfacique}
\label{x_reciprocidad:vii:subsec:conversion-informacion-area}

Nous fixons une même section orientée d'action et ses sections gravitationnelle et thermique, dans la branche positive. Les antécédents fournissent
\begin{equation}
 \begin{aligned}
 \ell_P^2&=\frac{\hbar G}{c^3},& t_P&=\frac{\ell_P}{c},\\
 T_{108}&=2\pi t_P,& h&=2\pi\hbar,\\
 E_P&=\frac{\hbar}{t_P}=k_B\Theta_{\rm clk}\log3.&&
 \end{aligned}
 \label{x_reciprocidad:vii:eq:unidad-accion-informacion}
\end{equation}
L'égalité thermique conserve le couple et la normalisation précédemment
construits ; lors du changement de leur base dimensionnelle, on transporte conjointement
\(k_B\) et \(\Theta_{\rm clk}\). L'incrément élémentaire d'aire est
\(\delta\mathcal A=\gamma_{\rm BI}a_0\ell_P^2\).

\begin{proposition}[Coefficient de conversion entre décision et aire]
\label{x_reciprocidad:vii:prop:tension-informacion-area}
Le quotient de l'énergie de décision par l'incrément surfacique satisfait
\begin{equation}
 \begin{aligned}
 \mathcal T_{I/A}:=\frac{E_P}{\delta\mathcal A}
 &=\frac{k_B\Theta_{\rm clk}\log3}{\gamma_{\rm BI}a_0\ell_P^2}\\
 &=\frac{c^3}{\gamma_{\rm BI}a_0G t_P}
  =\frac{c^4}{\gamma_{\rm BI}a_0G\ell_P}.
 \end{aligned}
 \label{x_reciprocidad:vii:eq:tension-informacion-area}
\end{equation}
Sa dimension est celle d'une énergie par aire.
\end{proposition}
\begin{proof}
La substitution de \(E_P=\hbar/t_P\) et de
\(\ell_P^2=\hbar G/c^3\) annule la section d'action commune.
L'identité \(\ell_P=ct_P\) donne la dernière expression. La positivité
des facteurs garantit que tous les quotients sont définis.
\end{proof}

La capacité \(81\) de la coupe cubique et la bipartition des trente-six
liens sectoriels correspondent à des domaines différents. Dans le
graphe cubique \(\{0,\ldots,N-1\}^3\), les \(N^2\) droites
parallèles entre deux faces opposées sont disjointes en liens : toute coupe
intercepte les \(N^2\), et une couche transversale atteint cette borne.
Pour \(N=9\), l'état produit uniforme des trits de cette couche
a \(s_{\rm corte}=81\log3\). L'attribution d'une décision du
même cycle à chaque lien publie
\begin{equation}
 S_{\rm corte}^{\rm fis}=81k_B\log3,\qquad
 E_{\rm corte}^{(1)}=81E_P
                  =\Theta_{\rm clk}S_{\rm corte}^{\rm fis}.
 \label{x_reciprocidad:vii:eq:corte-energia-informacion}
\end{equation}
La preuve combine la coupe minimale, l'additivité de l'entropie de l'état produit et~\eqref{x_reciprocidad:vii:eq:unidad-accion-informacion}. Dans l'état sectoriel, l'équation finie pertinente reste \eqref{x_reciprocidad:vii:eq:ley-modular-finita}, avec son entropie relative.

\section{Normalisation cosmologique de l'horizon}
\label{x_reciprocidad:vii:subsec:balance-horizonte}

Pour un rayon aréolaire \(R>0\), la réalisation cosmologique considérée utilise la triade
\begin{equation}
 E_\Lambda(R)=\frac{c^4R}{2G},\qquad
 S_{\rm H}(R)=\frac{k_B\pi R^2}{\ell_P^2},\qquad
 T_{\rm cos}(R)=\frac{\hbar c}{2\pi k_BR}.
 \label{x_reciprocidad:vii:eq:horizonte-cosmologico}
\end{equation}
Le facteur \(2\pi\) appartient à cette normalisation de température.
L'aire \(\mathcal A_{\rm H}=4\pi R^2\) exprime la même entropie
sous la forme \(S_{\rm H}/k_B=\mathcal A_{\rm H}/(4\ell_P^2)\).

\begin{theorem}[Cellule gravitationnelle de demi-action]
\label{x_reciprocidad:vii:thm:celda-media-accion}
Dans~\eqref{x_reciprocidad:vii:eq:horizonte-cosmologico},
\begin{align}
 T_{\rm cos}(R)S_{\rm H}(R)&=E_\Lambda(R),\\
 T_{\rm cos}(\ell_P)S_{\rm H}(\ell_P)
 &=\frac{E_P}{2}=\frac{k_B\Theta_{\rm clk}\log3}{2},\\
 E_\Lambda(\ell_P)t_P&=\frac{\hbar}{2},&
 E_\Lambda(\ell_P)T_{108}&=\frac h2.
 \label{x_reciprocidad:vii:eq:media-accion-horizonte}
\end{align}
En maintenant fixes les constantes de la carte et en faisant varier \(R\),
\begin{equation}
 T_{\rm cos}\,\mathrm dS_{\rm H}=2\,\mathrm dE_\Lambda,
 \qquad S_{\rm H}\,\mathrm dT_{\rm cos}=-\mathrm dE_\Lambda.
 \label{x_reciprocidad:vii:eq:diferencial-horizonte}
\end{equation}
\end{theorem}
\begin{proof}
Multiplier la température par l'entropie donne
\(\hbar cR/(2\ell_P^2)=c^4R/(2G)\).
En \(R=\ell_P\), le résultat est
\(\hbar/(2t_P)=E_P/2\). Les deux durées de
\eqref{x_reciprocidad:vii:eq:unidad-accion-informacion} donnent les relations d'action.
Pour les différentielles,
\(\mathrm dE_\Lambda=c^4\mathrm dR/(2G)\),
\(\mathrm dS_{\rm H}=2\pi k_BR\mathrm dR/\ell_P^2\) et
\(\mathrm dT_{\rm cos}=-T_{\rm cos}\mathrm dR/R\).
La substitution prouve les deux identités.
\end{proof}

Le produit \(E_\Lambda=T_{\rm cos}S_{\rm H}\) exprime une identité d'état. Ses variations radiales conservent les deux termes de~\eqref{x_reciprocidad:vii:eq:diferencial-horizonte} ; une interprétation thermodynamique complète spécifie en outre le générateur temporel, le travail et la réalisation matérielle correspondants. Dans cette carte, action et gravitation fixent l'unité \(\ell_P^2\), tandis que Boltzmann convertit l'information en entropie dimensionnelle. Le transport de sections avec \(G\hbar\) fixe conserve cette unité lorsque la section \(c\) reste fixe, et conserve \(S_{\rm H}/k_B\) à rayon fixe.

Dans la cellule \(R=\ell_P\), \(S_{\rm H}/k_B=\pi\), de sorte que
la multiplicité effective est \(\Omega_{\rm eff}=e^\pi\).
Pour l'involution hyperbolique tritique \(H_{\rm trit}\), de spectre
\(\{1,-1\}\), la diagonalisation donne
\(\operatorname{spec}(e^{\pi H_{\rm trit}})=\{e^\pi,e^{-\pi}\}\).
Cette relation spectrale spécifie la valeur expansive ; une identification
microscopique des deux espaces conserve en outre son application de
réalisation et l'état.

\section{Carte de Kruskal et échange des extérieurs}
\label{x_reciprocidad:vii:subsec:kruskal-dos-hojas}

Fixons un état publié \(x\), son rayon \(R=R(x)>0\) et les
sections positives \(c,G\). Dans la carte qui suit, \(x,R,c,G\)
sont des paramètres fixes ; \(t,r\) sont les coordonnées spatio-temporelles.
La réalisation extérieure des deux feuillets est spécifiée par
\begin{equation}
 f(r)=1-\frac Rr,\qquad
 \mathrm ds^2=-f(r)c^2\mathrm dt^2+f(r)^{-1}\mathrm dr^2
                 +r^2\mathrm d\Omega_2^2,\qquad r>R.
 \label{x_reciprocidad:vii:eq:metrica-exterior}
\end{equation}
Pour le lecteur de feuillet \(\sigma\in\{+1,-1\}\), nous définissons
\begin{align}
 r_*&=r+R\log\left(\frac rR-1\right),\\
 \mathcal U&=-\sigma\exp\left(-\frac{ct-r_*}{2R}\right),&
 \mathcal V&=\sigma\exp\left(\frac{ct+r_*}{2R}\right).
 \label{x_reciprocidad:vii:eq:carta-kruskal}
\end{align}
Les symboles \(\mathcal U,\mathcal V\) sont des coordonnées et se
distinguent des opérateurs de transport des routes.

\begin{theorem}[Réalisation géométrique et involution de feuillet]
\label{x_reciprocidad:vii:thm:kruskal-hojas}
La carte~\eqref{x_reciprocidad:vii:eq:carta-kruskal} réalise les deux feuillets comme les deux extérieurs \(\mathcal U\mathcal V<0\) de l'extension de Kruskal. On a
\begin{equation}
 -\mathcal U\mathcal V=\left(\frac rR-1\right)e^{r/R},\qquad
 \mathrm ds^2=-\frac{4R^3}{r}e^{-r/R}\,
                    \mathrm d\mathcal U\,\mathrm d\mathcal V
                    +r^2\mathrm d\Omega_2^2.
 \label{x_reciprocidad:vii:eq:metrica-kruskal}
\end{equation}
L'involution
\begin{equation}
 \mathcal C_{\rm H}:(\sigma,\mathcal U,\mathcal V)
 \longmapsto(-\sigma,-\mathcal U,-\mathcal V)
 \label{x_reciprocidad:vii:eq:involucion-horizonte}
\end{equation}
est isométrique, fixe le rayon aréolaire, échange les extérieurs et préserve
l'ensemble des horizons \(\mathcal U\mathcal V=0\).
La métrique est régulière en \(r=R\) et satisfait les équations du vide
dans le domaine régulier \(r>0\).
\end{theorem}
\begin{proof}
Multiplier les coordonnées donne la première identité. Comme
\(\mathrm dr_*/\mathrm dr=f^{-1}\),
\[
 \mathrm d\mathcal U=\frac{\mathcal U}{2R}
            (-c\,\mathrm dt+f^{-1}\mathrm dr),\qquad
 \mathrm d\mathcal V=\frac{\mathcal V}{2R}
            ( c\,\mathrm dt+f^{-1}\mathrm dr).
\]
Leur produit reproduit la partie radiale et temporelle de la métrique.
La fonction \(z\mapsto(z-1)e^z\) a pour dérivée \(ze^z>0\)
pour \(z>0\), ce qui détermine \(r\) à partir du produit.
À l'extérieur, on retrouve
\(\sigma=\operatorname{sign}\mathcal V\) et
\(t=(R/c)\log(-\mathcal V/\mathcal U)\).
En \(r=R\), le coefficient radial de la métrique étendue est
\(-4R^2/e\), fini et non nul. L'inversion conserve le
produit, le rayon et \(\mathrm d\mathcal U\mathrm d\mathcal V\),
et échange les deux choix de signe.

Pour vérifier le vide dans la carte extérieure, les composantes mixtes indépendantes du tenseur d'Einstein se réduisent à
\[
 \mathsf G^t{}_t=\mathsf G^r{}_r
 =\frac{rf'+f-1}{r^2},\qquad
 \mathsf G^\theta{}_\theta=\mathsf G^\phi{}_\phi
 =\frac{f'}r+\frac{f''}2.
\]
Avec \(f'=R/r^2\) et \(f''=-2R/r^3\), toutes s'annulent. L'expression régulière~\eqref{x_reciprocidad:vii:eq:metrica-kruskal} prolonge le résultat à travers les horizons dans le domaine \(r>0\).
\end{proof}

Cette réalisation conserve
\(E_{\rm H}=c^4R/(2G)\) et
\(S_{\rm H}=k_B\pi R^2/\ell_P^2\), qui dépendent du rayon aréolaire.
La température \(T_{\rm cos}\) de
\eqref{x_reciprocidad:vii:eq:horizonte-cosmologico} appartient à la réalisation
cosmologique qui y est déclarée. La carte extérieure
\eqref{x_reciprocidad:vii:eq:metrica-exterior} conserve sa propre normalisation temporelle.
L'application géométrique établit l'échange des extérieurs ; la comparaison
avec la formation par effondrement, les spectres de rayonnement ou les modes quasi normaux
utilise en outre leurs conditions physiques respectives.

\section{Transport d'états, d'observables et de raffinements}
\label{x_reciprocidad:vii:subsec:transporte-informacion-horizonte}

Désignons par \(F_x\) la carte des deux extérieurs qui vient d'être construite. Dans chaque carte angulaire régulière, son jacobien temporel et radial est
\begin{equation}
 \left|\frac{\partial(\mathcal U,\mathcal V)}{\partial(t,r)}\right|
 =\frac{cr}{2R^3}e^{r/R}>0.
 \label{x_reciprocidad:vii:eq:jacobiano-kruskal}
\end{equation}
L'identité résulte des deux différentielles de la preuve précédente.
Les autres registres de l'état sont transportés comme coordonnées de
fibre ; la carte conserve son identité et sa mémoire.

\begin{proposition}[Équivalence unitaire des lectures de feuillet et d'extérieur]
\label{x_reciprocidad:vii:prop:transporte-unitario-horizonte}
Soient \(\mu\) une mesure de feuillet et \(\nu=(F_x)_*\mu\) sa mesure
transportée. Le changement de variables définit un opérateur unitaire
\begin{equation}
 W_{\rm H}:L^2(\mu)\longrightarrow L^2(\nu),\qquad
 (W_{\rm H}\psi)(z)=\psi(F_x^{-1}z).
 \label{x_reciprocidad:vii:eq:unitaria-horizonte}
\end{equation}
Avec des densités par rapport à des mesures de coordonnées fixées, la même application incorpore le facteur jacobien correspondant. La conjugaison
\begin{equation}
 \mathfrak I_{\rm H}(B)=W_{\rm H}BW_{\rm H}^*
 \label{x_reciprocidad:vii:eq:mapa-observables-horizonte}
\end{equation}
est un \(*\)-isomorphisme unital complètement positif. Il transporte l'
état, ses espérances et son calcul fonctionnel modulaire. Pour la mesure
symétrique des deux feuillets, il entrelace aussi leurs involutions.
\end{proposition}
\begin{proof}
La définition de la mesure image donne \(\int|\psi\circ F_x^{-1}|^2\,\mathrm d\nu
=\int|\psi|^2\,\mathrm d\mu\). La carte est inversible dans les extérieurs, de sorte que l'isométrie est surjective. La conjugaison unitaire préserve le produit, l'adjoint et l'identité ; dans chaque amplification matricielle, elle conjugue par \(I_n\otimes W_{\rm H}\), ce qui prouve la positivité complète. Pour un état \(\rho\), la cyclicité de la trace donne \(\operatorname{Tr}(W_{\rm H}\rho W_{\rm H}^*
\,\mathfrak I_{\rm H}(B))=\operatorname{Tr}(\rho B)\). Le calcul fonctionnel implique \(g(W_{\rm H}\rho W_{\rm H}^*)=W_{\rm H}g(\rho)W_{\rm H}^*\) dans son domaine, en particulier pour le logarithme et le flot modulaire d'un état fidèle sur son support. L'égalité \(F_x\circ\mathcal C_{\rm hoja}=\mathcal C_{\rm H}\circ F_x\) prouve l'entrelacement des involutions ; la mesure symétrique les réalise unitairement.
\end{proof}

% REMATE_LOCAL_X_HORIZONTE
\Needspace{25\baselineskip}
Si les raffinements conservent \((\sigma,t,r,R)\) et utilisent la même
isométrie sur les registres de fibre, les inclusions
\(j_m^{\rm hoja}\) et \(j_m^{\rm ext}\) satisfont
\begin{equation}
 W_{{\rm H},m+1}j_m^{\rm hoja}
 =j_m^{\rm ext}W_{{\rm H},m}.
 \label{x_reciprocidad:vii:eq:naturalidad-horizonte}
\end{equation}
En effet, le changement de variables agit sur les coordonnées géométriques
conservées et l'inclusion sur les mêmes coordonnées de fibre ; évaluer
les deux compositions donne la même fonction et la même densité transportée.
L'égalité étend l'application au système raffiné sans réinitialiser les registres.

Pour une bipartition \(\mathcal H_A\otimes\mathcal H_B\), un
transport factorisé \(U_A\otimes U_B\) satisfait
\[
 \rho'_A=U_A\rho_AU_A^*,\qquad s(\rho'_A)=s(\rho_A).
\]
La première identité découle de la trace partielle et la seconde du spectre.
Pour une transformation unitaire générale, la formulation équivalente transporte également
l'algèbre du sous-système :
\(\mathcal B(\mathcal H_A)\otimes I\mapsto
U(\mathcal B(\mathcal H_A)\otimes I)U^*\).
On conserve ainsi la même question opérationnelle concernant l'information
accessible. Les égalités d'entropie s'appliquent aux états dont les entropies
pertinentes sont finies.

La relation entre information et gravitation s'exprime par des
opérations distinctes et composables : l'état et la bipartition fixent
l'information accessible ; l'incidence et Barbero--Immirzi fixent la résolution
surfacique ; l'action et la gravitation construisent l'unité \(\ell_P^2\) ;
Boltzmann réalise dimensionnellement l'entropie ; la carte d'horizon
transporte la géométrie et ses observables. Les lois finies conservent
l'entropie relative, et les réalisations géométriques conservent leurs mesures,
leur normalisation temporelle et leurs domaines.

% END OWNER


% BEGIN OWNER /Users/ruben/Documents/New project/output/EXTREMA_MEDIA_RAZON_RECIPROCIDAD_ES_EN_20260918/ARTICULO/ES/source/antecedentes_vii/80_observador_y_respuesta_relacional.tex
% Recuperación y articulación de los propietarios documentados en
% PROCEDENCIA_OBSERVADOR_RESPUESTA.md. Originales conservados íntegros.
% Antecedentes de integración: núcleo APP--TRIT--TPK, realización del
% transporte de borde, lectores areal/volumétrico y reloj energético.
% Las condiciones de representación se enuncian antes de cada uso.

\chapter{Observateur, mémoire et réponse relationnelle}
\label{x_reciprocidad:vii:sec:observador-respuesta}

L'opération de lecture agit sur le transport qui publie l'état
enrichi APP--TRIT--TPK. Son domaine conserve la route, le feuillet, l'orientation,
la retenue et la mémoire ; le résultat accessible à l'observateur est une application
de ce domaine. Cette distinction permet de décrire conjointement trois
opérations : le choix d'un représentant orienté du transport,
l'inscription de son registre dans un appareil et la réponse locale induite
par la frontière globale. Chaque opération transmet une information distincte.
La construction suivante spécifie leurs applications et les conditions qui
permettent de les composer.

\section{Transport de bord et lecture intrinsèque}
\label{x_reciprocidad:vii:sec:observador-ciclo}

La publication finie du transport fournit un cycle orienté
\(C_n\), avec \(n\geq1\), des sommets \(0,\ldots,n-1\) et des arêtes
\(e_j:j\longrightarrow j+1\pmod n\). Nous fixons la réalisation réticulaire
\(L\) du registre, son espace vectoriel
\(W=L\otimes_{\mathbb Z}\mathbb R\), un produit intérieur positif et
le représentant primitif non nul \(v\in L\) de la classe résiduelle
distinguée, dans la chambre déclarée. Les feuillets somme et produit d'APP
fournissent les polarisations du transport ; l'orientation du TRIT et la
mise à jour du TPK précèdent cette publication linéaire. Le cycle,
la réalisation réticulaire et la classe résiduelle sont les données que cette
section reçoit du transport construit précédemment.

Une \(0\)-cochaîne est une fonction sur les sommets à valeurs dans \(W\) ; une \(1\)-cochaîne attribue un vecteur à chaque arête orientée. Le cobord agit par
\[
 (d\Phi)(e_j)=\Phi(j+1)-\Phi(j).
\]
Nous désignons par \(\mathbf1_E\) la cochaîne scalaire constante sur les arêtes. Pour un transport \(\omega\in C^1(C_n;W)\), sa publication cumulée \(\Gamma_\omega\) s'obtient par sommes partielles et interpolation affine ; son défaut de clôture est
\begin{equation}
 B(\omega)=\Gamma_\omega(1)-\Gamma_\omega(0)
 =\sum_{j=0}^{n-1}\omega(e_j).
 \label{x_reciprocidad:vii:eq:observador-defecto}
\end{equation}

\begin{theorem}[Décomposition exacte et mode harmonique du cycle]
\label{x_reciprocidad:vii:thm:observador-descomposicion}
Tout transport possède une unique décomposition
\[
 \omega=d\Phi+u\mathbf1_E,\qquad \Phi(0)=0,
 \qquad u=\frac1n\sum_j\omega(e_j).
\]
En particulier, \(B(d\Phi)=0\) et \(B(\omega)=nu\).
\end{theorem}

\begin{proof}
Soit \(y_j=\omega(e_j)-u\). La somme des \(y_j\) est nulle. Nous définissons \(\Phi(k)=\sum_{j=0}^{k-1}y_j\) ; la valeur pour \(k=n\) coïncide avec \(\Phi(0)=0\), de sorte que la fonction est bien définie sur le cycle, y compris sur son arête finale. On a \(d\Phi(e_j)=y_j\). Si \(d\Phi+u\mathbf1_E=d\Psi+w\mathbf1_E\), la somme des arêtes donne \(nu=nw\). Alors \(d(\Phi-\Psi)=0\) ; la connexité du cycle et la normalisation au sommet initial impliquent \(\Phi=\Psi\). L'identité \(\sum_j(\Phi(j+1)-\Phi(j))=0\) prouve directement l'annulation télescopique.
\end{proof}

\begin{definition}[Transport admissible et lecture normalisée]
\label{x_reciprocidad:vii:def:observador-admisible}
Le transport est admissible lorsque son mode harmonique appartient à
\(\mathbb Rv\). Sur l'espace \(\mathcal T_{\rm adm}\) de ces
transports, nous définissons la lecture globale et sa densité par phase :
\begin{equation}
 \mathfrak a_v(\omega)
 =\frac{\langle B(\omega),v\rangle}{\langle v,v\rangle},
 \qquad
 \mu_v(\omega)=\frac{\mathfrak a_v(\omega)}n.
 \label{x_reciprocidad:vii:eq:observador-lectura}
\end{equation}
On a \(B(\omega)=\mathfrak a_v(\omega)v\).
\end{definition}

\begin{theorem}[Détermination de la lecture intrinsèque]
\label{x_reciprocidad:vii:thm:observador-unicidad}
Le quotient
\(\mathcal T_{\rm adm}/dC^0(C_n;W)\) est unidimensionnel, engendré
par la classe de \(v\mathbf1_E\). L'application \(\mu_v\) est l'unique
fonctionnelle linéaire sur \(\mathcal T_{\rm adm}\) qui satisfait
\[
 F(\omega+d\Phi)=F(\omega),\qquad F(v\mathbf1_E)=1.
\]
\end{theorem}

\begin{proof}
La décomposition précédente écrit chaque transport admissible sous la forme
\(d\Phi+\lambda v\mathbf1_E\). Sa classe modulo les cochaînes exactes
est déterminée par \(\lambda\). La classe de \(v\mathbf1_E\)
est non nulle parce que sa somme est \(nv\neq0\). Par linéarité et
invariance, toute fonctionnelle ayant les caractéristiques indiquées vaut
\(F(\omega)=\lambda F(v\mathbf1_E)=\lambda\). La formule
\eqref{x_reciprocidad:vii:eq:observador-lectura} donne précisément \(\mu_v=\lambda\).
La lecture globale est \(\mathfrak a_v=n\mu_v\).
\end{proof}

\begin{definition}[Observateur de bord]
\label{x_reciprocidad:vii:def:observador-borde}
Une fois fixée la polarisation du transport publié, un observateur de bord
est un couple \(\mathcal O=(\varepsilon,\Psi)\), où
\(\varepsilon\in\{+1,-1\}\) et \(\Psi\in C^0(C_n;W)\).
Son action de lecture est
\begin{equation}
 \omega^{\mathcal O}=\varepsilon\omega+d\Psi.
 \label{x_reciprocidad:vii:eq:observador-accion}
\end{equation}
Le signe fixe l'orientation et la cochaîne fixe la trivialisation. Deux choix de polarisation sont comparés au moyen de leur transport publié ; la relation \eqref{x_reciprocidad:vii:eq:observador-accion} décrit exactement les représentants qui diffèrent par l'orientation et le cobord.
\end{definition}

\begin{proposition}[Covariance de la lecture]
\label{x_reciprocidad:vii:prop:observador-covarianza}
L'action conserve l'admissibilité et vérifie
\[
 B(\omega^{\mathcal O})=\varepsilon B(\omega),\quad
 \mathfrak a_v(\omega^{\mathcal O})=\varepsilon\mathfrak a_v(\omega),
 \quad \mu_v(\omega^{\mathcal O})=\varepsilon\mu_v(\omega).
\]
Les trivialisations de même orientation publient la même clôture.
\end{proposition}

\begin{proof}
La somme du cobord \(d\Psi\) est nulle. Par conséquent, le mode harmonique
se transforme en \(\varepsilon u\), qui appartient toujours à
\(\mathbb Rv\), et le défaut se transforme en \(\varepsilon B\).
Les deux égalités scalaires s'obtiennent par la projection normalisée.
\end{proof}

L'unité de cette lecture est le représentant résiduel \(v\), avec sa normalisation interne. La liberté de trivialisation agit sur le représentant de la cochaîne ; l'information de clôture réside dans sa classe. La description permet de comparer des observateurs en maintenant séparés le transport généalogique, sa publication linéaire et le scalaire qui évalue la classe résiduelle.

\section{Identité de bord et intérieur cellulaire}
\label{x_reciprocidad:vii:sec:observador-stokes}

Soit \(K\) un complexe cellulaire fini orienté de dimension deux,
muni d'une \(2\)-chaîne fondamentale \(c_K\) telle que
\(\partial c_K=[C_n]\). Cette condition exprime que les incidences
intérieures se compensent avec leurs multiplicités et orientations.
Soient \(\Omega\in C^1(K;W)\) une extension de \(\omega\) et
\(F_\Omega=d\Omega\) sa courbure discrète additive.

\begin{theorem}[Identité holographique cellulaire de la lecture]
\label{x_reciprocidad:vii:thm:observador-stokes}
Avec les données précédentes,
\begin{equation}
 B(\omega)=\langle\Omega,\partial c_K\rangle
          =\langle d\Omega,c_K\rangle,
 \qquad
 \mathfrak a_v(\omega)
 =\frac{\langle\langle d\Omega,c_K\rangle,v\rangle}
        {\langle v,v\rangle}.
 \label{x_reciprocidad:vii:eq:observador-stokes}
\end{equation}
En particulier, \(F_\Omega=0\) implique \(\mathfrak a_v(\omega)=0\).
Si la fermeture est non nulle, toute extension du transport dans ce complexe
a une courbure discrète non nulle.
\end{theorem}

\begin{proof}
L'appariement entre chaînes et cochaînes satisfait
\(\langle d\Omega,c_K\rangle=\langle\Omega,\partial c_K\rangle\).
En le développant, chaque arête intérieure apparaît avec la somme des nombres
d'incidence des cellules contiguës ; la condition sur
\(\partial c_K\) laisse exactement le cycle orienté du bord.
La restriction de \(\Omega\) à ce cycle est \(\omega\), dont la somme
est \(B(\omega)\). La projection sur \(v\) donne la seconde formule
et l'implication de platitude.
\end{proof}

L'égalité compare deux évaluations de la même cochaîne étendue : le défaut cumulé du bord et le cobord intégré sur l'intérieur. Sa courbure prend ses valeurs dans \(W\) et utilise le cobord cellulaire additif. La réalisation par une connexion géométrique conserve en outre son propre espace de coefficients et sa loi de transport.

\section{Inscription isométrique du registre dans l'appareil}
\label{x_reciprocidad:vii:sec:observador-aparato}

Soit \(\mathfrak M_n\) un ensemble fini non vide de registres
publiés au niveau \(n\), avec résultat, route, retenue, feuillet et mémoire.
Son espace de registres est
\(\mathcal H_n=\ell^2(\mathfrak M_n)\), de base \(|m\rangle\).
La mise à jour visible \(F_n\) possède un relèvement de registre
\[
 \widehat F_n(m)=(F_n(m),m).
\]
Nous prenons comme registres du niveau suivant un ensemble contenant ces paires. La seconde coordonnée conserve l'antécédent et rend le relèvement injectif. Cette construction spécifie la mémoire retenue par l'inscription, y compris lorsque la mise à jour visible réunit des valeurs.

Dans un espace fini d'appareil \(\mathcal K_{A,n}\), nous fixons une
famille orthonormée \((|a_m\rangle)_{m\in\mathfrak M_n}\) et définissons
\[
 \iota_n:\mathcal H_n\longrightarrow\mathcal K_{A,n},
 \quad \iota_n|m\rangle=|a_m\rangle,
 \qquad
 V_n:\mathcal H_n\longrightarrow\mathcal H_{n+1},
 \quad V_n|m\rangle=|\widehat F_n(m)\rangle.
\]
Soit \(\mathcal P_n=\iota_n\mathcal H_n\) le sous-espace du pointeur.

\begin{theorem}[Compatibilité entre mise à jour et inscription]
\label{x_reciprocidad:vii:thm:observador-entrelazador}
Les applications \(\iota_n\) et \(V_n\) sont des isométries. Il existe une unique
isométrie \(W_n^0:\mathcal P_n\longrightarrow\mathcal K_{A,n+1}\)
qui vérifie
\begin{equation}
 W_n^0\iota_n=\iota_{n+1}V_n.
 \label{x_reciprocidad:vii:eq:observador-cuadrado}
\end{equation}
Si \(\dim\mathcal K_{A,n+1}\geq\dim\mathcal K_{A,n}\),
elle se prolonge en une isométrie de tout \(\mathcal K_{A,n}\).
Au moyen de compléments auxiliaires qui égalisent les dimensions, on obtient
un prolongement unitaire entre les espaces étendus.
\end{theorem}

\begin{proof}
Les images des bases par \(\iota_n\) sont orthonormales ; les
images par \(V_n\) le sont aussi par l'injectivité de
\(\widehat F_n\). Pour \(z=\iota_n\psi\), nous définissons
\(W_n^0z=\iota_{n+1}V_n\psi\). La représentation de \(z\) est unique,
et les trois applications conservent sa norme. La formule détermine l'application sur
une base de \(\mathcal P_n\) et prouve le carré.
Si la dimension d'arrivée est suffisante, le complément orthogonal de
\(W_n^0\mathcal P_n\) contient une famille orthonormale de la taille de
\(\mathcal P_n^\perp\) ; en envoyant une base de ce dernier sur cette famille,
on construit la prolongation. Enfin, on étend les deux espaces à
une même dimension et l'on complète les deux familles prescrites en bases
orthonormales. L'application entre ces bases est unitaire.
\end{proof}

Lorsque le registre inclut les étiquettes d'un instrument fini, cette même inscription réalise son environnement : si les opérateurs \(K_{y,a}:\mathcal H_S\to\mathcal H'_S\) satisfont \(\sum_{y,a}K_{y,a}^*K_{y,a}=I\), l'application
\[
 V_{\rm inst}\psi=\sum_{y,a}K_{y,a}\psi\otimes|y,a\rangle
\]
est isométrique, car sa norme au carré est
\(\sum_{y,a}\langle\psi,K_{y,a}^*K_{y,a}\psi\rangle=\|\psi\|^2\).
L'inscription \(\iota\) transporte chaque étiquette vers le pointeur qui lui
correspond. Le modèle fixe ainsi le registre de la prémesure ; un
dispositif matériel exige en outre son interaction et son étalonnage.


% BEGIN OWNER /Users/ruben/Documents/New project/output/EXTREMA_MEDIA_RAZON_RECIPROCIDAD_ES_EN_20260918/ARTICULO/ES/source/antecedentes_vii/80b_salida_y_condicionamiento.tex
\section{Application de sortie et conditionnement du registre}
\label{x_reciprocidad:vii:sec:salida-condicionamiento}

L’inscription précédente se compose avec le lecteur de l’appareil sur
l’état généalogique complet. Le couplage conserve l’opération qui
produit la corrélation ; l’application de sortie publie le résultat ; le
conditionnement restreint les prolongements compatibles avec ce registre.

Soient \((\Omega_S,\Sigma_S)\) et \((\Omega_M,\Sigma_M)\) les espaces
mesurables d’états généalogiques du système et de l’appareil, obtenus à partir de
leurs cylindres et registres compatibles. Un contexte de mesure s’écrit
\(\mathfrak M_a=(U_a,q_a,\mathcal C_a,m_0)\) : \(m_0\) est la
préparation de l’appareil, \(\mathcal C_a\) fixe les compatibilités et
l’horizon de prolongement, et
\[
 U_a:\Omega_S\times\Omega_M\longrightarrow\Omega_S\times\Omega_M,
 \qquad q_a:\Omega_M\longrightarrow Y_a
\]
sont, respectivement, le couplage mesurable et le lecteur mesurable vers l'espace des résultats \((Y_a,\Sigma_{Y_a})\). Dans un régime réversible, \(U_a\) est bijectif et bimesurable ; dans la réalisation de Hilbert, l'unitarité correspondante est également exigée. La sortie individuelle est la composition
\begin{equation}
 \operatorname{Out}_a(x)
 =q_a\bigl(\operatorname{pr}_M U_a(x,m_0)\bigr).
 \label{x_reciprocidad:vii:eq:out-genealogico}
\end{equation}
Elle est déterminée par l’état complet, la préparation de l’appareil et
le contexte. La description statistique d’une préparation constitue
une autre lecture de l’ensemble des états compatibles.

\begin{proposition}[Mesure des résultats et restriction à la fibre]
\label{x_reciprocidad:vii:prop:medida-salida}
Soit \(\Omega_a\in\Sigma_S\) l’ensemble des états compatibles avec la
préparation et le contexte, muni d’une probabilité \(\mu_a\) sur
\((\Omega_a,\Sigma_S|_{\Omega_a})\).
L’application \(\operatorname{Out}_a|_{\Omega_a}\) est mesurable et produit
la probabilité des résultats
\begin{equation}
 \nu_a(B)=\mu_a\bigl(\operatorname{Out}_a^{-1}(B)\cap\Omega_a\bigr),
 \qquad B\in\Sigma_{Y_a}.
 \label{x_reciprocidad:vii:eq:medida-imagen-salida}
\end{equation}
Pour un résultat \(y\) dont le singleton est mesurable, la fibre \(\Omega_{a,y}=\{x\in\Omega_a:\operatorname{Out}_a(x)=y\}\) est mesurable. Lorsque \(\mu_a(\Omega_{a,y})>0\),
\begin{equation}
 \mu_{a,y}(E)=
 \frac{\mu_a(E\cap\Omega_{a,y})}{\mu_a(\Omega_{a,y})}
 \label{x_reciprocidad:vii:eq:condicionamiento-salida}
\end{equation}
pour \(E\in\Sigma_S|_{\Omega_a}\) définit une probabilité concentrée
sur cette fibre.
\end{proposition}
\begin{proof}
L'inclusion \(x\mapsto(x,m_0)\), le couplage, la projection sur l'appareil et son lecteur sont mesurables ; leur composition l'est également. Les images réciproques conservent les réunions disjointes et l'ensemble total, de sorte que \eqref{x_reciprocidad:vii:eq:medida-imagen-salida} est dénombrablement additive et de masse un. La fibre est l'image réciproque du singleton \(\{y\}\) dans \(\Omega_a\). La restriction de \(\mu_a\) à cette fibre est une mesure positive et dénombrablement additive ; la division par sa masse positive la normalise.
\end{proof}

Pour un préfixe \(s\) de profondeur \(n\), les prolongements postérieurs
au registre forment
\[
 \operatorname{Fut}_{a,y}(s)
 =\{x\in\Omega_{a,y}:x|_n=s\}.
\]
La restriction conserve le préfixe déjà parcouru et le registre du
couplage ; elle modifie la famille des prolongements admis. La sortie
\eqref{x_reciprocidad:vii:eq:out-genealogico} agit sur les états complets ; la loi des
résultats \eqref{x_reciprocidad:vii:eq:medida-imagen-salida} s’obtient en transportant
la mesure de préparation par cette application. La représentation par une
matrice de densité exprime le niveau
statistique de la réalisation quantique et conserve ses conditions de
correspondance avec les fibres généalogiques.

% END OWNER


\section{Algèbre accessible et complément du pointeur}
\label{x_reciprocidad:vii:sec:observador-expectativa}

Nous définissons \(P_m=|a_m\rangle\langle a_m|\), \(P=\sum_mP_m\) et \(P_\perp=I-P\). L'algèbre du sous-espace du pointeur est \(P\mathcal B(\mathcal K_{A,n})P\), dont l'unité est \(P\). Dans celle-ci, le lecteur diagonal est
\[
 \mathcal D_P(A)=\sum_mP_mAP_m.
\]

\begin{proposition}[Espérance du registre et extension unitale]
\label{x_reciprocidad:vii:prop:observador-pinzamiento}
L'application \(\mathcal D_P\) est une espérance conditionnelle complètement positive, préservant la trace et idempotente sur l'algèbre du pointeur, avec \(\mathcal D_P(P)=P\). Son image est \(\bigoplus_m\mathbb CP_m\). Sur l'algèbre de l'appareil complet, l'application
\begin{equation}
 \mathcal E_P(A)=\sum_mP_mAP_m+P_\perp AP_\perp
 \label{x_reciprocidad:vii:eq:observador-complemento}
\end{equation}
est une espérance conditionnelle unitale, complètement positive et
préservant la trace. Son image est
\[
 \left(\bigoplus_m\mathbb CP_m\right)
 \oplus\mathcal B(\mathcal P_n^\perp).
\]
Les deux applications sont des projections orthogonales pour le produit de
Hilbert--Schmidt. Leurs spectres sont contenus dans \(\{0,1\}\).
Lorsqu'il existe des composantes éliminées, la séparation entre le
sous-espace fixe et le noyau vaut un.
\end{proposition}

\begin{proof}
Les projecteurs \(P_m\) sont mutuellement orthogonaux. Dans l'appareil complet, la famille à laquelle on ajoute \(P_\perp\) a pour somme \(I\). Chaque terme \(A\mapsto P_mAP_m\) est complètement positif : après toute amplification matricielle, il agit par congruence avec \(I\otimes P_m\). Il en va de même pour \(P_\perp\). La somme conserve cette propriété. La cyclicité de la trace et la somme des projecteurs prouvent la préservation de la trace dans chaque domaine. L'orthogonalité annule les termes croisés lors de l'itération de l'application, de sorte qu'elle est idempotente. Ses éléments fixes sont précisément les blocs annoncés et l'application est bimodulaire par rapport à cette sous-algèbre. En outre, \(\operatorname{Tr}(A^*P_mBP_m)
 =\operatorname{Tr}((P_mAP_m)^*B)\), d'où le caractère autoadjoint pour Hilbert--Schmidt. Une projection autoadjointe n'a que les valeurs spectrales zéro et un lorsque les deux sous-espaces sont non nuls. S'il ne reste que le sous-espace fixe, l'application est l'identité.
\end{proof}

La somme \(\sum_mP_mAP_m\), appliquée sans le complément à l'ensemble de
l'appareil, envoie \(I\) sur \(P\). Son unité naturelle appartient à l'algèbre
du pointeur. La formule \eqref{x_reciprocidad:vii:eq:observador-complemento} conserve
également le bloc auxiliaire ; une résolution orthonormée de ce bloc
fournit, si on le souhaite, la diagonalisation de l'appareil entier.

Si \(U\) est une transformation unitaire de l'appareil et
\(P'_m=UP_mU^*\), alors
\[
 \mathcal E_{P'}(UAU^*)=U\mathcal E_P(A)U^*.
\]
L'égalité s'obtient en substituant chaque projecteur dans la somme.
De même, pour \(A\in\mathcal B(\mathcal P_n)\), le carré
\eqref{x_reciprocidad:vii:eq:observador-cuadrado} implique
\[
 \mathcal D_{P_{n+1}}(W_n^0A(W_n^0)^*)
 =W_n^0\mathcal D_{P_n}(A)(W_n^0)^*.
\]
En effet, chaque unité matricielle \(|a_m\rangle\langle a_{m'}|\)
est transportée vers l'unité correspondant aux registres distincts
\(\widehat F_n(m),\widehat F_n(m')\) ; le lecteur diagonal conserve
exactement les termes avec \(m=m'\).

\section{Factorisation de la réponse par la frontière effective}
\label{x_reciprocidad:vii:sec:observador-factorizacion}

Soit \(\mathcal X_\infty^{\rm enr}\) le domaine des états compatibles
enrichis et soit
\(b:\mathcal X_\infty^{\rm enr}\longrightarrow B_\infty\) leur trace
globale de frontière. Nous notons \(B_{\rm ef}=\operatorname{im}b\)
la frontière effectivement publiée. Pour un espace de réponses
\(\mathcal V_v\), nous considérons
\(I_v:\mathcal X_\infty^{\rm enr}\longrightarrow\mathcal V_v\).

\begin{theorem}[Critère de réponse relationnelle]
\label{x_reciprocidad:vii:thm:observador-factorizacion}
Il existe une unique application \(\overline I_v:B_{\rm ef}\to\mathcal V_v\)
telle que \(I_v=\overline I_v\circ b\) si et seulement si
\begin{equation}
 b(x)=b(x')\quad\Longrightarrow\quad I_v(x)=I_v(x').
 \label{x_reciprocidad:vii:eq:observador-fibra}
\end{equation}
L'unicité s'étend à tout \(B_\infty\) lorsque \(b\) est
surjective.
\end{theorem}

\begin{proof}
Une composition par \(b\) est constante sur chaque fibre. Réciproquement,
pour \(\beta\in B_{\rm ef}\), on définit
\(\overline I_v(\beta)=I_v(x)\) avec \(b(x)=\beta\).
La condition \eqref{x_reciprocidad:vii:eq:observador-fibra} garantit l'indépendance
vis-à-vis du représentant ; toute valeur de \(\overline I_v\) est déterminée
par un état de cette fibre. Si \(b\) est surjective,
\(B_{\rm ef}=B_\infty\). Sinon, l'équation de composition
détermine uniquement les valeurs sur \(B_{\rm ef}\).
\end{proof}

Ce critère est ensembliste. Pour une factorisation continue, on
conserve en outre les topologies et la condition d'application quotient
de \(b:\mathcal X_\infty^{\rm enr}\to B_{\rm ef}\) ; avec cette condition, la
continuité de \(I_v\) équivaut à celle de son facteur. Cette précision
sépare la détermination d'une lecture de ses propriétés de régularité.

Soit \(\ell_v:\mathcal X_\infty^{\rm enr}\to\mathcal L_v\) la restriction accessible à un voisinage local. Le caractère global effectif se vérifie au moyen de deux états satisfaisant
\begin{equation}
 \ell_v(x)=\ell_v(x'),\qquad I_v(x)\neq I_v(x').
 \label{x_reciprocidad:vii:eq:observador-testigo-global}
\end{equation}
Si \(I_v\) se factorise par \(b\), ces deux états ont des frontières
distinctes. En outre, la réponse ne peut s'écrire comme fonction de
\(\ell_v\), car cette fonction attribuerait la même valeur aux deux états.
La différence de frontières acquiert ainsi un effet opératoire vérifiable.

\section{Poids d'ambivalence et réponse auto-inertielle}
\label{x_reciprocidad:vii:sec:observador-ambivalencia}

Pour la formule générale, nous considérons des lecteurs positifs
\(A_\partial,V_\partial\) et un lecteur énergétique \(Q\)
définis sur le domaine effectif de frontière déclaré.
Une réalisation sur
\(B_{\rm av}\subseteq B_{\rm ef}\) est explicitée ci-après. La mémoire demeure comme coordonnée
\(\mathsf M_\partial(\beta)\) de la frontière. La charnière
d'ambivalence fournit les deux orientations
\[
 r_{\rm av}=\frac{\pi_{\rm HMT}}{\varphi_{\rm HMT}^2},\qquad
 r_{\rm av}^{J}=\frac{\varphi_{\rm HMT}}{\pi_{\rm HMT}^2}.
\]
L'exposant \(J\) exprime l'échange de la paire génératrice.
\subsubsection{Réalisation normalisée des lecteurs de frontière.}
L'incidence modale APP--TRIT--TPK, développée dans la construction régionale, fournit l'enveloppe de chemins de matrice
\[
 F_{\rm av}=\begin{pmatrix}0&1\\1&1\end{pmatrix}.
\]
Le vecteur positif \(v=(1,\varphi_{\rm HMT})^{\mathsf T}\) satisfait \(F_{\rm av}v=\varphi_{\rm HMT}v\). Sa transition normalisée et ses poids stationnaires sont
\begin{equation}
 P_{ij}=\frac{(F_{\rm av})_{ij}v_j}{\varphi_{\rm HMT}v_i},
 \qquad
 P=\begin{pmatrix}0&1\\
 \varphi_{\rm HMT}^{-2}&\varphi_{\rm HMT}^{-1}\end{pmatrix},
 \qquad
 (\mu_{\rm ar},\mu_{\rm vol})
 =\frac{(1,\varphi_{\rm HMT}^{2})}{1+\varphi_{\rm HMT}^{2}}.
 \label{x_reciprocidad:vii:eq:lectores-parry}
\end{equation}
Les lignes de \(P\) ont pour somme un d'après
\(\varphi^{-2}+\varphi^{-1}=1\) ; l'équation stationnaire
\(\mu_{\rm ar}=\mu_{\rm vol}\varphi^{-2}\), jointe à la
normalisation, détermine les poids écrits. Cette mesure de Parry
pondère les chemins de l'enveloppe ; la fréquence chronologique
\((1/3,2/3)\) compte les instants du calendrier modal et conserve
sa fonction distincte.

La réalisation d'amplitude commune de mémoire est définie sur le secteur
\(\mathcal X_{\rm av}\subset\mathcal X_\infty^{\rm enr}\) dans lequel
les lecteurs des deux feuillets ont la factorisation
\begin{equation}
 \mathcal L_{\rm vol}(x)=\mu_{\rm vol}\mathcal M(x),\qquad
 \mathcal L_{\rm ar}(x)=\pi_{\rm HMT}\mu_{\rm ar}\mathcal M(x),
 \qquad \mathcal M(x)>0.
 \label{x_reciprocidad:vii:eq:lectores-amplitud-comun}
\end{equation}
Les deux valeurs appartiennent à la même droite positive d'amplitude de mémoire. La publication surfacique porte la marque régionale \(\pi_{\rm HMT}\) du retour. La lecture sans dimension s'obtient en divisant les deux sections par \(\mathcal L_{\rm vol}(x)\). Ainsi, les sommes des poids s'effectuent entre grandeurs de même type ; une éventuelle réalisation en aires et volumes physiques est ensuite composée avec leurs transducteurs dimensionnels respectifs.

\begin{proposition}[Descente des lecteurs normalisés et de l'horloge énergétique]
\label{x_reciprocidad:vii:prop:lectores-frontera-explicitos}
Sur \(B_{\rm av}=b(\mathcal X_{\rm av})\), les règles
\begin{equation}
 A_\partial(b(x))
 =\frac{\mathcal L_{\rm ar}(x)}{\mathcal L_{\rm vol}(x)}
 =r_{\rm av},\qquad V_\partial(b(x))=1
 \label{x_reciprocidad:vii:eq:lectores-normalizados}
\end{equation}
définissent deux lecteurs positifs indépendants du représentant.
Une orientation \(a\) de la section d'action étant fixée, l'énergie de
décision de \eqref{x_reciprocidad:v:ant:eq:ii-kb-aut-energia} détermine sur les états
\begin{equation}
 q_a(x)=\frac{h_a(x)}{108t_0(x)}
       =\frac{\hbar_a(x)}{t_P(x)}\in\mathcal L_E^+,
 \qquad t_P=\frac{54}{\pi_{\rm HMT}}t_0.
 \label{x_reciprocidad:vii:eq:lector-q-accion}
\end{equation}
Écrivons
\[
 b_{\rm av}:=b|_{\mathcal X_{\rm av}}:
 \mathcal X_{\rm av}\longrightarrow B_{\rm av}.
\]
Il existe un unique lecteur \(Q:B_{\rm av}\to\mathcal L_E^+\) tel que
\(q_a|_{\mathcal X_{\rm av}}=Q\circ b_{\rm av}\)
si et seulement si, pour tous \(x,x'\in\mathcal X_{\rm av}\),
\begin{equation}
 b(x)=b(x')\ \Longrightarrow\
 \frac{h_a(x)}{t_0(x)}=\frac{h_a(x')}{t_0(x')}.
 \label{x_reciprocidad:vii:eq:descenso-reloj}
\end{equation}
En particulier, une carte d'action et de durée fixes satisfait cette
condition et publie \(Q=h_a/(108t_0)\).
\end{proposition}
\begin{proof}
Le quotient de \eqref{x_reciprocidad:vii:eq:lectores-amplitud-comun} annule l'
amplitude commune et vaut
\(\pi_{\rm HMT}\mu_{\rm ar}/\mu_{\rm vol}
=\pi_{\rm HMT}/\varphi_{\rm HMT}^2\).
Le résultat est le même pour tous les représentants d'une fibre
de \(b_{\rm av}\) au sein de \(\mathcal X_{\rm av}\) et est positif.
Diviser les deux lecteurs par le même facteur positif
conserve les quotients pondérés dans lesquels ils interviennent.
Pour l'horloge, le critère de factorisation du
théorème~\ref{x_reciprocidad:vii:thm:observador-factorizacion}, appliqué à \(q_a\),
équivaut exactement à \eqref{x_reciprocidad:vii:eq:descenso-reloj}. Les sections
fixes le satisfont ; des sections variables dont le
quotient descend par \(b_{\rm av}\) le satisfont également. L'unicité résulte de la définition
de \(B_{\rm av}\) comme image effective.
\end{proof}

La valeur du lecteur \(Q\) appartient à la droite d'énergie \(\mathcal L_E=\mathcal L_{\rm act}\otimes\mathcal L_T^*\). Avec \(E_1,E_2\) dans cette même droite, \(E_1E_2/Q^2\) est sans dimension et indépendant de la base énergétique positive. Le symbole \(Q\) de cette section désigne exclusivement cette énergie d'horloge ; la quantité de chaleur de la réalisation de Landauer appartient au bilan thermique correspondant. La factorisation d'amplitude commune spécifie une réalisation concrète des lecteurs généraux. Chaque autre domaine conserve ses applications de frontière et sa propre vérification du passage au quotient. La mémoire complète demeure dans l'histoire enrichie et dans sa publication de frontière, bien qu'elle s'annule dans le quotient normalisé.


On définit
\[
 w_A(\beta)=
 \frac{A_\partial(\beta)r_{\rm av}}
 {A_\partial(\beta)r_{\rm av}+V_\partial(\beta)r_{\rm av}^{J}},
 \qquad w_V(\beta)=1-w_A(\beta).
\]
Dans la réalisation \eqref{x_reciprocidad:vii:eq:lectores-normalizados}, les poids
admettent l'évaluation exacte
\begin{equation}
 w_A=\frac{r_{\rm av}^{\,2}}{r_{\rm av}^{\,2}+r_{\rm av}^{J}}
     =\frac{\pi_{\rm HMT}^{4}}
            {\pi_{\rm HMT}^{4}+\varphi_{\rm HMT}^{5}},\qquad
 w_V=\frac{\varphi_{\rm HMT}^{5}}
            {\pi_{\rm HMT}^{4}+\varphi_{\rm HMT}^{5}}.
 \label{x_reciprocidad:vii:eq:pesos-amplitud-comun}
\end{equation}
La substitution conserve les deux facteurs de la définition : le lecteur
surfacique normalisé et sa pondération ultérieure par \(r_{\rm av}\).
La génération et le marquage régional du retour sont effectués une fois.
Multiplier le numérateur
et le dénominateur par \(\varphi_{\rm HMT}^4\pi_{\rm HMT}^2\)
donne la dernière expression. Ces poids restent constants dans le
secteur d'amplitude commune. À projecteurs et énergies fixés, une
variation de la réponse peut procéder du lecteur \(Q\) ;
le témoin de dépendance globale exige en outre les états de
même restriction locale et de réponses distinctes de
\eqref{x_reciprocidad:vii:eq:observador-testigo-global}.


Soient \(P_A,P_V\) des projections orthogonales complémentaires d'un
espace de réponse \(\mathcal H_v\), et soient \(E_1,E_2>0\)
deux énergies internes fixées pour la comparaison. La réponse est
\begin{equation}
 \overline I_v(\beta)=\frac{E_1E_2}{Q(\beta)^2}
       \bigl(w_A(\beta)P_A+w_V(\beta)P_V\bigr),
 \qquad I_v=\overline I_v\circ b.
 \label{x_reciprocidad:vii:eq:observador-respuesta-ponderada}
\end{equation}
On admet également des énergies qui soient des lecteurs explicites de
\(\beta\) ; dans ce cas, elles sont conservées comme telles dans la formule.

\begin{proposition}[Positivité et covariance de la réponse]
\label{x_reciprocidad:vii:prop:observador-respuesta-positiva}
La réponse \eqref{x_reciprocidad:vii:eq:observador-respuesta-ponderada} est
autoadjointe, définie positive et constante sur les fibres de \(b\).
Ses valeurs propres dans les secteurs non nuls sont
\[
 \lambda_A=E_1E_2Q^{-2}w_A,\qquad
 \lambda_V=E_1E_2Q^{-2}w_V.
\]
L'échange complet \((A_\partial,r_{\rm av},P_A)
 \leftrightarrow(V_\partial,r_{\rm av}^{J},P_V)\) permute les deux termes et conserve l'opérateur. Un transport unitaire de l'espace de réponse conjugue cet opérateur et conserve son spectre avec multiplicités.
\end{proposition}

\begin{proof}
La positivité des lecteurs donne \(0<w_A,w_V<1\). Pour
\(z=z_A+z_V\), avec \(z_A=P_Az\), \(z_V=P_Vz\),
\[
 \langle z,\overline I_vz\rangle
 =\frac{E_1E_2}{Q^2}
   (w_A\|z_A\|^2+w_V\|z_V\|^2)>0
\]
si \(z\neq0\). La décomposition orthogonale prouve les expressions
spectrales. La dépendance exclusive de \(\beta\) donne la constance
sur les fibres et la factorisation. L'échange des triplets
échange simultanément les numérateurs, les poids et les projecteurs, et
conserve donc la somme. Pour un transport unitaire \(U\), la substitution
\(P_A\mapsto UP_AU^*\), \(P_V\mapsto UP_VU^*\) produit
\(U\overline I_vU^*\), avec le même spectre.
\end{proof}

Un échange réalisé dans un même espace par une transformation unitaire
qui envoie \(P_A\) sur \(P_V\) requiert que les deux secteurs aient
la même dimension. L'échange des deux triplets comme étiquettes de la
formule est valable avec les dimensions propres de chaque secteur.

La condition \eqref{x_reciprocidad:vii:eq:observador-testigo-global} exige de comparer les réponses. Par exemple, multiplier simultanément \(A_\partial,V_\partial\) par un même facteur positif conserve les deux poids ; si \(E_1,E_2,Q\) sont également conservés, la réponse reste identique. En revanche, avec des poids et des énergies identiques, remplacer \(Q\) par \(qQ\), \(q>0\), multiplie la réponse par \(q^{-2}\). Si deux états de même restriction locale réalisent ce changement avec \(q\neq1\), ils fournissent le témoin global. Ce sont des conditions opératorielles de séparation ; l'existence des états se vérifie dans le domaine de frontière que l'on réalise.

\section{Connexion de l'état conjoint et lecture d'accélération}
\label{x_reciprocidad:vii:sec:observador-autoinercia}

Dans une réalisation différentiable de l'état conjoint
\(\Gamma=(x,\mathcal O,m)\), soient \(N\) son espace de réalisation,
\(\nabla^{\rm tot}\) sa connexion et \(\pi_S:N\to N_S\) une
projection différentiable vers le sous-système, muni de la connexion
\(\nabla^S\). La connexion et la projection conservent la provenance
des transports et des lecteurs qu'elles réalisent. Le cadre observé s'
obtient au moyen de
\[
 \mathfrak F_{\mathcal O}(\Gamma)
 =\operatorname{Res}_{\mathcal O}(\Gamma,\nabla^{\rm tot}).
\]
Pour une trajectoire deux fois différentiable \(\Gamma(t)\),
nous écrivons \(\gamma_S=\pi_S\circ\Gamma\).

\begin{definition}[Auto-inertie et défaut de projection]
\label{x_reciprocidad:vii:def:observador-autoinercia}
Une trajectoire est auto-inertielle par rapport à la connexion déclarée
lorsque \(\nabla^{\rm tot}_{\dot\Gamma}\dot\Gamma=0\).
Le défaut de sa projection est
\begin{equation}
 \mathcal K_{\pi_S}(\dot\Gamma,\dot\Gamma)
 =\nabla^S_{\dot\gamma_S}\dot\gamma_S
  -d\pi_S\bigl(\nabla^{\rm tot}_{\dot\Gamma}\dot\Gamma\bigr).
 \label{x_reciprocidad:vii:eq:observador-defecto-proyeccion}
\end{equation}
\end{definition}

\begin{proposition}[Accélération du sous-système]
\label{x_reciprocidad:vii:prop:observador-aceleracion}
Pour une trajectoire auto-inertielle,
\[
 a_S=\nabla^S_{\dot\gamma_S}\dot\gamma_S
     =\mathcal K_{\pi_S}(\dot\Gamma,\dot\Gamma).
\]
Dans les coordonnées \(x^B\) de \(N\) et \(y^a\) de \(N_S\), le
défaut est quadratique en la vitesse, avec pour composantes
\begin{equation}
 (\mathcal K_{\pi_S})^a{}_{BC}
 =\partial_B\partial_C\pi_S^a
  +(\Gamma^S)^a{}_{bc}\,
        \partial_B\pi_S^b\partial_C\pi_S^c
  -\partial_D\pi_S^a(\Gamma^{\rm tot})^D{}_{BC}.
 \label{x_reciprocidad:vii:eq:observador-defecto-coordenadas}
\end{equation}
\end{proposition}

\begin{proof}
La règle de dérivation en chaîne donne \(\ddot\gamma_S^a=\partial_B\pi_S^a\ddot\Gamma^B
 +\partial_B\partial_C\pi_S^a\dot\Gamma^B\dot\Gamma^C\). En ajoutant la contribution de \(\nabla^S\) et en soustrayant l'image de l'accélération totale, les termes contenant \(\ddot\Gamma\) s'annulent. Les termes restants sont \eqref{x_reciprocidad:vii:eq:observador-defecto-coordenadas} contractés avec \(\dot\Gamma^B\dot\Gamma^C\). Lorsque l'accélération totale s'annule, on obtient l'égalité annoncée.
\end{proof}

La réponse géométrique se factorise par \(b\) lorsque la connexion,
la projection et la donnée cinématique utilisée dans l'évaluation sont
déterminées par \(\beta=b(x)\), ou lorsque ces données cinématiques sont
maintenues fixes dans la comparaison. Sous ces conditions,
\eqref{x_reciprocidad:vii:eq:observador-defecto-coordenadas} donne le même résultat
pour tous les représentants d'une fibre ; le
théorème~\ref{x_reciprocidad:vii:thm:observador-factorizacion} construit son facteur
unique sur \(B_{\rm ef}\). Si la donnée cinématique varie, elle
fait partie du domaine de la réponse comparée.

La formulation auto-inertielle décrit une trajectoire géodésique de
l'état conjoint et une accélération effective de son sous-système due
au défaut de projection. La réalisation gravitationnelle identifie
ensuite cette connexion, le repère et la réponse à leurs observables
physiques. La covariance de l'observateur, l'inscription de la mémoire et
la factorisation par la frontière conservent ainsi des fonctions complémentaires :
la première détermine la classe de lecture, la deuxième retient la
généalogie et la troisième spécifie la dépendance relationnelle de la
réponse.

% END OWNER


% BEGIN OWNER /Users/ruben/Documents/New project/output/EXTREMA_MEDIA_RAZON_RECIPROCIDAD_ES_EN_20260918/ARTICULO/ES/source/antecedentes_vii/81_propagacion_bidireccional.tex
% Recuperación de 13_accion_interaccion_absorbedor y83_operador_absorbedor.
% Desarrollo del canal cronológico desde el transporte de rutas.
% Procedencia y condiciones en PROCEDENCIA_PROPAGACION_BIDIRECCIONAL.md.
% Los originales y la sección80 permanecen íntegros. Sin contenido de radión.

\chapter{Propagation bidirectionnelle et condition d'absorption à la frontière}
\label{x_reciprocidad:vii:sec:propagacion-bidireccional}

Le transport APP--TRIT--TPK conserve l'orientation de la route et la mémoire qui distingue son aller de son retour. Sur cet antécédent sont construites trois opérations successives : l'action d'une histoire enrichie, la propagation d'une source à travers les transports réalisés et l'imposition d'une condition commune à la frontière. L'inscription du registre dans un espace d'appareil, établie dans la section~\ref{x_reciprocidad:vii:sec:observador-aparato}, permet d'exprimer ces opérations au moyen d'applications linéaires sans confondre la publication observable avec l'état qui conserve ses antécédents.

La direction temporelle, la conjugaison des feuillets et la réflexion spatiale
ont des domaines distincts. Leurs relations sont fixées d'abord sur le registre ;
on spécifie ensuite les entrelacements utilisés par une réalisation
dynamique. Cet ordre détermine quelle information reste dans les propagateurs
et quelle condition introduit le choix d'une frontière absorbante.

\section{Action des routes et involutions du registre}
\label{x_reciprocidad:vii:subsec:accion-bidireccional}

Une route finie contient ses positions, ses directions cardinales
\((d_1,\ldots,d_N)\), ses feuillets \((\mu_1,\ldots,\mu_N)\), avec
\(\mu_t\in\{\Sigma,\Pi\}\), et les incréments de son registre
enrichi. Pour \(\lambda\geq0\), nous définissons
\begin{equation}
 \mathcal A_\lambda(\gamma)
 =\sum_{t=2}^N
 \bigl(1-\delta_{d_t,d_{t-1}}
       +\lambda(1-\delta_{\mu_t,\mu_{t-1}})\bigr).
 \label{x_reciprocidad:vii:eq:accion-ruta-bidireccional}
\end{equation}
Le premier terme compte les changements de direction ; le second compte les changements
de feuillet avec le poids déclaré. Pour \(N\leq1\), la somme est nulle.
L'action générale du secteur conserve en outre ses termes de frontière et de
mémoire :
\begin{equation}
 \mathcal S[\gamma]
 =\sum_{a\subset\gamma}\mathcal L_a
  +\mathcal B_{\rm in}(\partial_-\gamma)
  +\mathcal B_{\rm out}(\partial_+\gamma)
  +\mathcal M(\gamma).
 \label{x_reciprocidad:vii:eq:accion-enriquecida-frontera}
\end{equation}
Les termes de cette expression spécifient une classe variationnelle à
extrémités typées et à mémoire conservée. La dynamique du secteur détermine
\(\mathcal L_a\), les conditions admissibles et le terme \(\mathcal M\).

Sur les registres, nous définissons \(\mathcal C_{\rm r}\) comme l'échange global des feuillets ; \(\mathcal P_{\rm r}\) comme une réflexion de la carte spatiale et de ses directions ; et \(\mathcal T_{\rm r}\) comme l'inversion de l'ordre temporel, le remplacement de chaque direction par son opposée et l'échange des extrémités. La réflexion cardinale est une permutation involutive. L'inversion du registre signé est \(C_M(a_1,\ldots,a_N)=(-a_N,\ldots,-a_1)\). Les feuillets, les quotients et la mémoire sont transportés avec la route complète ; la lecture scalaire d'un mot constitue une opération ultérieure.

\begin{proposition}[Invariance de l'action locale de route]
\label{x_reciprocidad:vii:prop:invariancia-ruta-bidireccional}
Dans le domaine de registres précédent,
\[
 \mathcal A_\lambda(\gamma)
 =\mathcal A_\lambda(\mathcal C_{\rm r}\gamma)
 =\mathcal A_\lambda(\mathcal P_{\rm r}\gamma)
 =\mathcal A_\lambda(\mathcal T_{\rm r}\gamma)
 =\mathcal A_\lambda(\mathcal C_{\rm r}\mathcal P_{\rm r}
                              \mathcal T_{\rm r}\gamma).
\]
\end{proposition}
\begin{proof}
Une permutation globale de l'alphabet des directions conserve l'égalité
entre directions consécutives. L'échange des deux feuillets conserve
l'égalité entre feuillets consécutifs. L'inversion de l'ordre établit une
bijection entre paires adjacentes, et la substitution par les directions opposées
conserve également leurs égalités. Chaque opération conserve les deux dénombrements
de \eqref{x_reciprocidad:vii:eq:accion-ruta-bidireccional} ; leur composition les conserve.
\end{proof}

La proposition détermine la symétrie de la fonctionnelle locale. Pour qu'une
transformation agisse aussi sur un problème dynamique à frontières fixes,
elle doit transporter sa classe de prolongations admissibles et les trois autres
termes de \eqref{x_reciprocidad:vii:eq:accion-enriquecida-frontera}. En particulier, deux
routes ayant la même publication finale peuvent conserver des classes
variationnelles distinctes par leur mémoire.

Une interaction entre deux sections réalisées sur une frontière \(B\)
utilise une application de couplage
\[
 \mathcal I_B:E_B^{(1)}\otimes E_B^{(2)}\longrightarrow E_B^{(12)}
\]
et un terme \(\mathcal S_{\rm int}\) de la même action. Sa covariance s'exprime par
\[
 \rho_{12}(g)\mathcal I_B
 =\mathcal I_B\bigl(\rho_1(g)\otimes\rho_2(g)\bigr).
\]
Lorsqu'on compose les histoires, le registre de la frontière partagée est compté
une seule fois ; les mémoires des deux côtés sont transportées vers la section composée.
Le couplage et cette condition de covariance sont des données de l'interaction
concrète, distinctes des involutions du registre.

\section{Renversement de l'évolution réalisée}
\label{x_reciprocidad:vii:subsec:reversion-evolucion}

La réalisation du registre fournit un espace de Hilbert réel \(\mathcal H_{\mathbb R}\), une structure complexe orthogonale \(J_E\), avec \(J_E^2=-I\), et un hamiltonien autoadjoint \(H\). Son domaine \(\mathcal D(H)\) est invariant par \(J_E\) ; on suppose \(HJ_E=J_EH\) dans ce domaine. Soit \(\mathcal C\) une involution orthogonale qui conserve \(\mathcal D(H)\) et satisfait
\begin{equation}
 \mathcal C J_E\mathcal C^{-1}=-J_E,\qquad
 \mathcal C H\mathcal C^{-1}=H.
 \label{x_reciprocidad:vii:eq:dominio-reversion-evolucion}
\end{equation}
La constante interne d'action \(\hbar_{\rm int}>0\) conserve sa
génération antérieure ; elle normalise ici le groupe d'évolution.

\begin{theorem}[Conjugaison temporelle du groupe]
\label{x_reciprocidad:vii:thm:reversion-evolucion}
Le générateur \(A=-J_EH/\hbar_{\rm int}\), de domaine \(\mathcal D(H)\),
est antiadjoint. Son groupe orthogonal, unitaire par rapport à \(J_E\), satisfait
\begin{equation}
 U(t)=\exp(-J_EHt/\hbar_{\rm int}),\qquad
 \mathcal C U(t)\mathcal C^{-1}=U(-t).
 \label{x_reciprocidad:vii:eq:reversion-grupo-bidireccional}
\end{equation}
\end{theorem}
\begin{proof}
L'orthogonalité et \(J_E^2=-I\) impliquent \(J_E^*=-J_E\).
L'invariance du domaine et la commutation avec \(H\) permettent de traiter
\(H\) comme un opérateur autoadjoint complexe dans l'espace défini par
\(J_E\). Par calcul fonctionnel, \(-J_EH/\hbar_{\rm int}\) est
antiadjoint et engendre le groupe annoncé. En dimension finie, les mêmes
affirmations découlent de la série exponentielle et de \(A^*=-A\).
Les hypothèses sur \(\mathcal C\) donnent
\(\mathcal C A\mathcal C^{-1}=-A\). Conjuguer le groupe, ou son calcul
fonctionnel, produit \(\exp(-tA)=U(-t)\).
\end{proof}

L'application qui réalise l'involution signée \(C_M\) comme \(\mathcal C\)
conserve l'orientation temporelle et le domaine précédent. La conjugaison de
charge d'une représentation de champs conserve en outre ses propres
caractères et opérateurs. L'égalité
\eqref{x_reciprocidad:vii:eq:reversion-grupo-bidireccional} fixe le renversement temporel
du groupe sous les conditions exprimées.

\section{Propagateurs du canal chronologique d'une route}
\label{x_reciprocidad:vii:subsec:green-cronologico}

L'inscription isométrique du registre admet les prolongements unitaires du théorème~\ref{x_reciprocidad:vii:thm:observador-entrelazador}. Fixons une réalisation bilatérale d'un canal chronologique de fibres \(\mathcal H_n\), \(n\in\mathbb Z\), et de transports unitaires \(U_n:\mathcal H_n\to\mathcal H_{n+1}\). Le choix inclut les espaces auxiliaires lorsqu'ils sont nécessaires ; il conserve la route et ses registres précédents. Nous désignons par \(U(n,k)\) le transport de \(k\) à \(n\), avec
\[
 U(k,k)=I,\qquad U(n,k)U(k,l)=U(n,l),\qquad
 U(n,k)^*=U(k,n).
\]
Pour \(n>k\), c'est le produit \(U_{n-1}\cdots U_k\) ; pour \(n<k\),
c'est l'inverse du transport de \(n\) à \(k\).

L'espace des sources \(\mathcal J_c\) contient les sections à support
fini. L'espace des réponses \(\mathcal F\) contient toutes les
sections de ces fibres. L'action quadratique de la route, en unités
internes de cette réalisation, utilise
\[
 (D\psi)_n=\psi_{n+1}-U_n\psi_n,\qquad
 E[\psi]=\sum_n\|(D\psi)_n\|^2
\]
pour des sections à support fini. L'opérateur cinétique avec la convention
de signe \(L=-D^*D\) est
\begin{equation}
 (L\psi)_n=U_n^*\psi_{n+1}-2\psi_n+U_{n-1}\psi_{n-1}.
 \label{x_reciprocidad:vii:eq:diferencia-covariante-bidireccional}
\end{equation}
En effet, \((D^*z)_n=z_{n-1}-U_n^*z_n\), et la substitution produit cette formule. La variation de \(E\) est \(-2\operatorname{Re}\langle\delta\psi,L\psi\rangle\).

\begin{theorem}[Opérateurs de Green exacts de la différence covariante]
\label{x_reciprocidad:vii:thm:green-cronologico}
En écrivant \(x_+=\max(x,0)\), les applications
\begin{align}
 (G_{\rm ret}j)_n&=\sum_k(n-k)_+U(n,k)j_k,\label{x_reciprocidad:vii:eq:green-ret-cronologico}\\
 (G_{\rm adv}j)_n&=\sum_k(k-n)_+U(n,k)j_k\label{x_reciprocidad:vii:eq:green-adv-cronologico}
\end{align}
de \(\mathcal J_c\) dans \(\mathcal F\) satisfont
\[
 LG_{\rm ret}j=LG_{\rm adv}j=j.
\]
La réponse retardée s'annule avant la première position de la source et en celle-ci ; la réponse avancée s'annule après la dernière position et en celle-ci. Ce sont les seules solutions qui s'annulent suffisamment loin dans la direction respective. Pour \(\psi\) à support fini, on a également \(G_{\rm ret}L\psi=G_{\rm adv}L\psi=\psi\). Par rapport à l'accouplement avec des sources à support fini, \(G_{\rm adv}=G_{\rm ret}^{\dagger}\).
\end{theorem}
\begin{proof}
Chaque somme possède autant de termes que la source possède de positions.
La composition des transports réduit l'application de \(L\) à
l'identité scalaire
\[
 (n+1-k)_+-2(n-k)_++(n-1-k)_+=\delta_{nk}.
\]
La fonction \((k-n)_+\) satisfait la même identité. Cela prouve
les deux équations et les supports. Pour l'unicité, une solution
homogène qui s'annule en deux positions consécutives est déterminée
comme nulle à la suivante par \eqref{x_reciprocidad:vii:eq:diferencia-covariante-bidireccional} ;
la récurrence, utilisant des transports inversibles, se prolonge dans les deux
directions. La différence de deux solutions ayant le même support causal
s'annule suffisamment loin et est donc nulle.
Appliquer cet argument à \(\psi\) et à \(G_{\rm ret}L\psi\), ou au
cas avancé, donne les identités à droite sur le support fini.
Enfin, pour deux sources compactes, la somme double de l'appariement
est finie et
\[
 \bigl((k-n)_+U(k,n)\bigr)^*=(k-n)_+U(n,k).
\]
Échanger les indices démontre l'adjonction indiquée.
\end{proof}

Les réponses appartiennent à \(\mathcal F\) ; en général, elles croissent
linéairement hors de la source et n'appartiennent pas à \(\ell^2\).
Le domaine à support fini et son appariement spécifient l'adjonction
de Green utilisée ici. Le résultat construit la propagation du canal
chronologique à partir de son action ; une réalisation de champs spatiaux conserve
en outre son opérateur, sa géométrie et son domaine de propagation propres.

Soit \(C_0\) une involution isométrique, linéaire ou antilinéaire, de \(\mathcal H_0\). Le transport depuis la fibre de référence construit
\[
 C_n=U(-n,0)C_0U(0,n):\mathcal H_n\longrightarrow\mathcal H_{-n},
 \qquad (\mathcal C\psi)_n=C_{-n}\psi_{-n}.
\]
On a \(C_{-n}C_n=I\) et \(C_{n+1}U_n=U_{-n-1}^*C_n\), par composition. En particulier, \(\mathcal C^2=I\), et le changement \((n,k)\mapsto(-n,-k)\) dans les sommes précédentes prouve
\begin{equation}
 \mathcal C L=L\mathcal C,\qquad
 \mathcal C G_{\rm ret}\mathcal C^{-1}=G_{\rm adv}.
 \label{x_reciprocidad:vii:eq:intercambio-green-cronologico}
\end{equation}
Lorsque \(C_0\) est la réalisation déclarée de l'involution du registre, ces formules transportent cette involution à tout le canal, sans identifier le retour de phase à une réinitialisation de la mémoire.

\section{Échange causal et décomposition de la réponse}
\label{x_reciprocidad:vii:subsec:green-realizacion-campos}

Dans une réalisation de champs, soit
\(L:\mathcal D(L)\subset\mathcal H_\Phi\to\mathcal H_J\)
un opérateur possédant des opérateurs de Green retardé et avancé uniques sur les domaines de
sources admissibles. Nous notons \(J^+(S)\) et \(J^-(S)\) les
domaines causaux futur et passé d'un support \(S\). Les conditions sont
\[
 \begin{gathered}
 LG_{\rm ret}=LG_{\rm adv}=I,\\
 \operatorname{supp}(G_{\rm ret}j)\subset J^+(\operatorname{supp}j),\\
 \operatorname{supp}(G_{\rm adv}j)\subset J^-(\operatorname{supp}j).
 \end{gathered}
\]
Soient \(C_\Phi,C_J\) des involutions isométriques, toutes deux linéaires ou toutes deux
antilinéaires, qui transportent ces domaines,
inversent leurs orientations causales et satisfont
\(LC_\Phi=C_JL\). On conserve un appariement de Green pour lequel
\(G_{\rm adv}=G_{\rm ret}^{\dagger}\).

\begin{proposition}[Conjugaison des opérateurs de Green et parités]
\label{x_reciprocidad:vii:prop:green-paridades}
Sous les conditions précédentes,
\[
 C_\Phi G_{\rm ret}C_J^{-1}=G_{\rm adv}.
\]
Sur l'espace réel des opérateurs de réponse, définissons
\[
 \Theta(K)=C_\Phi K C_J^{-1},\qquad
 \mathsf P_\pm(K)=\tfrac12(K\pm\Theta(K)).
\]
Ce sont des projections complémentaires pour l'involution \(\Theta\), et
\begin{equation}
 G_{\rm sym}=\tfrac12(G_{\rm ret}+G_{\rm adv}),\qquad
 G_{\rm rad}=\tfrac12(G_{\rm ret}-G_{\rm adv})
 \label{x_reciprocidad:vii:eq:descomposicion-green}
\end{equation}
satisfont \(\Theta(G_{\rm sym})=G_{\rm sym}\),
\(\Theta(G_{\rm rad})=-G_{\rm rad}\),
\(G_{\rm sym}^{\dagger}=G_{\rm sym}\) et
\(LG_{\rm rad}=0\). Le support de la réponse symétrique est contenu
dans la réunion des deux domaines causaux.
\end{proposition}
\begin{proof}
L'opérateur conjugué résout l'équation de source, puisque
\[
 LC_\Phi G_{\rm ret}C_J^{-1}=C_JLG_{\rm ret}C_J^{-1}=I.
\]
L'inversion de l'orientation envoie son support dans le domaine avancé ;
son unicité l'identifie à \(G_{\rm adv}\). Les involutions donnent
\(\Theta^2=I\) ; développer \((I\pm\Theta)^2\) et
\((I+\Theta)(I-\Theta)\) démontre les affirmations sur les projections.
L'adjonction de Green démontre la symétrie et la différence des deux
équations de source donne \(LG_{\rm rad}=0\). La somme de deux réponses
a un support contenu dans l'union de leurs supports.
\end{proof}

Si les involutions sont antilinéaires, \(\Theta\) est linéaire sur les scalaires réels ; c'est l'espace d'opérateurs utilisé dans la proposition. La conjugaison du groupe temporel, la condition de support et l'adjonction de Green sont trois propriétés distinctes, transmises par les applications de réalisation correspondantes.

La réduction causale de l’actualisation bilatérale, développée dans la
section sur la conservation et la récupération, identifie une autre origine précise
du retard : l’élimination d’une composante de mémoire qui continue de
participer à l’évolution conjointe. Son noyau discret conserve
l’ordre des interactions et la contribution de la mémoire initiale.
Les opérateurs de Green construits ici résolvent, en revanche, une
équation de propagation avec des conditions de support. Les deux constructions
peuvent se composer lorsqu’on fixe l’application de réalisation qui identifie
leurs domaines et leurs opérateurs ; leurs conditions de causalité respectives
demeurent explicites.

\section{Projection des courants absorbés et variation de l'action}
\label{x_reciprocidad:vii:subsec:proyeccion-absorcion}

Soit \(\gamma_\partial:\mathcal H_\Phi\to\mathcal H_\partial\)
la lecture de frontière, distincte de la route \(\gamma\). Nous définissons
\[
 \mathcal A_\partial=\gamma_\partial G_{\rm rad}.
\]
La fenêtre finie fixe l'espace des sources et la lecture de frontière ;
les opérateurs de Green conservent leur domaine de réponses. À un niveau fini avec
des produits scalaires fixés, son inverse de
Moore--Penrose détermine
\begin{equation}
 P_{\rm abs}=I-\mathcal A_\partial^+\mathcal A_\partial.
 \label{x_reciprocidad:vii:eq:proyector-absorcion}
\end{equation}

\begin{theorem}[Condition d'absorption et réponse de bord]
\label{x_reciprocidad:vii:thm:proyeccion-absorcion}
\(P_{\rm abs}\) est la projection orthogonale sur \(\ker\mathcal A_\partial\). Pour toute source \(j\), la source \(j_{\rm abs}=P_{\rm abs}j\) satisfait
\begin{equation}
 \gamma_\partial G_{\rm ret}j_{\rm abs}
 =\gamma_\partial G_{\rm adv}j_{\rm abs}.
 \label{x_reciprocidad:vii:eq:igualdad-borde-absorbido}
\end{equation}
C'est la source absorbée la plus proche de \(j\) pour le produit intérieur
déclaré. Il existe des sources absorbées non nulles exactement lorsque
\(\operatorname{rank}\mathcal A_\partial<\dim\mathcal H_J\).
\end{theorem}
\begin{proof}
La décomposition orthogonale
\(\mathcal H_J=\ker\mathcal A_\partial\oplus
 \operatorname{im}\mathcal A_\partial^*\) montre que
\(\mathcal A_\partial^+\mathcal A_\partial\) est l'identité
sur le second terme et zéro sur le premier. Son complément est le
projecteur annoncé. Par définition,
\(2\mathcal A_\partial j_{\rm abs}
 =\gamma_\partial(G_{\rm ret}-G_{\rm adv})j_{\rm abs}=0\).
Pour \(k\in\ker\mathcal A_\partial\), la décomposition orthogonale donne
\[
 \|j-k\|^2=\|(I-P_{\rm abs})j\|^2+\|P_{\rm abs}j-k\|^2.
\]
L'unique minimum correspond à \(k=P_{\rm abs}j\). L'équivalence finale est le théorème du rang.
\end{proof}

Dans l'appariement de la source et du champ, l'action symétrique est
\(\mathscr S[j]=\tfrac12\langle j,G_{\rm sym}j\rangle\).
L'adjonction déclarée rend cette quantité réelle. Pour des variations réelles,
\[
 \mathrm d\mathscr S[j](h)
 =\operatorname{Re}\langle h,G_{\rm sym}j\rangle.
\]
L'action restreinte aux sources absorbées s'écrit sur l'espace ambiant comme
\begin{equation}
 \mathscr S_{\rm abs}[j]
 =\tfrac12\langle P_{\rm abs}j,G_{\rm sym}P_{\rm abs}j\rangle,
 \quad
 \mathrm d\mathscr S_{\rm abs}[j](h)
 =\operatorname{Re}\langle P_{\rm abs}h,G_{\rm sym}P_{\rm abs}j\rangle.
 \label{x_reciprocidad:vii:eq:variacion-restringida-absorcion}
\end{equation}
Lorsque l'accouplement identifie champ et source par Riesz, le gradient ambiant est \(P_{\rm abs}G_{\rm sym}P_{\rm abs}j\). La réponse de champ reste \(G_{\rm sym}j_{\rm abs}\) ; le projecteur du gradient exprime la restriction des variations. Les formules se démontrent en développant le terme quadratique de \(j+th\) et en utilisant le caractère autoadjoint. Un terme matériel \(\mathscr S_{\rm matter}\) ajoute sa propre différentielle.

\section{Condition absorbante explicite dans une fenêtre chronologique}
\label{x_reciprocidad:vii:subsec:absorcion-momentos}

Le canal de la sous-section~\ref{x_reciprocidad:vii:subsec:green-cronologico} permet
d'évaluer cette condition sans introduire de réponse cible. Soit une
source à support dans \(0,\ldots,N\), avec \(N\geq1\). Transportons
ses composantes vers la fibre de référence :
\[
 \widehat j_k=U(0,k)j_k,\qquad
 M_0(j)=\sum_{k=0}^N\widehat j_k,\qquad
 M_1(j)=\sum_{k=0}^Nk\widehat j_k.
\]
L'identité \((n-k)_+-(k-n)_+=n-k\) donne
\begin{equation}
 (G_{\rm ret}-G_{\rm adv})j\big|_n
 =U(n,0)\bigl(nM_0(j)-M_1(j)\bigr).
 \label{x_reciprocidad:vii:eq:diferencia-green-momentos}
\end{equation}

\begin{proposition}[Absorption et deux moments transportés]
\label{x_reciprocidad:vii:prop:absorcion-momentos}
Si la trace évalue le champ en deux positions distinctes \(a,b\),
l'égalité des traces avancée et retardée équivaut à
\(M_0(j)=M_1(j)=0\). Dans une fibre de dimension \(d\), l'espace
de telles sources est de dimension \((N-1)d\). Pour \(N\geq2\),
il contient les sources non nulles de composantes transportées
\((z,-2z,z,0,\ldots,0)\), \(z\neq0\).
\end{proposition}
\begin{proof}
L'inversibilité de \(U(a,0)\) et de \(U(b,0)\) réduit la condition
au bord à \(aM_0-M_1=0\) et \(bM_0-M_1=0\).
Soustraire les équations donne \((a-b)M_0=0\), puis \(M_1=0\).
L'application des deux moments est surjective : pour des valeurs prescrites
\((z_0,z_1)\), il suffit de prendre
\(\widehat j_0=z_0-z_1\), \(\widehat j_1=z_1\) et les autres
composantes nulles. Son rang est \(2d\) ; sa nullité est
\((N+1)d-2d\). La source indiquée possède deux moments nuls.
\end{proof}

Dans le repère transporté, écrivons
\[
 B_N=\begin{pmatrix}1&1&\cdots&1\\0&1&\cdots&N\end{pmatrix}.
\]
Le projecteur absorbant est évalué par
\begin{equation}
 P_N=\bigl[I-B_N^{\mathsf T}(B_NB_N^{\mathsf T})^{-1}B_N\bigr]
       \otimes I_{\mathcal H_0}.
 \label{x_reciprocidad:vii:eq:proyector-momentos-cronologicos}
\end{equation}
Les deux lignes de \(B_N\) sont indépendantes pour \(N\geq1\), de sorte que l'inverse existe. Le transport unitaire entre repères envoie ce projecteur sur les fibres d'origine. Pour \(N=2\),
\[
 P_2=\frac16
 \begin{pmatrix}1&-2&1\\-2&4&-2\\1&-2&1\end{pmatrix}
 \otimes I_{\mathcal H_0}.
\]
Cette construction produit une classe absorbée non nulle et son opération de
projection dans le canal chronologique. Son amplitude et son identification à un
courant matériel conservent les règles du secteur physique correspondant.

\section{Raffinement et continuité de la condition de bord}
\label{x_reciprocidad:vii:subsec:refinamiento-absorcion}

Soient \(R_m^\Phi,R_m^J,R_m^\partial\) les applications de raffinement des
champs, des sources et des frontières. Supposons qu'elles entrelacent les deux opérateurs de Green et la
trace. Par composition, on obtient
\begin{equation}
 R_m^\Phi G_{{\rm sym},m}=G_{{\rm sym},m+1}R_m^J,\qquad
 R_m^\partial\mathcal A_{\partial,m}
 =\mathcal A_{\partial,m+1}R_m^J.
 \label{x_reciprocidad:vii:eq:refinamiento-lectura-absorcion}
\end{equation}
La seconde égalité envoie une source absorbée sur une autre source absorbée.
La naturalité des projecteurs conserve en outre la structure
orthogonale des réalisations.

\begin{proposition}[Naturalité de la projection absorbante]
\label{x_reciprocidad:vii:prop:naturalidad-proyeccion-absorcion}
Pour une isométrie de sources \(R=R_m^J\), on a \(P_{{\rm abs},m+1}R=RP_{{\rm abs},m}\) exactement lorsque
\[
 R(\ker\mathcal A_{\partial,m})
       \subset\ker\mathcal A_{\partial,m+1},\qquad
 R((\ker\mathcal A_{\partial,m})^\perp)
       \subset(\ker\mathcal A_{\partial,m+1})^\perp.
\]
Avec ces conditions à tous les niveaux, les projecteurs définissent une projection orthogonale dans la complétion de la limite inductive isométrique des espaces de sources.
\end{proposition}
\begin{proof}
On décompose chaque source en ses composantes dans le noyau et dans son
complément orthogonal. Les deux inclusions impliquent qu'appliquer la
projection avant ou après le raffinement conserve respectivement la
première composante et annule la seconde. Réciproquement, l'égalité des
opérateurs appliquée à chaque composante donne les deux inclusions.
Dans la limite inductive, l'égalité permet de définir un opérateur sur l'
union des niveaux sans dépendre du représentant. Tous les projecteurs
ont une norme au plus égale à un, de sorte que l'opérateur s'étend à la
complétion. Ses identités d'autoadjonction et d'idempotence, valables sur
l'union dense, passent par continuité à la complétion.
\end{proof}

Pour prolonger également une réponse de Green à une complétion, on
conserve ses domaines et les estimations de continuité correspondantes.
La propagation chronologique à support fini et la factorisation de la
réponse par l'image du bord, démontrée dans
\ref{x_reciprocidad:vii:thm:observador-factorizacion}, maintiennent ainsi leurs domaines
respectifs pendant le raffinement.

La réalisation physique de l'émission et de l'absorption compose ces opérateurs avec les courants matériels, l'action et les conditions causales du secteur. L'entrelacement doit conserver l'involution de route, l'accouplement de la source et du champ, l'action transportée et les conditions de frontière à chaque niveau. La construction précédente apporte la réversion du groupe, les opérateurs de Green du canal chronologique, la décomposition temporelle et une projection absorbante explicite ; les réalisations spatiales emploient leurs propres applications et conditions de causalité.

% END OWNER

\chapter{Réalisation de l’action et comparaison métrologique}

% BEGIN OWNER /Users/ruben/Documents/New project/output/EXTREMA_MEDIA_RAZON_RECIPROCIDAD_ES_EN_20260918/ARTICULO/ES/source/antecedentes_i/sections/revision_planck_contraste.tex
\section{Confrontation de la constante de Planck réduite}
\label{x_reciprocidad:sec:revision-planck-contraste}

Le Système international en vigueur depuis le 20 mai 2019 fixe
exactement \(h=6.62607015\times10^{-34}\,\mathrm{J\,s}\)
\cite{x_reciprocidad:bipm2018}. Sa constante réduite est donc
\[
 \hbar_{\mathrm{SI}}=\frac{6.62607015}{2\pi}\,10^{-34}\,
 \mathrm{J\,s}
 =1.0545718176461563912\ldots\times10^{-34}\,\mathrm{J\,s}.
\]
Les points de suspension correspondent au développement d'une expression
exacte, et non à une incertitude expérimentale de \(h\) dans le SI actuel.
Cette valeur de référence est utilisée ici après l'évaluation des
lecteurs de HMT ; sa fonction est comparative.

Les deux coordonnées du retour de l'article sont
\begin{align*}
 H_5&=1.0545718176461544696\ldots,\\
 \eta_{\mathrm{ret}}&=1.0545718169150390611\ldots.
\end{align*}
Dans l'identification de \(\mathcal U_S\) à \(\mathrm{J\,s}\), la
différence relative de la section antérieure par rapport à
\(\hbar_{\mathrm{SI}}\) est d'environ \(-1.82\times10^{-15}\).
Celle de la section postérieure est d'environ \(-6.93\times10^{-10}\).
La seconde différence contient l'effet du retour utilisé par le
correcteur électronique ; elle n'équivaut pas à l'erreur d'évaluation de la
première section.

La précision de ces comparaisons permet de reconnaître la proximité
quantitative de la construction de \(H_5\) et, simultanément, de distinguer
proximité et identité exacte. Dans le présent manuscrit, les expressions
présentées ne fournissent pas une égalité exacte avec \(h/(2\pi)\).
Le résultat démontré sur \(\Sigma_H\) est sa décennie \(-34\) ;
le résultat analytique sur \(H_5\) est l'évaluation du caractère
déclaré. L'identification physique est examinée dans cette carte et avec les
différences explicites que nous venons de donner. La valeur SI n'intervient
ni dans la définition de la norme déterminantale ni dans la récurrence régionale.

% END OWNER

